% Options for packages loaded elsewhere
\PassOptionsToPackage{unicode}{hyperref}
\PassOptionsToPackage{hyphens}{url}
\PassOptionsToPackage{space}{xeCJK}
%
\documentclass[
  11pt,
]{ctexart}
\usepackage{amsmath,amssymb}
\usepackage{iftex}
\ifPDFTeX
  \usepackage[T1]{fontenc}
  \usepackage[utf8]{inputenc}
  \usepackage{textcomp} % provide euro and other symbols
\else % if luatex or xetex
  \usepackage{unicode-math} % this also loads fontspec
  \defaultfontfeatures{Scale=MatchLowercase}
  \defaultfontfeatures[\rmfamily]{Ligatures=TeX,Scale=1}
\fi
\usepackage{lmodern}
\ifPDFTeX\else
  % xetex/luatex font selection
  \ifXeTeX
    \usepackage{xeCJK}
    \setCJKmainfont[]{WenQuanYi Zen Hei}
          \fi
  \ifLuaTeX
    \usepackage[]{luatexja-fontspec}
    \setmainjfont[]{WenQuanYi Zen Hei}
  \fi
\fi
% Use upquote if available, for straight quotes in verbatim environments
\IfFileExists{upquote.sty}{\usepackage{upquote}}{}
\IfFileExists{microtype.sty}{% use microtype if available
  \usepackage[]{microtype}
  \UseMicrotypeSet[protrusion]{basicmath} % disable protrusion for tt fonts
}{}
\makeatletter
\@ifundefined{KOMAClassName}{% if non-KOMA class
  \IfFileExists{parskip.sty}{%
    \usepackage{parskip}
  }{% else
    \setlength{\parindent}{0pt}
    \setlength{\parskip}{6pt plus 2pt minus 1pt}}
}{% if KOMA class
  \KOMAoptions{parskip=half}}
\makeatother
\usepackage{xcolor}
\usepackage[margin=2.3cm]{geometry}
\usepackage{color}
\usepackage{fancyvrb}
\newcommand{\VerbBar}{|}
\newcommand{\VERB}{\Verb[commandchars=\\\{\}]}
\DefineVerbatimEnvironment{Highlighting}{Verbatim}{commandchars=\\\{\}}
% Add ',fontsize=\small' for more characters per line
\usepackage{framed}
\definecolor{shadecolor}{RGB}{248,248,248}
\newenvironment{Shaded}{\begin{snugshade}}{\end{snugshade}}
\newcommand{\AlertTok}[1]{\textcolor[rgb]{0.94,0.16,0.16}{#1}}
\newcommand{\AnnotationTok}[1]{\textcolor[rgb]{0.56,0.35,0.01}{\textbf{\textit{#1}}}}
\newcommand{\AttributeTok}[1]{\textcolor[rgb]{0.13,0.29,0.53}{#1}}
\newcommand{\BaseNTok}[1]{\textcolor[rgb]{0.00,0.00,0.81}{#1}}
\newcommand{\BuiltInTok}[1]{#1}
\newcommand{\CharTok}[1]{\textcolor[rgb]{0.31,0.60,0.02}{#1}}
\newcommand{\CommentTok}[1]{\textcolor[rgb]{0.56,0.35,0.01}{\textit{#1}}}
\newcommand{\CommentVarTok}[1]{\textcolor[rgb]{0.56,0.35,0.01}{\textbf{\textit{#1}}}}
\newcommand{\ConstantTok}[1]{\textcolor[rgb]{0.56,0.35,0.01}{#1}}
\newcommand{\ControlFlowTok}[1]{\textcolor[rgb]{0.13,0.29,0.53}{\textbf{#1}}}
\newcommand{\DataTypeTok}[1]{\textcolor[rgb]{0.13,0.29,0.53}{#1}}
\newcommand{\DecValTok}[1]{\textcolor[rgb]{0.00,0.00,0.81}{#1}}
\newcommand{\DocumentationTok}[1]{\textcolor[rgb]{0.56,0.35,0.01}{\textbf{\textit{#1}}}}
\newcommand{\ErrorTok}[1]{\textcolor[rgb]{0.64,0.00,0.00}{\textbf{#1}}}
\newcommand{\ExtensionTok}[1]{#1}
\newcommand{\FloatTok}[1]{\textcolor[rgb]{0.00,0.00,0.81}{#1}}
\newcommand{\FunctionTok}[1]{\textcolor[rgb]{0.13,0.29,0.53}{\textbf{#1}}}
\newcommand{\ImportTok}[1]{#1}
\newcommand{\InformationTok}[1]{\textcolor[rgb]{0.56,0.35,0.01}{\textbf{\textit{#1}}}}
\newcommand{\KeywordTok}[1]{\textcolor[rgb]{0.13,0.29,0.53}{\textbf{#1}}}
\newcommand{\NormalTok}[1]{#1}
\newcommand{\OperatorTok}[1]{\textcolor[rgb]{0.81,0.36,0.00}{\textbf{#1}}}
\newcommand{\OtherTok}[1]{\textcolor[rgb]{0.56,0.35,0.01}{#1}}
\newcommand{\PreprocessorTok}[1]{\textcolor[rgb]{0.56,0.35,0.01}{\textit{#1}}}
\newcommand{\RegionMarkerTok}[1]{#1}
\newcommand{\SpecialCharTok}[1]{\textcolor[rgb]{0.81,0.36,0.00}{\textbf{#1}}}
\newcommand{\SpecialStringTok}[1]{\textcolor[rgb]{0.31,0.60,0.02}{#1}}
\newcommand{\StringTok}[1]{\textcolor[rgb]{0.31,0.60,0.02}{#1}}
\newcommand{\VariableTok}[1]{\textcolor[rgb]{0.00,0.00,0.00}{#1}}
\newcommand{\VerbatimStringTok}[1]{\textcolor[rgb]{0.31,0.60,0.02}{#1}}
\newcommand{\WarningTok}[1]{\textcolor[rgb]{0.56,0.35,0.01}{\textbf{\textit{#1}}}}
\usepackage{longtable,booktabs,array}
\usepackage{calc} % for calculating minipage widths
% Correct order of tables after \paragraph or \subparagraph
\usepackage{etoolbox}
\makeatletter
\patchcmd\longtable{\par}{\if@noskipsec\mbox{}\fi\par}{}{}
\makeatother
% Allow footnotes in longtable head/foot
\IfFileExists{footnotehyper.sty}{\usepackage{footnotehyper}}{\usepackage{footnote}}
\makesavenoteenv{longtable}
\setlength{\emergencystretch}{3em} % prevent overfull lines
\providecommand{\tightlist}{%
  \setlength{\itemsep}{0pt}\setlength{\parskip}{0pt}}
\setcounter{secnumdepth}{-\maxdimen} % remove section numbering
% ===== Mini-OJ 教程 PDF 彩色样式(pandoc -H 注入;ctexart + xelatex/tectonic)=====
% ---- 配色 ----
\usepackage{xcolor}
\definecolor{primary}{HTML}{2C6E91}      % 主色:蓝
\definecolor{primarydark}{HTML}{1B4763}
\definecolor{primarylight}{HTML}{AFC8DA}
\definecolor{examgreen}{HTML}{2E8B57}    % 考试向:绿
\definecolor{engorange}{HTML}{C8772E}    % 工程向:橙
\definecolor{codebg}{HTML}{F6F8FA}
\definecolor{calloutbg}{HTML}{EEF4FA}
\definecolor{gold}{HTML}{D4A017}         % ★

% ---- 链接 ----
\usepackage{hyperref}
\hypersetup{colorlinks=true,linkcolor=primary,urlcolor=primary,citecolor=primary}

% ---- 章节着色 ----
\usepackage{titlesec}
\titleformat{\section}{\Large\bfseries\color{primary}}{\color{primarylight}\thesection}{0.6em}{}[{\color{primarylight}\titlerule[1.2pt]}]
\titleformat{\subsection}{\large\bfseries\color{primarydark}}{\color{primary}\thesubsection}{0.5em}{}
\titleformat{\subsubsection}{\normalsize\bfseries\color{primarydark}}{}{0em}{}
\titlespacing*{\section}{0pt}{16pt}{8pt}

% ---- 页眉页脚 ----
\usepackage{fancyhdr}
\pagestyle{fancy}\fancyhf{}
\renewcommand{\headrulewidth}{0.6pt}
\renewcommand{\footrulewidth}{0pt}
\fancyhead[L]{\small\color{primary}\nouppercase{\leftmark}}
\fancyhead[R]{\small\color{primarylight}\textbf{Mini-OJ 教程}}
\fancyfoot[C]{\small\color{primary}\thepage}

% ---- 代码:换行 + 圆角彩框 ----
\usepackage{fvextra}
\DefineVerbatimEnvironment{Highlighting}{Verbatim}{breaklines,breakanywhere,fontsize=\small,commandchars=\\\{\}}
\usepackage{tcolorbox}
\tcbuselibrary{skins,breakable}
\renewenvironment{Shaded}{%
  \begin{tcolorbox}[breakable,enhanced,boxrule=0pt,leftrule=2.5pt,arc=2pt,
    colback=codebg,colframe=primary,left=8pt,right=6pt,top=4pt,bottom=4pt]}{%
  \end{tcolorbox}}

% ---- 引用块 → 彩色提示框 ----
\renewenvironment{quote}{%
  \begin{tcolorbox}[breakable,enhanced,boxrule=0pt,leftrule=3pt,arc=1pt,
    colback=calloutbg,colframe=primary,left=8pt,right=8pt,top=5pt,bottom=5pt]}{%
  \end{tcolorbox}}

% ---- 表格 ----
\usepackage{booktabs}

% ---- 考试向 / 工程向 徽章 ----
\newcommand{\KS}{\,\colorbox{examgreen}{\textcolor{white}{\scriptsize\bfseries 考试向}}\,}
\newcommand{\GC}{\,\colorbox{engorange}{\textcolor{white}{\scriptsize\bfseries 工程向}}\,}

% ---- 符号字形(★☐①… 映射到 Dingbats/amssymb,并上色)----
\usepackage{pifont}\usepackage{amssymb}\usepackage{newunicodechar}
\newunicodechar{★}{\textcolor{gold}{\ding{72}}}
\newunicodechar{☆}{\textcolor{gold!60}{\ding{73}}}
\newunicodechar{✘}{\textcolor{red!65!black}{\ding{55}}}
\newunicodechar{✓}{\textcolor{examgreen}{\ding{51}}}
\newunicodechar{✔}{\textcolor{examgreen}{\ding{51}}}
\newunicodechar{☐}{$\square$}
\newunicodechar{①}{\ding{172}}\newunicodechar{②}{\ding{173}}\newunicodechar{③}{\ding{174}}
\newunicodechar{④}{\ding{175}}\newunicodechar{⑤}{\ding{176}}\newunicodechar{⑥}{\ding{177}}

% ---- 段落 ----
\setlength{\parindent}{0pt}\setlength{\parskip}{4pt}

% ---- 彩色标题页 ----
\renewcommand{\maketitle}{%
\begin{titlepage}
\thispagestyle{empty}
\vspace*{2.2cm}
{\setlength{\fboxsep}{0pt}\noindent\colorbox{primary}{\parbox{\textwidth}{\centering\color{white}%
  \vspace{22pt}{\Huge\bfseries 用 Mini-OJ 项目学 Java}\\[10pt]{\Large 判题机三代演进}\vspace{22pt}}}}
\vfill
\noindent{\large\bfseries\color{primarydark} winbeau · Mini-OJ 教学项目}\par\vspace{4pt}
\noindent{\large\color{primary} 古法手撸 · 教学优先 · 2026-06}\par\vspace{16pt}
\noindent{\color{examgreen}\rule{36pt}{8pt}}\ {\small 考试向(拟合试卷,先上手)}\par\vspace{4pt}
\noindent{\color{engorange}\rule{36pt}{8pt}}\ {\small 工程向(项目深度,选学)}\par\vspace{4pt}
\noindent{\color{gold}\ding{72}\ding{72}\ding{72}}\ {\small 考试相关度(历年真题)}
\vspace{2.2cm}
\end{titlepage}}
\ifLuaTeX
  \usepackage{selnolig}  % disable illegal ligatures
\fi
\IfFileExists{bookmark.sty}{\usepackage{bookmark}}{\usepackage{hyperref}}
\IfFileExists{xurl.sty}{\usepackage{xurl}}{} % add URL line breaks if available
\urlstyle{same}
\hypersetup{
  hidelinks,
  pdfcreator={LaTeX via pandoc}}

\title{用 Mini-OJ 项目学 Java}
\author{winbeau}
\date{2026-06}

\begin{document}
\maketitle

{
\setcounter{tocdepth}{2}
\tableofcontents
}
\hypertarget{mini-oj-ux5b66ux4e60ux8defux7ebfux56feux5224ux9898ux673aux4e09ux4ee3ux6f14ux8fdb}{%
\section{Mini-OJ
学习路线图(判题机三代演进)}\label{mini-oj-ux5b66ux4e60ux8defux7ebfux56feux5224ux9898ux673aux4e09ux4ee3ux6f14ux8fdb}}

\begin{quote}
环境:Ubuntu + nvim + JDK,\textbf{古法手撸}(纯 \texttt{javac/java} +
\texttt{Makefile}),简单、教学优先。 主线:用一个 Mini-OJ 串起 Java
课程,并贯穿一条\textbf{判题机三代演进}。
课程调整:\textbf{第13章网络(Socket)不考,已删除};\textbf{新增第15章泛型与集合}(并入
M5a)。
\end{quote}

\hypertarget{ux5224ux9898ux673aux4e09ux4ee3ux4e3bux7ebf}{%
\subsection{判题机三代主线}\label{ux5224ux9898ux673aux4e09ux4ee3ux4e3bux7ebf}}

\begin{longtable}[]{@{}
  >{\raggedright\arraybackslash}p{(\columnwidth - 6\tabcolsep) * \real{0.2500}}
  >{\raggedright\arraybackslash}p{(\columnwidth - 6\tabcolsep) * \real{0.2500}}
  >{\raggedright\arraybackslash}p{(\columnwidth - 6\tabcolsep) * \real{0.2500}}
  >{\raggedright\arraybackslash}p{(\columnwidth - 6\tabcolsep) * \real{0.2500}}@{}}
\toprule\noalign{}
\begin{minipage}[b]{\linewidth}\raggedright
代
\end{minipage} & \begin{minipage}[b]{\linewidth}\raggedright
里程碑
\end{minipage} & \begin{minipage}[b]{\linewidth}\raggedright
判题方式
\end{minipage} & \begin{minipage}[b]{\linewidth}\raggedright
关键技术
\end{minipage} \\
\midrule\noalign{}
\endhead
\bottomrule\noalign{}
\endlastfoot
\textbf{第一代} & M1--M3(已完成) & Java 对象在 JVM 内模拟 &
\texttt{Solution.solve()} + 多态 \\
\textbf{第二代} & M4 & 反射工厂 + 单文件配置 & \texttt{Class.forName} +
\texttt{config.txt} \\
\textbf{第三代} & M5a & 外部 C++ 判题机真编译真运行 &
\texttt{ProcessBuilder} + \texttt{setrlimit} \\
\end{longtable}

\hypertarget{ux8003ux8bd5ux76f8ux5173ux5ea6-ux5386ux5e74ux771fux9898ux5f3aux76f8ux5173}{%
\subsection{考试相关度(★ =
历年真题强相关)}\label{ux8003ux8bd5ux76f8ux5173ux5ea6-ux5386ux5e74ux771fux9898ux5f3aux76f8ux5173}}

依据新疆大学《面向对象程序设计 B / JAVA
程序设计》历年真题(2016-2017、2019-2020 A 卷 +
复习提纲、笔试模拟)。题型固定:\textbf{单选 + 判断 + 读程序 + 程序填空 +
编程},各约 20 分。

\begin{longtable}[]{@{}
  >{\raggedright\arraybackslash}p{(\columnwidth - 6\tabcolsep) * \real{0.2500}}
  >{\raggedright\arraybackslash}p{(\columnwidth - 6\tabcolsep) * \real{0.2500}}
  >{\raggedright\arraybackslash}p{(\columnwidth - 6\tabcolsep) * \real{0.2500}}
  >{\raggedright\arraybackslash}p{(\columnwidth - 6\tabcolsep) * \real{0.2500}}@{}}
\toprule\noalign{}
\begin{minipage}[b]{\linewidth}\raggedright
模块(章)
\end{minipage} & \begin{minipage}[b]{\linewidth}\raggedright
★相关度
\end{minipage} & \begin{minipage}[b]{\linewidth}\raggedright
真题题型(原题级)
\end{minipage} & \begin{minipage}[b]{\linewidth}\raggedright
里程碑
\end{minipage} \\
\midrule\noalign{}
\endhead
\bottomrule\noalign{}
\endlastfoot
类与对象/构造/封装/static/this(Ch4) & ★★★ & 编程大题(Vehicle
类)、读程序、选择/判断 & M2 \\
继承/重写/多态/抽象类/接口/异常(Ch5-7) & ★★★ & 编程大题(抽象类
Shape-\textgreater Circle/Rectangle)、读程序、选择 & M3 \\
Swing GUI + 事件(Ch9) & ★★★ & 编程(平方/按钮窗口)+ 填空(六按钮
GridLayout)双大题、选择 & M5c \\
集合 ArrayList/LinkedList/HashMap(Ch15) & ★★ &
读程序(Iterator)、填空(查询入 ArrayList) & M5a \\
JDBC 查询/更新(Ch11) & ★★ & 填空(查询入 List)、原题(score\_t / student
表) & M5b \\
多线程 Thread/Runnable/join/synchronized(Ch12) & ★★ & 填空(线程类)、选择
& M6a \\
String/StringBuffer/Random/Math(Ch8) & ★★ &
读程序(\texttt{==}/equals)、选择 & M4 \\
数组/运算符/流程(Ch2-3) & ★ & 选择、读程序 & M1 \\
异常 try-catch(Ch7) & ★ & 选择、判断 & M3 \\
文件 File/流/序列化(Ch10) & ★ &
少量选择/判断(File.mkdir、System.in.read、对象串行化) & M4/M5a \\
环境/编译运行(Ch1) & ☆ & 少量选择 & M0 \\
网络 Socket(Ch13) & ✘ 已删 & 历年仅 1 道判断 & --- \\
\end{longtable}

\begin{quote}
\textbf{工程化加料}(反射工厂、外部 C++
判题机、ProcessBuilder、ProblemService/泛型架构、线程池/队列、MVC
解耦)\textbf{不是考点},属项目深度。
\end{quote}

\hypertarget{ux4e24ux6b65ux8d70ux6bcfux4e2aux91ccux7a0bux7891ux5148ux62dfux5408ux8bd5ux5377ux518dux5de5ux7a0bux5316}{%
\subsection{两步走(每个里程碑先「拟合试卷」再「工程化」)}\label{ux4e24ux6b65ux8d70ux6bcfux4e2aux91ccux7a0bux7891ux5148ux62dfux5408ux8bd5ux5377ux518dux5de5ux7a0bux5316}}

本项目偏工程化。为让学生\textbf{先考得了、上得了手},每个里程碑拆成两步:

\begin{enumerate}
\def\labelenumi{\arabic{enumi}.}
\tightlist
\item
  \textbf{第一步 {[}preliminary ·
  拟合试卷{]}}:按真题难度做最小可运行版(教材/考试写法),快速有参与感、直接覆盖考点。
\item
  \textbf{第二步 {[}工程化{]}}:在此基础上引入反射 / 外部 C++ 判题机 /
  泛型架构 / 数据库 / 线程池 / MVC 解耦等工程方法。
\end{enumerate}

\begin{quote}
各章文档开头已标 \texttt{★相关度} 与 \texttt{{[}第一步{]}/{[}第二步{]}}
的具体内容。\textbf{考前只过第一步即可覆盖真题};想做工程深度再上第二步。
\end{quote}

\hypertarget{ux91ccux7a0bux7891ux603bux89c8}{%
\subsection{里程碑总览}\label{ux91ccux7a0bux7891ux603bux89c8}}

\begin{longtable}[]{@{}
  >{\raggedright\arraybackslash}p{(\columnwidth - 10\tabcolsep) * \real{0.1667}}
  >{\raggedright\arraybackslash}p{(\columnwidth - 10\tabcolsep) * \real{0.1667}}
  >{\raggedright\arraybackslash}p{(\columnwidth - 10\tabcolsep) * \real{0.1667}}
  >{\raggedright\arraybackslash}p{(\columnwidth - 10\tabcolsep) * \real{0.1667}}
  >{\raggedright\arraybackslash}p{(\columnwidth - 10\tabcolsep) * \real{0.1667}}
  >{\raggedright\arraybackslash}p{(\columnwidth - 10\tabcolsep) * \real{0.1667}}@{}}
\toprule\noalign{}
\begin{minipage}[b]{\linewidth}\raggedright
里程碑
\end{minipage} & \begin{minipage}[b]{\linewidth}\raggedright
章节
\end{minipage} & \begin{minipage}[b]{\linewidth}\raggedright
主题
\end{minipage} & \begin{minipage}[b]{\linewidth}\raggedright
★
\end{minipage} & \begin{minipage}[b]{\linewidth}\raggedright
代
\end{minipage} & \begin{minipage}[b]{\linewidth}\raggedright
产物
\end{minipage} \\
\midrule\noalign{}
\endhead
\bottomrule\noalign{}
\endlastfoot
M0 & Ch1 & 环境 + HelloWorld & ☆ & --- & 能编译运行的
\texttt{Main.java} \\
M1 & Ch2--3 & 单题判题器(硬编码) & ★ & 一 & 数组+流程控制判 A+B \\
M2 & Ch4 & 对象建模 & ★★★ & 一 & \texttt{oj.core} 包 + jar + javadoc \\
M3 & Ch5--7 & 接口/继承/多态/异常 & ★★★ & 一 & \texttt{Judge} 接口 +
\texttt{AbstractJudge} + 子类 \\
M4 & Ch8 + Ch10(单文件) & 第二代:反射 + 单文件配置 & ★★ Ch8 / 工程 反射
& 二 &
\texttt{ConfigFile}/\texttt{SingleFileProblemLoader}/\texttt{JudgeFactory} \\
M5a & Ch15 + Ch10(工业级) & 第三代:泛型集合 + C++ 判题机 & ★★ Ch15 /
工程 C++ & 三 &
\texttt{ProblemLoader(List)}/\texttt{ProblemRepository}/\texttt{MachineJudge} \\
M5b & Ch11 & 数据库衔接(FS/DB 分工) & ★★ & --- &
\texttt{Db}/\texttt{ProblemDao}/\texttt{SubmissionDao}/\texttt{ProblemService} \\
M5c & Ch9 & Swing 大前端(MVC,无网络) & ★★★ & --- &
\texttt{OjFrame}/\texttt{OjController} \\
M6a & Ch12 & 多线程并发判题(收官) & ★★ 线程 / 工程 队列 & 三 &
\texttt{JudgeTask}/\texttt{JudgeQueue}/\texttt{JudgeWorker} \\
\end{longtable}

\hypertarget{ux6700ux7ec8ux67b6ux6784ux65e0ux7f51ux7edcux672cux5730ux5927ux524dux7aef}{%
\subsection{最终架构(无网络,本地大前端)}\label{ux6700ux7ec8ux67b6ux6784ux65e0ux7f51ux7edcux672cux5730ux5927ux524dux7aef}}

\begin{verbatim}
        Swing 桌面客户端(Ch9)   ——MVC——
          | 选题/选语言/贴源码          |
          v                            v
   多线程判题队列(Ch12) --调用--> C++ 判题机(judge)
          |                       真编译/真运行/真 setrlimit
          v                            | 一行 JSON
     MySQL(Ch11)  <--元数据/历史---- ProblemService
          ^                            |
          +-- 文件系统:problems/<id>/{config,N.in,N.out} + submissions/ 源码
\end{verbatim}

\hypertarget{ux9879ux76eeux76eeux5f55ux6700ux7ec8ux5f62ux6001}{%
\subsection{项目目录(最终形态)}\label{ux9879ux76eeux76eeux5f55ux6700ux7ec8ux5f62ux6001}}

\begin{verbatim}
mini-oj/
+-- src/oj/
|   +-- core/       Status TestCase ProblemMeta Problem JudgeResult Submission JudgeTask
|   +-- judge/      Solution Judge AbstractJudge StandardJudge SpecialJudge
|   |   |           JudgeFactory(M4) MachineJudge(M5a)
|   |   +-- queue/  JudgeQueue JudgeWorker (M6a)
|   +-- exception/  JudgeException TimeLimitException (M3)
|   +-- io/         ConfigFile SingleFileProblemLoader (M4)
|   |               ProblemLoader ProblemRepository SubmissionStore (M5a)
|   +-- db/         Db ProblemDao SubmissionDao (M5b)
|   +-- service/    ProblemService (M5b)
|   +-- gui/        OjFrame OjController (M5c)
|   +-- Main.java
+-- problems/<id>/  config.txt + N.in/N.out
+-- submissions/    选手源码落地
+-- judge/          C++ 判题机(judge.cpp,编译型)
+-- lib/            mysql-connector-j.jar (M5b)
+-- Makefile
\end{verbatim}

\hypertarget{ux5de5ux5177ux94feux6f14ux8fdb}{%
\subsection{工具链演进}\label{ux5de5ux5177ux94feux6f14ux8fdb}}

\begin{itemize}
\tightlist
\item
  \textbf{M0--M1}:\texttt{nvim} + \texttt{javac} + \texttt{java} +
  \texttt{javap},单文件。
\item
  \textbf{M2}:多文件 +
  \texttt{package},\texttt{Makefile},\texttt{jar},\texttt{javadoc}。
\item
  \textbf{M4}:\texttt{File}/\texttt{BufferedReader} 读
  \texttt{config.txt};\texttt{g++} 预编译出 \texttt{judge} 二进制(M5a
  起调用)。
\item
  \textbf{M5b}:引入 MySQL 驱动 jar(\texttt{java\ -cp\ build:lib/*})。
\end{itemize}

\hypertarget{ux9010ux91ccux7a0bux7891ux8981ux70b9ux76f8ux5173ux5ea6-ux7b2cux4e00ux6b65ux62dfux5408ux8bd5ux5377-ux7b2cux4e8cux6b65ux5de5ux7a0bux5316}{%
\subsection{逐里程碑要点(★相关度 · 第一步拟合试卷 /
第二步工程化)}\label{ux9010ux91ccux7a0bux7891ux8981ux70b9ux76f8ux5173ux5ea6-ux7b2cux4e00ux6b65ux62dfux5408ux8bd5ux5377-ux7b2cux4e8cux6b65ux5de5ux7a0bux5316}}

\begin{itemize}
\tightlist
\item
  \textbf{M0 Ch1}
  ☆:走通编写-\textgreater 编译-\textgreater 运行;\texttt{Main}
  打印就绪。
\item
  \textbf{M1 Ch2--3} ★:数组+if/switch/循环,硬编码判 A+B 出
  AC/WA(真题:数组遍历、\texttt{++i}、运算符)。
\item
  \textbf{M2 Ch4} ★★★:\texttt{{[}第一步·考试向{]}} 用考试级写法把 OJ 的
  \texttt{TestCase/Problem/JudgeResult}
  写成最简普通类;\texttt{{[}第二步·工程向{]}} 收成
  \texttt{ProblemMeta}+\texttt{static}+包/jar。
\item
  \textbf{M3 Ch5--7} ★★★:\texttt{{[}第一步·考试向{]}}
  写一个简单判题(判出 AC/WA)+ try-catch;\texttt{{[}第二步·工程向{]}}
  抽成 \texttt{Judge} 接口/\texttt{AbstractJudge} 模板 + 多态 +
  自定义异常。
\item
  \textbf{M4 Ch8+Ch10} ★★:\texttt{{[}第一步·考试向{]}} 用
  \texttt{trim/equals/split}+\texttt{StringBuilder} 做 OJ
  输出比对与报告;\texttt{{[}第二步·工程向{]}}
  \texttt{config.txt}+反射工厂、单文件读写。
\item
  \textbf{M5a Ch15+Ch10} ★★:\texttt{{[}第一步·考试向{]}} 用
  \texttt{ArrayList/HashMap} 装 OJ
  的多组用例/题目(内存题库);\texttt{{[}第二步·工程向{]}}
  \texttt{ProblemLoader/Repository} + \texttt{ProcessBuilder} 调 C++
  判题机 + 序列化。
\item
  \textbf{M5b Ch11} ★★:\texttt{{[}第一步·考试向{]}} 用最简
  JDBC(\texttt{Statement})把 OJ 的提交/题目存取
  MySQL;\texttt{{[}第二步·工程向{]}}
  \texttt{ProblemDao}/事务/\texttt{PreparedStatement}/\texttt{ProblemService}
  FS/DB 分工。
\item
  \textbf{M5c Ch9} ★★★:\texttt{{[}第一步·考试向{]}} 写 OJ
  提交窗口最简版(选题+输入+提交+标签显示结果);\texttt{{[}第二步·工程向{]}}
  \texttt{OjFrame}/\texttt{OjController} MVC 解耦 + SwingWorker。
\item
  \textbf{M6a Ch12} ★★:\texttt{{[}第一步·考试向{]}} 给 OJ
  判题加最简并发(\texttt{Thread/Runnable}+\texttt{synchronized/join});\texttt{{[}第二步·工程向{]}}
  泛型阻塞队列 + 线程池 + 调 C++ 判题机(\textbf{收官})。
\end{itemize}

\hypertarget{ux5b66ux4e60ux65b9ux6cd5}{%
\subsection{学习方法}\label{ux5b66ux4e60ux65b9ux6cd5}}

\begin{itemize}
\tightlist
\item
  一里程碑一可运行产物,过了验收再进下一阶段。
\item
  每章新知识先用在 OJ 的真实需求上。
\item
  每个里程碑 \texttt{git} 提交一次,方便回看演进。
\end{itemize}

\begin{quote}
进度:M0 ☐ M1 ☐ M2 ☐ M3 ☐ M4 ☐ M5a ☐ M5b ☐ M5c ☐ M6a ☐
\end{quote}

\hypertarget{mini-oj-ux603bux8bbeux8ba1ux5168ux5c40ux5951ux7ea6}{%
\section{Mini-OJ
总设计(全局契约)}\label{mini-oj-ux603bux8bbeux8ba1ux5168ux5c40ux5951ux7ea6}}

\begin{quote}
这份是\textbf{地基}:定下全项目的包结构、类/接口清单、核心数据模型签名。各章文档都引用本文,\textbf{不得擅自改名或改签名}。
原则:\textbf{简单、教学优先};同时贯穿一条「判题机三代演进」主线。
\end{quote}

\hypertarget{ux5224ux9898ux673aux4e09ux4ee3ux4e3bux7ebfux5168ux9879ux76eeux7684ux9aa8ux67b6}{%
\subsection{1.
判题机三代主线(全项目的骨架)}\label{ux5224ux9898ux673aux4e09ux4ee3ux4e3bux7ebfux5168ux9879ux76eeux7684ux9aa8ux67b6}}

\begin{longtable}[]{@{}
  >{\raggedright\arraybackslash}p{(\columnwidth - 6\tabcolsep) * \real{0.2500}}
  >{\raggedright\arraybackslash}p{(\columnwidth - 6\tabcolsep) * \real{0.2500}}
  >{\raggedright\arraybackslash}p{(\columnwidth - 6\tabcolsep) * \real{0.2500}}
  >{\raggedright\arraybackslash}p{(\columnwidth - 6\tabcolsep) * \real{0.2500}}@{}}
\toprule\noalign{}
\begin{minipage}[b]{\linewidth}\raggedright
代
\end{minipage} & \begin{minipage}[b]{\linewidth}\raggedright
里程碑
\end{minipage} & \begin{minipage}[b]{\linewidth}\raggedright
判题方式
\end{minipage} & \begin{minipage}[b]{\linewidth}\raggedright
关键技术
\end{minipage} \\
\midrule\noalign{}
\endhead
\bottomrule\noalign{}
\endlastfoot
\textbf{第一代} & M1--M3(已完成) & Java 对象在 JVM 内模拟判题 &
\texttt{Solution.solve()} + \texttt{Judge} 多态 \\
\textbf{第二代} & M4 & 反射工厂 + 单文件配置 & \texttt{Class.forName} +
\texttt{config.txt} \\
\textbf{第三代} & M5a & 外部 C++ 判题机,真编译/真运行/真限资源 &
\texttt{ProcessBuilder} -\textgreater{} \texttt{judge} 二进制 +
\texttt{setrlimit} \\
\end{longtable}

第三代之后(M5b/M5c/M6a)都是围绕''第三代''做工程化:数据库供配置、Swing
做前端、多线程做并发。

\begin{quote}
课程调整:\textbf{第13章网络编程(Socket)不考,已删除}(原 M6b 取消,M6a
即收官);\textbf{新增第15章泛型与集合},并入 M5a。
考试相关度(★,基于历年真题)与「两步走(第一步拟合试卷 / 第二步工程化)」见
\texttt{mini-oj-plan.md} 的「考试相关度」「两步走」两节;各章开头也标了 ★
与第一步/第二步。
\end{quote}

\hypertarget{ux5173ux952eux8bbeux8ba1ux51b3ux7b56}{%
\subsection{2.
关键设计决策}\label{ux5173ux952eux8bbeux8ba1ux51b3ux7b56}}

\begin{longtable}[]{@{}
  >{\raggedright\arraybackslash}p{(\columnwidth - 4\tabcolsep) * \real{0.3333}}
  >{\raggedright\arraybackslash}p{(\columnwidth - 4\tabcolsep) * \real{0.3333}}
  >{\raggedright\arraybackslash}p{(\columnwidth - 4\tabcolsep) * \real{0.3333}}@{}}
\toprule\noalign{}
\begin{minipage}[b]{\linewidth}\raggedright
决策
\end{minipage} & \begin{minipage}[b]{\linewidth}\raggedright
选择
\end{minipage} & \begin{minipage}[b]{\linewidth}\raggedright
理由
\end{minipage} \\
\midrule\noalign{}
\endhead
\bottomrule\noalign{}
\endlastfoot
判题结果状态 & \texttt{enum\ Status} & 最干净(enum 不在大纲但极简单) \\
第一代用户解法 & \texttt{Solution} 接口:\texttt{String\ solve(String)} &
纯 Java,教接口+多态 \\
第二代换判题器 & \texttt{JudgeFactory.create(类名)} 反射 +
\texttt{config.txt} & 换题/换判法不改代码 \\
第三代真判题 & \texttt{MachineJudge} 用 \texttt{ProcessBuilder} 调外部
C++ 判题机 & 真限 CPU/内存、支持多语言 \\
题目元数据来源 & M4 \texttt{config.txt} -\textgreater{} M5b
MySQL(\texttt{ProblemMeta} 作接缝) & 数据源可换,\texttt{Problem} 不变 \\
FS/DB 分工 & 元数据/历史进 DB;\texttt{.in/.out} 与源码留文件系统 &
各取所长 \\
TLE/MLE & 由 C++ 判题机 \texttt{setrlimit}+墙钟掐断 & Java
侧不外套超时 \\
\end{longtable}

\hypertarget{ux6700ux7ec8ux5305ux7ed3ux6784ux9010ux7ae0ux957fux51faux6765}{%
\subsection{3.
最终包结构(逐章长出来)}\label{ux6700ux7ec8ux5305ux7ed3ux6784ux9010ux7ae0ux957fux51faux6765}}

\begin{verbatim}
mini-oj/
+-- src/oj/
|   +-- core/        Status, TestCase, ProblemMeta, Problem, JudgeResult, Submission, JudgeTask
|   +-- judge/       Solution, Judge, AbstractJudge, StandardJudge, SpecialJudge,
|   |   |            JudgeFactory(M4 反射), MachineJudge(M5a 第三代)
|   |   +-- queue/   JudgeQueue, JudgeWorker                              (M6a)
|   +-- exception/   JudgeException, TimeLimitException                   (M3)
|   +-- io/          ConfigFile, SingleFileProblemLoader (M4),
|   |                ProblemLoader, ProblemRepository, SubmissionStore   (M5a)
|   +-- db/          Db, ProblemDao, SubmissionDao                       (M5b)
|   +-- service/     ProblemService(DB 元数据 + FS 测试点合流)           (M5b)
|   +-- gui/         OjFrame, OjController                               (M5c)
|   +-- Main.java
+-- problems/        题库:<id>/config.txt, <id>/N.in, <id>/N.out
+-- submissions/     选手源码落地(<id>.cpp/.py)                          (M6a)
+-- judge/           C++ 判题机(单文件 judge.cpp,编译型)               (M5a 起调用)
+-- lib/             mysql-connector-j.jar                              (M5b)
+-- Makefile
\end{verbatim}

\begin{quote}
已\textbf{删除}
\texttt{net/}(JudgeServer/JudgeClient/Socket)------第13章网络不考。
\end{quote}

\hypertarget{ux5168ux91cfux7c7bux63a5ux53e3ux6e05ux5355}{%
\subsection{4.
全量类/接口清单}\label{ux5168ux91cfux7c7bux63a5ux53e3ux6e05ux5355}}

\begin{longtable}[]{@{}
  >{\raggedright\arraybackslash}p{(\columnwidth - 8\tabcolsep) * \real{0.2000}}
  >{\raggedright\arraybackslash}p{(\columnwidth - 8\tabcolsep) * \real{0.2000}}
  >{\raggedright\arraybackslash}p{(\columnwidth - 8\tabcolsep) * \real{0.2000}}
  >{\raggedright\arraybackslash}p{(\columnwidth - 8\tabcolsep) * \real{0.2000}}
  >{\raggedright\arraybackslash}p{(\columnwidth - 8\tabcolsep) * \real{0.2000}}@{}}
\toprule\noalign{}
\begin{minipage}[b]{\linewidth}\raggedright
类/接口
\end{minipage} & \begin{minipage}[b]{\linewidth}\raggedright
包
\end{minipage} & \begin{minipage}[b]{\linewidth}\raggedright
引入
\end{minipage} & \begin{minipage}[b]{\linewidth}\raggedright
代
\end{minipage} & \begin{minipage}[b]{\linewidth}\raggedright
一句话职责
\end{minipage} \\
\midrule\noalign{}
\endhead
\bottomrule\noalign{}
\endlastfoot
\texttt{Main} & (默认/oj) & M0 & --- & 程序入口、各阶段演示 \\
\texttt{Status}(enum) & core & M2 & --- &
判题状态:AC/WA/TLE/MLE/RE/CE/PE \\
\texttt{TestCase} & core & M2 & --- & 一组测试:输入 + 期望输出 \\
\texttt{ProblemMeta} & core & M4 & --- &
题目元数据:标题/判题类名/时限(数据源可换) \\
\texttt{Problem} & core & M2(重构) & --- & 题目:id +
\texttt{ProblemMeta} + 测试点 \\
\texttt{JudgeResult} & core & M2 & --- &
判题结果:状态/通过数/总数/详情/耗时 \\
\texttt{Submission} & core & M2 & --- &
一次提交:谁/哪题/语言/源码路径/结果 \\
\texttt{JudgeTask} & core & M6a & 三 & 在途工作项:源码字节 +
\texttt{complete/await} \\
\texttt{Solution} & judge & M2-\textgreater M3 & 一 &
用户解法:\texttt{solve(input)-\textgreater{}output} \\
\texttt{Judge}/\texttt{AbstractJudge}/\texttt{StandardJudge}/\texttt{SpecialJudge}
& judge & M3/M4 & 一/二 & Java 内判题器(接口+模板+子类) \\
\texttt{JudgeFactory} & judge & M4 & 二 & \texttt{Class.forName}
反射造判题器 \\
\texttt{MachineJudge} & judge & M5a & 三 & \texttt{ProcessBuilder}
调外部 C++ 判题机 \\
\texttt{JudgeException}/\texttt{TimeLimitException} & exception & M3 &
--- & 自定义异常 \\
\texttt{ConfigFile} & io & M4 & 二 & 读 \texttt{config.txt}
-\textgreater{} \texttt{ProblemMeta} \\
\texttt{SingleFileProblemLoader} & io & M4 & 二 & 单文件原型:读
input/output/config \\
\texttt{ProblemLoader} & io & M5a & 三 & 扫目录读多组用例
-\textgreater{} \texttt{List\textless{}TestCase\textgreater{}} \\
\texttt{ProblemRepository} & io & M5a & 三 &
\texttt{HashMap\textless{}Integer,Problem\textgreater{}} 缓存题目 \\
\texttt{SubmissionStore} & io & M5a & --- & 对象序列化(M5b 后由 DAO
取代) \\
\texttt{Db}/\texttt{ProblemDao}/\texttt{SubmissionDao} & db & M5b & ---
& JDBC 连接/元数据/提交历史(事务) \\
\texttt{ProblemService} & service & M5b & --- & DB 元数据 + FS
测试点合流入口 \\
\texttt{OjFrame}/\texttt{OjController} & gui & M5c & --- & Swing
大前端(MVC) \\
\texttt{JudgeQueue}/\texttt{JudgeWorker} & judge.queue & M6a & 三 &
泛型阻塞队列 + 并发判题工作线程 \\
\end{longtable}

\hypertarget{ux6838ux5fc3ux6570ux636eux6a21ux578bux5951ux7ea6ux7b7eux540dux9501ux5b9aux540eux7eedux7ae0ux8282ux7167ux6b64ux5b9eux73b0}{%
\subsection{5.
核心数据模型契约(签名锁定,后续章节照此实现)}\label{ux6838ux5fc3ux6570ux636eux6a21ux578bux5951ux7ea6ux7b7eux540dux9501ux5b9aux540eux7eedux7ae0ux8282ux7167ux6b64ux5b9eux73b0}}

\begin{Shaded}
\begin{Highlighting}[]
\CommentTok{// oj.core.Status  —— 全项目唯一权威定义}
\KeywordTok{enum}\NormalTok{ Status }\OperatorTok{\{}\NormalTok{ AC}\OperatorTok{,}\NormalTok{ WA}\OperatorTok{,}\NormalTok{ TLE}\OperatorTok{,}\NormalTok{ MLE}\OperatorTok{,}\NormalTok{ RE}\OperatorTok{,}\NormalTok{ CE}\OperatorTok{,}\NormalTok{ PE }\OperatorTok{\}}
\CommentTok{// 谁产生谁:第一代 Java 内判题(M3/M4)只产 AC/WA/RE(PE 由可选 FormatChecker 产);}
\CommentTok{//          第三代 C++ 判题机(M5a)产 AC/WA/TLE/MLE/RE/CE,其内部错 ERR 由 MachineJudge 映射成 RE。}

\CommentTok{// oj.core.TestCase}
\KeywordTok{class}\NormalTok{ TestCase }\OperatorTok{\{}
    \KeywordTok{private} \DataTypeTok{final} \BuiltInTok{String}\NormalTok{ input}\OperatorTok{,}\NormalTok{ expected}\OperatorTok{;}
    \FunctionTok{TestCase}\OperatorTok{(}\BuiltInTok{String}\NormalTok{ input}\OperatorTok{,} \BuiltInTok{String}\NormalTok{ expected}\OperatorTok{)}
    \BuiltInTok{String} \FunctionTok{getInput}\OperatorTok{();}  \BuiltInTok{String} \FunctionTok{getExpected}\OperatorTok{();}
\OperatorTok{\}}

\CommentTok{// oj.core.ProblemMeta  —— 元数据接缝:M4 来自 config.txt,M5b 来自 MySQL,Problem 不变}
\KeywordTok{class}\NormalTok{ ProblemMeta }\OperatorTok{\{}
    \KeywordTok{private} \DataTypeTok{final} \BuiltInTok{String}\NormalTok{ title}\OperatorTok{;}
    \KeywordTok{private} \DataTypeTok{final} \BuiltInTok{String}\NormalTok{ judgeClass}\OperatorTok{;}          \CommentTok{// 判题器全限定名,如 "oj.judge.StandardJudge"}
    \KeywordTok{private} \DataTypeTok{final} \DataTypeTok{long}\NormalTok{ timeLimitMs}\OperatorTok{;}
    \FunctionTok{ProblemMeta}\OperatorTok{(}\BuiltInTok{String}\NormalTok{ title}\OperatorTok{,} \BuiltInTok{String}\NormalTok{ judgeClass}\OperatorTok{,} \DataTypeTok{long}\NormalTok{ timeLimitMs}\OperatorTok{)}
    \BuiltInTok{String} \FunctionTok{getTitle}\OperatorTok{();} \BuiltInTok{String} \FunctionTok{getJudgeClass}\OperatorTok{();} \DataTypeTok{long} \FunctionTok{getTimeLimitMs}\OperatorTok{();}
\OperatorTok{\}}

\CommentTok{// oj.core.Problem  —— id + 元数据 + 测试点(M4 用 TestCase[];M5a 起 Ch15 重构为 List\textless{}TestCase\textgreater{})}
\KeywordTok{class}\NormalTok{ Problem }\OperatorTok{\{}
    \KeywordTok{private} \DataTypeTok{final} \DataTypeTok{int}\NormalTok{ id}\OperatorTok{;}
    \KeywordTok{private} \DataTypeTok{final}\NormalTok{ ProblemMeta meta}\OperatorTok{;}
    \KeywordTok{private} \DataTypeTok{final} \BuiltInTok{List}\OperatorTok{\textless{}}\NormalTok{TestCase}\OperatorTok{\textgreater{}}\NormalTok{ cases}\OperatorTok{;}       \CommentTok{// M4 阶段为 TestCase[]}
    \FunctionTok{Problem}\OperatorTok{(}\DataTypeTok{int}\NormalTok{ id}\OperatorTok{,}\NormalTok{ ProblemMeta meta}\OperatorTok{,} \BuiltInTok{List}\OperatorTok{\textless{}}\NormalTok{TestCase}\OperatorTok{\textgreater{}}\NormalTok{ cases}\OperatorTok{)}
    \DataTypeTok{int} \FunctionTok{getId}\OperatorTok{();} \BuiltInTok{String} \FunctionTok{getTitle}\OperatorTok{();} \BuiltInTok{String} \FunctionTok{getJudgeClass}\OperatorTok{();} \DataTypeTok{long} \FunctionTok{getTimeLimitMs}\OperatorTok{();}
    \BuiltInTok{List}\OperatorTok{\textless{}}\NormalTok{TestCase}\OperatorTok{\textgreater{}} \FunctionTok{getCases}\OperatorTok{();}
\OperatorTok{\}}

\CommentTok{// oj.core.JudgeResult}
\KeywordTok{class}\NormalTok{ JudgeResult }\OperatorTok{\{}
    \KeywordTok{private} \DataTypeTok{final}\NormalTok{ Status status}\OperatorTok{;}
    \KeywordTok{private} \DataTypeTok{final} \DataTypeTok{int}\NormalTok{ passed}\OperatorTok{,}\NormalTok{ total}\OperatorTok{;}
    \KeywordTok{private} \DataTypeTok{final} \BuiltInTok{String}\NormalTok{ detail}\OperatorTok{;}
    \KeywordTok{private} \DataTypeTok{final} \DataTypeTok{long}\NormalTok{ elapsedMs}\OperatorTok{;}
    \FunctionTok{JudgeResult}\OperatorTok{(}\NormalTok{Status status}\OperatorTok{,} \DataTypeTok{int}\NormalTok{ passed}\OperatorTok{,} \DataTypeTok{int}\NormalTok{ total}\OperatorTok{,} \BuiltInTok{String}\NormalTok{ detail}\OperatorTok{,} \DataTypeTok{long}\NormalTok{ elapsedMs}\OperatorTok{)}
\NormalTok{    Status }\FunctionTok{getStatus}\OperatorTok{();} \DataTypeTok{int} \FunctionTok{getPassed}\OperatorTok{();} \DataTypeTok{int} \FunctionTok{getTotal}\OperatorTok{();} \BuiltInTok{String} \FunctionTok{getDetail}\OperatorTok{();} \DataTypeTok{long} \FunctionTok{getElapsedMs}\OperatorTok{();}
    \DataTypeTok{boolean} \FunctionTok{isAccepted}\OperatorTok{();}
    \AttributeTok{@Override} \BuiltInTok{String} \FunctionTok{toString}\OperatorTok{();}              \CommentTok{// "AC 3/3 (12ms)"}
\OperatorTok{\}}   \CommentTok{// implements Serializable}

\CommentTok{// oj.core.Submission  —— 增加 lang + srcPath(用判题机/入库需要)}
\KeywordTok{class}\NormalTok{ Submission }\OperatorTok{\{}                            \CommentTok{// implements Serializable}
    \KeywordTok{private} \DataTypeTok{static} \DataTypeTok{int}\NormalTok{ counter }\OperatorTok{=} \DecValTok{0}\OperatorTok{;}
    \KeywordTok{private} \DataTypeTok{final} \DataTypeTok{int}\NormalTok{ id}\OperatorTok{;}                     \CommentTok{// counter 自增}
    \KeywordTok{private} \DataTypeTok{final} \DataTypeTok{int}\NormalTok{ problemId}\OperatorTok{;}
    \KeywordTok{private} \DataTypeTok{final} \BuiltInTok{String}\NormalTok{ userName}\OperatorTok{;}
    \KeywordTok{private} \DataTypeTok{final} \BuiltInTok{String}\NormalTok{ lang}\OperatorTok{;}                \CommentTok{// "cpp" | "python"}
    \KeywordTok{private} \DataTypeTok{final} \BuiltInTok{String}\NormalTok{ srcPath}\OperatorTok{;}             \CommentTok{// 选手源码在文件系统的路径}
    \KeywordTok{private}\NormalTok{ JudgeResult result}\OperatorTok{;}
    \FunctionTok{Submission}\OperatorTok{(}\DataTypeTok{int}\NormalTok{ problemId}\OperatorTok{,} \BuiltInTok{String}\NormalTok{ userName}\OperatorTok{,} \BuiltInTok{String}\NormalTok{ lang}\OperatorTok{,} \BuiltInTok{String}\NormalTok{ srcPath}\OperatorTok{)}
    \DataTypeTok{int} \FunctionTok{getId}\OperatorTok{();} \DataTypeTok{int} \FunctionTok{getProblemId}\OperatorTok{();} \BuiltInTok{String} \FunctionTok{getUserName}\OperatorTok{();} \BuiltInTok{String} \FunctionTok{getLang}\OperatorTok{();} \BuiltInTok{String} \FunctionTok{getSrcPath}\OperatorTok{();}
\NormalTok{    JudgeResult }\FunctionTok{getResult}\OperatorTok{();} \DataTypeTok{void} \FunctionTok{setResult}\OperatorTok{(}\NormalTok{JudgeResult r}\OperatorTok{);}
    \DataTypeTok{static} \DataTypeTok{int} \FunctionTok{count}\OperatorTok{();}
\OperatorTok{\}}

\CommentTok{// oj.judge.Solution(第一代,M3 起接口)/ Judge(M3)}
\KeywordTok{interface}\NormalTok{ Solution }\OperatorTok{\{} \BuiltInTok{String} \FunctionTok{solve}\OperatorTok{(}\BuiltInTok{String}\NormalTok{ input}\OperatorTok{);} \OperatorTok{\}}
\KeywordTok{interface}\NormalTok{ Judge }\OperatorTok{\{}\NormalTok{ JudgeResult }\FunctionTok{judge}\OperatorTok{(}\NormalTok{Problem problem}\OperatorTok{,}\NormalTok{ Solution solution}\OperatorTok{);} \OperatorTok{\}}

\CommentTok{// oj.judge.MachineJudge(第三代,M5a)}
\KeywordTok{class}\NormalTok{ MachineJudge }\OperatorTok{\{}
    \FunctionTok{MachineJudge}\OperatorTok{(}\BuiltInTok{String}\NormalTok{ judgeBinary}\OperatorTok{)}
\NormalTok{    JudgeResult }\FunctionTok{judge}\OperatorTok{(}\BuiltInTok{String}\NormalTok{ problemDir}\OperatorTok{,} \BuiltInTok{String}\NormalTok{ srcFile}\OperatorTok{,} \BuiltInTok{String}\NormalTok{ lang}\OperatorTok{,} \DataTypeTok{long}\NormalTok{ timeMs}\OperatorTok{,} \DataTypeTok{int}\NormalTok{ memMb}\OperatorTok{)}
\OperatorTok{\}}
\end{Highlighting}
\end{Shaded}

\hypertarget{ux547dux540d-ux98ceux683cux7ea6ux5b9a}{%
\subsection{6. 命名 /
风格约定}\label{ux547dux540d-ux98ceux683cux7ea6ux5b9a}}

\begin{itemize}
\tightlist
\item
  包名小写 \texttt{oj.xxx};类名大驼峰;方法/变量小驼峰;常量全大写。
\item
  缩进 4 空格;采用 K\&R(左括号不换行)。每个 public 类/方法写一句
  Javadoc。
\item
  一个文件一个 public 类;\texttt{main} 只放在 \texttt{Main}。
\end{itemize}

\hypertarget{ux6587ux6863ux6e05ux5355ux672cux7cfbux5217}{%
\subsection{7.
文档清单(本系列)}\label{ux6587ux6863ux6e05ux5355ux672cux7cfbux5217}}

\begin{longtable}[]{@{}
  >{\raggedright\arraybackslash}p{(\columnwidth - 6\tabcolsep) * \real{0.2500}}
  >{\raggedright\arraybackslash}p{(\columnwidth - 6\tabcolsep) * \real{0.2500}}
  >{\raggedright\arraybackslash}p{(\columnwidth - 6\tabcolsep) * \real{0.2500}}
  >{\raggedright\arraybackslash}p{(\columnwidth - 6\tabcolsep) * \real{0.2500}}@{}}
\toprule\noalign{}
\begin{minipage}[b]{\linewidth}\raggedright
文件
\end{minipage} & \begin{minipage}[b]{\linewidth}\raggedright
对应
\end{minipage} & \begin{minipage}[b]{\linewidth}\raggedright
里程碑
\end{minipage} & \begin{minipage}[b]{\linewidth}\raggedright
代
\end{minipage} \\
\midrule\noalign{}
\endhead
\bottomrule\noalign{}
\endlastfoot
\texttt{00-项目总设计.md} & 全局契约 & --- & --- \\
\texttt{01-Ch1-环境与HelloWorld.md} & 第1章 & M0 & --- \\
\texttt{02-Ch2-3-单题判题器.md} & 第2-3章 & M1 & 一 \\
\texttt{03-Ch4-对象建模.md} & 第4章 & M2 & 一 \\
\texttt{04-Ch5-7-判题器与异常.md} & 第5-7章 & M3 & 一 \\
\texttt{05-M4-第二代判题机.md} & 第8章 + 第10章(单文件) & M4 & 二 \\
\texttt{06-M5a-第三代判题机.md} & \textbf{第15章} + 第10章(工业级) + C++
& M5a & 三 \\
\texttt{07-M5b-数据库衔接.md} & 第11章 & M5b & --- \\
\texttt{08-M5c-Swing客户端.md} & 第9章 & M5c & --- \\
\texttt{09-M6a-多线程收官.md} & 第12章(收官) & M6a & 三 \\
\end{longtable}

\begin{quote}
\textbf{判题机}(真实编译运行 C++/Python 提交)是独立组件,代码在仓库
\texttt{judge/}(单文件 \texttt{judge.cpp})。 第一代 Java
内判题(\texttt{Judge}/\texttt{Solution})在 M3/M4 用于教 OOP;第三代(M5a
起)真实判题走 \texttt{MachineJudge} 调 \texttt{judge/}。
第13章网络(Socket/\texttt{JudgeServer}/\texttt{JudgeClient})已删除,不出现在任何章节。
\end{quote}

\hypertarget{ux7b2c1ux7ae0-ux73afux5883ux4e0e-helloworldux91ccux7a0bux7891-m0}{%
\section{01 · 第1章 环境与 HelloWorld(里程碑
M0)}\label{ux7b2c1ux7ae0-ux73afux5883ux4e0e-helloworldux91ccux7a0bux7891-m0}}

\begin{quote}
\textbf{考试相关度 ☆} ·
真题考点参考:编译运行、字节码、\texttt{.class}(少量选择)。本章即「第一步·上手」,纯命令行
\texttt{javac}/\texttt{java}/\texttt{javap},无两步拆分。
\end{quote}

\begin{quote}
目标:装好
JDK,理解''编写-\textgreater 编译-\textgreater 运行'',跑出第一个程序。\textbf{本章只有一个文件、一个方法。}
\end{quote}

\hypertarget{ux77e5ux8bc6ux70b9ux5927ux7eb2ux5bf9ux7167}{%
\subsection{知识点(大纲对照)}\label{ux77e5ux8bc6ux70b9ux5927ux7eb2ux5bf9ux7167}}

\begin{longtable}[]{@{}ll@{}}
\toprule\noalign{}
要点 & 处理 \\
\midrule\noalign{}
\endhead
\bottomrule\noalign{}
\endlastfoot
Java 地位/特点/James Gosling & \textbf{了解即可},不实现 \\
JDK 安装与环境变量 & \textbf{动手} \\
编写-\textgreater 编译-\textgreater 运行 三步 & \textbf{动手}(核心) \\
\texttt{javap} 反编译 & \textbf{动手}(看一眼字节码即可) \\
编程风格 Allman/K\&R、注释 & \textbf{了解},本项目统一用 K\&R \\
\end{longtable}

\hypertarget{ux8981ux5199ux7684ux4ee3ux7801}{%
\subsection{要写的代码}\label{ux8981ux5199ux7684ux4ee3ux7801}}

\textbf{文件}:\texttt{Main.java}(默认包,先不分包)

\begin{verbatim}
类 Main
+-- public static void main(String[] args)
\end{verbatim}

\begin{longtable}[]{@{}
  >{\raggedright\arraybackslash}p{(\columnwidth - 6\tabcolsep) * \real{0.2500}}
  >{\raggedright\arraybackslash}p{(\columnwidth - 6\tabcolsep) * \real{0.2500}}
  >{\raggedright\arraybackslash}p{(\columnwidth - 6\tabcolsep) * \real{0.2500}}
  >{\raggedright\arraybackslash}p{(\columnwidth - 6\tabcolsep) * \real{0.2500}}@{}}
\toprule\noalign{}
\begin{minipage}[b]{\linewidth}\raggedright
方法
\end{minipage} & \begin{minipage}[b]{\linewidth}\raggedright
签名
\end{minipage} & \begin{minipage}[b]{\linewidth}\raggedright
职责
\end{minipage} & \begin{minipage}[b]{\linewidth}\raggedright
关键逻辑
\end{minipage} \\
\midrule\noalign{}
\endhead
\bottomrule\noalign{}
\endlastfoot
入口 & \texttt{public\ static\ void\ main(String{[}{]}\ args)} &
打印就绪信息 & 一行
\texttt{System.out.println("Mini-OJ\ judge\ ready");} \\
\end{longtable}

\begin{quote}
就这么简单。重点不在代码,在于走通工具链。
\end{quote}

\hypertarget{ux52a8ux624bux6b65ux9aa4}{%
\subsection{动手步骤}\label{ux52a8ux624bux6b65ux9aa4}}

\begin{Shaded}
\begin{Highlighting}[]
\CommentTok{\# 1. 装 JDK(二选一)}
\FunctionTok{sudo}\NormalTok{ apt install openjdk{-}21{-}jdk}
\ExtensionTok{java} \AttributeTok{{-}version} \KeywordTok{\&\&} \ExtensionTok{javac} \AttributeTok{{-}version}

\CommentTok{\# 2. 建目录、写代码}
\FunctionTok{mkdir} \AttributeTok{{-}p}\NormalTok{ \textasciitilde{}/mini{-}oj }\KeywordTok{\&\&} \BuiltInTok{cd}\NormalTok{ \textasciitilde{}/mini{-}oj}
\ExtensionTok{nvim}\NormalTok{ Main.java          }\CommentTok{\# 写上面的 Main 类}

\CommentTok{\# 3. 编译 {-}\textgreater{} 运行}
\ExtensionTok{javac}\NormalTok{ Main.java         }\CommentTok{\# 生成 Main.class}
\ExtensionTok{java}\NormalTok{ Main               }\CommentTok{\# 输出: Mini{-}OJ judge ready}

\CommentTok{\# 4. 反编译看字节码(理解 .class)}
\ExtensionTok{javap} \AttributeTok{{-}c}\NormalTok{ Main}
\end{Highlighting}
\end{Shaded}

\hypertarget{ux9a8cux6536ux6807ux51c6}{%
\subsection{验收标准}\label{ux9a8cux6536ux6807ux51c6}}

\begin{itemize}
\tightlist
\item[$\square$]
  \texttt{java\ -version} 正常,\texttt{echo\ \$JAVA\_HOME}
  有值(如配了)。
\item[$\square$]
  终端打印 \texttt{Mini-OJ\ judge\ ready}。
\item[$\square$]
  能用自己的话说清:\texttt{Main.java} --javac--\textgreater{}
  \texttt{Main.class} --java(JVM)--\textgreater{} 运行。
\item[$\square$]
  \texttt{javap\ -c\ Main} 能看到 \texttt{main} 方法的字节码(认出
  \texttt{getstatic}/\texttt{invokevirtual} 即可,不求全懂)。
\end{itemize}

\hypertarget{ux672cux7ae0ux4ea7ux7269}{%
\subsection{本章产物}\label{ux672cux7ae0ux4ea7ux7269}}

一个能编译运行的 \texttt{Main.java}。下一章在它里面加判题逻辑。

\hypertarget{ux7b2c2-3ux7ae0-ux5355ux9898ux5224ux9898ux5668ux91ccux7a0bux7891-m1}{%
\section{02 · 第2-3章 单题判题器(里程碑
M1)}\label{ux7b2c2-3ux7ae0-ux5355ux9898ux5224ux9898ux5668ux91ccux7a0bux7891-m1}}

\begin{quote}
\textbf{考试相关度 ★} ·
真题以选择/读程序考数组遍历、运算符、\texttt{++i}、流程控制。本章即基础语法,无需两步拆分。
\end{quote}

\begin{quote}
目标:\textbf{还不用类},只用基本类型、数组、循环、分支,写出能判 A+B
的判题器,跑出 \texttt{AC/WA}。 全部代码仍在单文件 \texttt{Main.java},用
\texttt{static} 方法 + 数组组织。
\end{quote}

\hypertarget{ux77e5ux8bc6ux70b9ux5927ux7eb2ux5bf9ux7167-1}{%
\subsection{知识点(大纲对照)}\label{ux77e5ux8bc6ux70b9ux5927ux7eb2ux5bf9ux7167-1}}

\begin{longtable}[]{@{}
  >{\raggedright\arraybackslash}p{(\columnwidth - 2\tabcolsep) * \real{0.5000}}
  >{\raggedright\arraybackslash}p{(\columnwidth - 2\tabcolsep) * \real{0.5000}}@{}}
\toprule\noalign{}
\begin{minipage}[b]{\linewidth}\raggedright
要点
\end{minipage} & \begin{minipage}[b]{\linewidth}\raggedright
处理
\end{minipage} \\
\midrule\noalign{}
\endhead
\bottomrule\noalign{}
\endlastfoot
标识符/关键字、基本数据类型、类型转换 &
\textbf{动手}(int/double/String/类型转换) \\
Scanner 输入 / println 输出 & \textbf{动手} \\
数组:声明/分配/length/初始化/遍历 & \textbf{动手}(用例存数组) \\
算术/关系/逻辑/赋值运算符 & \textbf{动手} \\
位运算符、instanceof & \textbf{了解即可},不实现 \\
if / if-else / switch & \textbf{动手} \\
for / while / do-while、break/continue & \textbf{动手} \\
\end{longtable}

\hypertarget{ux8981ux5199ux7684ux4ee3ux7801-1}{%
\subsection{要写的代码}\label{ux8981ux5199ux7684ux4ee3ux7801-1}}

\textbf{文件}:\texttt{Main.java}(默认包,纯 \texttt{static} 方法)

\begin{verbatim}
类 Main
+-- main                  程序入口,依次跑三道题的判题
+-- solve                 题1:模拟用户解法 a+b
+-- judgeAplusB           题1:遍历用例,统计 AC/WA 并打印
+-- grade                 题2:switch 求成绩等级(练 switch)
+-- isPrime               题3:循环判素数(练 for/while/break)
+-- compareInt            通用:比对期望值与实际值
\end{verbatim}

\hypertarget{ux65b9ux6cd5ux6e05ux5355ux9010ux4e2a}{%
\subsubsection{方法清单(逐个)}\label{ux65b9ux6cd5ux6e05ux5355ux9010ux4e2a}}

\begin{longtable}[]{@{}
  >{\raggedright\arraybackslash}p{(\columnwidth - 6\tabcolsep) * \real{0.2500}}
  >{\raggedright\arraybackslash}p{(\columnwidth - 6\tabcolsep) * \real{0.2500}}
  >{\raggedright\arraybackslash}p{(\columnwidth - 6\tabcolsep) * \real{0.2500}}
  >{\raggedright\arraybackslash}p{(\columnwidth - 6\tabcolsep) * \real{0.2500}}@{}}
\toprule\noalign{}
\begin{minipage}[b]{\linewidth}\raggedright
方法
\end{minipage} & \begin{minipage}[b]{\linewidth}\raggedright
签名
\end{minipage} & \begin{minipage}[b]{\linewidth}\raggedright
职责
\end{minipage} & \begin{minipage}[b]{\linewidth}\raggedright
关键逻辑
\end{minipage} \\
\midrule\noalign{}
\endhead
\bottomrule\noalign{}
\endlastfoot
入口 & \texttt{public\ static\ void\ main(String{[}{]}\ a)} & 跑三道题 &
顺序调 \texttt{judgeAplusB()} 等,打印分隔 \\
解法1 & \texttt{static\ int\ solve(int\ a,\ int\ b)} & 模拟用户 A+B 解法
& \texttt{return\ a\ +\ b;} \\
判题1 & \texttt{static\ void\ judgeAplusB()} &
内置用例-\textgreater 逐组判-\textgreater 统计 & 见下方算法 \\
比对 & \texttt{static\ boolean\ compareInt(int\ expected,\ int\ actual)}
& 判断一组是否通过 & \texttt{return\ expected\ ==\ actual;} \\
解法2 & \texttt{static\ char\ grade(int\ score)} &
分数-\textgreater 等级(练 switch) &
\texttt{switch(score/10)\{case\ 10,9-\textgreater{}\textquotesingle{}A\textquotesingle{};\ ...\}} \\
解法3 & \texttt{static\ boolean\ isPrime(int\ n)} & 判素数(练循环) &
\texttt{for(i=2;i*i\textless{}=n;i++)\ if(n\%i==0)\ return\ false;} \\
\end{longtable}

\begin{quote}
题2/题3 只需各写一个简版 \texttt{judgeXxx()} 仿照
\texttt{judgeAplusB},确认 switch / 循环
都用上即可,不必三题都做满------\textbf{够覆盖知识点就停}。
\end{quote}

\hypertarget{judgeaplusb-ux7b97ux6cd5ux6838ux5fc3ux8bb2ux6e05ux6570ux7ec4ux5faaux73afux5206ux652f}{%
\subsubsection{\texorpdfstring{\texttt{judgeAplusB}
算法(核心,讲清数组+循环+分支)}{judgeAplusB 算法(核心,讲清数组+循环+分支)}}\label{judgeaplusb-ux7b97ux6cd5ux6838ux5fc3ux8bb2ux6e05ux6570ux7ec4ux5faaux73afux5206ux652f}}

\begin{verbatim}
// 用例:每行 {a, b};期望:对应 a+b
int[][] inputs   = { {1,2}, {10,20}, {0,0} };
int[] expected   = {   3,     30,      0  };

int passed = 0;
for (int i = 0; i < inputs.length; i++) {          // 数组 length + for 遍历
    int actual = solve(inputs[i][0], inputs[i][1]);
    if (compareInt(expected[i], actual)) {          // 分支
        passed++;
    } else {
        System.out.println("case#" + i + " 期望 " + expected[i] + " 实际 " + actual);
        // 可 continue / break,体会两者区别
    }
}
String status = (passed == inputs.length) ? "AC" : "WA";   // 三元 + 关系运算
System.out.println("A+B: " + status + " " + passed + "/" + inputs.length);
\end{verbatim}

\begin{quote}
故意制造一组错的期望值,观察输出从 \texttt{AC} 变 \texttt{WA}
并打出失败行------这就是''判题''的内核。
\end{quote}

\hypertarget{ux52a8ux624bux6b65ux9aa4-1}{%
\subsection{动手步骤}\label{ux52a8ux624bux6b65ux9aa4-1}}

\begin{Shaded}
\begin{Highlighting}[]
\BuiltInTok{cd}\NormalTok{ \textasciitilde{}/mini{-}oj}
\ExtensionTok{nvim}\NormalTok{ Main.java            }\CommentTok{\# 按上面方法清单补全}
\ExtensionTok{javac}\NormalTok{ Main.java }\KeywordTok{\&\&} \ExtensionTok{java}\NormalTok{ Main}
\CommentTok{\# 期望输出形如:}
\CommentTok{\#   A+B: AC 3/3}
\CommentTok{\#   Grade(85): B}
\CommentTok{\#   isPrime(7): true}
\end{Highlighting}
\end{Shaded}

\hypertarget{ux9a8cux6536ux6807ux51c6-1}{%
\subsection{验收标准}\label{ux9a8cux6536ux6807ux51c6-1}}

\begin{itemize}
\tightlist
\item[$\square$]
  三道题都能跑,A+B 打印 \texttt{AC\ 3/3}。
\item[$\square$]
  把某组 \texttt{expected} 改错 -\textgreater{} 输出变 \texttt{WA\ 2/3}
  且打出失败的 case 行。
\item[$\square$]
  代码里确实用到了:数组遍历、\texttt{if}/三元、\texttt{switch}、\texttt{for}(或
  \texttt{while})。
\item[$\square$]
  能说清''为什么现在代码很难扩展到多道题''------为下一章引入''类/对象''埋伏笔。
\end{itemize}

\hypertarget{ux672cux7ae0ux4ea7ux7269---ux4e0bux4e00ux7ae0}{%
\subsection{本章产物 -\textgreater{}
下一章}\label{ux672cux7ae0ux4ea7ux7269---ux4e0bux4e00ux7ae0}}

一个单文件、面向过程的判题器。\textbf{痛点}:题目、用例、结果全是散落的数组和变量,加一道题要改一堆地方。第3章(M2)用\textbf{类与对象}把它们封装成
\texttt{Problem}/\texttt{TestCase}/\texttt{JudgeResult}。

\hypertarget{ux7b2c4ux7ae0-ux5bf9ux8c61ux5efaux6a21ux91ccux7a0bux7891-m2}{%
\section{03 · 第4章 对象建模(里程碑
M2)}\label{ux7b2c4ux7ae0-ux5bf9ux8c61ux5efaux6a21ux91ccux7a0bux7891-m2}}

\begin{quote}
\textbf{考试相关度 ★★★(编程大题核心)} ·
真题考点参考:定义类(属性/构造/get-set/封装)。\textbf{{[}第一步{]}}
\KS 用考试级写法把 OJ 的 \texttt{TestCase/Problem/JudgeResult}
写成最简普通类(字段+构造+get/set),能 new 能打印;\textbf{{[}第二步{]}}
\GC 收成 \texttt{ProblemMeta} + \texttt{static} 计数 + 包/jar。
\end{quote}

\begin{quote}
目标:把 M1 散落的数组/变量,封装成 \texttt{oj.core}
的几个类;判题逻辑收进一个
\texttt{SimpleJudge}(还没学接口,先用具体类)。学会包、jar、javadoc。
\end{quote}

\hypertarget{ux77e5ux8bc6ux70b9ux5927ux7eb2ux5bf9ux7167-2}{%
\subsection{知识点(大纲对照)}\label{ux77e5ux8bc6ux70b9ux5927ux7eb2ux5bf9ux7167-2}}

\begin{longtable}[]{@{}ll@{}}
\toprule\noalign{}
要点 & 处理 \\
\midrule\noalign{}
\endhead
\bottomrule\noalign{}
\endlastfoot
类/成员变量/方法/构造方法/\texttt{this} & \textbf{动手}(贯穿) \\
封装(private + getter) & \textbf{动手} \\
实例成员 vs 类成员(\texttt{static}) &
\textbf{动手}(\texttt{Submission.counter}) \\
方法重载、可变参数 & \textbf{动手}(构造重载 + \texttt{TestCase...}) \\
对象数组 & \textbf{动手}(\texttt{Problem.cases}) \\
包 \texttt{package} / \texttt{import} / 访问权限 & \textbf{动手}(迁进
\texttt{oj.core}) \\
可运行 jar、javadoc & \textbf{动手} \\
基本类型的类封装、\texttt{var} 局部变量 &
\textbf{了解即可},可不专门写 \\
\end{longtable}

\hypertarget{ux8981ux5199ux7684ux7c7bux672cux7ae0ux5171-6-ux4e2a}{%
\subsection{要写的类(本章共 6
个)}\label{ux8981ux5199ux7684ux7c7bux672cux7ae0ux5171-6-ux4e2a}}

\begin{verbatim}
oj.core:   Status(enum)  TestCase  Problem  JudgeResult  Submission
oj.judge:  Solution(本章先做成具体类)  SimpleJudge
\end{verbatim}

\hypertarget{statusenum-ojcorestatus.java}{%
\subsubsection{\texorpdfstring{Status(enum) ---
\texttt{oj/core/Status.java}}{Status(enum) --- oj/core/Status.java}}\label{statusenum-ojcorestatus.java}}

\begin{Shaded}
\begin{Highlighting}[]
\KeywordTok{package}\ImportTok{ oj}\OperatorTok{.}\ImportTok{core}\OperatorTok{;}
\KeywordTok{public} \KeywordTok{enum}\NormalTok{ Status }\OperatorTok{\{}\NormalTok{ AC}\OperatorTok{,}\NormalTok{ WA}\OperatorTok{,}\NormalTok{ TLE}\OperatorTok{,}\NormalTok{ MLE}\OperatorTok{,}\NormalTok{ RE}\OperatorTok{,}\NormalTok{ CE}\OperatorTok{,}\NormalTok{ PE }\OperatorTok{\}}   \CommentTok{// 权威定义见 00 §5;MLE/CE 主要由判题机(第11章)产生,Java 内判题用不到}
\end{Highlighting}
\end{Shaded}

\hypertarget{testcase-ojcoretestcase.java}{%
\subsubsection{\texorpdfstring{TestCase ---
\texttt{oj/core/TestCase.java}}{TestCase --- oj/core/TestCase.java}}\label{testcase-ojcoretestcase.java}}

\begin{Shaded}
\begin{Highlighting}[]
\KeywordTok{package}\ImportTok{ oj}\OperatorTok{.}\ImportTok{core}\OperatorTok{;}
\KeywordTok{public} \KeywordTok{class}\NormalTok{ TestCase }\OperatorTok{\{}
    \KeywordTok{private} \DataTypeTok{final} \BuiltInTok{String}\NormalTok{ input}\OperatorTok{;}
    \KeywordTok{private} \DataTypeTok{final} \BuiltInTok{String}\NormalTok{ expected}\OperatorTok{;}
    \KeywordTok{public} \FunctionTok{TestCase}\OperatorTok{(}\BuiltInTok{String}\NormalTok{ input}\OperatorTok{,} \BuiltInTok{String}\NormalTok{ expected}\OperatorTok{)} \OperatorTok{\{}
        \KeywordTok{this}\OperatorTok{.}\FunctionTok{input} \OperatorTok{=}\NormalTok{ input}\OperatorTok{;}            \CommentTok{// this 区分参数与字段}
        \KeywordTok{this}\OperatorTok{.}\FunctionTok{expected} \OperatorTok{=}\NormalTok{ expected}\OperatorTok{;}
    \OperatorTok{\}}
    \KeywordTok{public} \BuiltInTok{String} \FunctionTok{getInput}\OperatorTok{()}    \OperatorTok{\{} \ControlFlowTok{return}\NormalTok{ input}\OperatorTok{;} \OperatorTok{\}}
    \KeywordTok{public} \BuiltInTok{String} \FunctionTok{getExpected}\OperatorTok{()} \OperatorTok{\{} \ControlFlowTok{return}\NormalTok{ expected}\OperatorTok{;} \OperatorTok{\}}
\OperatorTok{\}}
\end{Highlighting}
\end{Shaded}

\hypertarget{problem-ojcoreproblem.javaux5bf9ux8c61ux6570ux7ec4-ux6784ux9020ux91cdux8f7d-ux53efux53d8ux53c2ux6570}{%
\subsubsection{\texorpdfstring{Problem ---
\texttt{oj/core/Problem.java}(对象数组 + 构造重载 +
可变参数)}{Problem --- oj/core/Problem.java(对象数组 + 构造重载 + 可变参数)}}\label{problem-ojcoreproblem.javaux5bf9ux8c61ux6570ux7ec4-ux6784ux9020ux91cdux8f7d-ux53efux53d8ux53c2ux6570}}

\begin{Shaded}
\begin{Highlighting}[]
\KeywordTok{package}\ImportTok{ oj}\OperatorTok{.}\ImportTok{core}\OperatorTok{;}
\KeywordTok{public} \KeywordTok{class}\NormalTok{ Problem }\OperatorTok{\{}
    \KeywordTok{private} \DataTypeTok{final} \DataTypeTok{int}\NormalTok{ id}\OperatorTok{;}
    \KeywordTok{private} \DataTypeTok{final} \BuiltInTok{String}\NormalTok{ title}\OperatorTok{;}
    \KeywordTok{private} \DataTypeTok{final}\NormalTok{ TestCase}\OperatorTok{[]}\NormalTok{ cases}\OperatorTok{;}        \CommentTok{// 对象数组}
    \KeywordTok{private} \DataTypeTok{final} \DataTypeTok{long}\NormalTok{ timeLimitMs}\OperatorTok{;}
    \CommentTok{// 构造重载①:默认时限 1000ms;可变参数 TestCase... 本质就是 TestCase[]}
    \KeywordTok{public} \FunctionTok{Problem}\OperatorTok{(}\DataTypeTok{int}\NormalTok{ id}\OperatorTok{,} \BuiltInTok{String}\NormalTok{ title}\OperatorTok{,}\NormalTok{ TestCase}\KeywordTok{...}\NormalTok{ cases}\OperatorTok{)} \OperatorTok{\{}
        \KeywordTok{this}\OperatorTok{(}\NormalTok{id}\OperatorTok{,}\NormalTok{ title}\OperatorTok{,} \DecValTok{1000}\OperatorTok{,}\NormalTok{ cases}\OperatorTok{);}      \CommentTok{// this(...) 调另一个构造方法}
    \OperatorTok{\}}
    \CommentTok{// 构造重载②:显式时限}
    \KeywordTok{public} \FunctionTok{Problem}\OperatorTok{(}\DataTypeTok{int}\NormalTok{ id}\OperatorTok{,} \BuiltInTok{String}\NormalTok{ title}\OperatorTok{,} \DataTypeTok{long}\NormalTok{ timeLimitMs}\OperatorTok{,}\NormalTok{ TestCase}\KeywordTok{...}\NormalTok{ cases}\OperatorTok{)} \OperatorTok{\{}
        \KeywordTok{this}\OperatorTok{.}\FunctionTok{id} \OperatorTok{=}\NormalTok{ id}\OperatorTok{;} \KeywordTok{this}\OperatorTok{.}\FunctionTok{title} \OperatorTok{=}\NormalTok{ title}\OperatorTok{;}
        \KeywordTok{this}\OperatorTok{.}\FunctionTok{timeLimitMs} \OperatorTok{=}\NormalTok{ timeLimitMs}\OperatorTok{;} \KeywordTok{this}\OperatorTok{.}\FunctionTok{cases} \OperatorTok{=}\NormalTok{ cases}\OperatorTok{;}
    \OperatorTok{\}}
    \KeywordTok{public} \DataTypeTok{int} \FunctionTok{getId}\OperatorTok{()}           \OperatorTok{\{} \ControlFlowTok{return}\NormalTok{ id}\OperatorTok{;} \OperatorTok{\}}
    \KeywordTok{public} \BuiltInTok{String} \FunctionTok{getTitle}\OperatorTok{()}     \OperatorTok{\{} \ControlFlowTok{return}\NormalTok{ title}\OperatorTok{;} \OperatorTok{\}}
    \KeywordTok{public}\NormalTok{ TestCase}\OperatorTok{[]} \FunctionTok{getCases}\OperatorTok{()} \OperatorTok{\{} \ControlFlowTok{return}\NormalTok{ cases}\OperatorTok{;} \OperatorTok{\}}
    \KeywordTok{public} \DataTypeTok{long} \FunctionTok{getTimeLimitMs}\OperatorTok{()} \OperatorTok{\{} \ControlFlowTok{return}\NormalTok{ timeLimitMs}\OperatorTok{;} \OperatorTok{\}}
\OperatorTok{\}}
\end{Highlighting}
\end{Shaded}

\begin{quote}
注:为可变参数好用,把 \texttt{timeLimitMs} 放在 \texttt{cases} 前;签名与
00 契约等价(\texttt{TestCase...} 即 \texttt{TestCase{[}{]}})。
\end{quote}

\hypertarget{judgeresult-ojcorejudgeresult.java}{%
\subsubsection{\texorpdfstring{JudgeResult ---
\texttt{oj/core/JudgeResult.java}}{JudgeResult --- oj/core/JudgeResult.java}}\label{judgeresult-ojcorejudgeresult.java}}

\begin{Shaded}
\begin{Highlighting}[]
\KeywordTok{package}\ImportTok{ oj}\OperatorTok{.}\ImportTok{core}\OperatorTok{;}
\KeywordTok{public} \KeywordTok{class}\NormalTok{ JudgeResult }\OperatorTok{\{}
    \KeywordTok{private} \DataTypeTok{final}\NormalTok{ Status status}\OperatorTok{;}
    \KeywordTok{private} \DataTypeTok{final} \DataTypeTok{int}\NormalTok{ passed}\OperatorTok{,}\NormalTok{ total}\OperatorTok{;}
    \KeywordTok{private} \DataTypeTok{final} \BuiltInTok{String}\NormalTok{ detail}\OperatorTok{;}
    \KeywordTok{private} \DataTypeTok{final} \DataTypeTok{long}\NormalTok{ elapsedMs}\OperatorTok{;}
    \KeywordTok{public} \FunctionTok{JudgeResult}\OperatorTok{(}\NormalTok{Status status}\OperatorTok{,} \DataTypeTok{int}\NormalTok{ passed}\OperatorTok{,} \DataTypeTok{int}\NormalTok{ total}\OperatorTok{,} \BuiltInTok{String}\NormalTok{ detail}\OperatorTok{,} \DataTypeTok{long}\NormalTok{ elapsedMs}\OperatorTok{)} \OperatorTok{\{}
        \KeywordTok{this}\OperatorTok{.}\FunctionTok{status}\OperatorTok{=}\NormalTok{status}\OperatorTok{;} \KeywordTok{this}\OperatorTok{.}\FunctionTok{passed}\OperatorTok{=}\NormalTok{passed}\OperatorTok{;} \KeywordTok{this}\OperatorTok{.}\FunctionTok{total}\OperatorTok{=}\NormalTok{total}\OperatorTok{;}
        \KeywordTok{this}\OperatorTok{.}\FunctionTok{detail}\OperatorTok{=}\NormalTok{detail}\OperatorTok{;} \KeywordTok{this}\OperatorTok{.}\FunctionTok{elapsedMs}\OperatorTok{=}\NormalTok{elapsedMs}\OperatorTok{;}
    \OperatorTok{\}}
    \KeywordTok{public}\NormalTok{ Status }\FunctionTok{getStatus}\OperatorTok{()}  \OperatorTok{\{} \ControlFlowTok{return}\NormalTok{ status}\OperatorTok{;} \OperatorTok{\}}
    \KeywordTok{public} \DataTypeTok{int} \FunctionTok{getPassed}\OperatorTok{()}     \OperatorTok{\{} \ControlFlowTok{return}\NormalTok{ passed}\OperatorTok{;} \OperatorTok{\}}
    \KeywordTok{public} \DataTypeTok{int} \FunctionTok{getTotal}\OperatorTok{()}      \OperatorTok{\{} \ControlFlowTok{return}\NormalTok{ total}\OperatorTok{;} \OperatorTok{\}}
    \KeywordTok{public} \BuiltInTok{String} \FunctionTok{getDetail}\OperatorTok{()}  \OperatorTok{\{} \ControlFlowTok{return}\NormalTok{ detail}\OperatorTok{;} \OperatorTok{\}}
    \KeywordTok{public} \DataTypeTok{long} \FunctionTok{getElapsedMs}\OperatorTok{()} \OperatorTok{\{} \ControlFlowTok{return}\NormalTok{ elapsedMs}\OperatorTok{;} \OperatorTok{\}}
    \KeywordTok{public} \DataTypeTok{boolean} \FunctionTok{isAccepted}\OperatorTok{()\{} \ControlFlowTok{return}\NormalTok{ status }\OperatorTok{==}\NormalTok{ Status}\OperatorTok{.}\FunctionTok{AC}\OperatorTok{;} \OperatorTok{\}}
    \AttributeTok{@Override} \KeywordTok{public} \BuiltInTok{String} \FunctionTok{toString}\OperatorTok{()} \OperatorTok{\{}
        \ControlFlowTok{return}\NormalTok{ status }\OperatorTok{+} \StringTok{" "} \OperatorTok{+}\NormalTok{ passed }\OperatorTok{+} \StringTok{"/"} \OperatorTok{+}\NormalTok{ total }\OperatorTok{+} \StringTok{" ("} \OperatorTok{+}\NormalTok{ elapsedMs }\OperatorTok{+} \StringTok{"ms)"}\OperatorTok{;}
    \OperatorTok{\}}
\OperatorTok{\}}
\end{Highlighting}
\end{Shaded}

\hypertarget{submission-ojcoresubmission.javaux7c7bux53d8ux91cf-static}{%
\subsubsection{\texorpdfstring{Submission ---
\texttt{oj/core/Submission.java}(类变量
static)}{Submission --- oj/core/Submission.java(类变量 static)}}\label{submission-ojcoresubmission.javaux7c7bux53d8ux91cf-static}}

\begin{Shaded}
\begin{Highlighting}[]
\KeywordTok{package}\ImportTok{ oj}\OperatorTok{.}\ImportTok{core}\OperatorTok{;}
\KeywordTok{public} \KeywordTok{class}\NormalTok{ Submission }\OperatorTok{\{}
    \KeywordTok{private} \DataTypeTok{static} \DataTypeTok{int}\NormalTok{ counter }\OperatorTok{=} \DecValTok{0}\OperatorTok{;}     \CommentTok{// 类变量:所有提交共享}
    \KeywordTok{private} \DataTypeTok{final} \DataTypeTok{int}\NormalTok{ id}\OperatorTok{;}
    \KeywordTok{private} \DataTypeTok{final} \DataTypeTok{int}\NormalTok{ problemId}\OperatorTok{;}
    \KeywordTok{private} \DataTypeTok{final} \BuiltInTok{String}\NormalTok{ userName}\OperatorTok{;}
    \KeywordTok{private}\NormalTok{ JudgeResult result}\OperatorTok{;}         \CommentTok{// 判完回填}
    \KeywordTok{public} \FunctionTok{Submission}\OperatorTok{(}\DataTypeTok{int}\NormalTok{ problemId}\OperatorTok{,} \BuiltInTok{String}\NormalTok{ userName}\OperatorTok{)} \OperatorTok{\{}
        \KeywordTok{this}\OperatorTok{.}\FunctionTok{id} \OperatorTok{=} \OperatorTok{++}\NormalTok{counter}\OperatorTok{;}           \CommentTok{// 自增得到唯一 id}
        \KeywordTok{this}\OperatorTok{.}\FunctionTok{problemId} \OperatorTok{=}\NormalTok{ problemId}\OperatorTok{;}
        \KeywordTok{this}\OperatorTok{.}\FunctionTok{userName} \OperatorTok{=}\NormalTok{ userName}\OperatorTok{;}
    \OperatorTok{\}}
    \KeywordTok{public} \DataTypeTok{int} \FunctionTok{getId}\OperatorTok{()}            \OperatorTok{\{} \ControlFlowTok{return}\NormalTok{ id}\OperatorTok{;} \OperatorTok{\}}
    \KeywordTok{public} \DataTypeTok{int} \FunctionTok{getProblemId}\OperatorTok{()}     \OperatorTok{\{} \ControlFlowTok{return}\NormalTok{ problemId}\OperatorTok{;} \OperatorTok{\}}
    \KeywordTok{public} \BuiltInTok{String} \FunctionTok{getUserName}\OperatorTok{()}   \OperatorTok{\{} \ControlFlowTok{return}\NormalTok{ userName}\OperatorTok{;} \OperatorTok{\}}
    \KeywordTok{public}\NormalTok{ JudgeResult }\FunctionTok{getResult}\OperatorTok{()\{} \ControlFlowTok{return}\NormalTok{ result}\OperatorTok{;} \OperatorTok{\}}
    \KeywordTok{public} \DataTypeTok{void} \FunctionTok{setResult}\OperatorTok{(}\NormalTok{JudgeResult r}\OperatorTok{)\{} \KeywordTok{this}\OperatorTok{.}\FunctionTok{result} \OperatorTok{=}\NormalTok{ r}\OperatorTok{;} \OperatorTok{\}}
    \KeywordTok{public} \DataTypeTok{static} \DataTypeTok{int} \FunctionTok{count}\OperatorTok{()}     \OperatorTok{\{} \ControlFlowTok{return}\NormalTok{ counter}\OperatorTok{;} \OperatorTok{\}}   \CommentTok{// 类方法}
\OperatorTok{\}}
\end{Highlighting}
\end{Shaded}

\begin{quote}
\texttt{time} 字段(LocalDateTime)第5章(M4)再加,避免本章引入日期类。
\end{quote}

\hypertarget{solutionux672cux7ae0ux5177ux4f53ux7c7b-ojjudgeaplusb.java}{%
\subsubsection{\texorpdfstring{Solution(本章=具体类) ---
\texttt{oj/judge/AplusB.java}}{Solution(本章=具体类) --- oj/judge/AplusB.java}}\label{solutionux672cux7ae0ux5177ux4f53ux7c7b-ojjudgeaplusb.java}}

\begin{Shaded}
\begin{Highlighting}[]
\KeywordTok{package}\ImportTok{ oj}\OperatorTok{.}\ImportTok{judge}\OperatorTok{;}
\KeywordTok{import} \ImportTok{java}\OperatorTok{.}\ImportTok{util}\OperatorTok{.}\ImportTok{Scanner}\OperatorTok{;}
\KeywordTok{public} \KeywordTok{class}\NormalTok{ AplusB }\OperatorTok{\{}                  \CommentTok{// 第6章会抽象成 Solution 接口}
    \KeywordTok{public} \BuiltInTok{String} \FunctionTok{solve}\OperatorTok{(}\BuiltInTok{String}\NormalTok{ input}\OperatorTok{)} \OperatorTok{\{}
        \BuiltInTok{Scanner}\NormalTok{ sc }\OperatorTok{=} \KeywordTok{new} \BuiltInTok{Scanner}\OperatorTok{(}\NormalTok{input}\OperatorTok{);}
        \ControlFlowTok{return} \BuiltInTok{String}\OperatorTok{.}\FunctionTok{valueOf}\OperatorTok{(}\NormalTok{sc}\OperatorTok{.}\FunctionTok{nextInt}\OperatorTok{()} \OperatorTok{+}\NormalTok{ sc}\OperatorTok{.}\FunctionTok{nextInt}\OperatorTok{());}
    \OperatorTok{\}}
\OperatorTok{\}}
\end{Highlighting}
\end{Shaded}

\hypertarget{simplejudge-ojjudgesimplejudge.javaux672cux7ae0ux6838ux5fc3}{%
\subsubsection{\texorpdfstring{SimpleJudge ---
\texttt{oj/judge/SimpleJudge.java}(本章核心)}{SimpleJudge --- oj/judge/SimpleJudge.java(本章核心)}}\label{simplejudge-ojjudgesimplejudge.javaux672cux7ae0ux6838ux5fc3}}

\begin{Shaded}
\begin{Highlighting}[]
\KeywordTok{package}\ImportTok{ oj}\OperatorTok{.}\ImportTok{judge}\OperatorTok{;}
\KeywordTok{import} \ImportTok{oj}\OperatorTok{.}\ImportTok{core}\OperatorTok{.*;}
\KeywordTok{public} \KeywordTok{class}\NormalTok{ SimpleJudge }\OperatorTok{\{}
    \CommentTok{// 方法重载示例:可指定/不指定用户名}
    \KeywordTok{public}\NormalTok{ JudgeResult }\FunctionTok{judge}\OperatorTok{(}\NormalTok{Problem p}\OperatorTok{,}\NormalTok{ AplusB solution}\OperatorTok{)} \OperatorTok{\{}
        \ControlFlowTok{return} \FunctionTok{judge}\OperatorTok{(}\NormalTok{p}\OperatorTok{,}\NormalTok{ solution}\OperatorTok{,} \StringTok{"anonymous"}\OperatorTok{);}
    \OperatorTok{\}}
    \KeywordTok{public}\NormalTok{ JudgeResult }\FunctionTok{judge}\OperatorTok{(}\NormalTok{Problem p}\OperatorTok{,}\NormalTok{ AplusB solution}\OperatorTok{,} \BuiltInTok{String}\NormalTok{ user}\OperatorTok{)} \OperatorTok{\{}
\NormalTok{        TestCase}\OperatorTok{[]}\NormalTok{ cases }\OperatorTok{=}\NormalTok{ p}\OperatorTok{.}\FunctionTok{getCases}\OperatorTok{();}
        \DataTypeTok{int}\NormalTok{ passed }\OperatorTok{=} \DecValTok{0}\OperatorTok{;}
        \BuiltInTok{String}\NormalTok{ detail }\OperatorTok{=} \StringTok{""}\OperatorTok{;}
        \DataTypeTok{long}\NormalTok{ start }\OperatorTok{=} \BuiltInTok{System}\OperatorTok{.}\FunctionTok{currentTimeMillis}\OperatorTok{();}
        \ControlFlowTok{for} \OperatorTok{(}\DataTypeTok{int}\NormalTok{ i }\OperatorTok{=} \DecValTok{0}\OperatorTok{;}\NormalTok{ i }\OperatorTok{\textless{}}\NormalTok{ cases}\OperatorTok{.}\FunctionTok{length}\OperatorTok{;}\NormalTok{ i}\OperatorTok{++)} \OperatorTok{\{}
            \BuiltInTok{String}\NormalTok{ actual   }\OperatorTok{=}\NormalTok{ solution}\OperatorTok{.}\FunctionTok{solve}\OperatorTok{(}\NormalTok{cases}\OperatorTok{[}\NormalTok{i}\OperatorTok{].}\FunctionTok{getInput}\OperatorTok{()).}\FunctionTok{trim}\OperatorTok{();}
            \BuiltInTok{String}\NormalTok{ expected }\OperatorTok{=}\NormalTok{ cases}\OperatorTok{[}\NormalTok{i}\OperatorTok{].}\FunctionTok{getExpected}\OperatorTok{().}\FunctionTok{trim}\OperatorTok{();}
            \ControlFlowTok{if} \OperatorTok{(}\NormalTok{actual}\OperatorTok{.}\FunctionTok{equals}\OperatorTok{(}\NormalTok{expected}\OperatorTok{))} \OperatorTok{\{}
\NormalTok{                passed}\OperatorTok{++;}
            \OperatorTok{\}} \ControlFlowTok{else} \ControlFlowTok{if} \OperatorTok{(}\NormalTok{detail}\OperatorTok{.}\FunctionTok{isEmpty}\OperatorTok{())} \OperatorTok{\{}
\NormalTok{                detail }\OperatorTok{=} \StringTok{"case\#"} \OperatorTok{+}\NormalTok{ i }\OperatorTok{+} \StringTok{" 期望["} \OperatorTok{+}\NormalTok{ expected }\OperatorTok{+} \StringTok{"] 实际["} \OperatorTok{+}\NormalTok{ actual }\OperatorTok{+} \StringTok{"]"}\OperatorTok{;}
            \OperatorTok{\}}
        \OperatorTok{\}}
        \DataTypeTok{long}\NormalTok{ elapsed }\OperatorTok{=} \BuiltInTok{System}\OperatorTok{.}\FunctionTok{currentTimeMillis}\OperatorTok{()} \OperatorTok{{-}}\NormalTok{ start}\OperatorTok{;}
\NormalTok{        Status st }\OperatorTok{=} \OperatorTok{(}\NormalTok{passed }\OperatorTok{==}\NormalTok{ cases}\OperatorTok{.}\FunctionTok{length}\OperatorTok{)} \OperatorTok{?}\NormalTok{ Status}\OperatorTok{.}\FunctionTok{AC} \OperatorTok{:}\NormalTok{ Status}\OperatorTok{.}\FunctionTok{WA}\OperatorTok{;}
        \ControlFlowTok{return} \KeywordTok{new} \FunctionTok{JudgeResult}\OperatorTok{(}\NormalTok{st}\OperatorTok{,}\NormalTok{ passed}\OperatorTok{,}\NormalTok{ cases}\OperatorTok{.}\FunctionTok{length}\OperatorTok{,}\NormalTok{ detail}\OperatorTok{,}\NormalTok{ elapsed}\OperatorTok{);}
    \OperatorTok{\}}
\OperatorTok{\}}
\end{Highlighting}
\end{Shaded}

\hypertarget{mainux6f14ux793a-ojmain.java}{%
\subsubsection{\texorpdfstring{Main(演示) ---
\texttt{oj/Main.java}}{Main(演示) --- oj/Main.java}}\label{mainux6f14ux793a-ojmain.java}}

\begin{Shaded}
\begin{Highlighting}[]
\KeywordTok{package}\ImportTok{ oj}\OperatorTok{;}
\KeywordTok{import} \ImportTok{oj}\OperatorTok{.}\ImportTok{core}\OperatorTok{.*;}
\KeywordTok{import} \ImportTok{oj}\OperatorTok{.}\ImportTok{judge}\OperatorTok{.*;}
\KeywordTok{public} \KeywordTok{class}\NormalTok{ Main }\OperatorTok{\{}
    \KeywordTok{public} \DataTypeTok{static} \DataTypeTok{void} \FunctionTok{main}\OperatorTok{(}\BuiltInTok{String}\OperatorTok{[]}\NormalTok{ args}\OperatorTok{)} \OperatorTok{\{}
\NormalTok{        Problem p }\OperatorTok{=} \KeywordTok{new} \FunctionTok{Problem}\OperatorTok{(}\DecValTok{1}\OperatorTok{,} \StringTok{"A+B"}\OperatorTok{,}
            \KeywordTok{new} \FunctionTok{TestCase}\OperatorTok{(}\StringTok{"1 2"}\OperatorTok{,} \StringTok{"3"}\OperatorTok{),}
            \KeywordTok{new} \FunctionTok{TestCase}\OperatorTok{(}\StringTok{"10 20"}\OperatorTok{,} \StringTok{"30"}\OperatorTok{));}
\NormalTok{        JudgeResult r }\OperatorTok{=} \KeywordTok{new} \FunctionTok{SimpleJudge}\OperatorTok{().}\FunctionTok{judge}\OperatorTok{(}\NormalTok{p}\OperatorTok{,} \KeywordTok{new} \FunctionTok{AplusB}\OperatorTok{(),} \StringTok{"alice"}\OperatorTok{);}
        \BuiltInTok{System}\OperatorTok{.}\FunctionTok{out}\OperatorTok{.}\FunctionTok{println}\OperatorTok{(}\NormalTok{p}\OperatorTok{.}\FunctionTok{getTitle}\OperatorTok{()} \OperatorTok{+} \StringTok{" {-}\textgreater{} "} \OperatorTok{+}\NormalTok{ r}\OperatorTok{);}   \CommentTok{// 用到 toString}
        \BuiltInTok{System}\OperatorTok{.}\FunctionTok{out}\OperatorTok{.}\FunctionTok{println}\OperatorTok{(}\StringTok{"总提交数: "} \OperatorTok{+}\NormalTok{ Submission}\OperatorTok{.}\FunctionTok{count}\OperatorTok{());}
    \OperatorTok{\}}
\OperatorTok{\}}
\end{Highlighting}
\end{Shaded}

\hypertarget{ux52a8ux624bux6b65ux9aa4ux7f16ux8bd1jarjavadoc}{%
\subsection{动手步骤(编译/jar/javadoc)}\label{ux52a8ux624bux6b65ux9aa4ux7f16ux8bd1jarjavadoc}}

目录:\texttt{src/oj/...}。写个 \texttt{Makefile}:

\begin{Shaded}
\begin{Highlighting}[]
\NormalTok{build:}
\NormalTok{    javac {-}d build $(shell find src {-}name \textquotesingle{}*.java\textquotesingle{})}
\NormalTok{run: build}
\NormalTok{    java {-}cp build oj.Main}
\NormalTok{jar: build}
\NormalTok{    jar cfe mini{-}oj.jar oj.Main {-}C build .}
\NormalTok{doc:}
\NormalTok{    javadoc {-}d doc {-}sourcepath src {-}subpackages oj}
\end{Highlighting}
\end{Shaded}

\begin{Shaded}
\begin{Highlighting}[]
\FunctionTok{make}\NormalTok{ run          }\CommentTok{\# 跑演示}
\FunctionTok{make}\NormalTok{ jar }\KeywordTok{\&\&} \ExtensionTok{java} \AttributeTok{{-}jar}\NormalTok{ mini{-}oj.jar}
\FunctionTok{make}\NormalTok{ doc          }\CommentTok{\# 生成 API 文档到 doc/}
\end{Highlighting}
\end{Shaded}

\hypertarget{ux52a8ux624bux5b9eux73b0ux4e24ux6b65ux8d70javac-java-ux5b9eux64cd}{%
\subsection{动手实现:两步走(javac / java
实操)}\label{ux52a8ux624bux5b9eux73b0ux4e24ux6b65ux8d70javac-java-ux5b9eux64cd}}

\begin{quote}
约定:源码放 \texttt{src/oj/...},编译到 \texttt{build/},全程
\texttt{javac}/\texttt{java},不用 make。
\end{quote}

\hypertarget{ux7b2cux4e00ux6b65-ux628a-oj-ux7684ux6570ux636eux5199ux6210ux6700ux7b80ux666eux901aux7c7b}{%
\subsubsection{\texorpdfstring{第一步 \KS --- 把 OJ
的数据写成最简普通类}{第一步 --- 把 OJ 的数据写成最简普通类}}\label{ux7b2cux4e00ux6b65-ux628a-oj-ux7684ux6570ux636eux5199ux6210ux6700ux7b80ux666eux901aux7c7b}}

\textbf{\texttt{src/oj/core/TestCase.java}}

\begin{Shaded}
\begin{Highlighting}[]
\KeywordTok{package}\ImportTok{ oj}\OperatorTok{.}\ImportTok{core}\OperatorTok{;}
\KeywordTok{public} \KeywordTok{class}\NormalTok{ TestCase }\OperatorTok{\{}
    \KeywordTok{private} \BuiltInTok{String}\NormalTok{ input}\OperatorTok{,}\NormalTok{ expected}\OperatorTok{;}
    \KeywordTok{public} \FunctionTok{TestCase}\OperatorTok{(}\BuiltInTok{String}\NormalTok{ input}\OperatorTok{,} \BuiltInTok{String}\NormalTok{ expected}\OperatorTok{)} \OperatorTok{\{} \KeywordTok{this}\OperatorTok{.}\FunctionTok{input} \OperatorTok{=}\NormalTok{ input}\OperatorTok{;} \KeywordTok{this}\OperatorTok{.}\FunctionTok{expected} \OperatorTok{=}\NormalTok{ expected}\OperatorTok{;} \OperatorTok{\}}
    \KeywordTok{public} \BuiltInTok{String} \FunctionTok{getInput}\OperatorTok{()}    \OperatorTok{\{} \ControlFlowTok{return}\NormalTok{ input}\OperatorTok{;} \OperatorTok{\}}
    \KeywordTok{public} \BuiltInTok{String} \FunctionTok{getExpected}\OperatorTok{()} \OperatorTok{\{} \ControlFlowTok{return}\NormalTok{ expected}\OperatorTok{;} \OperatorTok{\}}
\OperatorTok{\}}
\end{Highlighting}
\end{Shaded}

讲解:一组测试 = 输入 + 期望输出。考试编程题的''普通类''写法:私有字段 +
构造方法 + getter。

\textbf{\texttt{src/oj/core/Problem.java}}

\begin{Shaded}
\begin{Highlighting}[]
\KeywordTok{package}\ImportTok{ oj}\OperatorTok{.}\ImportTok{core}\OperatorTok{;}
\KeywordTok{public} \KeywordTok{class}\NormalTok{ Problem }\OperatorTok{\{}
    \KeywordTok{private} \DataTypeTok{int}\NormalTok{ id}\OperatorTok{;} \KeywordTok{private} \BuiltInTok{String}\NormalTok{ title}\OperatorTok{;} \KeywordTok{private}\NormalTok{ TestCase}\OperatorTok{[]}\NormalTok{ cases}\OperatorTok{;}
    \KeywordTok{public} \FunctionTok{Problem}\OperatorTok{(}\DataTypeTok{int}\NormalTok{ id}\OperatorTok{,} \BuiltInTok{String}\NormalTok{ title}\OperatorTok{,}\NormalTok{ TestCase}\OperatorTok{[]}\NormalTok{ cases}\OperatorTok{)} \OperatorTok{\{} \KeywordTok{this}\OperatorTok{.}\FunctionTok{id} \OperatorTok{=}\NormalTok{ id}\OperatorTok{;} \KeywordTok{this}\OperatorTok{.}\FunctionTok{title} \OperatorTok{=}\NormalTok{ title}\OperatorTok{;} \KeywordTok{this}\OperatorTok{.}\FunctionTok{cases} \OperatorTok{=}\NormalTok{ cases}\OperatorTok{;} \OperatorTok{\}}
    \KeywordTok{public} \DataTypeTok{int} \FunctionTok{getId}\OperatorTok{()}           \OperatorTok{\{} \ControlFlowTok{return}\NormalTok{ id}\OperatorTok{;} \OperatorTok{\}}
    \KeywordTok{public} \BuiltInTok{String} \FunctionTok{getTitle}\OperatorTok{()}     \OperatorTok{\{} \ControlFlowTok{return}\NormalTok{ title}\OperatorTok{;} \OperatorTok{\}}
    \KeywordTok{public}\NormalTok{ TestCase}\OperatorTok{[]} \FunctionTok{getCases}\OperatorTok{()} \OperatorTok{\{} \ControlFlowTok{return}\NormalTok{ cases}\OperatorTok{;} \OperatorTok{\}}   \CommentTok{// 对象数组}
\OperatorTok{\}}
\end{Highlighting}
\end{Shaded}

讲解:题目 = 编号 + 标题 + 一组测试点(对象数组)。

\textbf{\texttt{src/oj/core/JudgeResult.java}}

\begin{Shaded}
\begin{Highlighting}[]
\KeywordTok{package}\ImportTok{ oj}\OperatorTok{.}\ImportTok{core}\OperatorTok{;}
\KeywordTok{public} \KeywordTok{class}\NormalTok{ JudgeResult }\OperatorTok{\{}
    \KeywordTok{private} \BuiltInTok{String}\NormalTok{ status}\OperatorTok{;} \KeywordTok{private} \DataTypeTok{int}\NormalTok{ passed}\OperatorTok{,}\NormalTok{ total}\OperatorTok{;}
    \KeywordTok{public} \FunctionTok{JudgeResult}\OperatorTok{(}\BuiltInTok{String}\NormalTok{ status}\OperatorTok{,} \DataTypeTok{int}\NormalTok{ passed}\OperatorTok{,} \DataTypeTok{int}\NormalTok{ total}\OperatorTok{)} \OperatorTok{\{} \KeywordTok{this}\OperatorTok{.}\FunctionTok{status} \OperatorTok{=}\NormalTok{ status}\OperatorTok{;} \KeywordTok{this}\OperatorTok{.}\FunctionTok{passed} \OperatorTok{=}\NormalTok{ passed}\OperatorTok{;} \KeywordTok{this}\OperatorTok{.}\FunctionTok{total} \OperatorTok{=}\NormalTok{ total}\OperatorTok{;} \OperatorTok{\}}
    \AttributeTok{@Override} \KeywordTok{public} \BuiltInTok{String} \FunctionTok{toString}\OperatorTok{()} \OperatorTok{\{} \ControlFlowTok{return}\NormalTok{ status }\OperatorTok{+} \StringTok{" "} \OperatorTok{+}\NormalTok{ passed }\OperatorTok{+} \StringTok{"/"} \OperatorTok{+}\NormalTok{ total}\OperatorTok{;} \OperatorTok{\}}
\OperatorTok{\}}
\end{Highlighting}
\end{Shaded}

\textbf{\texttt{src/oj/Demo.java}}

\begin{Shaded}
\begin{Highlighting}[]
\KeywordTok{package}\ImportTok{ oj}\OperatorTok{;}
\KeywordTok{import} \ImportTok{oj}\OperatorTok{.}\ImportTok{core}\OperatorTok{.*;}
\KeywordTok{public} \KeywordTok{class}\NormalTok{ Demo }\OperatorTok{\{}
    \KeywordTok{public} \DataTypeTok{static} \DataTypeTok{void} \FunctionTok{main}\OperatorTok{(}\BuiltInTok{String}\OperatorTok{[]}\NormalTok{ args}\OperatorTok{)} \OperatorTok{\{}
\NormalTok{        Problem p }\OperatorTok{=} \KeywordTok{new} \FunctionTok{Problem}\OperatorTok{(}\DecValTok{1}\OperatorTok{,} \StringTok{"A+B"}\OperatorTok{,} \KeywordTok{new}\NormalTok{ TestCase}\OperatorTok{[]\{} \KeywordTok{new} \FunctionTok{TestCase}\OperatorTok{(}\StringTok{"1 2"}\OperatorTok{,}\StringTok{"3"}\OperatorTok{),} \KeywordTok{new} \FunctionTok{TestCase}\OperatorTok{(}\StringTok{"10 20"}\OperatorTok{,}\StringTok{"30"}\OperatorTok{)} \OperatorTok{\});}
        \BuiltInTok{System}\OperatorTok{.}\FunctionTok{out}\OperatorTok{.}\FunctionTok{println}\OperatorTok{(}\NormalTok{p}\OperatorTok{.}\FunctionTok{getTitle}\OperatorTok{()} \OperatorTok{+} \StringTok{" 有 "} \OperatorTok{+}\NormalTok{ p}\OperatorTok{.}\FunctionTok{getCases}\OperatorTok{().}\FunctionTok{length} \OperatorTok{+} \StringTok{" 组用例"}\OperatorTok{);}
        \BuiltInTok{System}\OperatorTok{.}\FunctionTok{out}\OperatorTok{.}\FunctionTok{println}\OperatorTok{(}\KeywordTok{new} \FunctionTok{JudgeResult}\OperatorTok{(}\StringTok{"AC"}\OperatorTok{,} \DecValTok{2}\OperatorTok{,} \DecValTok{2}\OperatorTok{));}
    \OperatorTok{\}}
\OperatorTok{\}}
\end{Highlighting}
\end{Shaded}

编译运行:

\begin{Shaded}
\begin{Highlighting}[]
\ExtensionTok{javac} \AttributeTok{{-}d}\NormalTok{ build src/oj/core/TestCase.java src/oj/core/Problem.java src/oj/core/JudgeResult.java src/oj/Demo.java}
\ExtensionTok{java} \AttributeTok{{-}cp}\NormalTok{ build oj.Demo}
\CommentTok{\# A+B 有 2 组用例}
\CommentTok{\# AC 2/2}
\end{Highlighting}
\end{Shaded}

\hypertarget{ux7b2cux4e8cux6b65-ux5347ux7ea7ux4e3aux5951ux7ea6ux7248}{%
\subsubsection{\texorpdfstring{第二步 \GC ---
升级为契约版}{第二步 --- 升级为契约版}}\label{ux7b2cux4e8cux6b65-ux5347ux7ea7ux4e3aux5951ux7ea6ux7248}}

把 \texttt{Status} 改 \texttt{enum}、\texttt{JudgeResult} 补
\texttt{detail/elapsedMs}、抽出
\texttt{ProblemMeta}(标题/判题类名/时限)让
\texttt{Problem\ =\ id\ +\ ProblemMeta\ +\ 测试点}、\texttt{Submission}
加 \texttt{static} 计数(代码见上文【新产物架构】)。整项目一条命令编译:

\begin{Shaded}
\begin{Highlighting}[]
\ExtensionTok{javac} \AttributeTok{{-}d}\NormalTok{ build }\VariableTok{$(}\FunctionTok{find}\NormalTok{ src }\AttributeTok{{-}name} \StringTok{\textquotesingle{}*.java\textquotesingle{}}\VariableTok{)}
\ExtensionTok{java} \AttributeTok{{-}cp}\NormalTok{ build oj.Main}
\end{Highlighting}
\end{Shaded}

\hypertarget{ux9a8cux6536ux6807ux51c6-2}{%
\subsection{验收标准}\label{ux9a8cux6536ux6807ux51c6-2}}

\begin{itemize}
\tightlist
\item[$\square$]
  M1 的判题逻辑现在全部由对象承载,\texttt{make\ run} 输出
  \texttt{A+B\ -\textgreater{}\ AC\ 2/2\ (...)}。
\item[$\square$]
  用到了:\texttt{private}+getter、构造重载、\texttt{this}、\texttt{static}
  计数、对象数组、\texttt{TestCase...} 可变参数、包与 \texttt{import}。
\item[$\square$]
  \texttt{java\ -jar\ mini-oj.jar} 能跑;\texttt{doc/index.html} 能打开。
\item[$\square$]
  能说清''加一道新题''现在只需 \texttt{new\ Problem(...)},比 M1
  改数组干净多了。
\end{itemize}

\hypertarget{ux672cux7ae0ux4ea7ux7269---ux4e0bux4e00ux7ae0-1}{%
\subsection{本章产物 -\textgreater{}
下一章}\label{ux672cux7ae0ux4ea7ux7269---ux4e0bux4e00ux7ae0-1}}

\texttt{SimpleJudge}
只会''精确比对'',且写死在一个类里。第5-7章(M3)学\textbf{接口/抽象类/多态}后,把它重构成
\texttt{Judge} 接口 + \texttt{AbstractJudge} +
多种判题器,并用\textbf{异常}处理用户解法崩溃。

\hypertarget{ux7b2c5-7ux7ae0-ux5224ux9898ux5668ux63a5ux53e3ux7ee7ux627fux591aux6001ux5f02ux5e38ux91ccux7a0bux7891-m3}{%
\section{04 · 第5-7章 判题器:接口/继承/多态/异常(里程碑
M3)}\label{ux7b2c5-7ux7ae0-ux5224ux9898ux5668ux63a5ux53e3ux7ee7ux627fux591aux6001ux5f02ux5e38ux91ccux7a0bux7891-m3}}

\begin{quote}
\textbf{考试相关度 ★★★(编程大题核心)} ·
真题考点参考:抽象类/继承/重写/多态/异常。\textbf{{[}第一步{]}}
\KS 用考试级写法写一个简单判题(把一题判出 AC/WA)+ try-catch
兜异常;\textbf{{[}第二步{]}} \GC 抽成 \texttt{Judge} 接口 /
\texttt{AbstractJudge} 模板 + 多态 + 自定义异常。
\end{quote}

\begin{quote}
\textbf{本章是面向对象的重头戏。} 我们把 M2 写死的 \texttt{SimpleJudge}
重构成一套可扩展的判题器: \texttt{Judge} 接口 -\textgreater{}
\texttt{AbstractJudge} 抽象类(模板)-\textgreater{}
\texttt{StandardJudge}
子类(继承+重写);并用\textbf{自定义异常}处理用户解法崩溃。
主线一句话:\textbf{抽象出''判题器''这个概念,以后加新判法只加类、不改老代码(开闭原则)。}
\end{quote}

\hypertarget{ux77e5ux8bc6ux70b9ux5927ux7eb2ux5bf9ux7167ux672cux7ae0ux8986ux76d6ux6700ux5bc6}{%
\subsection{知识点(大纲对照)------本章覆盖最密}\label{ux77e5ux8bc6ux70b9ux5927ux7eb2ux5bf9ux7167ux672cux7ae0ux8986ux76d6ux6700ux5bc6}}

\begin{longtable}[]{@{}
  >{\raggedright\arraybackslash}p{(\columnwidth - 4\tabcolsep) * \real{0.3333}}
  >{\raggedright\arraybackslash}p{(\columnwidth - 4\tabcolsep) * \real{0.3333}}
  >{\raggedright\arraybackslash}p{(\columnwidth - 4\tabcolsep) * \real{0.3333}}@{}}
\toprule\noalign{}
\begin{minipage}[b]{\linewidth}\raggedright
主题
\end{minipage} & \begin{minipage}[b]{\linewidth}\raggedright
要点
\end{minipage} & \begin{minipage}[b]{\linewidth}\raggedright
处理
\end{minipage} \\
\midrule\noalign{}
\endhead
\bottomrule\noalign{}
\endlastfoot
接口(Ch6) &
\texttt{interface}/实现/接口回调/函数接口+Lambda/接口参数/接口多态 &
\textbf{动手}(\texttt{Solution}、\texttt{Judge}) \\
继承(Ch5) & 子类 \texttt{extends}/方法重写
\texttt{@Override}/\texttt{super}/\texttt{final}/上转型/多态 &
\textbf{动手}(\texttt{AbstractJudge}-\textgreater 子类) \\
抽象(Ch5) & \texttt{abstract} 类与方法/面向抽象/开闭原则 &
\textbf{动手}(模板方法) \\
内部类(Ch7) & 内部类/匿名类/Lambda 代替匿名类 &
\textbf{动手}(比较器回调) \\
异常(Ch7) & \texttt{try-catch}/自定义异常类/断言 & \textbf{动手}(RE 处理
+ \texttt{JudgeException}) \\
成员变量隐藏 & (Ch5) & \textbf{了解即可},不刻意制造 \\
\end{longtable}

\hypertarget{ux7b2cux4e00ux6b65ux628a-solution-ux53d8ux6210ux63a5ux53e3ch6}{%
\subsection{\texorpdfstring{第一步:把 \texttt{Solution}
变成接口(Ch6)}{第一步:把 Solution 变成接口(Ch6)}}\label{ux7b2cux4e00ux6b65ux628a-solution-ux53d8ux6210ux63a5ux53e3ch6}}

\texttt{oj/judge/Solution.java}

\begin{Shaded}
\begin{Highlighting}[]
\KeywordTok{package}\ImportTok{ oj}\OperatorTok{.}\ImportTok{judge}\OperatorTok{;}
\KeywordTok{public} \KeywordTok{interface}\NormalTok{ Solution }\OperatorTok{\{}            \CommentTok{// 函数式接口(只一个抽象方法){-}\textgreater{} 可用 Lambda}
    \BuiltInTok{String} \FunctionTok{solve}\OperatorTok{(}\BuiltInTok{String}\NormalTok{ input}\OperatorTok{);}
\OperatorTok{\}}
\end{Highlighting}
\end{Shaded}

原来的 \texttt{AplusB} 改成实现接口:

\begin{Shaded}
\begin{Highlighting}[]
\KeywordTok{package}\ImportTok{ oj}\OperatorTok{.}\ImportTok{judge}\OperatorTok{;}
\KeywordTok{import} \ImportTok{java}\OperatorTok{.}\ImportTok{util}\OperatorTok{.}\ImportTok{Scanner}\OperatorTok{;}
\KeywordTok{public} \KeywordTok{class}\NormalTok{ AplusB }\KeywordTok{implements}\NormalTok{ Solution }\OperatorTok{\{}
    \AttributeTok{@Override} \KeywordTok{public} \BuiltInTok{String} \FunctionTok{solve}\OperatorTok{(}\BuiltInTok{String}\NormalTok{ input}\OperatorTok{)} \OperatorTok{\{}
        \BuiltInTok{Scanner}\NormalTok{ sc }\OperatorTok{=} \KeywordTok{new} \BuiltInTok{Scanner}\OperatorTok{(}\NormalTok{input}\OperatorTok{);}
        \ControlFlowTok{return} \BuiltInTok{String}\OperatorTok{.}\FunctionTok{valueOf}\OperatorTok{(}\NormalTok{sc}\OperatorTok{.}\FunctionTok{nextInt}\OperatorTok{()} \OperatorTok{+}\NormalTok{ sc}\OperatorTok{.}\FunctionTok{nextInt}\OperatorTok{());}
    \OperatorTok{\}}
\OperatorTok{\}}
\end{Highlighting}
\end{Shaded}

\textbf{接口=函数式接口的好处}:用户解法可以直接用 Lambda 写,不必建类:

\begin{Shaded}
\begin{Highlighting}[]
\NormalTok{Solution s }\OperatorTok{=}\NormalTok{ input }\OperatorTok{{-}\textgreater{}} \OperatorTok{\{}                 \CommentTok{// Lambda 实现 Solution}
    \BuiltInTok{Scanner}\NormalTok{ sc }\OperatorTok{=} \KeywordTok{new} \BuiltInTok{Scanner}\OperatorTok{(}\NormalTok{input}\OperatorTok{);}
    \ControlFlowTok{return} \BuiltInTok{String}\OperatorTok{.}\FunctionTok{valueOf}\OperatorTok{(}\NormalTok{sc}\OperatorTok{.}\FunctionTok{nextInt}\OperatorTok{()} \OperatorTok{+}\NormalTok{ sc}\OperatorTok{.}\FunctionTok{nextInt}\OperatorTok{());}
\OperatorTok{\};}
\end{Highlighting}
\end{Shaded}

\hypertarget{ux7b2cux4e8cux6b65judge-ux63a5ux53e3-abstractjudge-ux62bdux8c61ux7c7bch5ch6-ux5408ux4f53}{%
\subsection{\texorpdfstring{第二步:\texttt{Judge} 接口 +
\texttt{AbstractJudge} 抽象类(Ch5+Ch6
合体)}{第二步:Judge 接口 + AbstractJudge 抽象类(Ch5+Ch6 合体)}}\label{ux7b2cux4e8cux6b65judge-ux63a5ux53e3-abstractjudge-ux62bdux8c61ux7c7bch5ch6-ux5408ux4f53}}

\hypertarget{ux63a5ux53e3-ojjudgejudge.java}{%
\subsubsection{\texorpdfstring{接口
\texttt{oj/judge/Judge.java}}{接口 oj/judge/Judge.java}}\label{ux63a5ux53e3-ojjudgejudge.java}}

\begin{Shaded}
\begin{Highlighting}[]
\KeywordTok{package}\ImportTok{ oj}\OperatorTok{.}\ImportTok{judge}\OperatorTok{;}
\KeywordTok{import} \ImportTok{oj}\OperatorTok{.}\ImportTok{core}\OperatorTok{.*;}
\KeywordTok{public} \KeywordTok{interface}\NormalTok{ Judge }\OperatorTok{\{}
\NormalTok{    JudgeResult }\FunctionTok{judge}\OperatorTok{(}\NormalTok{Problem problem}\OperatorTok{,}\NormalTok{ Solution solution}\OperatorTok{);}
\OperatorTok{\}}
\end{Highlighting}
\end{Shaded}

\hypertarget{ux62bdux8c61ux7c7bux6a21ux677fux65b9ux6cd5ojjudgeabstractjudge.javaux672cux7ae0ux6838ux5fc3}{%
\subsubsection{\texorpdfstring{抽象类(模板方法)\texttt{oj/judge/AbstractJudge.java}------\textbf{本章核心}}{抽象类(模板方法)oj/judge/AbstractJudge.java------本章核心}}\label{ux62bdux8c61ux7c7bux6a21ux677fux65b9ux6cd5ojjudgeabstractjudge.javaux672cux7ae0ux6838ux5fc3}}

\begin{Shaded}
\begin{Highlighting}[]
\KeywordTok{package}\ImportTok{ oj}\OperatorTok{.}\ImportTok{judge}\OperatorTok{;}
\KeywordTok{import} \ImportTok{oj}\OperatorTok{.}\ImportTok{core}\OperatorTok{.*;}
\KeywordTok{public} \KeywordTok{abstract} \KeywordTok{class}\NormalTok{ AbstractJudge }\KeywordTok{implements}\NormalTok{ Judge }\OperatorTok{\{}
    \CommentTok{// final:模板流程固定,子类不许改;子类只能定制 compare(开闭原则的"闭")}
    \AttributeTok{@Override}
    \KeywordTok{public} \DataTypeTok{final}\NormalTok{ JudgeResult }\FunctionTok{judge}\OperatorTok{(}\NormalTok{Problem p}\OperatorTok{,}\NormalTok{ Solution s}\OperatorTok{)} \OperatorTok{\{}
\NormalTok{        TestCase}\OperatorTok{[]}\NormalTok{ cases }\OperatorTok{=}\NormalTok{ p}\OperatorTok{.}\FunctionTok{getCases}\OperatorTok{();}
        \DataTypeTok{int}\NormalTok{ passed }\OperatorTok{=} \DecValTok{0}\OperatorTok{;}
        \BuiltInTok{String}\NormalTok{ detail }\OperatorTok{=} \StringTok{""}\OperatorTok{;}
        \DataTypeTok{long}\NormalTok{ start }\OperatorTok{=} \BuiltInTok{System}\OperatorTok{.}\FunctionTok{currentTimeMillis}\OperatorTok{();}
        \ControlFlowTok{try} \OperatorTok{\{}
            \ControlFlowTok{for} \OperatorTok{(}\DataTypeTok{int}\NormalTok{ i }\OperatorTok{=} \DecValTok{0}\OperatorTok{;}\NormalTok{ i }\OperatorTok{\textless{}}\NormalTok{ cases}\OperatorTok{.}\FunctionTok{length}\OperatorTok{;}\NormalTok{ i}\OperatorTok{++)} \OperatorTok{\{}
                \BuiltInTok{String}\NormalTok{ actual }\OperatorTok{=}\NormalTok{ s}\OperatorTok{.}\FunctionTok{solve}\OperatorTok{(}\NormalTok{cases}\OperatorTok{[}\NormalTok{i}\OperatorTok{].}\FunctionTok{getInput}\OperatorTok{());}   \CommentTok{// 用户解法,可能抛异常}
                \ControlFlowTok{if} \OperatorTok{(}\FunctionTok{compare}\OperatorTok{(}\NormalTok{cases}\OperatorTok{[}\NormalTok{i}\OperatorTok{].}\FunctionTok{getExpected}\OperatorTok{(),}\NormalTok{ actual}\OperatorTok{))} \OperatorTok{\{}  \CommentTok{// 调子类实现(多态)}
\NormalTok{                    passed}\OperatorTok{++;}
                \OperatorTok{\}} \ControlFlowTok{else} \ControlFlowTok{if} \OperatorTok{(}\NormalTok{detail}\OperatorTok{.}\FunctionTok{isEmpty}\OperatorTok{())} \OperatorTok{\{}
\NormalTok{                    detail }\OperatorTok{=} \StringTok{"case\#"} \OperatorTok{+}\NormalTok{ i }\OperatorTok{+} \StringTok{" 期望["} \OperatorTok{+}\NormalTok{ cases}\OperatorTok{[}\NormalTok{i}\OperatorTok{].}\FunctionTok{getExpected}\OperatorTok{().}\FunctionTok{trim}\OperatorTok{()}
                           \OperatorTok{+} \StringTok{"] 实际["} \OperatorTok{+} \OperatorTok{(}\NormalTok{actual }\OperatorTok{==} \KeywordTok{null} \OperatorTok{?} \StringTok{"null"} \OperatorTok{:}\NormalTok{ actual}\OperatorTok{.}\FunctionTok{trim}\OperatorTok{())} \OperatorTok{+} \StringTok{"]"}\OperatorTok{;}
                \OperatorTok{\}}
            \OperatorTok{\}}
        \OperatorTok{\}} \ControlFlowTok{catch} \OperatorTok{(}\BuiltInTok{RuntimeException}\NormalTok{ e}\OperatorTok{)} \OperatorTok{\{}                          \CommentTok{// 捕获运行期异常 {-}\textgreater{} RE}
            \DataTypeTok{long}\NormalTok{ t }\OperatorTok{=} \BuiltInTok{System}\OperatorTok{.}\FunctionTok{currentTimeMillis}\OperatorTok{()} \OperatorTok{{-}}\NormalTok{ start}\OperatorTok{;}
            \ControlFlowTok{return} \KeywordTok{new} \FunctionTok{JudgeResult}\OperatorTok{(}\NormalTok{Status}\OperatorTok{.}\FunctionTok{RE}\OperatorTok{,}\NormalTok{ passed}\OperatorTok{,}\NormalTok{ cases}\OperatorTok{.}\FunctionTok{length}\OperatorTok{,}
                                   \StringTok{"运行异常: "} \OperatorTok{+}\NormalTok{ e}\OperatorTok{.}\FunctionTok{getMessage}\OperatorTok{(),}\NormalTok{ t}\OperatorTok{);}
        \OperatorTok{\}}
        \DataTypeTok{long}\NormalTok{ elapsed }\OperatorTok{=} \BuiltInTok{System}\OperatorTok{.}\FunctionTok{currentTimeMillis}\OperatorTok{()} \OperatorTok{{-}}\NormalTok{ start}\OperatorTok{;}
        \ControlFlowTok{assert}\NormalTok{ passed }\OperatorTok{\textless{}=}\NormalTok{ cases}\OperatorTok{.}\FunctionTok{length} \OperatorTok{:} \StringTok{"passed 不应超过 total"}\OperatorTok{;}   \CommentTok{// 断言:内部不变量}
\NormalTok{        Status st }\OperatorTok{=} \OperatorTok{(}\NormalTok{passed }\OperatorTok{==}\NormalTok{ cases}\OperatorTok{.}\FunctionTok{length}\OperatorTok{)} \OperatorTok{?}\NormalTok{ Status}\OperatorTok{.}\FunctionTok{AC} \OperatorTok{:}\NormalTok{ Status}\OperatorTok{.}\FunctionTok{WA}\OperatorTok{;}
        \ControlFlowTok{return} \KeywordTok{new} \FunctionTok{JudgeResult}\OperatorTok{(}\NormalTok{st}\OperatorTok{,}\NormalTok{ passed}\OperatorTok{,}\NormalTok{ cases}\OperatorTok{.}\FunctionTok{length}\OperatorTok{,}\NormalTok{ detail}\OperatorTok{,}\NormalTok{ elapsed}\OperatorTok{);}
    \OperatorTok{\}}

    \CommentTok{// 抽象方法:"怎么比对一组输出"留给子类 {-}\textgreater{} 这就是可扩展点}
    \KeywordTok{protected} \KeywordTok{abstract} \DataTypeTok{boolean} \FunctionTok{compare}\OperatorTok{(}\BuiltInTok{String}\NormalTok{ expected}\OperatorTok{,} \BuiltInTok{String}\NormalTok{ actual}\OperatorTok{);}
\OperatorTok{\}}
\end{Highlighting}
\end{Shaded}

\begin{quote}
这一个抽象方法 \texttt{compare} 是整套设计的''活接头'':换判法 =
写个新子类重写 \texttt{compare},模板流程一行不动。
\end{quote}

\hypertarget{ux5b50ux7c7bux7cbeux786eux5224ux9898ux5668-ojjudgestandardjudge.javaux7ee7ux627f-ux91cdux5199-super}{%
\subsubsection{\texorpdfstring{子类:精确判题器
\texttt{oj/judge/StandardJudge.java}(继承 + 重写 +
super)}{子类:精确判题器 oj/judge/StandardJudge.java(继承 + 重写 + super)}}\label{ux5b50ux7c7bux7cbeux786eux5224ux9898ux5668-ojjudgestandardjudge.javaux7ee7ux627f-ux91cdux5199-super}}

\begin{Shaded}
\begin{Highlighting}[]
\KeywordTok{package}\ImportTok{ oj}\OperatorTok{.}\ImportTok{judge}\OperatorTok{;}
\KeywordTok{public} \KeywordTok{class}\NormalTok{ StandardJudge }\KeywordTok{extends}\NormalTok{ AbstractJudge }\OperatorTok{\{}
    \AttributeTok{@Override}
    \KeywordTok{protected} \DataTypeTok{boolean} \FunctionTok{compare}\OperatorTok{(}\BuiltInTok{String}\NormalTok{ expected}\OperatorTok{,} \BuiltInTok{String}\NormalTok{ actual}\OperatorTok{)} \OperatorTok{\{}
        \ControlFlowTok{if} \OperatorTok{(}\NormalTok{actual }\OperatorTok{==} \KeywordTok{null}\OperatorTok{)} \ControlFlowTok{return} \KeywordTok{false}\OperatorTok{;}
        \ControlFlowTok{return}\NormalTok{ expected}\OperatorTok{.}\FunctionTok{trim}\OperatorTok{().}\FunctionTok{equals}\OperatorTok{(}\NormalTok{actual}\OperatorTok{.}\FunctionTok{trim}\OperatorTok{());}      \CommentTok{// 去首尾空白后精确比对}
    \OperatorTok{\}}
\OperatorTok{\}}
\end{Highlighting}
\end{Shaded}

\begin{quote}
\texttt{super} 的典型用法(可选演示):子类构造里 \texttt{super(...)}
调父类构造;本章父类无显式构造,故从略,留到有字段的子类时再讲。
\end{quote}

\hypertarget{ux7b2cux4e09ux6b65ux591aux6001-ux5f00ux95edux539fux5219ch5}{%
\subsection{第三步:多态 +
开闭原则(Ch5)}\label{ux7b2cux4e09ux6b65ux591aux6001-ux5f00ux95edux539fux5219ch5}}

\begin{Shaded}
\begin{Highlighting}[]
\NormalTok{Judge judge}\OperatorTok{;}
\ControlFlowTok{if} \OperatorTok{(}\NormalTok{useFloat}\OperatorTok{)}\NormalTok{ judge }\OperatorTok{=} \KeywordTok{new} \FunctionTok{SpecialJudge}\OperatorTok{();}   \CommentTok{// 第8章(M4)实现的浮点判题器}
\ControlFlowTok{else}\NormalTok{          judge }\OperatorTok{=} \KeywordTok{new} \FunctionTok{StandardJudge}\OperatorTok{();}  \CommentTok{// 上转型:父类型引用指向子类对象}
\NormalTok{JudgeResult r }\OperatorTok{=}\NormalTok{ judge}\OperatorTok{.}\FunctionTok{judge}\OperatorTok{(}\NormalTok{problem}\OperatorTok{,}\NormalTok{ solution}\OperatorTok{);}   \CommentTok{// 运行时多态:调到具体子类的 compare}
\end{Highlighting}
\end{Shaded}

\textbf{开闭原则}:将来要支持''忽略空白''``按行无序''等新判法,只新增
\texttt{XxxJudge\ extends\ AbstractJudge},
\texttt{Main}/调用方代码完全不用改。

\hypertarget{ux7b2cux56dbux6b65ux5185ux90e8ux7c7b-ux533fux540dux7c7b-lambdach7}{%
\subsection{第四步:内部类 / 匿名类 /
Lambda(Ch7)}\label{ux7b2cux56dbux6b65ux5185ux90e8ux7c7b-ux533fux540dux7c7b-lambdach7}}

判题里常要一个''输出比较器''回调。三种写法演示\textbf{演化关系}:

\begin{Shaded}
\begin{Highlighting}[]
\CommentTok{// 1) 命名内部类}
\KeywordTok{class}\NormalTok{ IgnoreCaseComparator }\KeywordTok{implements}\NormalTok{ OutputComparator }\OperatorTok{\{}
    \KeywordTok{public} \DataTypeTok{boolean} \FunctionTok{same}\OperatorTok{(}\BuiltInTok{String}\NormalTok{ a}\OperatorTok{,} \BuiltInTok{String}\NormalTok{ b}\OperatorTok{)} \OperatorTok{\{} \ControlFlowTok{return}\NormalTok{ a}\OperatorTok{.}\FunctionTok{trim}\OperatorTok{().}\FunctionTok{equalsIgnoreCase}\OperatorTok{(}\NormalTok{b}\OperatorTok{.}\FunctionTok{trim}\OperatorTok{());} \OperatorTok{\}}
\OperatorTok{\}}
\CommentTok{// 2) 匿名类(就地实现)}
\NormalTok{OutputComparator c2 }\OperatorTok{=} \KeywordTok{new} \FunctionTok{OutputComparator}\OperatorTok{()} \OperatorTok{\{}
    \KeywordTok{public} \DataTypeTok{boolean} \FunctionTok{same}\OperatorTok{(}\BuiltInTok{String}\NormalTok{ a}\OperatorTok{,} \BuiltInTok{String}\NormalTok{ b}\OperatorTok{)} \OperatorTok{\{} \ControlFlowTok{return}\NormalTok{ a}\OperatorTok{.}\FunctionTok{trim}\OperatorTok{().}\FunctionTok{equals}\OperatorTok{(}\NormalTok{b}\OperatorTok{.}\FunctionTok{trim}\OperatorTok{());} \OperatorTok{\}}
\OperatorTok{\};}
\CommentTok{// 3) Lambda 代替匿名类(函数式接口才行)}
\NormalTok{OutputComparator c3 }\OperatorTok{=} \OperatorTok{(}\NormalTok{a}\OperatorTok{,}\NormalTok{ b}\OperatorTok{)} \OperatorTok{{-}\textgreater{}}\NormalTok{ a}\OperatorTok{.}\FunctionTok{trim}\OperatorTok{().}\FunctionTok{equals}\OperatorTok{(}\NormalTok{b}\OperatorTok{.}\FunctionTok{trim}\OperatorTok{());}
\end{Highlighting}
\end{Shaded}

其中:

\begin{Shaded}
\begin{Highlighting}[]
\KeywordTok{package}\ImportTok{ oj}\OperatorTok{.}\ImportTok{judge}\OperatorTok{;}
\KeywordTok{public} \KeywordTok{interface}\NormalTok{ OutputComparator }\OperatorTok{\{}     \CommentTok{// 函数式接口}
    \DataTypeTok{boolean} \FunctionTok{same}\OperatorTok{(}\BuiltInTok{String}\NormalTok{ a}\OperatorTok{,} \BuiltInTok{String}\NormalTok{ b}\OperatorTok{);}
\OperatorTok{\}}
\end{Highlighting}
\end{Shaded}

\begin{quote}
体会:匿名类 -\textgreater{} Lambda
是同一件事的简写。\texttt{OutputComparator} 属''可选增强'',核心判题用
\texttt{compare} 即可,不强求全用上。
\end{quote}

\hypertarget{ux7b2cux4e94ux6b65ux81eaux5b9aux4e49ux5f02ux5e38ch7}{%
\subsection{第五步:自定义异常(Ch7)}\label{ux7b2cux4e94ux6b65ux81eaux5b9aux4e49ux5f02ux5e38ch7}}

\texttt{oj/exception/JudgeException.java}

\begin{Shaded}
\begin{Highlighting}[]
\KeywordTok{package}\ImportTok{ oj}\OperatorTok{.}\ImportTok{exception}\OperatorTok{;}
\KeywordTok{public} \KeywordTok{class}\NormalTok{ JudgeException }\KeywordTok{extends} \BuiltInTok{Exception} \OperatorTok{\{}     \CommentTok{// 受检异常:判题流程的异常基类}
    \KeywordTok{public} \FunctionTok{JudgeException}\OperatorTok{(}\BuiltInTok{String}\NormalTok{ message}\OperatorTok{)} \OperatorTok{\{} \KeywordTok{super}\OperatorTok{(}\NormalTok{message}\OperatorTok{);} \OperatorTok{\}}
\OperatorTok{\}}
\end{Highlighting}
\end{Shaded}

\texttt{oj/exception/TimeLimitException.java}

\begin{Shaded}
\begin{Highlighting}[]
\KeywordTok{package}\ImportTok{ oj}\OperatorTok{.}\ImportTok{exception}\OperatorTok{;}
\KeywordTok{public} \KeywordTok{class}\NormalTok{ TimeLimitException }\KeywordTok{extends}\NormalTok{ JudgeException }\OperatorTok{\{}   \CommentTok{// 继承自定义异常}
    \KeywordTok{public} \FunctionTok{TimeLimitException}\OperatorTok{(}\BuiltInTok{String}\NormalTok{ message}\OperatorTok{)} \OperatorTok{\{} \KeywordTok{super}\OperatorTok{(}\NormalTok{message}\OperatorTok{);} \OperatorTok{\}}
\OperatorTok{\}}
\end{Highlighting}
\end{Shaded}

\begin{quote}
本章先\textbf{定义}这两个异常类(练''怎么写自定义异常、怎么继承异常'')。\texttt{TimeLimitException}
的\textbf{真正触发}在第12章(M6a)用线程做超时控制时;本章 RE 用捕获
\texttt{RuntimeException} 即可。
\end{quote}

\hypertarget{ux7b2cux516dux6b65ux66f4ux65b0-main-ux6f14ux793a}{%
\subsection{第六步:更新 Main
演示}\label{ux7b2cux516dux6b65ux66f4ux65b0-main-ux6f14ux793a}}

\begin{Shaded}
\begin{Highlighting}[]
\KeywordTok{package}\ImportTok{ oj}\OperatorTok{;}
\KeywordTok{import} \ImportTok{oj}\OperatorTok{.}\ImportTok{core}\OperatorTok{.*;}
\KeywordTok{import} \ImportTok{oj}\OperatorTok{.}\ImportTok{judge}\OperatorTok{.*;}
\KeywordTok{public} \KeywordTok{class}\NormalTok{ Main }\OperatorTok{\{}
    \KeywordTok{public} \DataTypeTok{static} \DataTypeTok{void} \FunctionTok{main}\OperatorTok{(}\BuiltInTok{String}\OperatorTok{[]}\NormalTok{ args}\OperatorTok{)} \OperatorTok{\{}
\NormalTok{        Problem p }\OperatorTok{=} \KeywordTok{new} \FunctionTok{Problem}\OperatorTok{(}\DecValTok{1}\OperatorTok{,} \StringTok{"A+B"}\OperatorTok{,}
            \KeywordTok{new} \FunctionTok{TestCase}\OperatorTok{(}\StringTok{"1 2"}\OperatorTok{,} \StringTok{"3"}\OperatorTok{),} \KeywordTok{new} \FunctionTok{TestCase}\OperatorTok{(}\StringTok{"10 20"}\OperatorTok{,} \StringTok{"30"}\OperatorTok{));}

\NormalTok{        Judge judge }\OperatorTok{=} \KeywordTok{new} \FunctionTok{StandardJudge}\OperatorTok{();}          \CommentTok{// 多态:用接口类型接收}
\NormalTok{        Solution good }\OperatorTok{=} \KeywordTok{new} \FunctionTok{AplusB}\OperatorTok{();}
        \BuiltInTok{System}\OperatorTok{.}\FunctionTok{out}\OperatorTok{.}\FunctionTok{println}\OperatorTok{(}\StringTok{"正常: "} \OperatorTok{+}\NormalTok{ judge}\OperatorTok{.}\FunctionTok{judge}\OperatorTok{(}\NormalTok{p}\OperatorTok{,}\NormalTok{ good}\OperatorTok{));}        \CommentTok{// AC 2/2}

\NormalTok{        Solution bad }\OperatorTok{=}\NormalTok{ input }\OperatorTok{{-}\textgreater{}} \OperatorTok{\{} \ControlFlowTok{throw} \KeywordTok{new} \BuiltInTok{RuntimeException}\OperatorTok{(}\StringTok{"boom"}\OperatorTok{);} \OperatorTok{\};}  \CommentTok{// Lambda 造一个会崩的解法}
        \BuiltInTok{System}\OperatorTok{.}\FunctionTok{out}\OperatorTok{.}\FunctionTok{println}\OperatorTok{(}\StringTok{"崩溃: "} \OperatorTok{+}\NormalTok{ judge}\OperatorTok{.}\FunctionTok{judge}\OperatorTok{(}\NormalTok{p}\OperatorTok{,}\NormalTok{ bad}\OperatorTok{));}         \CommentTok{// RE}
    \OperatorTok{\}}
\OperatorTok{\}}
\end{Highlighting}
\end{Shaded}

\begin{quote}
跑断言需加 \texttt{-ea}:\texttt{java\ -ea\ -cp\ build\ oj.Main}。
\end{quote}

\hypertarget{ux52a8ux624bux5b9eux73b0ux4e24ux6b65ux8d70javac-java-ux5b9eux64cd-1}{%
\subsection{动手实现:两步走(javac / java
实操)}\label{ux52a8ux624bux5b9eux73b0ux4e24ux6b65ux8d70javac-java-ux5b9eux64cd-1}}

\begin{quote}
接上一章 M2 的 \texttt{core} 类(\texttt{Problem/TestCase/JudgeResult})。
\end{quote}

\hypertarget{ux7b2cux4e00ux6b65-ux4e00ux4e2aux7b80ux5355ux5224ux9898ux5668ux5224ux51fa-acwaux5e76ux7528-try-catch-ux515cux5f02ux5e38}{%
\subsubsection{\texorpdfstring{第一步 \KS --- 一个简单判题器,判出
AC/WA,并用 try-catch
兜异常}{第一步 --- 一个简单判题器,判出 AC/WA,并用 try-catch 兜异常}}\label{ux7b2cux4e00ux6b65-ux4e00ux4e2aux7b80ux5355ux5224ux9898ux5668ux5224ux51fa-acwaux5e76ux7528-try-catch-ux515cux5f02ux5e38}}

\textbf{\texttt{src/oj/judge/Solution.java}}(函数式接口,考点 Ch6)

\begin{Shaded}
\begin{Highlighting}[]
\KeywordTok{package}\ImportTok{ oj}\OperatorTok{.}\ImportTok{judge}\OperatorTok{;}
\KeywordTok{public} \KeywordTok{interface}\NormalTok{ Solution }\OperatorTok{\{} \BuiltInTok{String} \FunctionTok{solve}\OperatorTok{(}\BuiltInTok{String}\NormalTok{ input}\OperatorTok{);} \OperatorTok{\}}   \CommentTok{// 输入一组数据 {-}\textgreater{} 产生输出}
\end{Highlighting}
\end{Shaded}

\textbf{\texttt{src/oj/judge/SimpleJudge.java}}

\begin{Shaded}
\begin{Highlighting}[]
\KeywordTok{package}\ImportTok{ oj}\OperatorTok{.}\ImportTok{judge}\OperatorTok{;}
\KeywordTok{import} \ImportTok{oj}\OperatorTok{.}\ImportTok{core}\OperatorTok{.*;}
\KeywordTok{public} \KeywordTok{class}\NormalTok{ SimpleJudge }\OperatorTok{\{}
    \KeywordTok{public}\NormalTok{ JudgeResult }\FunctionTok{judge}\OperatorTok{(}\NormalTok{Problem p}\OperatorTok{,}\NormalTok{ Solution s}\OperatorTok{)} \OperatorTok{\{}
\NormalTok{        TestCase}\OperatorTok{[]}\NormalTok{ cs }\OperatorTok{=}\NormalTok{ p}\OperatorTok{.}\FunctionTok{getCases}\OperatorTok{();}
        \DataTypeTok{int}\NormalTok{ passed }\OperatorTok{=} \DecValTok{0}\OperatorTok{;}
        \ControlFlowTok{try} \OperatorTok{\{}
            \ControlFlowTok{for} \OperatorTok{(}\NormalTok{TestCase c }\OperatorTok{:}\NormalTok{ cs}\OperatorTok{)}
                \ControlFlowTok{if} \OperatorTok{(}\NormalTok{s}\OperatorTok{.}\FunctionTok{solve}\OperatorTok{(}\NormalTok{c}\OperatorTok{.}\FunctionTok{getInput}\OperatorTok{()).}\FunctionTok{trim}\OperatorTok{().}\FunctionTok{equals}\OperatorTok{(}\NormalTok{c}\OperatorTok{.}\FunctionTok{getExpected}\OperatorTok{().}\FunctionTok{trim}\OperatorTok{()))}\NormalTok{ passed}\OperatorTok{++;}
        \OperatorTok{\}} \ControlFlowTok{catch} \OperatorTok{(}\BuiltInTok{RuntimeException}\NormalTok{ e}\OperatorTok{)} \OperatorTok{\{}
            \ControlFlowTok{return} \KeywordTok{new} \FunctionTok{JudgeResult}\OperatorTok{(}\StringTok{"RE"}\OperatorTok{,}\NormalTok{ passed}\OperatorTok{,}\NormalTok{ cs}\OperatorTok{.}\FunctionTok{length}\OperatorTok{);}   \CommentTok{// 用户解法崩溃 {-}\textgreater{} RE}
        \OperatorTok{\}}
        \ControlFlowTok{return} \KeywordTok{new} \FunctionTok{JudgeResult}\OperatorTok{(}\NormalTok{passed }\OperatorTok{==}\NormalTok{ cs}\OperatorTok{.}\FunctionTok{length} \OperatorTok{?} \StringTok{"AC"} \OperatorTok{:} \StringTok{"WA"}\OperatorTok{,}\NormalTok{ passed}\OperatorTok{,}\NormalTok{ cs}\OperatorTok{.}\FunctionTok{length}\OperatorTok{);}
    \OperatorTok{\}}
\OperatorTok{\}}
\end{Highlighting}
\end{Shaded}

讲解:遍历用例、\texttt{trim().equals()} 比对、\texttt{try-catch}
把崩溃兜成 RE------考试读程序/编程题的常见写法。

\textbf{\texttt{src/oj/Demo3.java}}

\begin{Shaded}
\begin{Highlighting}[]
\KeywordTok{package}\ImportTok{ oj}\OperatorTok{;}
\KeywordTok{import} \ImportTok{oj}\OperatorTok{.}\ImportTok{core}\OperatorTok{.*;}
\KeywordTok{import} \ImportTok{oj}\OperatorTok{.}\ImportTok{judge}\OperatorTok{.*;}
\KeywordTok{public} \KeywordTok{class}\NormalTok{ Demo3 }\OperatorTok{\{}
    \KeywordTok{public} \DataTypeTok{static} \DataTypeTok{void} \FunctionTok{main}\OperatorTok{(}\BuiltInTok{String}\OperatorTok{[]}\NormalTok{ args}\OperatorTok{)} \OperatorTok{\{}
\NormalTok{        Problem p }\OperatorTok{=} \KeywordTok{new} \FunctionTok{Problem}\OperatorTok{(}\DecValTok{1}\OperatorTok{,} \StringTok{"A+B"}\OperatorTok{,} \KeywordTok{new}\NormalTok{ TestCase}\OperatorTok{[]\{} \KeywordTok{new} \FunctionTok{TestCase}\OperatorTok{(}\StringTok{"1 2"}\OperatorTok{,}\StringTok{"3"}\OperatorTok{),} \KeywordTok{new} \FunctionTok{TestCase}\OperatorTok{(}\StringTok{"10 20"}\OperatorTok{,}\StringTok{"30"}\OperatorTok{)} \OperatorTok{\});}
\NormalTok{        Solution good }\OperatorTok{=}\NormalTok{ in }\OperatorTok{{-}\textgreater{}} \OperatorTok{\{} \BuiltInTok{String}\OperatorTok{[]}\NormalTok{ t }\OperatorTok{=}\NormalTok{ in}\OperatorTok{.}\FunctionTok{split}\OperatorTok{(}\StringTok{"}\SpecialCharTok{\textbackslash{}\textbackslash{}}\StringTok{s+"}\OperatorTok{);} \ControlFlowTok{return} \StringTok{""} \OperatorTok{+} \OperatorTok{(}\BuiltInTok{Integer}\OperatorTok{.}\FunctionTok{parseInt}\OperatorTok{(}\NormalTok{t}\OperatorTok{[}\DecValTok{0}\OperatorTok{])} \OperatorTok{+} \BuiltInTok{Integer}\OperatorTok{.}\FunctionTok{parseInt}\OperatorTok{(}\NormalTok{t}\OperatorTok{[}\DecValTok{1}\OperatorTok{]));} \OperatorTok{\};}
\NormalTok{        Solution bad  }\OperatorTok{=}\NormalTok{ in }\OperatorTok{{-}\textgreater{}} \OperatorTok{\{} \ControlFlowTok{throw} \KeywordTok{new} \BuiltInTok{RuntimeException}\OperatorTok{(}\StringTok{"boom"}\OperatorTok{);} \OperatorTok{\};}   \CommentTok{// Lambda 实现接口}
        \BuiltInTok{System}\OperatorTok{.}\FunctionTok{out}\OperatorTok{.}\FunctionTok{println}\OperatorTok{(}\KeywordTok{new} \FunctionTok{SimpleJudge}\OperatorTok{().}\FunctionTok{judge}\OperatorTok{(}\NormalTok{p}\OperatorTok{,}\NormalTok{ good}\OperatorTok{));}   \CommentTok{// AC 2/2}
        \BuiltInTok{System}\OperatorTok{.}\FunctionTok{out}\OperatorTok{.}\FunctionTok{println}\OperatorTok{(}\KeywordTok{new} \FunctionTok{SimpleJudge}\OperatorTok{().}\FunctionTok{judge}\OperatorTok{(}\NormalTok{p}\OperatorTok{,}\NormalTok{ bad}\OperatorTok{));}    \CommentTok{// RE 0/2}
    \OperatorTok{\}}
\OperatorTok{\}}
\end{Highlighting}
\end{Shaded}

编译运行:

\begin{Shaded}
\begin{Highlighting}[]
\ExtensionTok{javac} \AttributeTok{{-}d}\NormalTok{ build src/oj/core/}\PreprocessorTok{*}\NormalTok{.java src/oj/judge/Solution.java src/oj/judge/SimpleJudge.java src/oj/Demo3.java}
\ExtensionTok{java} \AttributeTok{{-}cp}\NormalTok{ build oj.Demo3}
\CommentTok{\# AC 2/2}
\CommentTok{\# RE 0/2}
\end{Highlighting}
\end{Shaded}

\hypertarget{ux7b2cux4e8cux6b65-ux62bdux8c61ux5316ux63a5ux53e3-ux6a21ux677f-ux591aux6001-ux81eaux5b9aux4e49ux5f02ux5e38}{%
\subsubsection{\texorpdfstring{第二步 \GC --- 抽象化:接口 + 模板 + 多态
+
自定义异常}{第二步 --- 抽象化:接口 + 模板 + 多态 + 自定义异常}}\label{ux7b2cux4e8cux6b65-ux62bdux8c61ux5316ux63a5ux53e3-ux6a21ux677f-ux591aux6001-ux81eaux5b9aux4e49ux5f02ux5e38}}

把 \texttt{SimpleJudge} 重构成 \texttt{Judge} 接口 +
\texttt{AbstractJudge}(模板方法,\texttt{final\ judge} + 抽象
\texttt{compare})+ \texttt{StandardJudge}/\texttt{SpecialJudge} 子类,加
\texttt{oj.exception.JudgeException}(代码见上文【新产物架构】)。多态调用
+ 开闭原则;断言要加 \texttt{-ea}:

\begin{Shaded}
\begin{Highlighting}[]
\ExtensionTok{javac} \AttributeTok{{-}d}\NormalTok{ build }\VariableTok{$(}\FunctionTok{find}\NormalTok{ src }\AttributeTok{{-}name} \StringTok{\textquotesingle{}*.java\textquotesingle{}}\VariableTok{)}
\ExtensionTok{java} \AttributeTok{{-}ea} \AttributeTok{{-}cp}\NormalTok{ build oj.Main}
\end{Highlighting}
\end{Shaded}

\hypertarget{ux9a8cux6536ux6807ux51c6-3}{%
\subsection{验收标准}\label{ux9a8cux6536ux6807ux51c6-3}}

\begin{itemize}
\tightlist
\item[$\square$]
  \texttt{Judge}/\texttt{Solution}
  是\textbf{接口};\texttt{AbstractJudge} 是\textbf{抽象类}且实现了
  \texttt{Judge};\texttt{StandardJudge} \textbf{继承}
  \texttt{AbstractJudge} 并 \texttt{@Override\ compare}。
\item[$\square$]
  用\textbf{接口类型变量}接收子类对象,跑出多态(\texttt{judge.judge(...)}
  调到子类逻辑)。
\item[$\square$]
  会崩溃的解法被判 \textbf{RE} 而不是让程序挂掉。
\item[$\square$]
  至少演示了''匿名类 -\textgreater{} Lambda''两种等价写法之一。
\item[$\square$]
  定义了自定义异常
  \texttt{JudgeException}/\texttt{TimeLimitException}(\texttt{extends}),并能说清它们的继承关系。
\item[$\square$]
  \textbf{能讲清开闭原则}:新增一种判法时,你只新增了类、没动
  \texttt{AbstractJudge} 和 \texttt{Main}。
\end{itemize}

\hypertarget{ux672cux7ae0ux4ea7ux7269---ux4e0bux4e00ux7ae0-2}{%
\subsection{本章产物 -\textgreater{}
下一章}\label{ux672cux7ae0ux4ea7ux7269---ux4e0bux4e00ux7ae0-2}}

判题器框架成型,但只能''精确/浮点''地比。第8章(M4)用
\textbf{String/正则/反射} 实现
\texttt{SpecialJudge}、做格式校验(PE),并用\textbf{反射}按类名动态加载判题器,让判法彻底''可插拔''。

\hypertarget{m4-ux7b2cux4e8cux4ee3ux5224ux9898ux673ach8-ux5b57ux7b26ux4e32ux53cdux5c04-ch10-ux5355ux6587ux4ef6}{%
\section{05 · M4 第二代判题机（Ch8 字符串/反射 + Ch10
单文件）}\label{m4-ux7b2cux4e8cux4ee3ux5224ux9898ux673ach8-ux5b57ux7b26ux4e32ux53cdux5c04-ch10-ux5355ux6587ux4ef6}}

\begin{quote}
\textbf{考试相关度 ★★(Ch8 String)· 反射/单文件属工程} ·
真题考点参考:String
\texttt{==}/equals、StringBuffer、Random/Math。\textbf{{[}第一步{]}}
\KS 用 \texttt{trim/equals/split}+\texttt{StringBuilder} 做 OJ
的输出比对与判题报告;\textbf{{[}第二步{]}} \GC \texttt{config.txt} +
\texttt{Class.forName} 反射工厂、单文件读写。
\end{quote}

\begin{quote}
第一代判题机能跑、但题目是硬编码的，每新增一道题都要改 Java
源码重新编译。第二代的目标是：\textbf{读一份配置文件，自动决定用哪个判题器、跑哪组测试点}------``配置驱动，反射解耦''。
\end{quote}

\begin{center}\rule{0.5\linewidth}{0.5pt}\end{center}

\hypertarget{ux6838ux5fc3ux75dbux70b9}{%
\subsection{【核心痛点】}\label{ux6838ux5fc3ux75dbux70b9}}

M1--M3 的代码大致如下：

\begin{Shaded}
\begin{Highlighting}[]
\CommentTok{// 硬编码时代 —— 全是 magic number 和写死的类名}
\NormalTok{Problem p }\OperatorTok{=} \KeywordTok{new} \FunctionTok{Problem}\OperatorTok{(}\DecValTok{1}\OperatorTok{,} \StringTok{"A+B"}\OperatorTok{,} \KeywordTok{new}\NormalTok{ TestCase}\OperatorTok{[]\{}
    \KeywordTok{new} \FunctionTok{TestCase}\OperatorTok{(}\StringTok{"1 2"}\OperatorTok{,} \StringTok{"3"}\OperatorTok{),}
    \KeywordTok{new} \FunctionTok{TestCase}\OperatorTok{(}\StringTok{"10 20"}\OperatorTok{,} \StringTok{"30"}\OperatorTok{)}
\OperatorTok{\});}
\NormalTok{Judge judge }\OperatorTok{=} \KeywordTok{new} \FunctionTok{StandardJudge}\OperatorTok{();}   \CommentTok{// 换题换判题器？改代码、重新编译}
\NormalTok{JudgeResult r }\OperatorTok{=}\NormalTok{ judge}\OperatorTok{.}\FunctionTok{judge}\OperatorTok{(}\NormalTok{p}\OperatorTok{,} \KeywordTok{new} \FunctionTok{MySolution}\OperatorTok{());}
\end{Highlighting}
\end{Shaded}

三个直接缺陷：

\begin{enumerate}
\def\labelenumi{\arabic{enumi}.}
\tightlist
\item
  \textbf{题目元数据硬编码}：标题、时间限制、判题器类名全散在
  \texttt{main} 里，加一道题就改一处代码。
\item
  \textbf{判题器写死}：\texttt{new\ StandardJudge()} 或
  \texttt{new\ SpecialJudge()} 在源码里固定，无法按题目切换。
\item
  \textbf{测试点无处管理}：\texttt{TestCase}
  对象在代码里手打，完全无法复用，也无法离线维护。
\end{enumerate}

M4 的答案：把元数据写进 \texttt{config.txt}，把测试点写进
\texttt{input.txt}/\texttt{output.txt}，用 \texttt{Class.forName}
反射造判题器------代码不动，换题只换文件。

\begin{center}\rule{0.5\linewidth}{0.5pt}\end{center}

\hypertarget{ux5f15ux5165ux8bfeux672cux77e5ux8bc6ux70b9}{%
\subsection{【引入课本知识点】}\label{ux5f15ux5165ux8bfeux672cux77e5ux8bc6ux70b9}}

\hypertarget{ch8-ux5b57ux7b26ux4e32ux5904ux7406}{%
\subsubsection{Ch8 ·
字符串处理}\label{ch8-ux5b57ux7b26ux4e32ux5904ux7406}}

\textbf{\texttt{String.split} 切分 key=value}

\texttt{config.txt} 每行形如 \texttt{title=A+B\ Problem}，读出来之后：

\begin{Shaded}
\begin{Highlighting}[]
\BuiltInTok{String}\OperatorTok{[]}\NormalTok{ parts }\OperatorTok{=}\NormalTok{ line}\OperatorTok{.}\FunctionTok{split}\OperatorTok{(}\StringTok{"="}\OperatorTok{,} \DecValTok{2}\OperatorTok{);}   \CommentTok{// 限制 2 份，防止 value 里含 "="}
\BuiltInTok{String}\NormalTok{ key   }\OperatorTok{=}\NormalTok{ parts}\OperatorTok{[}\DecValTok{0}\OperatorTok{].}\FunctionTok{trim}\OperatorTok{();}
\BuiltInTok{String}\NormalTok{ value }\OperatorTok{=}\NormalTok{ parts}\OperatorTok{[}\DecValTok{1}\OperatorTok{].}\FunctionTok{trim}\OperatorTok{();}
\end{Highlighting}
\end{Shaded}

\texttt{split("=",\ 2)} 是本节最小可用的正则应用：\texttt{=}
是字面量分隔符，\texttt{2} 是 limit 参数------只切第一个等号，value
部分原封不动保留。

\textbf{\texttt{Math.abs} 浮点容差}（\texttt{SpecialJudge} 核心逻辑）

精确比对对浮点题不公平，标准做法：

\begin{Shaded}
\begin{Highlighting}[]
\DataTypeTok{double}\NormalTok{ expected }\OperatorTok{=} \BuiltInTok{Double}\OperatorTok{.}\FunctionTok{parseDouble}\OperatorTok{(}\NormalTok{exp}\OperatorTok{.}\FunctionTok{trim}\OperatorTok{());}
\DataTypeTok{double}\NormalTok{ actual   }\OperatorTok{=} \BuiltInTok{Double}\OperatorTok{.}\FunctionTok{parseDouble}\OperatorTok{(}\NormalTok{out}\OperatorTok{.}\FunctionTok{trim}\OperatorTok{());}
\DataTypeTok{boolean}\NormalTok{ ok }\OperatorTok{=} \BuiltInTok{Math}\OperatorTok{.}\FunctionTok{abs}\OperatorTok{(}\NormalTok{expected }\OperatorTok{{-}}\NormalTok{ actual}\OperatorTok{)} \OperatorTok{\textless{}} \FloatTok{1e{-}6}\OperatorTok{;}
\end{Highlighting}
\end{Shaded}

\texttt{Math.abs} 返回绝对值，\texttt{1e-6} 是科学计数法字面量------两个
Ch8 知识点同时落地。

\hypertarget{ch8-ux53cdux5c04ux57faux7840}{%
\subsubsection{Ch8 · 反射基础}\label{ch8-ux53cdux5c04ux57faux7840}}

\texttt{Class.forName(String\ name)}
在运行期按全限定名查找并加载类；\texttt{.getDeclaredConstructor().newInstance()}
无参构造出对象------这是 Java 反射的最小实践：

\begin{Shaded}
\begin{Highlighting}[]
\BuiltInTok{Class}\OperatorTok{\textless{}?\textgreater{}}\NormalTok{ clazz }\OperatorTok{=} \BuiltInTok{Class}\OperatorTok{.}\FunctionTok{forName}\OperatorTok{(}\StringTok{"oj.judge.StandardJudge"}\OperatorTok{);}
\NormalTok{Judge judge }\OperatorTok{=} \OperatorTok{(}\NormalTok{Judge}\OperatorTok{)}\NormalTok{ clazz}\OperatorTok{.}\FunctionTok{getDeclaredConstructor}\OperatorTok{().}\FunctionTok{newInstance}\OperatorTok{();}
\end{Highlighting}
\end{Shaded}

不需要提前
\texttt{import}，不需要写死类型------判题器类名从配置文件读进来就能用。

\hypertarget{ch10-ux6587ux4ef6ux8bfbux5199ux57faux7840}{%
\subsubsection{Ch10 ·
文件读写基础}\label{ch10-ux6587ux4ef6ux8bfbux5199ux57faux7840}}

\begin{verbatim}
File dir = new File("problems/1/");
File cfg = new File(dir, "config.txt");   // 拼路径，平台无关
BufferedReader br = new BufferedReader(new FileReader(cfg));
String line;
while ((line = br.readLine()) != null) {
    // 逐行处理
}
br.close();
\end{verbatim}

\texttt{File} 是路径抽象，\texttt{BufferedReader} 带缓冲逐行读------Ch10
两个核心 API 的教学入口。\texttt{readLine()} 返回 \texttt{null} 表示
EOF，这是初学者最常犯错的边界条件，此处正好强调。

\begin{center}\rule{0.5\linewidth}{0.5pt}\end{center}

\hypertarget{ux4e09ux4ee3ux6f14ux8fdbux5b9aux4f4d}{%
\subsection{【三代演进定位】}\label{ux4e09ux4ee3ux6f14ux8fdbux5b9aux4f4d}}

\begin{longtable}[]{@{}
  >{\raggedright\arraybackslash}p{(\columnwidth - 8\tabcolsep) * \real{0.1200}}
  >{\raggedright\arraybackslash}p{(\columnwidth - 8\tabcolsep) * \real{0.1200}}
  >{\raggedright\arraybackslash}p{(\columnwidth - 8\tabcolsep) * \real{0.2000}}
  >{\raggedright\arraybackslash}p{(\columnwidth - 8\tabcolsep) * \real{0.3200}}
  >{\raggedright\arraybackslash}p{(\columnwidth - 8\tabcolsep) * \real{0.2400}}@{}}
\toprule\noalign{}
\begin{minipage}[b]{\linewidth}\raggedright
代次
\end{minipage} & \begin{minipage}[b]{\linewidth}\raggedright
阶段
\end{minipage} & \begin{minipage}[b]{\linewidth}\raggedright
题目来源
\end{minipage} & \begin{minipage}[b]{\linewidth}\raggedright
判题器决定方式
\end{minipage} & \begin{minipage}[b]{\linewidth}\raggedright
测试点来源
\end{minipage} \\
\midrule\noalign{}
\endhead
\bottomrule\noalign{}
\endlastfoot
第一代（M1--M3） & 已完成 & 硬编码 \texttt{new\ Problem(...)} &
\texttt{new\ StandardJudge()} 写死 & \texttt{new\ TestCase(...)} 手打 \\
\textbf{第二代（M4，本章）} & \textbf{本章} &
\textbf{\texttt{config.txt} 读元数据} & \textbf{\texttt{Class.forName}
反射造判题器} & \textbf{\texttt{input.txt}/\texttt{output.txt}
文件读取} \\
第三代（M5a） & 下一章 & 多组 \texttt{N.in}/\texttt{N.out} 文件 & 调外部
C++ 判题机二进制 & 文件系统扫目录 \\
\end{longtable}

第二代的核心贡献是\textbf{``配置驱动''闭环}：

\begin{verbatim}
problems/1/
+-- config.txt        <- title / judgeClass / timeLimitMs
+-- input.txt         <- 单组输入
+-- output.txt        <- 单组期望输出
\end{verbatim}

\texttt{SingleFileProblemLoader} 读这三个文件，组装出 \texttt{Problem}
对象；\texttt{JudgeFactory} 用 \texttt{config.txt} 里的
\texttt{judgeClass} 字段反射造出对应
\texttt{Judge}------从此换题不改代码。

\begin{quote}
第三代会把''单文件单组''扩展成''多文件多组''，并把 JVM
内判题替换为外部二进制------但那是 M5a
的事，本章只做第二代的''最小可用闭环''。
\end{quote}

\begin{center}\rule{0.5\linewidth}{0.5pt}\end{center}

\hypertarget{ux65b0ux4ea7ux7269ux67b6ux6784}{%
\subsection{【新产物架构】}\label{ux65b0ux4ea7ux7269ux67b6ux6784}}

\hypertarget{ux6d89ux53caux7c7bux63a5ux53e3ux603bux89c8}{%
\subsubsection{涉及类/接口总览}\label{ux6d89ux53caux7c7bux63a5ux53e3ux603bux89c8}}

\begin{longtable}[]{@{}
  >{\raggedright\arraybackslash}p{(\columnwidth - 4\tabcolsep) * \real{0.5238}}
  >{\raggedright\arraybackslash}p{(\columnwidth - 4\tabcolsep) * \real{0.1905}}
  >{\raggedright\arraybackslash}p{(\columnwidth - 4\tabcolsep) * \real{0.2857}}@{}}
\toprule\noalign{}
\begin{minipage}[b]{\linewidth}\raggedright
类 / 接口
\end{minipage} & \begin{minipage}[b]{\linewidth}\raggedright
包
\end{minipage} & \begin{minipage}[b]{\linewidth}\raggedright
职责
\end{minipage} \\
\midrule\noalign{}
\endhead
\bottomrule\noalign{}
\endlastfoot
\texttt{ProblemMeta} & \texttt{oj.core} &
不可变值对象，持有题目元数据 \\
\texttt{ConfigFile} & \texttt{oj.io} & 静态工具类，解析
\texttt{config.txt} -\textgreater{} \texttt{ProblemMeta} \\
\texttt{SingleFileProblemLoader} & \texttt{oj.io} &
静态工具类，聚合三文件 -\textgreater{} \texttt{Problem} \\
\texttt{JudgeFactory} & \texttt{oj.judge} &
静态工厂，\texttt{Class.forName} 反射造 \texttt{Judge} \\
\texttt{Problem}（更新） & \texttt{oj.core} & 新增 \texttt{meta}
字段，\texttt{judgeClass} 移入 \texttt{ProblemMeta} \\
\end{longtable}

以下沿用（不改动）：\texttt{TestCase}、\texttt{JudgeResult}、\texttt{Status}、\texttt{Judge}、\texttt{AbstractJudge}、\texttt{StandardJudge}、\texttt{SpecialJudge}、\texttt{Solution}。

\begin{center}\rule{0.5\linewidth}{0.5pt}\end{center}

\hypertarget{oj.core.problemmeta}{%
\subsubsection{\texorpdfstring{\texttt{oj.core.ProblemMeta}}{oj.core.ProblemMeta}}\label{oj.core.problemmeta}}

元数据值对象。字段全 \texttt{final}，只有 getter，不可变。

\begin{Shaded}
\begin{Highlighting}[]
\KeywordTok{package}\ImportTok{ oj}\OperatorTok{.}\ImportTok{core}\OperatorTok{;}

\KeywordTok{public} \DataTypeTok{final} \KeywordTok{class}\NormalTok{ ProblemMeta }\OperatorTok{\{}
    \KeywordTok{private} \DataTypeTok{final} \BuiltInTok{String}\NormalTok{ title}\OperatorTok{;}
    \KeywordTok{private} \DataTypeTok{final} \BuiltInTok{String}\NormalTok{ judgeClass}\OperatorTok{;}   \CommentTok{// 判题器全限定名，如 "oj.judge.StandardJudge"}
    \KeywordTok{private} \DataTypeTok{final} \DataTypeTok{long}\NormalTok{ timeLimitMs}\OperatorTok{;}

    \KeywordTok{public} \FunctionTok{ProblemMeta}\OperatorTok{(}\BuiltInTok{String}\NormalTok{ title}\OperatorTok{,} \BuiltInTok{String}\NormalTok{ judgeClass}\OperatorTok{,} \DataTypeTok{long}\NormalTok{ timeLimitMs}\OperatorTok{)} \OperatorTok{\{}
        \KeywordTok{this}\OperatorTok{.}\FunctionTok{title} \OperatorTok{=}\NormalTok{ title}\OperatorTok{;}
        \KeywordTok{this}\OperatorTok{.}\FunctionTok{judgeClass} \OperatorTok{=}\NormalTok{ judgeClass}\OperatorTok{;}
        \KeywordTok{this}\OperatorTok{.}\FunctionTok{timeLimitMs} \OperatorTok{=}\NormalTok{ timeLimitMs}\OperatorTok{;}
    \OperatorTok{\}}

    \KeywordTok{public} \BuiltInTok{String} \FunctionTok{getTitle}\OperatorTok{()} \OperatorTok{\{} \ControlFlowTok{return}\NormalTok{ title}\OperatorTok{;} \OperatorTok{\}}
    \KeywordTok{public} \BuiltInTok{String} \FunctionTok{getJudgeClass}\OperatorTok{()} \OperatorTok{\{} \ControlFlowTok{return}\NormalTok{ judgeClass}\OperatorTok{;} \OperatorTok{\}}
    \KeywordTok{public} \DataTypeTok{long} \FunctionTok{getTimeLimitMs}\OperatorTok{()} \OperatorTok{\{} \ControlFlowTok{return}\NormalTok{ timeLimitMs}\OperatorTok{;} \OperatorTok{\}}
\OperatorTok{\}}
\end{Highlighting}
\end{Shaded}

\begin{center}\rule{0.5\linewidth}{0.5pt}\end{center}

\hypertarget{oj.core.problemm4-ux7248ux672c}{%
\subsubsection{\texorpdfstring{\texttt{oj.core.Problem}（M4
版本）}{oj.core.Problem（M4 版本）}}\label{oj.core.problemm4-ux7248ux672c}}

相比 M1--M3 版本，新增 \texttt{meta}
字段，\texttt{getTitle()}/\texttt{getJudgeClass()}/\texttt{getTimeLimitMs()}
委托给 \texttt{meta}。测试点仍用 \texttt{TestCase{[}{]}}（M5a 才迁移到
\texttt{List\textless{}TestCase\textgreater{}}）。

\begin{Shaded}
\begin{Highlighting}[]
\KeywordTok{package}\ImportTok{ oj}\OperatorTok{.}\ImportTok{core}\OperatorTok{;}

\KeywordTok{public} \KeywordTok{class}\NormalTok{ Problem }\OperatorTok{\{}
    \KeywordTok{private} \DataTypeTok{final} \DataTypeTok{int}\NormalTok{ id}\OperatorTok{;}
    \KeywordTok{private} \DataTypeTok{final}\NormalTok{ ProblemMeta meta}\OperatorTok{;}
    \KeywordTok{private} \DataTypeTok{final}\NormalTok{ TestCase}\OperatorTok{[]}\NormalTok{ cases}\OperatorTok{;}

    \KeywordTok{public} \FunctionTok{Problem}\OperatorTok{(}\DataTypeTok{int}\NormalTok{ id}\OperatorTok{,}\NormalTok{ ProblemMeta meta}\OperatorTok{,}\NormalTok{ TestCase}\OperatorTok{[]}\NormalTok{ cases}\OperatorTok{)} \OperatorTok{\{}
        \KeywordTok{this}\OperatorTok{.}\FunctionTok{id} \OperatorTok{=}\NormalTok{ id}\OperatorTok{;}
        \KeywordTok{this}\OperatorTok{.}\FunctionTok{meta} \OperatorTok{=}\NormalTok{ meta}\OperatorTok{;}
        \KeywordTok{this}\OperatorTok{.}\FunctionTok{cases} \OperatorTok{=}\NormalTok{ cases}\OperatorTok{;}
    \OperatorTok{\}}

    \KeywordTok{public} \DataTypeTok{int} \FunctionTok{getId}\OperatorTok{()} \OperatorTok{\{} \ControlFlowTok{return}\NormalTok{ id}\OperatorTok{;} \OperatorTok{\}}
    \KeywordTok{public}\NormalTok{ ProblemMeta }\FunctionTok{getMeta}\OperatorTok{()} \OperatorTok{\{} \ControlFlowTok{return}\NormalTok{ meta}\OperatorTok{;} \OperatorTok{\}}

    \CommentTok{// 便捷代理，让调用方无需先拿 meta}
    \KeywordTok{public} \BuiltInTok{String} \FunctionTok{getTitle}\OperatorTok{()} \OperatorTok{\{} \ControlFlowTok{return}\NormalTok{ meta}\OperatorTok{.}\FunctionTok{getTitle}\OperatorTok{();} \OperatorTok{\}}
    \KeywordTok{public} \BuiltInTok{String} \FunctionTok{getJudgeClass}\OperatorTok{()} \OperatorTok{\{} \ControlFlowTok{return}\NormalTok{ meta}\OperatorTok{.}\FunctionTok{getJudgeClass}\OperatorTok{();} \OperatorTok{\}}
    \KeywordTok{public} \DataTypeTok{long} \FunctionTok{getTimeLimitMs}\OperatorTok{()} \OperatorTok{\{} \ControlFlowTok{return}\NormalTok{ meta}\OperatorTok{.}\FunctionTok{getTimeLimitMs}\OperatorTok{();} \OperatorTok{\}}

    \KeywordTok{public}\NormalTok{ TestCase}\OperatorTok{[]} \FunctionTok{getCases}\OperatorTok{()} \OperatorTok{\{} \ControlFlowTok{return}\NormalTok{ cases}\OperatorTok{;} \OperatorTok{\}}
\OperatorTok{\}}
\end{Highlighting}
\end{Shaded}

\begin{center}\rule{0.5\linewidth}{0.5pt}\end{center}

\hypertarget{oj.io.configfile}{%
\subsubsection{\texorpdfstring{\texttt{oj.io.ConfigFile}}{oj.io.ConfigFile}}\label{oj.io.configfile}}

解析 \texttt{config.txt}，返回 \texttt{ProblemMeta}。

\texttt{config.txt} 格式约定（顺序不限，\texttt{\#} 开头为注释行）：

\begin{verbatim}
title=A+B Problem
judge=oj.judge.StandardJudge
timeLimitMs=1000
\end{verbatim}

\begin{Shaded}
\begin{Highlighting}[]
\KeywordTok{package}\ImportTok{ oj}\OperatorTok{.}\ImportTok{io}\OperatorTok{;}

\KeywordTok{import} \ImportTok{oj}\OperatorTok{.}\ImportTok{core}\OperatorTok{.}\ImportTok{ProblemMeta}\OperatorTok{;}
\KeywordTok{import} \ImportTok{java}\OperatorTok{.}\ImportTok{io}\OperatorTok{.*;}

\KeywordTok{public} \KeywordTok{class}\NormalTok{ ConfigFile }\OperatorTok{\{}

    \CommentTok{/**}
     \CommentTok{*} \CommentTok{解析}\NormalTok{ dir }\CommentTok{下的}\NormalTok{ config}\CommentTok{.}\NormalTok{txt}\CommentTok{，返回}\NormalTok{ ProblemMeta}\CommentTok{。}
     \CommentTok{*} \CommentTok{未识别的}\NormalTok{ key }\CommentTok{直接跳过，不抛异常。}
     \CommentTok{*/}
    \KeywordTok{public} \DataTypeTok{static}\NormalTok{ ProblemMeta }\FunctionTok{read}\OperatorTok{(}\BuiltInTok{File}\NormalTok{ dir}\OperatorTok{)} \KeywordTok{throws} \BuiltInTok{IOException} \OperatorTok{\{}
        \BuiltInTok{File}\NormalTok{ cfg }\OperatorTok{=} \KeywordTok{new} \BuiltInTok{File}\OperatorTok{(}\NormalTok{dir}\OperatorTok{,} \StringTok{"config.txt"}\OperatorTok{);}
        \BuiltInTok{String}\NormalTok{ title }\OperatorTok{=} \StringTok{""}\OperatorTok{;}
        \BuiltInTok{String}\NormalTok{ judgeClass }\OperatorTok{=} \StringTok{"oj.judge.StandardJudge"}\OperatorTok{;}   \CommentTok{// 默认值}
        \DataTypeTok{long}\NormalTok{ timeLimitMs }\OperatorTok{=} \DecValTok{1000L}\OperatorTok{;}

        \BuiltInTok{BufferedReader}\NormalTok{ br }\OperatorTok{=} \KeywordTok{new} \BuiltInTok{BufferedReader}\OperatorTok{(}\KeywordTok{new} \BuiltInTok{FileReader}\OperatorTok{(}\NormalTok{cfg}\OperatorTok{));}
        \BuiltInTok{String}\NormalTok{ line}\OperatorTok{;}
        \ControlFlowTok{while} \OperatorTok{((}\NormalTok{line }\OperatorTok{=}\NormalTok{ br}\OperatorTok{.}\FunctionTok{readLine}\OperatorTok{())} \OperatorTok{!=} \KeywordTok{null}\OperatorTok{)} \OperatorTok{\{}
\NormalTok{            line }\OperatorTok{=}\NormalTok{ line}\OperatorTok{.}\FunctionTok{trim}\OperatorTok{();}
            \ControlFlowTok{if} \OperatorTok{(}\NormalTok{line}\OperatorTok{.}\FunctionTok{isEmpty}\OperatorTok{()} \OperatorTok{||}\NormalTok{ line}\OperatorTok{.}\FunctionTok{startsWith}\OperatorTok{(}\StringTok{"\#"}\OperatorTok{))} \OperatorTok{\{}
                \ControlFlowTok{continue}\OperatorTok{;}   \CommentTok{// 跳过空行和注释}
            \OperatorTok{\}}
            \BuiltInTok{String}\OperatorTok{[]}\NormalTok{ parts }\OperatorTok{=}\NormalTok{ line}\OperatorTok{.}\FunctionTok{split}\OperatorTok{(}\StringTok{"="}\OperatorTok{,} \DecValTok{2}\OperatorTok{);}
            \ControlFlowTok{if} \OperatorTok{(}\NormalTok{parts}\OperatorTok{.}\FunctionTok{length} \OperatorTok{\textless{}} \DecValTok{2}\OperatorTok{)} \OperatorTok{\{}
                \ControlFlowTok{continue}\OperatorTok{;}   \CommentTok{// 格式不合法，跳过}
            \OperatorTok{\}}
            \BuiltInTok{String}\NormalTok{ key   }\OperatorTok{=}\NormalTok{ parts}\OperatorTok{[}\DecValTok{0}\OperatorTok{].}\FunctionTok{trim}\OperatorTok{();}
            \BuiltInTok{String}\NormalTok{ value }\OperatorTok{=}\NormalTok{ parts}\OperatorTok{[}\DecValTok{1}\OperatorTok{].}\FunctionTok{trim}\OperatorTok{();}
            \ControlFlowTok{switch} \OperatorTok{(}\NormalTok{key}\OperatorTok{)} \OperatorTok{\{}
                \ControlFlowTok{case} \StringTok{"title"}\OperatorTok{:}
\NormalTok{                    title }\OperatorTok{=}\NormalTok{ value}\OperatorTok{;}
                    \ControlFlowTok{break}\OperatorTok{;}
                \ControlFlowTok{case} \StringTok{"judge"}\OperatorTok{:}
\NormalTok{                    judgeClass }\OperatorTok{=}\NormalTok{ value}\OperatorTok{;}
                    \ControlFlowTok{break}\OperatorTok{;}
                \ControlFlowTok{case} \StringTok{"timeLimitMs"}\OperatorTok{:}
\NormalTok{                    timeLimitMs }\OperatorTok{=} \BuiltInTok{Long}\OperatorTok{.}\FunctionTok{parseLong}\OperatorTok{(}\NormalTok{value}\OperatorTok{);}
                    \ControlFlowTok{break}\OperatorTok{;}
                \KeywordTok{default}\OperatorTok{:}
                    \CommentTok{// 未知 key 忽略，方便未来扩展}
                    \ControlFlowTok{break}\OperatorTok{;}
            \OperatorTok{\}}
        \OperatorTok{\}}
\NormalTok{        br}\OperatorTok{.}\FunctionTok{close}\OperatorTok{();}
        \ControlFlowTok{return} \KeywordTok{new} \FunctionTok{ProblemMeta}\OperatorTok{(}\NormalTok{title}\OperatorTok{,}\NormalTok{ judgeClass}\OperatorTok{,}\NormalTok{ timeLimitMs}\OperatorTok{);}
    \OperatorTok{\}}
\OperatorTok{\}}
\end{Highlighting}
\end{Shaded}

\textbf{教学要点}：\texttt{split("=",\ 2)} 的 limit 参数防止
\texttt{value} 里含 \texttt{=} 被误切；\texttt{switch} 上的字符串（Java
7+）比 \texttt{if-else\ if} 链更易扩展。

\begin{center}\rule{0.5\linewidth}{0.5pt}\end{center}

\hypertarget{oj.io.singlefileproblemloader}{%
\subsubsection{\texorpdfstring{\texttt{oj.io.SingleFileProblemLoader}}{oj.io.SingleFileProblemLoader}}\label{oj.io.singlefileproblemloader}}

聚合三个文件，组装 \texttt{Problem}。单组测试点，\texttt{TestCase{[}{]}}
长度始终为 1。

\begin{Shaded}
\begin{Highlighting}[]
\KeywordTok{package}\ImportTok{ oj}\OperatorTok{.}\ImportTok{io}\OperatorTok{;}

\KeywordTok{import} \ImportTok{oj}\OperatorTok{.}\ImportTok{core}\OperatorTok{.*;}
\KeywordTok{import} \ImportTok{java}\OperatorTok{.}\ImportTok{io}\OperatorTok{.*;}
\KeywordTok{import} \ImportTok{java}\OperatorTok{.}\ImportTok{nio}\OperatorTok{.}\ImportTok{file}\OperatorTok{.*;}

\KeywordTok{public} \KeywordTok{class}\NormalTok{ SingleFileProblemLoader }\OperatorTok{\{}

    \CommentTok{/**}
     \CommentTok{*} \CommentTok{读}\NormalTok{ dir }\CommentTok{下的}\NormalTok{ config}\CommentTok{.}\NormalTok{txt }\CommentTok{/}\NormalTok{ input}\CommentTok{.}\NormalTok{txt }\CommentTok{/}\NormalTok{ output}\CommentTok{.}\NormalTok{txt}\CommentTok{，}
     \CommentTok{*} \CommentTok{返回含单组测试点的}\NormalTok{ Problem}\CommentTok{。}
     \CommentTok{*/}
    \KeywordTok{public} \DataTypeTok{static}\NormalTok{ Problem }\FunctionTok{load}\OperatorTok{(}\DataTypeTok{int}\NormalTok{ id}\OperatorTok{,} \BuiltInTok{File}\NormalTok{ dir}\OperatorTok{)} \KeywordTok{throws} \BuiltInTok{IOException} \OperatorTok{\{}
\NormalTok{        ProblemMeta meta }\OperatorTok{=}\NormalTok{ ConfigFile}\OperatorTok{.}\FunctionTok{read}\OperatorTok{(}\NormalTok{dir}\OperatorTok{);}

        \BuiltInTok{String}\NormalTok{ input    }\OperatorTok{=} \FunctionTok{readAll}\OperatorTok{(}\KeywordTok{new} \BuiltInTok{File}\OperatorTok{(}\NormalTok{dir}\OperatorTok{,} \StringTok{"input.txt"}\OperatorTok{));}
        \BuiltInTok{String}\NormalTok{ expected }\OperatorTok{=} \FunctionTok{readAll}\OperatorTok{(}\KeywordTok{new} \BuiltInTok{File}\OperatorTok{(}\NormalTok{dir}\OperatorTok{,} \StringTok{"output.txt"}\OperatorTok{));}

\NormalTok{        TestCase}\OperatorTok{[]}\NormalTok{ cases }\OperatorTok{=} \OperatorTok{\{} \KeywordTok{new} \FunctionTok{TestCase}\OperatorTok{(}\NormalTok{input}\OperatorTok{,}\NormalTok{ expected}\OperatorTok{)} \OperatorTok{\};}
        \ControlFlowTok{return} \KeywordTok{new} \FunctionTok{Problem}\OperatorTok{(}\NormalTok{id}\OperatorTok{,}\NormalTok{ meta}\OperatorTok{,}\NormalTok{ cases}\OperatorTok{);}
    \OperatorTok{\}}

    \CommentTok{/**} \CommentTok{读整个文件为字符串，行尾统一用} \CommentTok{\textbackslash{}}\NormalTok{n}\CommentTok{。} \CommentTok{*/}
    \KeywordTok{private} \DataTypeTok{static} \BuiltInTok{String} \FunctionTok{readAll}\OperatorTok{(}\BuiltInTok{File}\NormalTok{ f}\OperatorTok{)} \KeywordTok{throws} \BuiltInTok{IOException} \OperatorTok{\{}
        \BuiltInTok{StringBuilder}\NormalTok{ sb }\OperatorTok{=} \KeywordTok{new} \BuiltInTok{StringBuilder}\OperatorTok{();}
        \BuiltInTok{BufferedReader}\NormalTok{ br }\OperatorTok{=} \KeywordTok{new} \BuiltInTok{BufferedReader}\OperatorTok{(}\KeywordTok{new} \BuiltInTok{FileReader}\OperatorTok{(}\NormalTok{f}\OperatorTok{));}
        \BuiltInTok{String}\NormalTok{ line}\OperatorTok{;}
        \ControlFlowTok{while} \OperatorTok{((}\NormalTok{line }\OperatorTok{=}\NormalTok{ br}\OperatorTok{.}\FunctionTok{readLine}\OperatorTok{())} \OperatorTok{!=} \KeywordTok{null}\OperatorTok{)} \OperatorTok{\{}
\NormalTok{            sb}\OperatorTok{.}\FunctionTok{append}\OperatorTok{(}\NormalTok{line}\OperatorTok{).}\FunctionTok{append}\OperatorTok{(}\CharTok{\textquotesingle{}\textbackslash{}n\textquotesingle{}}\OperatorTok{);}
        \OperatorTok{\}}
\NormalTok{        br}\OperatorTok{.}\FunctionTok{close}\OperatorTok{();}
        \ControlFlowTok{return}\NormalTok{ sb}\OperatorTok{.}\FunctionTok{toString}\OperatorTok{().}\FunctionTok{trim}\OperatorTok{();}   \CommentTok{// 去掉末尾多余空行}
    \OperatorTok{\}}
\OperatorTok{\}}
\end{Highlighting}
\end{Shaded}

\textbf{教学要点}：\texttt{private\ static\ readAll}
体现''封装小工具''原则------重复的文件读逻辑只写一次，两处调用（input /
output）共享同一实现。\texttt{trim()}
消除末尾换行符，避免比对时因空白差异误判为 WA。

\begin{center}\rule{0.5\linewidth}{0.5pt}\end{center}

\hypertarget{oj.judge.judgefactory}{%
\subsubsection{\texorpdfstring{\texttt{oj.judge.JudgeFactory}}{oj.judge.JudgeFactory}}\label{oj.judge.judgefactory}}

反射工厂。核心只有一行：\texttt{Class.forName(className).getDeclaredConstructor().newInstance()}。

\begin{Shaded}
\begin{Highlighting}[]
\KeywordTok{package}\ImportTok{ oj}\OperatorTok{.}\ImportTok{judge}\OperatorTok{;}

\KeywordTok{public} \KeywordTok{class}\NormalTok{ JudgeFactory }\OperatorTok{\{}

    \CommentTok{/**}
     \CommentTok{*} \CommentTok{按全限定类名反射造}\NormalTok{ Judge }\CommentTok{实例。}
     \CommentTok{*} \CommentTok{要求目标类：①} \CommentTok{实现}\NormalTok{ Judge }\CommentTok{接口；②} \CommentTok{有无参构造器。}
     \CommentTok{*}
\CommentTok{     * @}\NormalTok{param className }\CommentTok{如} \CommentTok{"}\NormalTok{oj}\CommentTok{.}\NormalTok{judge}\CommentTok{.}\NormalTok{StandardJudge}\CommentTok{"} \CommentTok{或} \CommentTok{"}\NormalTok{oj}\CommentTok{.}\NormalTok{judge}\CommentTok{.}\NormalTok{SpecialJudge}\CommentTok{"}
     \CommentTok{*} \CommentTok{@}\NormalTok{return }\CommentTok{对应的}\NormalTok{ Judge }\CommentTok{实例}
     \CommentTok{*} \CommentTok{@}\NormalTok{throws ReflectiveOperationException }\CommentTok{类不存在、无无参构造、类型不匹配时抛出}
     \CommentTok{*/}
    \KeywordTok{public} \DataTypeTok{static}\NormalTok{ Judge }\FunctionTok{create}\OperatorTok{(}\BuiltInTok{String}\NormalTok{ className}\OperatorTok{)} \KeywordTok{throws}\NormalTok{ ReflectiveOperationException }\OperatorTok{\{}
        \BuiltInTok{Class}\OperatorTok{\textless{}?\textgreater{}}\NormalTok{ clazz }\OperatorTok{=} \BuiltInTok{Class}\OperatorTok{.}\FunctionTok{forName}\OperatorTok{(}\NormalTok{className}\OperatorTok{);}
        \ControlFlowTok{return} \OperatorTok{(}\NormalTok{Judge}\OperatorTok{)}\NormalTok{ clazz}\OperatorTok{.}\FunctionTok{getDeclaredConstructor}\OperatorTok{().}\FunctionTok{newInstance}\OperatorTok{();}
    \OperatorTok{\}}
\OperatorTok{\}}
\end{Highlighting}
\end{Shaded}

\textbf{教学要点}： - \texttt{Class\textless{}?\textgreater{}}
用通配符，因为编译期不知道具体类型。 - \texttt{getDeclaredConstructor()}
比 \texttt{newInstance()}（已废弃）更规范。 - 强制转型 \texttt{(Judge)}
若配置文件写错类名会在运行期抛
\texttt{ClassCastException}------这正是向学生演示''运行期类型检查''的绝佳时机。

\begin{center}\rule{0.5\linewidth}{0.5pt}\end{center}

\hypertarget{ux6587ux4ef6---ux53cdux5c04---ux5224ux9898ux95edux73afux793aux4f8b}{%
\subsubsection{``文件 -\textgreater{} 反射 -\textgreater{}
判题''闭环示例}\label{ux6587ux4ef6---ux53cdux5c04---ux5224ux9898ux95edux73afux793aux4f8b}}

以下是 \texttt{main} 里的完整使用示例，串起所有新产物：

\begin{Shaded}
\begin{Highlighting}[]
\KeywordTok{import} \ImportTok{oj}\OperatorTok{.}\ImportTok{core}\OperatorTok{.*;}
\KeywordTok{import} \ImportTok{oj}\OperatorTok{.}\ImportTok{io}\OperatorTok{.*;}
\KeywordTok{import} \ImportTok{oj}\OperatorTok{.}\ImportTok{judge}\OperatorTok{.*;}

\KeywordTok{public} \KeywordTok{class}\NormalTok{ M4Demo }\OperatorTok{\{}
    \KeywordTok{public} \DataTypeTok{static} \DataTypeTok{void} \FunctionTok{main}\OperatorTok{(}\BuiltInTok{String}\OperatorTok{[]}\NormalTok{ args}\OperatorTok{)} \KeywordTok{throws} \BuiltInTok{Exception} \OperatorTok{\{}
        \CommentTok{// 1. 从文件系统加载题目（读三个文件）}
        \BuiltInTok{File}\NormalTok{ problemDir }\OperatorTok{=} \KeywordTok{new} \BuiltInTok{File}\OperatorTok{(}\StringTok{"problems/1/"}\OperatorTok{);}
\NormalTok{        Problem problem }\OperatorTok{=}\NormalTok{ SingleFileProblemLoader}\OperatorTok{.}\FunctionTok{load}\OperatorTok{(}\DecValTok{1}\OperatorTok{,}\NormalTok{ problemDir}\OperatorTok{);}

        \BuiltInTok{System}\OperatorTok{.}\FunctionTok{out}\OperatorTok{.}\FunctionTok{println}\OperatorTok{(}\StringTok{"题目: "} \OperatorTok{+}\NormalTok{ problem}\OperatorTok{.}\FunctionTok{getTitle}\OperatorTok{());}
        \BuiltInTok{System}\OperatorTok{.}\FunctionTok{out}\OperatorTok{.}\FunctionTok{println}\OperatorTok{(}\StringTok{"判题器: "} \OperatorTok{+}\NormalTok{ problem}\OperatorTok{.}\FunctionTok{getJudgeClass}\OperatorTok{());}
        \BuiltInTok{System}\OperatorTok{.}\FunctionTok{out}\OperatorTok{.}\FunctionTok{println}\OperatorTok{(}\StringTok{"时间限制: "} \OperatorTok{+}\NormalTok{ problem}\OperatorTok{.}\FunctionTok{getTimeLimitMs}\OperatorTok{()} \OperatorTok{+} \StringTok{"ms"}\OperatorTok{);}

        \CommentTok{// 2. 反射造判题器（判题器类名来自 config.txt，不写死）}
\NormalTok{        Judge judge }\OperatorTok{=}\NormalTok{ JudgeFactory}\OperatorTok{.}\FunctionTok{create}\OperatorTok{(}\NormalTok{problem}\OperatorTok{.}\FunctionTok{getJudgeClass}\OperatorTok{());}

        \CommentTok{// 3. 准备一个 Solution（教学占位，仍在 JVM 内运行）}
\NormalTok{        Solution solution }\OperatorTok{=}\NormalTok{ input }\OperatorTok{{-}\textgreater{}} \OperatorTok{\{}
            \BuiltInTok{String}\OperatorTok{[]}\NormalTok{ nums }\OperatorTok{=}\NormalTok{ input}\OperatorTok{.}\FunctionTok{trim}\OperatorTok{().}\FunctionTok{split}\OperatorTok{(}\StringTok{"}\SpecialCharTok{\textbackslash{}\textbackslash{}}\StringTok{s+"}\OperatorTok{);}
            \DataTypeTok{int}\NormalTok{ a }\OperatorTok{=} \BuiltInTok{Integer}\OperatorTok{.}\FunctionTok{parseInt}\OperatorTok{(}\NormalTok{nums}\OperatorTok{[}\DecValTok{0}\OperatorTok{]);}
            \DataTypeTok{int}\NormalTok{ b }\OperatorTok{=} \BuiltInTok{Integer}\OperatorTok{.}\FunctionTok{parseInt}\OperatorTok{(}\NormalTok{nums}\OperatorTok{[}\DecValTok{1}\OperatorTok{]);}
            \ControlFlowTok{return} \BuiltInTok{String}\OperatorTok{.}\FunctionTok{valueOf}\OperatorTok{(}\NormalTok{a }\OperatorTok{+}\NormalTok{ b}\OperatorTok{);}
        \OperatorTok{\};}

        \CommentTok{// 4. 判题，打印结果}
\NormalTok{        JudgeResult result }\OperatorTok{=}\NormalTok{ judge}\OperatorTok{.}\FunctionTok{judge}\OperatorTok{(}\NormalTok{problem}\OperatorTok{,}\NormalTok{ solution}\OperatorTok{);}
        \BuiltInTok{System}\OperatorTok{.}\FunctionTok{out}\OperatorTok{.}\FunctionTok{println}\OperatorTok{(}\NormalTok{result}\OperatorTok{);}   \CommentTok{// 输出如: AC 1/1 (3ms)}
    \OperatorTok{\}}
\OperatorTok{\}}
\end{Highlighting}
\end{Shaded}

运行时目录结构：

\begin{verbatim}
problems/
+-- 1/
    +-- config.txt    ->  title=A+B Problem
    |                    judge=oj.judge.StandardJudge
    |                    timeLimitMs=1000
    +-- input.txt     ->  1 2
    +-- output.txt    ->  3
\end{verbatim}

将 \texttt{config.txt} 里的 \texttt{judge} 改为
\texttt{oj.judge.SpecialJudge}，无需改代码，重新运行即切换为浮点判题器------这就是''配置驱动''的含义。

\begin{center}\rule{0.5\linewidth}{0.5pt}\end{center}

\hypertarget{ux52a8ux624bux5b9eux73b0ux4e24ux6b65ux8d70javac-java-ux5b9eux64cd-2}{%
\subsection{动手实现:两步走(javac / java
实操)}\label{ux52a8ux624bux5b9eux73b0ux4e24ux6b65ux8d70javac-java-ux5b9eux64cd-2}}

\hypertarget{ux7b2cux4e00ux6b65-ux7528-string-ux5199-oj-ux7684ux8f93ux51faux6bd4ux5bf9-ux62a5ux544a}{%
\subsubsection{\texorpdfstring{第一步 \KS --- 用 String 写 OJ 的输出比对
+
报告}{第一步 --- 用 String 写 OJ 的输出比对 + 报告}}\label{ux7b2cux4e00ux6b65-ux7528-string-ux5199-oj-ux7684ux8f93ux51faux6bd4ux5bf9-ux62a5ux544a}}

\textbf{\texttt{src/oj/judge/OutputChecker.java}}

\begin{Shaded}
\begin{Highlighting}[]
\KeywordTok{package}\ImportTok{ oj}\OperatorTok{.}\ImportTok{judge}\OperatorTok{;}
\KeywordTok{public} \KeywordTok{class}\NormalTok{ OutputChecker }\OperatorTok{\{}
    \CommentTok{/**} \CommentTok{逐行去首尾空白、忽略末尾空行后比较(}\NormalTok{OJ }\CommentTok{最常用的"宽松精确"比对)} \CommentTok{*/}
    \KeywordTok{public} \DataTypeTok{static} \DataTypeTok{boolean} \FunctionTok{same}\OperatorTok{(}\BuiltInTok{String}\NormalTok{ expected}\OperatorTok{,} \BuiltInTok{String}\NormalTok{ actual}\OperatorTok{)} \OperatorTok{\{}
        \ControlFlowTok{return} \FunctionTok{norm}\OperatorTok{(}\NormalTok{expected}\OperatorTok{).}\FunctionTok{equals}\OperatorTok{(}\FunctionTok{norm}\OperatorTok{(}\NormalTok{actual}\OperatorTok{));}
    \OperatorTok{\}}
    \KeywordTok{private} \DataTypeTok{static} \BuiltInTok{String} \FunctionTok{norm}\OperatorTok{(}\BuiltInTok{String}\NormalTok{ s}\OperatorTok{)} \OperatorTok{\{}
        \BuiltInTok{StringBuilder}\NormalTok{ sb }\OperatorTok{=} \KeywordTok{new} \BuiltInTok{StringBuilder}\OperatorTok{();}
        \ControlFlowTok{for} \OperatorTok{(}\BuiltInTok{String}\NormalTok{ line }\OperatorTok{:}\NormalTok{ s}\OperatorTok{.}\FunctionTok{split}\OperatorTok{(}\StringTok{"}\SpecialCharTok{\textbackslash{}n}\StringTok{"}\OperatorTok{,} \OperatorTok{{-}}\DecValTok{1}\OperatorTok{))}\NormalTok{ sb}\OperatorTok{.}\FunctionTok{append}\OperatorTok{(}\NormalTok{line}\OperatorTok{.}\FunctionTok{trim}\OperatorTok{()).}\FunctionTok{append}\OperatorTok{(}\CharTok{\textquotesingle{}\textbackslash{}n\textquotesingle{}}\OperatorTok{);}
        \ControlFlowTok{return}\NormalTok{ sb}\OperatorTok{.}\FunctionTok{toString}\OperatorTok{().}\FunctionTok{trim}\OperatorTok{();}
    \OperatorTok{\}}
\OperatorTok{\}}
\end{Highlighting}
\end{Shaded}

讲解:\texttt{split} 切行、\texttt{trim()} 去空白、\texttt{StringBuilder}
拼接、\texttt{equals} 比较------全是 Ch8 考点,用在 OJ 自己的判题上。

\textbf{\texttt{src/oj/Demo4.java}}

\begin{Shaded}
\begin{Highlighting}[]
\KeywordTok{package}\ImportTok{ oj}\OperatorTok{;}
\KeywordTok{import} \ImportTok{oj}\OperatorTok{.}\ImportTok{judge}\OperatorTok{.}\ImportTok{OutputChecker}\OperatorTok{;}
\KeywordTok{public} \KeywordTok{class}\NormalTok{ Demo4 }\OperatorTok{\{}
    \KeywordTok{public} \DataTypeTok{static} \DataTypeTok{void} \FunctionTok{main}\OperatorTok{(}\BuiltInTok{String}\OperatorTok{[]}\NormalTok{ args}\OperatorTok{)} \OperatorTok{\{}
        \BuiltInTok{System}\OperatorTok{.}\FunctionTok{out}\OperatorTok{.}\FunctionTok{println}\OperatorTok{(}\NormalTok{OutputChecker}\OperatorTok{.}\FunctionTok{same}\OperatorTok{(}\StringTok{"3}\SpecialCharTok{\textbackslash{}n}\StringTok{"}\OperatorTok{,} \StringTok{"  3  "}\OperatorTok{));}        \CommentTok{// true(空白无关)}
        \BuiltInTok{System}\OperatorTok{.}\FunctionTok{out}\OperatorTok{.}\FunctionTok{println}\OperatorTok{(}\NormalTok{OutputChecker}\OperatorTok{.}\FunctionTok{same}\OperatorTok{(}\StringTok{"1 2 3"}\OperatorTok{,} \StringTok{"1 2 4"}\OperatorTok{));}      \CommentTok{// false}
        \BuiltInTok{StringBuilder}\NormalTok{ report }\OperatorTok{=} \KeywordTok{new} \BuiltInTok{StringBuilder}\OperatorTok{(}\StringTok{"判题报告:}\SpecialCharTok{\textbackslash{}n}\StringTok{"}\OperatorTok{);}
\NormalTok{        report}\OperatorTok{.}\FunctionTok{append}\OperatorTok{(}\StringTok{"case\#1: "}\OperatorTok{).}\FunctionTok{append}\OperatorTok{(}\NormalTok{OutputChecker}\OperatorTok{.}\FunctionTok{same}\OperatorTok{(}\StringTok{"3"}\OperatorTok{,} \StringTok{"3"}\OperatorTok{)} \OperatorTok{?} \StringTok{"AC"} \OperatorTok{:} \StringTok{"WA"}\OperatorTok{).}\FunctionTok{append}\OperatorTok{(}\CharTok{\textquotesingle{}\textbackslash{}n\textquotesingle{}}\OperatorTok{);}
        \BuiltInTok{System}\OperatorTok{.}\FunctionTok{out}\OperatorTok{.}\FunctionTok{print}\OperatorTok{(}\NormalTok{report}\OperatorTok{);}
    \OperatorTok{\}}
\OperatorTok{\}}
\end{Highlighting}
\end{Shaded}

编译运行:

\begin{Shaded}
\begin{Highlighting}[]
\ExtensionTok{javac} \AttributeTok{{-}d}\NormalTok{ build src/oj/judge/OutputChecker.java src/oj/Demo4.java}
\ExtensionTok{java} \AttributeTok{{-}cp}\NormalTok{ build oj.Demo4}
\CommentTok{\# true}
\CommentTok{\# false}
\CommentTok{\# 判题报告:}
\CommentTok{\# case\#1: AC}
\end{Highlighting}
\end{Shaded}

\hypertarget{ux7b2cux4e8cux6b65-ux914dux7f6eux9a71ux52a8-ux53cdux5c04ux5de5ux5382}{%
\subsubsection{\texorpdfstring{第二步 \GC --- 配置驱动 +
反射工厂}{第二步 --- 配置驱动 + 反射工厂}}\label{ux7b2cux4e8cux6b65-ux914dux7f6eux9a71ux52a8-ux53cdux5c04ux5de5ux5382}}

加 \texttt{ConfigFile.read(dir)} 解析
\texttt{config.txt}、\texttt{JudgeFactory.create(类名)} 用
\texttt{Class.forName} 反射造判题器、\texttt{SingleFileProblemLoader}
读单组用例(代码见上文【新产物架构】)。换判题器只改配置:

\begin{Shaded}
\begin{Highlighting}[]
\ExtensionTok{javac} \AttributeTok{{-}d}\NormalTok{ build }\VariableTok{$(}\FunctionTok{find}\NormalTok{ src }\AttributeTok{{-}name} \StringTok{\textquotesingle{}*.java\textquotesingle{}}\VariableTok{)}
\ExtensionTok{java} \AttributeTok{{-}cp}\NormalTok{ build oj.Main}
\end{Highlighting}
\end{Shaded}

\hypertarget{ux9a8cux6536ux6807ux51c6-4}{%
\subsection{验收标准}\label{ux9a8cux6536ux6807ux51c6-4}}

\begin{enumerate}
\def\labelenumi{\arabic{enumi}.}
\tightlist
\item
  \textbf{\texttt{ConfigFile.read}}：给定合法
  \texttt{config.txt}，能正确解析
  \texttt{title}、\texttt{judge}、\texttt{timeLimitMs}
  三个字段；遇到空行、\texttt{\#} 注释行、未知 key 不抛异常。
\item
  \textbf{\texttt{SingleFileProblemLoader.load}}：读取
  \texttt{problems/1/} 下三个文件后，返回的 \texttt{Problem} 对象
  \texttt{getTitle()} / \texttt{getJudgeClass()} /
  \texttt{getCases().length\ ==\ 1} 均符合预期。
\item
  \textbf{\texttt{JudgeFactory.create("oj.judge.StandardJudge")}}：返回的对象
  \texttt{instanceof\ StandardJudge} 为
  \texttt{true}；传入不存在的类名抛
  \texttt{ClassNotFoundException}（不是静默失败）。
\item
  \textbf{\texttt{JudgeFactory.create("oj.judge.SpecialJudge")}}：仅修改
  \texttt{config.txt}，不改任何 Java 源码，判题器自动切换。
\item
  \textbf{闭环冒烟测试}：\texttt{M4Demo.main} 在 \texttt{problems/1/}
  目录存在时输出 \texttt{AC\ 1/1\ (Xms)}；将 \texttt{output.txt}
  改为错误答案后输出 \texttt{WA\ 0/1\ (Xms)}。
\item
  \textbf{\texttt{Problem.getJudgeClass()}}：调用链
  \texttt{problem.getJudgeClass()} 实际委托给
  \texttt{problem.getMeta().getJudgeClass()}，两者返回值相同。
\end{enumerate}

\hypertarget{m5a-ux7b2cux4e09ux4ee3ux5224ux9898ux673a}{%
\section{06 · M5a
第三代判题机}\label{m5a-ux7b2cux4e09ux4ee3ux5224ux9898ux673a}}

\begin{quote}
\textbf{考试相关度 ★★(Ch15 集合)· C++ 判题机属工程} ·
真题考点参考:ArrayList/LinkedList/HashMap/Iterator。\textbf{{[}第一步{]}}
\KS 用 \texttt{ArrayList/HashMap} 把 OJ
的多组用例/题目装成内存题库;\textbf{{[}第二步{]}}
\GC \texttt{ProblemLoader/Repository} + \texttt{ProcessBuilder} 调外部
C++ 判题机 + 序列化。
\end{quote}

\begin{quote}
到目前为止，Java 在 JVM 里模拟判题------既无法真正卡住 CPU
时间，也无法强制限制内存。是时候让外部 C++ 沙箱接管真正的资源控制了。
\end{quote}

\hypertarget{ux4e0bux8f7dux4e0eux4f7fux7528-c-ux5224ux9898ux673a}{%
\subsection{下载与使用 C++
判题机}\label{ux4e0bux8f7dux4e0eux4f7fux7528-c-ux5224ux9898ux673a}}

第三代判题机是一个\textbf{独立的 C++ 单文件程序}
\texttt{judge.cpp}------与 Java 解耦,Java 只负责''调用 +
读结果''。它支持 \texttt{cpp/c/python/java},判
\texttt{AC/WA/PE/TLE/MLE/RE/CE},带 \texttt{setrlimit} + 进程组掐断 +
编译超时,健壮性不影响教学推进。

v 下载 judge.cpp(单文件源码) v 下载完整判题机(含 Makefile + 示例)

\textbf{① 放置位置} ------ 放到项目根目录下的 \texttt{judge/}:

\begin{verbatim}
mini-oj/
+-- src/oj/...           # Java 代码(MachineJudge 在 oj.judge)
+-- problems/<id>/       # 题目:config.txt + N.in/N.out
+-- judge/
    +-- judge.cpp
    +-- judge            # ② 编译出来的二进制
\end{verbatim}

\textbf{② 编译}(一次即可,需要 g++):

\begin{Shaded}
\begin{Highlighting}[]
\BuiltInTok{cd}\NormalTok{ judge }\KeywordTok{\&\&} \ExtensionTok{g++} \AttributeTok{{-}O2} \AttributeTok{{-}std}\OperatorTok{=}\NormalTok{c++17 }\AttributeTok{{-}o}\NormalTok{ judge judge.cpp   }\CommentTok{\# 或 make}
\end{Highlighting}
\end{Shaded}

\textbf{③ 命令行用法}:

\begin{Shaded}
\begin{Highlighting}[]
\ExtensionTok{./judge/judge} \AttributeTok{{-}{-}problem}\NormalTok{ problems/1 }\AttributeTok{{-}{-}src}\NormalTok{ sub.cpp }\AttributeTok{{-}{-}lang}\NormalTok{ cpp }\AttributeTok{{-}{-}time{-}ms}\NormalTok{ 1000 }\AttributeTok{{-}{-}mem{-}mb}\NormalTok{ 256 }\PreprocessorTok{[{-}{-}}\SpecialStringTok{special}\PreprocessorTok{]}
\CommentTok{\# {-}\textgreater{} \{"status":"AC","passed":3,"total":3,"time\_ms":5,"mem\_kb":1840,"detail":""\}}
\end{Highlighting}
\end{Shaded}

\textbf{④ Java 怎么调} ------ 就是下面 \texttt{MachineJudge}
干的事(\texttt{ProcessBuilder} 启动 + 正则解析 JSON):

\begin{Shaded}
\begin{Highlighting}[]
\NormalTok{MachineJudge mj }\OperatorTok{=} \KeywordTok{new} \FunctionTok{MachineJudge}\OperatorTok{(}\StringTok{"judge/judge"}\OperatorTok{);}
\NormalTok{JudgeResult r }\OperatorTok{=}\NormalTok{ mj}\OperatorTok{.}\FunctionTok{judge}\OperatorTok{(}\StringTok{"problems/1"}\OperatorTok{,} \StringTok{"sub.cpp"}\OperatorTok{,} \StringTok{"cpp"}\OperatorTok{,} \DecValTok{1000}\OperatorTok{,} \DecValTok{256}\OperatorTok{);}
\end{Highlighting}
\end{Shaded}

\begin{quote}
判题机完整说明见仓库
\texttt{judge/README.md};它在本仓库已编译验证(AC/WA/TLE/RE/CE/MLE/PE +
python + 浮点 special 全通过)。
\end{quote}

\hypertarget{ux6838ux5fc3ux75dbux70b9-1}{%
\subsection{【核心痛点】}\label{ux6838ux5fc3ux75dbux70b9-1}}

\textbf{第一代、第二代判题的天花板}

M1--M4 的判题链路始终在 JVM 内部：\texttt{Solution.solve()}
直接返回字符串，\texttt{StandardJudge}
做字符串比对。这种做法在教学早期够用，但有三个根本缺陷：

\begin{enumerate}
\def\labelenumi{\arabic{enumi}.}
\tightlist
\item
  \textbf{无法真正限制 CPU}：Java 没有
  \texttt{setrlimit(RLIMIT\_CPU,\ ...)}，\texttt{Thread.interrupt()}
  只能请求终止，无法强杀死循环。
\item
  \textbf{无法真正限制内存}：JVM 堆本身就是几百 MB
  的大容器，\texttt{-Xmx} 限制的是 JVM 堆，不是选手代码的实际内存分配。
\item
  \textbf{无法支持多语言}：第二代只能判 Java，而 OJ 需要支持
  C++/Python/Java 等语言，每种语言的编译/运行命令完全不同。
\end{enumerate}

\textbf{破局方案}：Java
只做''调度员''，真正的编译、运行、资源限制全部委托给 C++ 编写的外部
judge 二进制。Java 用 \texttt{ProcessBuilder}
启动这个二进制，传入源文件路径和时间/内存限制，读取其标准输出的一行
JSON，正则解析后回填到 \texttt{JudgeResult}。

\begin{center}\rule{0.5\linewidth}{0.5pt}\end{center}

\hypertarget{ux5f15ux5165ux8bfeux672cux77e5ux8bc6ux70b9-1}{%
\subsection{【引入课本知识点】}\label{ux5f15ux5165ux8bfeux672cux77e5ux8bc6ux70b9-1}}

\hypertarget{ch15-ux6cdbux578bux4e0eux96c6ux5408}{%
\subsubsection{Ch15 ·
泛型与集合}\label{ch15-ux6cdbux578bux4e0eux96c6ux5408}}

\textbf{泛型类与接口声明}

泛型让容器在编译期就确定元素类型，消除强制转型：

\begin{Shaded}
\begin{Highlighting}[]
\CommentTok{// 编译器保证 list 里只有 TestCase，取出无需 (TestCase) 强转}
\BuiltInTok{List}\OperatorTok{\textless{}}\NormalTok{TestCase}\OperatorTok{\textgreater{}}\NormalTok{ cases }\OperatorTok{=} \KeywordTok{new} \BuiltInTok{LinkedList}\OperatorTok{\textless{}\textgreater{}();}
\NormalTok{cases}\OperatorTok{.}\FunctionTok{add}\OperatorTok{(}\KeywordTok{new} \FunctionTok{TestCase}\OperatorTok{(}\StringTok{"1 2"}\OperatorTok{,} \StringTok{"3"}\OperatorTok{));}
\NormalTok{TestCase tc }\OperatorTok{=}\NormalTok{ cases}\OperatorTok{.}\FunctionTok{get}\OperatorTok{(}\DecValTok{0}\OperatorTok{);}  \CommentTok{// 类型安全}
\end{Highlighting}
\end{Shaded}

\texttt{ProblemLoader} 用 \texttt{List\textless{}TestCase\textgreater{}}
替代 M4 的 \texttt{TestCase{[}{]}}，原因有二： -
测试点数量未知，扫目录时逐条 \texttt{add()}，不需要预先
\texttt{new\ TestCase{[}n{]}}。 - \texttt{List} 提供迭代器和
Stream，后续统计更方便。

\textbf{\texttt{HashMap\textless{}Integer,\ Problem\textgreater{}}
题目映射}

\begin{Shaded}
\begin{Highlighting}[]
\BuiltInTok{Map}\OperatorTok{\textless{}}\BuiltInTok{Integer}\OperatorTok{,}\NormalTok{ Problem}\OperatorTok{\textgreater{}}\NormalTok{ cache }\OperatorTok{=} \KeywordTok{new} \BuiltInTok{HashMap}\OperatorTok{\textless{}\textgreater{}();}
\NormalTok{cache}\OperatorTok{.}\FunctionTok{put}\OperatorTok{(}\DecValTok{1001}\OperatorTok{,}\NormalTok{ problem}\OperatorTok{);}
\NormalTok{Problem p }\OperatorTok{=}\NormalTok{ cache}\OperatorTok{.}\FunctionTok{get}\OperatorTok{(}\DecValTok{1001}\OperatorTok{);}  \CommentTok{// O(1) 查找}
\end{Highlighting}
\end{Shaded}

\texttt{ProblemRepository}
用此结构做内存缓存：第一次从磁盘加载，之后命中缓存直接返回。

\textbf{Stream 流式统计示例}

\begin{Shaded}
\begin{Highlighting}[]
\CommentTok{// 统计一批提交里各 Status 出现次数}
\BuiltInTok{Map}\OperatorTok{\textless{}}\NormalTok{Status}\OperatorTok{,} \BuiltInTok{Long}\OperatorTok{\textgreater{}}\NormalTok{ stat }\OperatorTok{=}\NormalTok{ submissions}\OperatorTok{.}\FunctionTok{stream}\OperatorTok{()}
    \OperatorTok{.}\FunctionTok{map}\OperatorTok{(}\NormalTok{s }\OperatorTok{{-}\textgreater{}}\NormalTok{ s}\OperatorTok{.}\FunctionTok{getResult}\OperatorTok{().}\FunctionTok{getStatus}\OperatorTok{())}
    \OperatorTok{.}\FunctionTok{collect}\OperatorTok{(}\NormalTok{Collectors}\OperatorTok{.}\FunctionTok{groupingBy}\OperatorTok{(}\NormalTok{s }\OperatorTok{{-}\textgreater{}}\NormalTok{ s}\OperatorTok{,}\NormalTok{ Collectors}\OperatorTok{.}\FunctionTok{counting}\OperatorTok{()));}
\end{Highlighting}
\end{Shaded}

这是 M5a 结尾的统计演示，让学生感受 Stream 的链式风格。

\hypertarget{ch10-ux76eeux5f55ux904dux5386ux4e0eux5bf9ux8c61ux5e8fux5217ux5316}{%
\subsubsection{Ch10 ·
目录遍历与对象序列化}\label{ch10-ux76eeux5f55ux904dux5386ux4e0eux5bf9ux8c61ux5e8fux5217ux5316}}

\textbf{目录遍历}

\begin{Shaded}
\begin{Highlighting}[]
\BuiltInTok{File}\NormalTok{ dir }\OperatorTok{=} \KeywordTok{new} \BuiltInTok{File}\OperatorTok{(}\StringTok{"problems/1001"}\OperatorTok{);}
\BuiltInTok{File}\OperatorTok{[]}\NormalTok{ files }\OperatorTok{=}\NormalTok{ dir}\OperatorTok{.}\FunctionTok{listFiles}\OperatorTok{((}\NormalTok{d}\OperatorTok{,}\NormalTok{ name}\OperatorTok{)} \OperatorTok{{-}\textgreater{}}\NormalTok{ name}\OperatorTok{.}\FunctionTok{endsWith}\OperatorTok{(}\StringTok{".in"}\OperatorTok{));}
\end{Highlighting}
\end{Shaded}

\texttt{ProblemLoader} 用此方式扫
\texttt{problems/\textless{}id\textgreater{}/} 下所有
\texttt{N.in}/\texttt{N.out} 文件对。

\textbf{对象序列化（\texttt{Serializable}）}

实现 \texttt{Serializable} 的对象可以用 \texttt{ObjectOutputStream}
直接写入文件，用 \texttt{ObjectInputStream} 原样读回，不需要手写 CSV
格式。\texttt{SubmissionStore} 正是用这种方式把 \texttt{Submission}
对象持久化到 \texttt{submissions.dat}。

\begin{center}\rule{0.5\linewidth}{0.5pt}\end{center}

\hypertarget{ux4e09ux4ee3ux6f14ux8fdbux5b9aux4f4d-1}{%
\subsection{【三代演进定位】}\label{ux4e09ux4ee3ux6f14ux8fdbux5b9aux4f4d-1}}

\begin{longtable}[]{@{}
  >{\raggedright\arraybackslash}p{(\columnwidth - 8\tabcolsep) * \real{0.2000}}
  >{\raggedright\arraybackslash}p{(\columnwidth - 8\tabcolsep) * \real{0.2000}}
  >{\raggedright\arraybackslash}p{(\columnwidth - 8\tabcolsep) * \real{0.2000}}
  >{\raggedright\arraybackslash}p{(\columnwidth - 8\tabcolsep) * \real{0.2000}}
  >{\raggedright\arraybackslash}p{(\columnwidth - 8\tabcolsep) * \real{0.2000}}@{}}
\toprule\noalign{}
\begin{minipage}[b]{\linewidth}\raggedright
代际
\end{minipage} & \begin{minipage}[b]{\linewidth}\raggedright
判题入口
\end{minipage} & \begin{minipage}[b]{\linewidth}\raggedright
资源限制
\end{minipage} & \begin{minipage}[b]{\linewidth}\raggedright
多语言
\end{minipage} & \begin{minipage}[b]{\linewidth}\raggedright
定位
\end{minipage} \\
\midrule\noalign{}
\endhead
\bottomrule\noalign{}
\endlastfoot
第一代（M1--M3） & \texttt{Solution.solve()} & 无 & 仅 Java & JVM
内模拟，教学启蒙 \\
第二代（M4） & 反射 \texttt{JudgeFactory.create()} & 无 & 仅 Java &
反射解耦，配置驱动 \\
\textbf{第三代（M5a）} & \textbf{\texttt{MachineJudge} + ProcessBuilder}
& \textbf{setrlimit（C++ 端）} & \textbf{C++/Python/Java} &
\textbf{跨界真判题，工业级沙箱} \\
\end{longtable}

\textbf{第三代的关键跨越}：Java 进程通过 \texttt{ProcessBuilder}
启动一个预编译好的 C++
二进制（\texttt{judge}），后者接收命令行参数（源文件、语言、时间/内存限制、题目目录），在
Linux 内核层面用 \texttt{setrlimit}
施加资源上限，运行完毕后向标准输出打印一行 JSON：

\begin{Shaded}
\begin{Highlighting}[]
\FunctionTok{\{}\DataTypeTok{"status"}\FunctionTok{:}\StringTok{"AC"}\FunctionTok{,}\DataTypeTok{"passed"}\FunctionTok{:}\DecValTok{3}\FunctionTok{,}\DataTypeTok{"total"}\FunctionTok{:}\DecValTok{3}\FunctionTok{,}\DataTypeTok{"time\_ms"}\FunctionTok{:}\DecValTok{45}\FunctionTok{,}\DataTypeTok{"mem\_kb"}\FunctionTok{:}\DecValTok{2048}\FunctionTok{,}\DataTypeTok{"detail"}\FunctionTok{:}\StringTok{""}\FunctionTok{\}}
\end{Highlighting}
\end{Shaded}

Java 用正则表达式解析这行 JSON，构造 \texttt{JudgeResult} 返回。C++
侧代码不在本章展开，本章聚焦 \textbf{Java 侧的调用接口与数据流}。

\begin{center}\rule{0.5\linewidth}{0.5pt}\end{center}

\hypertarget{ux65b0ux4ea7ux7269ux67b6ux6784-1}{%
\subsection{【新产物架构】}\label{ux65b0ux4ea7ux7269ux67b6ux6784-1}}

本节给出 M5a 四个新产物的完整实现。

\hypertarget{oj.io.problemloader-ux626bux76eeux5f55ux8fd4ux56de-listtestcase}{%
\subsubsection{\texorpdfstring{1. \texttt{oj.io.ProblemLoader} ---
扫目录，返回
\texttt{List\textless{}TestCase\textgreater{}}}{1. oj.io.ProblemLoader --- 扫目录，返回 List\textless TestCase\textgreater{}}}\label{oj.io.problemloader-ux626bux76eeux5f55ux8fd4ux56de-listtestcase}}

\textbf{职责}：扫描 \texttt{problems/\textless{}id\textgreater{}/}
目录，按编号顺序读取
\texttt{1.in}/\texttt{1.out}、\texttt{2.in}/\texttt{2.out}\ldots\ldots 构造
\texttt{List\textless{}TestCase\textgreater{}}。

\begin{Shaded}
\begin{Highlighting}[]
\KeywordTok{package}\ImportTok{ oj}\OperatorTok{.}\ImportTok{io}\OperatorTok{;}

\KeywordTok{import} \ImportTok{oj}\OperatorTok{.}\ImportTok{core}\OperatorTok{.}\ImportTok{TestCase}\OperatorTok{;}

\KeywordTok{import} \ImportTok{java}\OperatorTok{.}\ImportTok{io}\OperatorTok{.}\ImportTok{File}\OperatorTok{;}
\KeywordTok{import} \ImportTok{java}\OperatorTok{.}\ImportTok{io}\OperatorTok{.}\ImportTok{IOException}\OperatorTok{;}
\KeywordTok{import} \ImportTok{java}\OperatorTok{.}\ImportTok{nio}\OperatorTok{.}\ImportTok{file}\OperatorTok{.}\ImportTok{Files}\OperatorTok{;}
\KeywordTok{import} \ImportTok{java}\OperatorTok{.}\ImportTok{util}\OperatorTok{.}\ImportTok{ArrayList}\OperatorTok{;}
\KeywordTok{import} \ImportTok{java}\OperatorTok{.}\ImportTok{util}\OperatorTok{.}\ImportTok{Arrays}\OperatorTok{;}
\KeywordTok{import} \ImportTok{java}\OperatorTok{.}\ImportTok{util}\OperatorTok{.}\ImportTok{List}\OperatorTok{;}

\KeywordTok{public} \KeywordTok{class}\NormalTok{ ProblemLoader }\OperatorTok{\{}

    \KeywordTok{private} \DataTypeTok{final} \BuiltInTok{String}\NormalTok{ problemsRoot}\OperatorTok{;}

    \KeywordTok{public} \FunctionTok{ProblemLoader}\OperatorTok{(}\BuiltInTok{String}\NormalTok{ problemsRoot}\OperatorTok{)} \OperatorTok{\{}
        \KeywordTok{this}\OperatorTok{.}\FunctionTok{problemsRoot} \OperatorTok{=}\NormalTok{ problemsRoot}\OperatorTok{;}
    \OperatorTok{\}}

    \CommentTok{/**}
     \CommentTok{*} \CommentTok{扫描}\NormalTok{ problems}\CommentTok{/}\DataTypeTok{\textless{}}\KeywordTok{id}\DataTypeTok{\textgreater{}}\CommentTok{/} \CommentTok{目录，返回按序号排列的测试点列表。}
     \CommentTok{*} \CommentTok{每对}\NormalTok{ N}\CommentTok{.}\NormalTok{in }\CommentTok{/}\NormalTok{ N}\CommentTok{.}\NormalTok{out }\CommentTok{构成一个}\NormalTok{ TestCase}\CommentTok{。}
     \CommentTok{*}
\CommentTok{     * @}\NormalTok{param id }\CommentTok{题目}\NormalTok{ ID}
     \CommentTok{*} \CommentTok{@}\NormalTok{return }\CommentTok{有序的}\NormalTok{ TestCase }\CommentTok{列表，若目录不存在则返回空列表}
     \CommentTok{*/}
    \KeywordTok{public} \BuiltInTok{List}\OperatorTok{\textless{}}\NormalTok{TestCase}\OperatorTok{\textgreater{}} \FunctionTok{cases}\OperatorTok{(}\DataTypeTok{int}\NormalTok{ id}\OperatorTok{)} \OperatorTok{\{}
        \BuiltInTok{File}\NormalTok{ dir }\OperatorTok{=} \KeywordTok{new} \BuiltInTok{File}\OperatorTok{(}\NormalTok{problemsRoot}\OperatorTok{,} \BuiltInTok{String}\OperatorTok{.}\FunctionTok{valueOf}\OperatorTok{(}\NormalTok{id}\OperatorTok{));}
        \BuiltInTok{List}\OperatorTok{\textless{}}\NormalTok{TestCase}\OperatorTok{\textgreater{}}\NormalTok{ result }\OperatorTok{=} \KeywordTok{new} \BuiltInTok{ArrayList}\OperatorTok{\textless{}\textgreater{}();}
        \ControlFlowTok{if} \OperatorTok{(!}\NormalTok{dir}\OperatorTok{.}\FunctionTok{isDirectory}\OperatorTok{())} \OperatorTok{\{}
            \ControlFlowTok{return}\NormalTok{ result}\OperatorTok{;}
        \OperatorTok{\}}

        \CommentTok{// 收集所有 .in 文件，按文件名数字排序}
        \BuiltInTok{File}\OperatorTok{[]}\NormalTok{ inFiles }\OperatorTok{=}\NormalTok{ dir}\OperatorTok{.}\FunctionTok{listFiles}\OperatorTok{((}\NormalTok{d}\OperatorTok{,}\NormalTok{ name}\OperatorTok{)} \OperatorTok{{-}\textgreater{}}\NormalTok{ name}\OperatorTok{.}\FunctionTok{endsWith}\OperatorTok{(}\StringTok{".in"}\OperatorTok{));}
        \ControlFlowTok{if} \OperatorTok{(}\NormalTok{inFiles }\OperatorTok{==} \KeywordTok{null} \OperatorTok{||}\NormalTok{ inFiles}\OperatorTok{.}\FunctionTok{length} \OperatorTok{==} \DecValTok{0}\OperatorTok{)} \OperatorTok{\{}
            \ControlFlowTok{return}\NormalTok{ result}\OperatorTok{;}
        \OperatorTok{\}}
        \BuiltInTok{Arrays}\OperatorTok{.}\FunctionTok{sort}\OperatorTok{(}\NormalTok{inFiles}\OperatorTok{,} \OperatorTok{(}\NormalTok{a}\OperatorTok{,}\NormalTok{ b}\OperatorTok{)} \OperatorTok{{-}\textgreater{}} \OperatorTok{\{}
            \DataTypeTok{int}\NormalTok{ na }\OperatorTok{=} \FunctionTok{indexOf}\OperatorTok{(}\NormalTok{a}\OperatorTok{.}\FunctionTok{getName}\OperatorTok{());}
            \DataTypeTok{int}\NormalTok{ nb }\OperatorTok{=} \FunctionTok{indexOf}\OperatorTok{(}\NormalTok{b}\OperatorTok{.}\FunctionTok{getName}\OperatorTok{());}
            \ControlFlowTok{return} \BuiltInTok{Integer}\OperatorTok{.}\FunctionTok{compare}\OperatorTok{(}\NormalTok{na}\OperatorTok{,}\NormalTok{ nb}\OperatorTok{);}
        \OperatorTok{\});}

        \ControlFlowTok{for} \OperatorTok{(}\BuiltInTok{File}\NormalTok{ inFile }\OperatorTok{:}\NormalTok{ inFiles}\OperatorTok{)} \OperatorTok{\{}
            \BuiltInTok{String}\NormalTok{ baseName }\OperatorTok{=}\NormalTok{ inFile}\OperatorTok{.}\FunctionTok{getName}\OperatorTok{().}\FunctionTok{replace}\OperatorTok{(}\StringTok{".in"}\OperatorTok{,} \StringTok{""}\OperatorTok{);}
            \BuiltInTok{File}\NormalTok{ outFile }\OperatorTok{=} \KeywordTok{new} \BuiltInTok{File}\OperatorTok{(}\NormalTok{dir}\OperatorTok{,}\NormalTok{ baseName }\OperatorTok{+} \StringTok{".out"}\OperatorTok{);}
            \ControlFlowTok{if} \OperatorTok{(!}\NormalTok{outFile}\OperatorTok{.}\FunctionTok{exists}\OperatorTok{())} \OperatorTok{\{}
                \ControlFlowTok{continue}\OperatorTok{;}
            \OperatorTok{\}}
            \ControlFlowTok{try} \OperatorTok{\{}
                \BuiltInTok{String}\NormalTok{ input }\OperatorTok{=} \KeywordTok{new} \BuiltInTok{String}\OperatorTok{(}\NormalTok{Files}\OperatorTok{.}\FunctionTok{readAllBytes}\OperatorTok{(}\NormalTok{inFile}\OperatorTok{.}\FunctionTok{toPath}\OperatorTok{())).}\FunctionTok{trim}\OperatorTok{();}
                \BuiltInTok{String}\NormalTok{ expected }\OperatorTok{=} \KeywordTok{new} \BuiltInTok{String}\OperatorTok{(}\NormalTok{Files}\OperatorTok{.}\FunctionTok{readAllBytes}\OperatorTok{(}\NormalTok{outFile}\OperatorTok{.}\FunctionTok{toPath}\OperatorTok{())).}\FunctionTok{trim}\OperatorTok{();}
\NormalTok{                result}\OperatorTok{.}\FunctionTok{add}\OperatorTok{(}\KeywordTok{new} \FunctionTok{TestCase}\OperatorTok{(}\NormalTok{input}\OperatorTok{,}\NormalTok{ expected}\OperatorTok{));}
            \OperatorTok{\}} \ControlFlowTok{catch} \OperatorTok{(}\BuiltInTok{IOException}\NormalTok{ e}\OperatorTok{)} \OperatorTok{\{}
                \BuiltInTok{System}\OperatorTok{.}\FunctionTok{err}\OperatorTok{.}\FunctionTok{println}\OperatorTok{(}\StringTok{"[ProblemLoader] 读取测试点失败: "} \OperatorTok{+}\NormalTok{ inFile}\OperatorTok{.}\FunctionTok{getName}\OperatorTok{());}
            \OperatorTok{\}}
        \OperatorTok{\}}
        \ControlFlowTok{return}\NormalTok{ result}\OperatorTok{;}
    \OperatorTok{\}}

    \CommentTok{/**} \CommentTok{从} \CommentTok{"3.}\NormalTok{in}\CommentTok{"} \CommentTok{这样的文件名中提取数字} \CommentTok{3} \CommentTok{*/}
    \KeywordTok{private} \DataTypeTok{int} \FunctionTok{indexOf}\OperatorTok{(}\BuiltInTok{String}\NormalTok{ filename}\OperatorTok{)} \OperatorTok{\{}
        \ControlFlowTok{try} \OperatorTok{\{}
            \ControlFlowTok{return} \BuiltInTok{Integer}\OperatorTok{.}\FunctionTok{parseInt}\OperatorTok{(}\NormalTok{filename}\OperatorTok{.}\FunctionTok{replaceAll}\OperatorTok{(}\StringTok{"[\^{}0{-}9]"}\OperatorTok{,} \StringTok{""}\OperatorTok{));}
        \OperatorTok{\}} \ControlFlowTok{catch} \OperatorTok{(}\BuiltInTok{NumberFormatException}\NormalTok{ e}\OperatorTok{)} \OperatorTok{\{}
            \ControlFlowTok{return} \BuiltInTok{Integer}\OperatorTok{.}\FunctionTok{MAX\_VALUE}\OperatorTok{;}
        \OperatorTok{\}}
    \OperatorTok{\}}
\OperatorTok{\}}
\end{Highlighting}
\end{Shaded}

\textbf{目录结构约定}：

\begin{verbatim}
problems/
  1001/
    config.txt        <- 题目元数据（M4 沿用）
    1.in  1.out
    2.in  2.out
    3.in  3.out
\end{verbatim}

\begin{center}\rule{0.5\linewidth}{0.5pt}\end{center}

\hypertarget{oj.io.problemrepository-hashmap-ux7f13ux5b58ux9898ux76ee}{%
\subsubsection{\texorpdfstring{2. \texttt{oj.io.ProblemRepository} ---
\texttt{HashMap}
缓存题目}{2. oj.io.ProblemRepository --- HashMap 缓存题目}}\label{oj.io.problemrepository-hashmap-ux7f13ux5b58ux9898ux76ee}}

\textbf{职责}：组合 \texttt{ConfigFile}（M4 产物，读
\texttt{config.txt}）与 \texttt{ProblemLoader}，对外提供
\texttt{get(id)} 接口；内部用
\texttt{HashMap\textless{}Integer,\ Problem\textgreater{}}
缓存，同一题目只从磁盘加载一次。

\begin{Shaded}
\begin{Highlighting}[]
\KeywordTok{package}\ImportTok{ oj}\OperatorTok{.}\ImportTok{io}\OperatorTok{;}

\KeywordTok{import} \ImportTok{oj}\OperatorTok{.}\ImportTok{core}\OperatorTok{.}\ImportTok{Problem}\OperatorTok{;}
\KeywordTok{import} \ImportTok{oj}\OperatorTok{.}\ImportTok{core}\OperatorTok{.}\ImportTok{ProblemMeta}\OperatorTok{;}
\KeywordTok{import} \ImportTok{oj}\OperatorTok{.}\ImportTok{core}\OperatorTok{.}\ImportTok{TestCase}\OperatorTok{;}

\KeywordTok{import} \ImportTok{java}\OperatorTok{.}\ImportTok{io}\OperatorTok{.}\ImportTok{File}\OperatorTok{;}
\KeywordTok{import} \ImportTok{java}\OperatorTok{.}\ImportTok{util}\OperatorTok{.}\ImportTok{HashMap}\OperatorTok{;}
\KeywordTok{import} \ImportTok{java}\OperatorTok{.}\ImportTok{util}\OperatorTok{.}\ImportTok{List}\OperatorTok{;}
\KeywordTok{import} \ImportTok{java}\OperatorTok{.}\ImportTok{util}\OperatorTok{.}\ImportTok{Map}\OperatorTok{;}

\KeywordTok{public} \KeywordTok{class}\NormalTok{ ProblemRepository }\OperatorTok{\{}

    \KeywordTok{private} \DataTypeTok{final} \BuiltInTok{String}\NormalTok{ problemsRoot}\OperatorTok{;}
    \KeywordTok{private} \DataTypeTok{final}\NormalTok{ ProblemLoader loader}\OperatorTok{;}
    \KeywordTok{private} \DataTypeTok{final} \BuiltInTok{Map}\OperatorTok{\textless{}}\BuiltInTok{Integer}\OperatorTok{,}\NormalTok{ Problem}\OperatorTok{\textgreater{}}\NormalTok{ cache }\OperatorTok{=} \KeywordTok{new} \BuiltInTok{HashMap}\OperatorTok{\textless{}\textgreater{}();}

    \KeywordTok{public} \FunctionTok{ProblemRepository}\OperatorTok{(}\BuiltInTok{String}\NormalTok{ problemsRoot}\OperatorTok{)} \OperatorTok{\{}
        \KeywordTok{this}\OperatorTok{.}\FunctionTok{problemsRoot} \OperatorTok{=}\NormalTok{ problemsRoot}\OperatorTok{;}
        \KeywordTok{this}\OperatorTok{.}\FunctionTok{loader} \OperatorTok{=} \KeywordTok{new} \FunctionTok{ProblemLoader}\OperatorTok{(}\NormalTok{problemsRoot}\OperatorTok{);}
    \OperatorTok{\}}

    \CommentTok{/**}
     \CommentTok{*} \CommentTok{返回指定}\NormalTok{ ID }\CommentTok{的题目，优先命中缓存。}
     \CommentTok{*} \CommentTok{若缓存未命中则从磁盘加载}\NormalTok{ config}\CommentTok{.}\NormalTok{txt }\CommentTok{+} \CommentTok{测试点，存入缓存后返回。}
     \CommentTok{*}
\CommentTok{     * @}\NormalTok{param id }\CommentTok{题目}\NormalTok{ ID}
     \CommentTok{*} \CommentTok{@}\NormalTok{return Problem }\CommentTok{对象，若不存在则返回}\NormalTok{ null}
     \CommentTok{*/}
    \KeywordTok{public}\NormalTok{ Problem }\FunctionTok{get}\OperatorTok{(}\DataTypeTok{int}\NormalTok{ id}\OperatorTok{)} \OperatorTok{\{}
        \ControlFlowTok{if} \OperatorTok{(}\NormalTok{cache}\OperatorTok{.}\FunctionTok{containsKey}\OperatorTok{(}\NormalTok{id}\OperatorTok{))} \OperatorTok{\{}
            \ControlFlowTok{return}\NormalTok{ cache}\OperatorTok{.}\FunctionTok{get}\OperatorTok{(}\NormalTok{id}\OperatorTok{);}
        \OperatorTok{\}}
\NormalTok{        Problem p }\OperatorTok{=} \FunctionTok{load}\OperatorTok{(}\NormalTok{id}\OperatorTok{);}
        \ControlFlowTok{if} \OperatorTok{(}\NormalTok{p }\OperatorTok{!=} \KeywordTok{null}\OperatorTok{)} \OperatorTok{\{}
\NormalTok{            cache}\OperatorTok{.}\FunctionTok{put}\OperatorTok{(}\NormalTok{id}\OperatorTok{,}\NormalTok{ p}\OperatorTok{);}
        \OperatorTok{\}}
        \ControlFlowTok{return}\NormalTok{ p}\OperatorTok{;}
    \OperatorTok{\}}

    \KeywordTok{private}\NormalTok{ Problem }\FunctionTok{load}\OperatorTok{(}\DataTypeTok{int}\NormalTok{ id}\OperatorTok{)} \OperatorTok{\{}
        \BuiltInTok{File}\NormalTok{ dir }\OperatorTok{=} \KeywordTok{new} \BuiltInTok{File}\OperatorTok{(}\NormalTok{problemsRoot}\OperatorTok{,} \BuiltInTok{String}\OperatorTok{.}\FunctionTok{valueOf}\OperatorTok{(}\NormalTok{id}\OperatorTok{));}
        \BuiltInTok{File}\NormalTok{ configFile }\OperatorTok{=} \KeywordTok{new} \BuiltInTok{File}\OperatorTok{(}\NormalTok{dir}\OperatorTok{,} \StringTok{"config.txt"}\OperatorTok{);}
        \ControlFlowTok{if} \OperatorTok{(!}\NormalTok{configFile}\OperatorTok{.}\FunctionTok{exists}\OperatorTok{())} \OperatorTok{\{}
            \ControlFlowTok{return} \KeywordTok{null}\OperatorTok{;}
        \OperatorTok{\}}
\NormalTok{        ProblemMeta meta }\OperatorTok{=}\NormalTok{ ConfigFile}\OperatorTok{.}\FunctionTok{read}\OperatorTok{(}\NormalTok{dir}\OperatorTok{);}  \CommentTok{// M4 契约:传目录,内部读 config.txt}
        \BuiltInTok{List}\OperatorTok{\textless{}}\NormalTok{TestCase}\OperatorTok{\textgreater{}}\NormalTok{ cases }\OperatorTok{=}\NormalTok{ loader}\OperatorTok{.}\FunctionTok{cases}\OperatorTok{(}\NormalTok{id}\OperatorTok{);}
        \ControlFlowTok{return} \KeywordTok{new} \FunctionTok{Problem}\OperatorTok{(}\NormalTok{id}\OperatorTok{,}\NormalTok{ meta}\OperatorTok{,}\NormalTok{ cases}\OperatorTok{);}
        \CommentTok{// Problem 构造器在 M5a 重构为接受 List\textless{}TestCase\textgreater{}}
    \OperatorTok{\}}

    \CommentTok{/**} \CommentTok{清除缓存（测试用途）} \CommentTok{*/}
    \KeywordTok{public} \DataTypeTok{void} \FunctionTok{evict}\OperatorTok{(}\DataTypeTok{int}\NormalTok{ id}\OperatorTok{)} \OperatorTok{\{}
\NormalTok{        cache}\OperatorTok{.}\FunctionTok{remove}\OperatorTok{(}\NormalTok{id}\OperatorTok{);}
    \OperatorTok{\}}

    \CommentTok{/**} \CommentTok{已缓存的题目数量} \CommentTok{*/}
    \KeywordTok{public} \DataTypeTok{int} \FunctionTok{cacheSize}\OperatorTok{()} \OperatorTok{\{}
        \ControlFlowTok{return}\NormalTok{ cache}\OperatorTok{.}\FunctionTok{size}\OperatorTok{();}
    \OperatorTok{\}}
\OperatorTok{\}}
\end{Highlighting}
\end{Shaded}

\begin{center}\rule{0.5\linewidth}{0.5pt}\end{center}

\hypertarget{oj.io.submissionstore-ux5bf9ux8c61ux5e8fux5217ux5316ux6301ux4e45ux5316}{%
\subsubsection{\texorpdfstring{3. \texttt{oj.io.SubmissionStore} ---
对象序列化持久化}{3. oj.io.SubmissionStore --- 对象序列化持久化}}\label{oj.io.submissionstore-ux5bf9ux8c61ux5e8fux5217ux5316ux6301ux4e45ux5316}}

\textbf{职责}：把 \texttt{Submission} 对象列表追加写入
\texttt{submissions.dat}，支持全量读回。利用 Java
对象序列化（\texttt{ObjectOutputStream}/\texttt{ObjectInputStream}）完成持久化，不需要手写文本格式。

\begin{Shaded}
\begin{Highlighting}[]
\KeywordTok{package}\ImportTok{ oj}\OperatorTok{.}\ImportTok{io}\OperatorTok{;}

\KeywordTok{import} \ImportTok{oj}\OperatorTok{.}\ImportTok{core}\OperatorTok{.}\ImportTok{Submission}\OperatorTok{;}

\KeywordTok{import} \ImportTok{java}\OperatorTok{.}\ImportTok{io}\OperatorTok{.}\ImportTok{EOFException}\OperatorTok{;}
\KeywordTok{import} \ImportTok{java}\OperatorTok{.}\ImportTok{io}\OperatorTok{.}\ImportTok{File}\OperatorTok{;}
\KeywordTok{import} \ImportTok{java}\OperatorTok{.}\ImportTok{io}\OperatorTok{.}\ImportTok{FileInputStream}\OperatorTok{;}
\KeywordTok{import} \ImportTok{java}\OperatorTok{.}\ImportTok{io}\OperatorTok{.}\ImportTok{FileOutputStream}\OperatorTok{;}
\KeywordTok{import} \ImportTok{java}\OperatorTok{.}\ImportTok{io}\OperatorTok{.}\ImportTok{IOException}\OperatorTok{;}
\KeywordTok{import} \ImportTok{java}\OperatorTok{.}\ImportTok{io}\OperatorTok{.}\ImportTok{ObjectInputStream}\OperatorTok{;}
\KeywordTok{import} \ImportTok{java}\OperatorTok{.}\ImportTok{io}\OperatorTok{.}\ImportTok{ObjectOutputStream}\OperatorTok{;}
\KeywordTok{import} \ImportTok{java}\OperatorTok{.}\ImportTok{util}\OperatorTok{.}\ImportTok{ArrayList}\OperatorTok{;}
\KeywordTok{import} \ImportTok{java}\OperatorTok{.}\ImportTok{util}\OperatorTok{.}\ImportTok{List}\OperatorTok{;}

\KeywordTok{public} \KeywordTok{class}\NormalTok{ SubmissionStore }\OperatorTok{\{}

    \KeywordTok{private} \DataTypeTok{final} \BuiltInTok{File}\NormalTok{ storeFile}\OperatorTok{;}

    \KeywordTok{public} \FunctionTok{SubmissionStore}\OperatorTok{(}\BuiltInTok{String}\NormalTok{ path}\OperatorTok{)} \OperatorTok{\{}
        \KeywordTok{this}\OperatorTok{.}\FunctionTok{storeFile} \OperatorTok{=} \KeywordTok{new} \BuiltInTok{File}\OperatorTok{(}\NormalTok{path}\OperatorTok{);}
    \OperatorTok{\}}

    \CommentTok{/**}
     \CommentTok{*} \CommentTok{将一条提交追加写入持久化文件。}
     \CommentTok{*} \CommentTok{注意：}\NormalTok{ObjectOutputStream }\CommentTok{写多个对象到同一文件时，}
     \CommentTok{*} \CommentTok{后续追加必须跳过文件头（}\NormalTok{magic }\CommentTok{+}\NormalTok{ version}\CommentTok{），用内部类}\NormalTok{ AppendObjectOutputStream }\CommentTok{实现。}
     \CommentTok{*}
\CommentTok{     * @}\NormalTok{param submission }\CommentTok{要持久化的提交对象}
     \CommentTok{*/}
    \KeywordTok{public} \DataTypeTok{void} \FunctionTok{save}\OperatorTok{(}\NormalTok{Submission submission}\OperatorTok{)} \KeywordTok{throws} \BuiltInTok{IOException} \OperatorTok{\{}
        \ControlFlowTok{if} \OperatorTok{(}\NormalTok{storeFile}\OperatorTok{.}\FunctionTok{exists}\OperatorTok{()} \OperatorTok{\&\&}\NormalTok{ storeFile}\OperatorTok{.}\FunctionTok{length}\OperatorTok{()} \OperatorTok{\textgreater{}} \DecValTok{0}\OperatorTok{)} \OperatorTok{\{}
            \ControlFlowTok{try} \OperatorTok{(}\NormalTok{AppendObjectOutputStream out }\OperatorTok{=}
                     \KeywordTok{new} \FunctionTok{AppendObjectOutputStream}\OperatorTok{(}\KeywordTok{new} \BuiltInTok{FileOutputStream}\OperatorTok{(}\NormalTok{storeFile}\OperatorTok{,} \KeywordTok{true}\OperatorTok{)))} \OperatorTok{\{}
\NormalTok{                out}\OperatorTok{.}\FunctionTok{writeObject}\OperatorTok{(}\NormalTok{submission}\OperatorTok{);}
            \OperatorTok{\}}
        \OperatorTok{\}} \ControlFlowTok{else} \OperatorTok{\{}
            \ControlFlowTok{try} \OperatorTok{(}\BuiltInTok{ObjectOutputStream}\NormalTok{ out }\OperatorTok{=}
                     \KeywordTok{new} \BuiltInTok{ObjectOutputStream}\OperatorTok{(}\KeywordTok{new} \BuiltInTok{FileOutputStream}\OperatorTok{(}\NormalTok{storeFile}\OperatorTok{,} \KeywordTok{true}\OperatorTok{)))} \OperatorTok{\{}
\NormalTok{                out}\OperatorTok{.}\FunctionTok{writeObject}\OperatorTok{(}\NormalTok{submission}\OperatorTok{);}
            \OperatorTok{\}}
        \OperatorTok{\}}
    \OperatorTok{\}}

    \CommentTok{/**}
     \CommentTok{*} \CommentTok{读取文件中所有提交记录。}
     \CommentTok{*}
\CommentTok{     * @}\NormalTok{return }\CommentTok{提交列表，若文件不存在则返回空列表}
     \CommentTok{*/}
    \KeywordTok{public} \BuiltInTok{List}\OperatorTok{\textless{}}\NormalTok{Submission}\OperatorTok{\textgreater{}} \FunctionTok{loadAll}\OperatorTok{()} \KeywordTok{throws} \BuiltInTok{IOException}\OperatorTok{,} \BuiltInTok{ClassNotFoundException} \OperatorTok{\{}
        \BuiltInTok{List}\OperatorTok{\textless{}}\NormalTok{Submission}\OperatorTok{\textgreater{}}\NormalTok{ list }\OperatorTok{=} \KeywordTok{new} \BuiltInTok{ArrayList}\OperatorTok{\textless{}\textgreater{}();}
        \ControlFlowTok{if} \OperatorTok{(!}\NormalTok{storeFile}\OperatorTok{.}\FunctionTok{exists}\OperatorTok{()} \OperatorTok{||}\NormalTok{ storeFile}\OperatorTok{.}\FunctionTok{length}\OperatorTok{()} \OperatorTok{==} \DecValTok{0}\OperatorTok{)} \OperatorTok{\{}
            \ControlFlowTok{return}\NormalTok{ list}\OperatorTok{;}
        \OperatorTok{\}}
        \ControlFlowTok{try} \OperatorTok{(}\BuiltInTok{ObjectInputStream}\NormalTok{ in }\OperatorTok{=}
                 \KeywordTok{new} \BuiltInTok{ObjectInputStream}\OperatorTok{(}\KeywordTok{new} \BuiltInTok{FileInputStream}\OperatorTok{(}\NormalTok{storeFile}\OperatorTok{)))} \OperatorTok{\{}
            \ControlFlowTok{while} \OperatorTok{(}\KeywordTok{true}\OperatorTok{)} \OperatorTok{\{}
                \ControlFlowTok{try} \OperatorTok{\{}
\NormalTok{                    list}\OperatorTok{.}\FunctionTok{add}\OperatorTok{((}\NormalTok{Submission}\OperatorTok{)}\NormalTok{ in}\OperatorTok{.}\FunctionTok{readObject}\OperatorTok{());}
                \OperatorTok{\}} \ControlFlowTok{catch} \OperatorTok{(}\BuiltInTok{EOFException}\NormalTok{ e}\OperatorTok{)} \OperatorTok{\{}
                    \ControlFlowTok{break}\OperatorTok{;}  \CommentTok{// 正常结束，文件读完}
                \OperatorTok{\}}
            \OperatorTok{\}}
        \OperatorTok{\}}
        \ControlFlowTok{return}\NormalTok{ list}\OperatorTok{;}
    \OperatorTok{\}}

    \CommentTok{/**}
     \CommentTok{*} \CommentTok{追加写模式的}\NormalTok{ ObjectOutputStream}\CommentTok{：覆盖}\NormalTok{ writeStreamHeader}\CommentTok{()，}
     \CommentTok{*} \CommentTok{避免在文件中段再次写入}\NormalTok{ magic bytes}\CommentTok{，防止读回时报}\NormalTok{ StreamCorruptedException}\CommentTok{。}
     \CommentTok{*/}
    \KeywordTok{private} \DataTypeTok{static} \KeywordTok{class}\NormalTok{ AppendObjectOutputStream }\KeywordTok{extends} \BuiltInTok{ObjectOutputStream} \OperatorTok{\{}
        \FunctionTok{AppendObjectOutputStream}\OperatorTok{(}\BuiltInTok{FileOutputStream}\NormalTok{ out}\OperatorTok{)} \KeywordTok{throws} \BuiltInTok{IOException} \OperatorTok{\{}
            \KeywordTok{super}\OperatorTok{(}\NormalTok{out}\OperatorTok{);}
        \OperatorTok{\}}

        \AttributeTok{@Override}
        \KeywordTok{protected} \DataTypeTok{void} \FunctionTok{writeStreamHeader}\OperatorTok{()} \KeywordTok{throws} \BuiltInTok{IOException} \OperatorTok{\{}
            \FunctionTok{reset}\OperatorTok{();}  \CommentTok{// 只重置状态，不写 magic + version}
        \OperatorTok{\}}
    \OperatorTok{\}}
\OperatorTok{\}}
\end{Highlighting}
\end{Shaded}

\begin{quote}
\textbf{教学提示}：\texttt{AppendObjectOutputStream}
是一个常见陷阱。直接用
\texttt{new\ ObjectOutputStream(new\ FileOutputStream(file,\ true))}
追加写，第二条对象写入时会再次写入 4 字节 magic（\texttt{0xACED}）和 2
字节版本（\texttt{0x0005}），导致读回时
\texttt{StreamCorruptedException}。覆盖 \texttt{writeStreamHeader()}
改为 \texttt{reset()} 是标准解法。
\end{quote}

\begin{center}\rule{0.5\linewidth}{0.5pt}\end{center}

\hypertarget{oj.judge.machinejudge-ux7b2cux4e09ux4ee3ux5224ux9898ux673aux6838ux5fc3}{%
\subsubsection{\texorpdfstring{4. \texttt{oj.judge.MachineJudge} ---
第三代判题机核心}{4. oj.judge.MachineJudge --- 第三代判题机核心}}\label{oj.judge.machinejudge-ux7b2cux4e09ux4ee3ux5224ux9898ux673aux6838ux5fc3}}

\textbf{职责}：通过 \texttt{ProcessBuilder} 启动外部 C++ judge
二进制，传入判题参数，读取其标准输出的 JSON 行，用正则解析后构造
\texttt{JudgeResult} 返回。

\begin{Shaded}
\begin{Highlighting}[]
\KeywordTok{package}\ImportTok{ oj}\OperatorTok{.}\ImportTok{judge}\OperatorTok{;}

\KeywordTok{import} \ImportTok{oj}\OperatorTok{.}\ImportTok{core}\OperatorTok{.}\ImportTok{JudgeResult}\OperatorTok{;}
\KeywordTok{import} \ImportTok{oj}\OperatorTok{.}\ImportTok{core}\OperatorTok{.}\ImportTok{Status}\OperatorTok{;}

\KeywordTok{import} \ImportTok{java}\OperatorTok{.}\ImportTok{io}\OperatorTok{.}\ImportTok{BufferedReader}\OperatorTok{;}
\KeywordTok{import} \ImportTok{java}\OperatorTok{.}\ImportTok{io}\OperatorTok{.}\ImportTok{File}\OperatorTok{;}
\KeywordTok{import} \ImportTok{java}\OperatorTok{.}\ImportTok{io}\OperatorTok{.}\ImportTok{IOException}\OperatorTok{;}
\KeywordTok{import} \ImportTok{java}\OperatorTok{.}\ImportTok{io}\OperatorTok{.}\ImportTok{InputStreamReader}\OperatorTok{;}
\KeywordTok{import} \ImportTok{java}\OperatorTok{.}\ImportTok{util}\OperatorTok{.}\ImportTok{concurrent}\OperatorTok{.}\ImportTok{TimeUnit}\OperatorTok{;}
\KeywordTok{import} \ImportTok{java}\OperatorTok{.}\ImportTok{util}\OperatorTok{.}\ImportTok{regex}\OperatorTok{.}\ImportTok{Matcher}\OperatorTok{;}
\KeywordTok{import} \ImportTok{java}\OperatorTok{.}\ImportTok{util}\OperatorTok{.}\ImportTok{regex}\OperatorTok{.}\ImportTok{Pattern}\OperatorTok{;}

\KeywordTok{public} \KeywordTok{class}\NormalTok{ MachineJudge }\OperatorTok{\{}

    \CommentTok{/**}
     \CommentTok{*}\NormalTok{ C}\CommentTok{++}\NormalTok{ judge }\CommentTok{二进制输出的}\NormalTok{ JSON }\CommentTok{正则。}
     \CommentTok{*} \CommentTok{示例：\{"}\NormalTok{status}\CommentTok{":"}\NormalTok{AC}\CommentTok{","}\NormalTok{passed}\CommentTok{":3,"}\NormalTok{total}\CommentTok{":3,"}\NormalTok{time\_ms}\CommentTok{":45,"}\NormalTok{mem\_kb}\CommentTok{":2048,"}\NormalTok{detail}\CommentTok{":""\}}
     \CommentTok{*/}
    \KeywordTok{private} \DataTypeTok{static} \DataTypeTok{final} \BuiltInTok{Pattern}\NormalTok{ JSON\_PATTERN }\OperatorTok{=} \BuiltInTok{Pattern}\OperatorTok{.}\FunctionTok{compile}\OperatorTok{(}
        \StringTok{"}\SpecialCharTok{\textbackslash{}\textbackslash{}}\StringTok{\{}\SpecialCharTok{\textbackslash{}"}\StringTok{status}\SpecialCharTok{\textbackslash{}"}\StringTok{:}\SpecialCharTok{\textbackslash{}"}\StringTok{(}\SpecialCharTok{\textbackslash{}\textbackslash{}}\StringTok{w+)}\SpecialCharTok{\textbackslash{}"}\StringTok{,"} \OperatorTok{+}
        \StringTok{"}\SpecialCharTok{\textbackslash{}"}\StringTok{passed}\SpecialCharTok{\textbackslash{}"}\StringTok{:(}\SpecialCharTok{\textbackslash{}\textbackslash{}}\StringTok{d+),"} \OperatorTok{+}
        \StringTok{"}\SpecialCharTok{\textbackslash{}"}\StringTok{total}\SpecialCharTok{\textbackslash{}"}\StringTok{:(}\SpecialCharTok{\textbackslash{}\textbackslash{}}\StringTok{d+),"} \OperatorTok{+}
        \StringTok{"}\SpecialCharTok{\textbackslash{}"}\StringTok{time\_ms}\SpecialCharTok{\textbackslash{}"}\StringTok{:(}\SpecialCharTok{\textbackslash{}\textbackslash{}}\StringTok{d+),"} \OperatorTok{+}
        \StringTok{"}\SpecialCharTok{\textbackslash{}"}\StringTok{mem\_kb}\SpecialCharTok{\textbackslash{}"}\StringTok{:(}\SpecialCharTok{\textbackslash{}\textbackslash{}}\StringTok{d+),"} \OperatorTok{+}
        \StringTok{"}\SpecialCharTok{\textbackslash{}"}\StringTok{detail}\SpecialCharTok{\textbackslash{}"}\StringTok{:}\SpecialCharTok{\textbackslash{}"}\StringTok{([\^{}}\SpecialCharTok{\textbackslash{}"}\StringTok{]*)}\SpecialCharTok{\textbackslash{}"\textbackslash{}\textbackslash{}}\StringTok{\}"}
    \OperatorTok{);}

    \KeywordTok{private} \DataTypeTok{final} \BuiltInTok{String}\NormalTok{ judgeBinary}\OperatorTok{;}  \CommentTok{// C++ judge 二进制的绝对路径}

    \KeywordTok{public} \FunctionTok{MachineJudge}\OperatorTok{(}\BuiltInTok{String}\NormalTok{ judgeBinary}\OperatorTok{)} \OperatorTok{\{}
        \KeywordTok{this}\OperatorTok{.}\FunctionTok{judgeBinary} \OperatorTok{=}\NormalTok{ judgeBinary}\OperatorTok{;}
    \OperatorTok{\}}

    \CommentTok{/**}
     \CommentTok{*} \CommentTok{调用外部}\NormalTok{ C}\CommentTok{++} \CommentTok{判题机，返回判题结果。}
     \CommentTok{*}
\CommentTok{     * @}\NormalTok{param problemDir }\CommentTok{题目目录（含}\NormalTok{ N}\CommentTok{.}\NormalTok{in}\CommentTok{/}\NormalTok{N}\CommentTok{.}\NormalTok{out}\CommentTok{）}
     \CommentTok{*} \CommentTok{@}\NormalTok{param srcFile    }\CommentTok{选手源文件路径（.}\NormalTok{cpp }\CommentTok{/} \CommentTok{.}\NormalTok{py }\CommentTok{/} \CommentTok{.}\NormalTok{java}\CommentTok{）}
     \CommentTok{*} \CommentTok{@}\NormalTok{param lang       }\CommentTok{语言标识："}\NormalTok{cpp}\CommentTok{"} \CommentTok{/} \CommentTok{"}\NormalTok{python}\CommentTok{"} \CommentTok{/} \CommentTok{"}\NormalTok{java}\CommentTok{"}
     \CommentTok{*} \CommentTok{@}\NormalTok{param timeMs     }\CommentTok{时间限制（毫秒）}
     \CommentTok{*} \CommentTok{@}\NormalTok{param memMb      }\CommentTok{内存限制（}\NormalTok{MB}\CommentTok{）}
     \CommentTok{*} \CommentTok{@}\NormalTok{return JudgeResult}\CommentTok{，若}\NormalTok{ judge }\CommentTok{进程异常则返回}\NormalTok{ RE}
     \CommentTok{*/}
    \KeywordTok{public}\NormalTok{ JudgeResult }\FunctionTok{judge}\OperatorTok{(}\BuiltInTok{String}\NormalTok{ problemDir}\OperatorTok{,} \BuiltInTok{String}\NormalTok{ srcFile}\OperatorTok{,}
                             \BuiltInTok{String}\NormalTok{ lang}\OperatorTok{,} \DataTypeTok{long}\NormalTok{ timeMs}\OperatorTok{,} \DataTypeTok{int}\NormalTok{ memMb}\OperatorTok{)} \OperatorTok{\{}
        \BuiltInTok{ProcessBuilder}\NormalTok{ pb }\OperatorTok{=} \KeywordTok{new} \BuiltInTok{ProcessBuilder}\OperatorTok{(}
\NormalTok{            judgeBinary}\OperatorTok{,}
            \StringTok{"{-}{-}problem"}\OperatorTok{,}\NormalTok{ problemDir}\OperatorTok{,}
            \StringTok{"{-}{-}src"}\OperatorTok{,}\NormalTok{     srcFile}\OperatorTok{,}
            \StringTok{"{-}{-}lang"}\OperatorTok{,}\NormalTok{    lang}\OperatorTok{,}
            \StringTok{"{-}{-}time"}\OperatorTok{,}    \BuiltInTok{String}\OperatorTok{.}\FunctionTok{valueOf}\OperatorTok{(}\NormalTok{timeMs}\OperatorTok{),}
            \StringTok{"{-}{-}mem"}\OperatorTok{,}     \BuiltInTok{String}\OperatorTok{.}\FunctionTok{valueOf}\OperatorTok{(}\NormalTok{memMb}\OperatorTok{)}
        \OperatorTok{);}
\NormalTok{        pb}\OperatorTok{.}\FunctionTok{redirectErrorStream}\OperatorTok{(}\KeywordTok{false}\OperatorTok{);}  \CommentTok{// 不合并 stderr，只读 stdout}
\NormalTok{        pb}\OperatorTok{.}\FunctionTok{directory}\OperatorTok{(}\KeywordTok{new} \BuiltInTok{File}\OperatorTok{(}\StringTok{"."}\OperatorTok{));}

        \DataTypeTok{long}\NormalTok{ wallStart }\OperatorTok{=} \BuiltInTok{System}\OperatorTok{.}\FunctionTok{currentTimeMillis}\OperatorTok{();}
        \BuiltInTok{Process}\NormalTok{ proc}\OperatorTok{;}
        \ControlFlowTok{try} \OperatorTok{\{}
\NormalTok{            proc }\OperatorTok{=}\NormalTok{ pb}\OperatorTok{.}\FunctionTok{start}\OperatorTok{();}
        \OperatorTok{\}} \ControlFlowTok{catch} \OperatorTok{(}\BuiltInTok{IOException}\NormalTok{ e}\OperatorTok{)} \OperatorTok{\{}
            \ControlFlowTok{return} \FunctionTok{errorResult}\OperatorTok{(}\StringTok{"启动 judge 进程失败: "} \OperatorTok{+}\NormalTok{ e}\OperatorTok{.}\FunctionTok{getMessage}\OperatorTok{());}
        \OperatorTok{\}}

        \CommentTok{// 读取 judge 标准输出（预期一行 JSON）}
        \BuiltInTok{String}\NormalTok{ jsonLine }\OperatorTok{=} \StringTok{""}\OperatorTok{;}
        \ControlFlowTok{try} \OperatorTok{(}\BuiltInTok{BufferedReader}\NormalTok{ reader }\OperatorTok{=}
                 \KeywordTok{new} \BuiltInTok{BufferedReader}\OperatorTok{(}\KeywordTok{new} \BuiltInTok{InputStreamReader}\OperatorTok{(}\NormalTok{proc}\OperatorTok{.}\FunctionTok{getInputStream}\OperatorTok{())))} \OperatorTok{\{}
\NormalTok{            jsonLine }\OperatorTok{=}\NormalTok{ reader}\OperatorTok{.}\FunctionTok{readLine}\OperatorTok{();}
        \OperatorTok{\}} \ControlFlowTok{catch} \OperatorTok{(}\BuiltInTok{IOException}\NormalTok{ e}\OperatorTok{)} \OperatorTok{\{}
            \ControlFlowTok{return} \FunctionTok{errorResult}\OperatorTok{(}\StringTok{"读取 judge 输出失败: "} \OperatorTok{+}\NormalTok{ e}\OperatorTok{.}\FunctionTok{getMessage}\OperatorTok{());}
        \OperatorTok{\}}

        \CommentTok{// 等待进程退出，最多 timeMs + 5000ms 宽限}
        \DataTypeTok{boolean}\NormalTok{ finished}\OperatorTok{;}
        \ControlFlowTok{try} \OperatorTok{\{}
\NormalTok{            finished }\OperatorTok{=}\NormalTok{ proc}\OperatorTok{.}\FunctionTok{waitFor}\OperatorTok{(}\NormalTok{timeMs }\OperatorTok{+} \DecValTok{5000}\OperatorTok{,} \BuiltInTok{TimeUnit}\OperatorTok{.}\FunctionTok{MILLISECONDS}\OperatorTok{);}
        \OperatorTok{\}} \ControlFlowTok{catch} \OperatorTok{(}\BuiltInTok{InterruptedException}\NormalTok{ e}\OperatorTok{)} \OperatorTok{\{}
            \BuiltInTok{Thread}\OperatorTok{.}\FunctionTok{currentThread}\OperatorTok{().}\FunctionTok{interrupt}\OperatorTok{();}
\NormalTok{            proc}\OperatorTok{.}\FunctionTok{destroyForcibly}\OperatorTok{();}
            \ControlFlowTok{return} \FunctionTok{errorResult}\OperatorTok{(}\StringTok{"等待 judge 进程被中断"}\OperatorTok{);}
        \OperatorTok{\}}
        \ControlFlowTok{if} \OperatorTok{(!}\NormalTok{finished}\OperatorTok{)} \OperatorTok{\{}
\NormalTok{            proc}\OperatorTok{.}\FunctionTok{destroyForcibly}\OperatorTok{();}
            \ControlFlowTok{return} \FunctionTok{errorResult}\OperatorTok{(}\StringTok{"judge 进程超时未退出"}\OperatorTok{);}
        \OperatorTok{\}}

        \DataTypeTok{long}\NormalTok{ elapsed }\OperatorTok{=} \BuiltInTok{System}\OperatorTok{.}\FunctionTok{currentTimeMillis}\OperatorTok{()} \OperatorTok{{-}}\NormalTok{ wallStart}\OperatorTok{;}

        \ControlFlowTok{if} \OperatorTok{(}\NormalTok{jsonLine }\OperatorTok{==} \KeywordTok{null} \OperatorTok{||}\NormalTok{ jsonLine}\OperatorTok{.}\FunctionTok{isBlank}\OperatorTok{())} \OperatorTok{\{}
            \ControlFlowTok{return} \FunctionTok{errorResult}\OperatorTok{(}\StringTok{"judge 无输出"}\OperatorTok{);}
        \OperatorTok{\}}
        \ControlFlowTok{return} \FunctionTok{parseJson}\OperatorTok{(}\NormalTok{jsonLine}\OperatorTok{.}\FunctionTok{trim}\OperatorTok{(),}\NormalTok{ elapsed}\OperatorTok{);}
    \OperatorTok{\}}

    \CommentTok{/**}
     \CommentTok{*} \CommentTok{用正则解析}\NormalTok{ C}\CommentTok{++}\NormalTok{ judge }\CommentTok{输出的}\NormalTok{ JSON }\CommentTok{行，构造}\NormalTok{ JudgeResult}\CommentTok{。}
     \CommentTok{*}\NormalTok{ C}\CommentTok{++} \CommentTok{侧若报}\NormalTok{ ERR}\CommentTok{（沙箱内部错误），映射为}\NormalTok{ RE}\CommentTok{。}
     \CommentTok{*/}
    \KeywordTok{private}\NormalTok{ JudgeResult }\FunctionTok{parseJson}\OperatorTok{(}\BuiltInTok{String}\NormalTok{ json}\OperatorTok{,} \DataTypeTok{long}\NormalTok{ elapsed}\OperatorTok{)} \OperatorTok{\{}
        \BuiltInTok{Matcher}\NormalTok{ m }\OperatorTok{=}\NormalTok{ JSON\_PATTERN}\OperatorTok{.}\FunctionTok{matcher}\OperatorTok{(}\NormalTok{json}\OperatorTok{);}
        \ControlFlowTok{if} \OperatorTok{(!}\NormalTok{m}\OperatorTok{.}\FunctionTok{find}\OperatorTok{())} \OperatorTok{\{}
            \ControlFlowTok{return} \FunctionTok{errorResult}\OperatorTok{(}\StringTok{"JSON 格式不匹配: "} \OperatorTok{+}\NormalTok{ json}\OperatorTok{);}
        \OperatorTok{\}}

        \BuiltInTok{String}\NormalTok{ rawStatus }\OperatorTok{=}\NormalTok{ m}\OperatorTok{.}\FunctionTok{group}\OperatorTok{(}\DecValTok{1}\OperatorTok{);}
        \DataTypeTok{int}\NormalTok{ passed  }\OperatorTok{=} \BuiltInTok{Integer}\OperatorTok{.}\FunctionTok{parseInt}\OperatorTok{(}\NormalTok{m}\OperatorTok{.}\FunctionTok{group}\OperatorTok{(}\DecValTok{2}\OperatorTok{));}
        \DataTypeTok{int}\NormalTok{ total   }\OperatorTok{=} \BuiltInTok{Integer}\OperatorTok{.}\FunctionTok{parseInt}\OperatorTok{(}\NormalTok{m}\OperatorTok{.}\FunctionTok{group}\OperatorTok{(}\DecValTok{3}\OperatorTok{));}
        \DataTypeTok{long}\NormalTok{ timeMs }\OperatorTok{=} \BuiltInTok{Long}\OperatorTok{.}\FunctionTok{parseLong}\OperatorTok{(}\NormalTok{m}\OperatorTok{.}\FunctionTok{group}\OperatorTok{(}\DecValTok{4}\OperatorTok{));}
        \CommentTok{// mem\_kb 目前记录在 detail，供日志使用}
        \DataTypeTok{int}\NormalTok{ memKb   }\OperatorTok{=} \BuiltInTok{Integer}\OperatorTok{.}\FunctionTok{parseInt}\OperatorTok{(}\NormalTok{m}\OperatorTok{.}\FunctionTok{group}\OperatorTok{(}\DecValTok{5}\OperatorTok{));}
        \BuiltInTok{String}\NormalTok{ detail }\OperatorTok{=}\NormalTok{ m}\OperatorTok{.}\FunctionTok{group}\OperatorTok{(}\DecValTok{6}\OperatorTok{);}

\NormalTok{        Status status }\OperatorTok{=} \FunctionTok{mapStatus}\OperatorTok{(}\NormalTok{rawStatus}\OperatorTok{);}
        \BuiltInTok{String}\NormalTok{ fullDetail }\OperatorTok{=}\NormalTok{ detail}\OperatorTok{.}\FunctionTok{isBlank}\OperatorTok{()}
            \OperatorTok{?} \BuiltInTok{String}\OperatorTok{.}\FunctionTok{format}\OperatorTok{(}\StringTok{"mem=}\SpecialCharTok{\%d}\StringTok{KB"}\OperatorTok{,}\NormalTok{ memKb}\OperatorTok{)}
            \OperatorTok{:} \BuiltInTok{String}\OperatorTok{.}\FunctionTok{format}\OperatorTok{(}\StringTok{"}\SpecialCharTok{\%s}\StringTok{ mem=}\SpecialCharTok{\%d}\StringTok{KB"}\OperatorTok{,}\NormalTok{ detail}\OperatorTok{,}\NormalTok{ memKb}\OperatorTok{);}

        \ControlFlowTok{return} \KeywordTok{new} \FunctionTok{JudgeResult}\OperatorTok{(}\NormalTok{status}\OperatorTok{,}\NormalTok{ passed}\OperatorTok{,}\NormalTok{ total}\OperatorTok{,}\NormalTok{ fullDetail}\OperatorTok{,}\NormalTok{ timeMs}\OperatorTok{);}
    \OperatorTok{\}}

    \CommentTok{/**}
     \CommentTok{*} \CommentTok{将}\NormalTok{ C}\CommentTok{++}\NormalTok{ judge }\CommentTok{返回的状态字符串映射到}\NormalTok{ oj}\CommentTok{.}\NormalTok{core}\CommentTok{.}\NormalTok{Status}\CommentTok{。}
     \CommentTok{*}\NormalTok{ ERR}\CommentTok{（沙箱内部异常）统一映射为}\NormalTok{ RE}\CommentTok{。}
     \CommentTok{*/}
    \KeywordTok{private}\NormalTok{ Status }\FunctionTok{mapStatus}\OperatorTok{(}\BuiltInTok{String}\NormalTok{ raw}\OperatorTok{)} \OperatorTok{\{}
        \ControlFlowTok{switch} \OperatorTok{(}\NormalTok{raw}\OperatorTok{.}\FunctionTok{toUpperCase}\OperatorTok{())} \OperatorTok{\{}
            \ControlFlowTok{case} \StringTok{"AC"}\OperatorTok{:}  \ControlFlowTok{return}\NormalTok{ Status}\OperatorTok{.}\FunctionTok{AC}\OperatorTok{;}
            \ControlFlowTok{case} \StringTok{"WA"}\OperatorTok{:}  \ControlFlowTok{return}\NormalTok{ Status}\OperatorTok{.}\FunctionTok{WA}\OperatorTok{;}
            \ControlFlowTok{case} \StringTok{"TLE"}\OperatorTok{:} \ControlFlowTok{return}\NormalTok{ Status}\OperatorTok{.}\FunctionTok{TLE}\OperatorTok{;}
            \ControlFlowTok{case} \StringTok{"MLE"}\OperatorTok{:} \ControlFlowTok{return}\NormalTok{ Status}\OperatorTok{.}\FunctionTok{MLE}\OperatorTok{;}
            \ControlFlowTok{case} \StringTok{"RE"}\OperatorTok{:}  \ControlFlowTok{return}\NormalTok{ Status}\OperatorTok{.}\FunctionTok{RE}\OperatorTok{;}
            \ControlFlowTok{case} \StringTok{"CE"}\OperatorTok{:}  \ControlFlowTok{return}\NormalTok{ Status}\OperatorTok{.}\FunctionTok{CE}\OperatorTok{;}
            \ControlFlowTok{case} \StringTok{"PE"}\OperatorTok{:}  \ControlFlowTok{return}\NormalTok{ Status}\OperatorTok{.}\FunctionTok{PE}\OperatorTok{;}
            \ControlFlowTok{case} \StringTok{"ERR"}\OperatorTok{:} \ControlFlowTok{return}\NormalTok{ Status}\OperatorTok{.}\FunctionTok{RE}\OperatorTok{;}   \CommentTok{// 沙箱内部错误 {-}\textgreater{} RE}
            \KeywordTok{default}\OperatorTok{:}    \ControlFlowTok{return}\NormalTok{ Status}\OperatorTok{.}\FunctionTok{RE}\OperatorTok{;}
        \OperatorTok{\}}
    \OperatorTok{\}}

    \CommentTok{/**} \CommentTok{构造一个表示}\NormalTok{ Java }\CommentTok{侧调用失败（}\NormalTok{RE}\CommentTok{）的}\NormalTok{ JudgeResult }\CommentTok{*/}
    \KeywordTok{private}\NormalTok{ JudgeResult }\FunctionTok{errorResult}\OperatorTok{(}\BuiltInTok{String}\NormalTok{ detail}\OperatorTok{)} \OperatorTok{\{}
        \ControlFlowTok{return} \KeywordTok{new} \FunctionTok{JudgeResult}\OperatorTok{(}\NormalTok{Status}\OperatorTok{.}\FunctionTok{RE}\OperatorTok{,} \DecValTok{0}\OperatorTok{,} \DecValTok{0}\OperatorTok{,}\NormalTok{ detail}\OperatorTok{,} \DecValTok{0L}\OperatorTok{);}
    \OperatorTok{\}}
\OperatorTok{\}}
\end{Highlighting}
\end{Shaded}

\textbf{为什么换 C++？}

\begin{longtable}[]{@{}
  >{\raggedright\arraybackslash}p{(\columnwidth - 4\tabcolsep) * \real{0.3333}}
  >{\raggedright\arraybackslash}p{(\columnwidth - 4\tabcolsep) * \real{0.3333}}
  >{\raggedright\arraybackslash}p{(\columnwidth - 4\tabcolsep) * \real{0.3333}}@{}}
\toprule\noalign{}
\begin{minipage}[b]{\linewidth}\raggedright
能力
\end{minipage} & \begin{minipage}[b]{\linewidth}\raggedright
Java（第一/二代）
\end{minipage} & \begin{minipage}[b]{\linewidth}\raggedright
C++ judge（第三代）
\end{minipage} \\
\midrule\noalign{}
\endhead
\bottomrule\noalign{}
\endlastfoot
CPU 时间限制 & 不可靠（Thread.interrupt） &
\texttt{setrlimit(RLIMIT\_CPU,\ t)} 内核强制 \\
内存限制 & 不可靠（JVM 堆外无法控制） &
\texttt{setrlimit(RLIMIT\_AS,\ m)} 地址空间硬限 \\
多语言编译 & 仅 Java & 调
\texttt{g++}/\texttt{python3}/\texttt{javac} \\
隔离性 & 无 & chroot / seccomp（可扩展） \\
性能 & JVM 启动 \textasciitilde200ms & native 进程 \textasciitilde5ms \\
\end{longtable}

\begin{center}\rule{0.5\linewidth}{0.5pt}\end{center}

\hypertarget{stream-ux7edfux8ba1ux6f14ux793aux6574ux5408ux793aux4f8b}{%
\subsubsection{5. Stream
统计演示（整合示例）}\label{stream-ux7edfux8ba1ux6f14ux793aux6574ux5408ux793aux4f8b}}

以下代码展示如何用 Stream 对一批提交做状态统计，通常放在 \texttt{Main}
或测试类中：

\begin{Shaded}
\begin{Highlighting}[]
\KeywordTok{import} \ImportTok{oj}\OperatorTok{.}\ImportTok{core}\OperatorTok{.}\ImportTok{Status}\OperatorTok{;}
\KeywordTok{import} \ImportTok{oj}\OperatorTok{.}\ImportTok{core}\OperatorTok{.}\ImportTok{Submission}\OperatorTok{;}
\KeywordTok{import} \ImportTok{oj}\OperatorTok{.}\ImportTok{io}\OperatorTok{.}\ImportTok{SubmissionStore}\OperatorTok{;}

\KeywordTok{import} \ImportTok{java}\OperatorTok{.}\ImportTok{util}\OperatorTok{.}\ImportTok{List}\OperatorTok{;}
\KeywordTok{import} \ImportTok{java}\OperatorTok{.}\ImportTok{util}\OperatorTok{.}\ImportTok{Map}\OperatorTok{;}
\KeywordTok{import} \ImportTok{java}\OperatorTok{.}\ImportTok{util}\OperatorTok{.}\ImportTok{stream}\OperatorTok{.}\ImportTok{Collectors}\OperatorTok{;}

\KeywordTok{public} \KeywordTok{class}\NormalTok{ StatsDemo }\OperatorTok{\{}
    \KeywordTok{public} \DataTypeTok{static} \DataTypeTok{void} \FunctionTok{main}\OperatorTok{(}\BuiltInTok{String}\OperatorTok{[]}\NormalTok{ args}\OperatorTok{)} \KeywordTok{throws} \BuiltInTok{Exception} \OperatorTok{\{}
\NormalTok{        SubmissionStore store }\OperatorTok{=} \KeywordTok{new} \FunctionTok{SubmissionStore}\OperatorTok{(}\StringTok{"submissions.dat"}\OperatorTok{);}
        \BuiltInTok{List}\OperatorTok{\textless{}}\NormalTok{Submission}\OperatorTok{\textgreater{}}\NormalTok{ all }\OperatorTok{=}\NormalTok{ store}\OperatorTok{.}\FunctionTok{loadAll}\OperatorTok{();}

        \CommentTok{// Stream：按 Status 分组计数}
        \BuiltInTok{Map}\OperatorTok{\textless{}}\NormalTok{Status}\OperatorTok{,} \BuiltInTok{Long}\OperatorTok{\textgreater{}}\NormalTok{ stat }\OperatorTok{=}\NormalTok{ all}\OperatorTok{.}\FunctionTok{stream}\OperatorTok{()}
            \OperatorTok{.}\FunctionTok{map}\OperatorTok{(}\NormalTok{s }\OperatorTok{{-}\textgreater{}}\NormalTok{ s}\OperatorTok{.}\FunctionTok{getResult}\OperatorTok{().}\FunctionTok{getStatus}\OperatorTok{())}
            \OperatorTok{.}\FunctionTok{collect}\OperatorTok{(}\NormalTok{Collectors}\OperatorTok{.}\FunctionTok{groupingBy}\OperatorTok{(}\NormalTok{s }\OperatorTok{{-}\textgreater{}}\NormalTok{ s}\OperatorTok{,}\NormalTok{ Collectors}\OperatorTok{.}\FunctionTok{counting}\OperatorTok{()));}

\NormalTok{        stat}\OperatorTok{.}\FunctionTok{forEach}\OperatorTok{((}\NormalTok{status}\OperatorTok{,}\NormalTok{ count}\OperatorTok{)} \OperatorTok{{-}\textgreater{}}
            \BuiltInTok{System}\OperatorTok{.}\FunctionTok{out}\OperatorTok{.}\FunctionTok{printf}\OperatorTok{(}\StringTok{"}\SpecialCharTok{\%{-}4s}\StringTok{ : }\SpecialCharTok{\%d}\StringTok{ 次}\SpecialCharTok{\%n}\StringTok{"}\OperatorTok{,}\NormalTok{ status}\OperatorTok{,}\NormalTok{ count}\OperatorTok{));}

        \CommentTok{// Stream：筛选 AC 提交，打印用户名}
        \BuiltInTok{System}\OperatorTok{.}\FunctionTok{out}\OperatorTok{.}\FunctionTok{println}\OperatorTok{(}\StringTok{"{-}{-}{-} AC 提交者 {-}{-}{-}"}\OperatorTok{);}
\NormalTok{        all}\OperatorTok{.}\FunctionTok{stream}\OperatorTok{()}
            \OperatorTok{.}\FunctionTok{filter}\OperatorTok{(}\NormalTok{s }\OperatorTok{{-}\textgreater{}}\NormalTok{ s}\OperatorTok{.}\FunctionTok{getResult}\OperatorTok{().}\FunctionTok{isAccepted}\OperatorTok{())}
            \OperatorTok{.}\FunctionTok{map}\OperatorTok{(}\NormalTok{Submission}\OperatorTok{::}\NormalTok{getUserName}\OperatorTok{)}
            \OperatorTok{.}\FunctionTok{distinct}\OperatorTok{()}
            \OperatorTok{.}\FunctionTok{sorted}\OperatorTok{()}
            \OperatorTok{.}\FunctionTok{forEach}\OperatorTok{(}\BuiltInTok{System}\OperatorTok{.}\FunctionTok{out}\OperatorTok{::}\NormalTok{println}\OperatorTok{);}
    \OperatorTok{\}}
\OperatorTok{\}}
\end{Highlighting}
\end{Shaded}

\begin{center}\rule{0.5\linewidth}{0.5pt}\end{center}

\hypertarget{ux7c7bux5173ux7cfbux603bux89c8}{%
\subsubsection{类关系总览}\label{ux7c7bux5173ux7cfbux603bux89c8}}

\begin{verbatim}
ProblemRepository
  +- uses ConfigFile          (M4 产物，读 config.txt)
  +- uses ProblemLoader
           +- returns List<TestCase>   <- Ch15 泛型集合

MachineJudge
  +- ProcessBuilder -> [C++ judge binary]
  +- Pattern.compile -> JudgeResult

SubmissionStore
  +- ObjectOutputStream (写)  <- Ch10 序列化
  +- ObjectInputStream  (读)
\end{verbatim}

\begin{center}\rule{0.5\linewidth}{0.5pt}\end{center}

\hypertarget{ux52a8ux624bux5b9eux73b0ux4e24ux6b65ux8d70javac-java-ux5b9eux64cd-3}{%
\subsection{动手实现:两步走(javac / java
实操)}\label{ux52a8ux624bux5b9eux73b0ux4e24ux6b65ux8d70javac-java-ux5b9eux64cd-3}}

\hypertarget{ux7b2cux4e00ux6b65-ux7528ux96c6ux5408ux628a-oj-ux9898ux5e93ux88c5ux8fdbux5185ux5b58}{%
\subsubsection{\texorpdfstring{第一步 \KS --- 用集合把 OJ
题库装进内存}{第一步 --- 用集合把 OJ 题库装进内存}}\label{ux7b2cux4e00ux6b65-ux7528ux96c6ux5408ux628a-oj-ux9898ux5e93ux88c5ux8fdbux5185ux5b58}}

\textbf{\texttt{src/oj/io/Bank.java}}

\begin{Shaded}
\begin{Highlighting}[]
\KeywordTok{package}\ImportTok{ oj}\OperatorTok{.}\ImportTok{io}\OperatorTok{;}
\KeywordTok{import} \ImportTok{oj}\OperatorTok{.}\ImportTok{core}\OperatorTok{.*;}
\KeywordTok{import} \ImportTok{java}\OperatorTok{.}\ImportTok{util}\OperatorTok{.*;}
\KeywordTok{public} \KeywordTok{class}\NormalTok{ Bank }\OperatorTok{\{}
    \KeywordTok{private} \BuiltInTok{Map}\OperatorTok{\textless{}}\BuiltInTok{Integer}\OperatorTok{,}\NormalTok{ Problem}\OperatorTok{\textgreater{}}\NormalTok{ problems }\OperatorTok{=} \KeywordTok{new} \BuiltInTok{HashMap}\OperatorTok{\textless{}\textgreater{}();}   \CommentTok{// 题号 {-}\textgreater{} 题目}
    \KeywordTok{public} \DataTypeTok{void} \FunctionTok{add}\OperatorTok{(}\NormalTok{Problem p}\OperatorTok{)}   \OperatorTok{\{}\NormalTok{ problems}\OperatorTok{.}\FunctionTok{put}\OperatorTok{(}\NormalTok{p}\OperatorTok{.}\FunctionTok{getId}\OperatorTok{(),}\NormalTok{ p}\OperatorTok{);} \OperatorTok{\}}
    \KeywordTok{public}\NormalTok{ Problem }\FunctionTok{get}\OperatorTok{(}\DataTypeTok{int}\NormalTok{ id}\OperatorTok{)}   \OperatorTok{\{} \ControlFlowTok{return}\NormalTok{ problems}\OperatorTok{.}\FunctionTok{get}\OperatorTok{(}\NormalTok{id}\OperatorTok{);} \OperatorTok{\}}
    \KeywordTok{public} \DataTypeTok{int} \FunctionTok{size}\OperatorTok{()}            \OperatorTok{\{} \ControlFlowTok{return}\NormalTok{ problems}\OperatorTok{.}\FunctionTok{size}\OperatorTok{();} \OperatorTok{\}}
    \KeywordTok{public} \BuiltInTok{List}\OperatorTok{\textless{}}\BuiltInTok{Integer}\OperatorTok{\textgreater{}} \FunctionTok{ids}\OperatorTok{()}   \OperatorTok{\{} \ControlFlowTok{return} \KeywordTok{new} \BuiltInTok{ArrayList}\OperatorTok{\textless{}\textgreater{}(}\NormalTok{problems}\OperatorTok{.}\FunctionTok{keySet}\OperatorTok{());} \OperatorTok{\}}
\OperatorTok{\}}
\end{Highlighting}
\end{Shaded}

讲解:\texttt{HashMap\textless{}Integer,Problem\textgreater{}}
做题号映射、\texttt{ArrayList} 列题号------Ch15 集合考点,用在 OJ
题库管理上。

\textbf{\texttt{src/oj/Demo5a.java}}

\begin{Shaded}
\begin{Highlighting}[]
\KeywordTok{package}\ImportTok{ oj}\OperatorTok{;}
\KeywordTok{import} \ImportTok{oj}\OperatorTok{.}\ImportTok{core}\OperatorTok{.*;}
\KeywordTok{import} \ImportTok{oj}\OperatorTok{.}\ImportTok{io}\OperatorTok{.}\ImportTok{Bank}\OperatorTok{;}
\KeywordTok{import} \ImportTok{java}\OperatorTok{.}\ImportTok{util}\OperatorTok{.*;}
\KeywordTok{public} \KeywordTok{class}\NormalTok{ Demo5a }\OperatorTok{\{}
    \KeywordTok{public} \DataTypeTok{static} \DataTypeTok{void} \FunctionTok{main}\OperatorTok{(}\BuiltInTok{String}\OperatorTok{[]}\NormalTok{ args}\OperatorTok{)} \OperatorTok{\{}
\NormalTok{        Bank bank }\OperatorTok{=} \KeywordTok{new} \FunctionTok{Bank}\OperatorTok{();}
\NormalTok{        bank}\OperatorTok{.}\FunctionTok{add}\OperatorTok{(}\KeywordTok{new} \FunctionTok{Problem}\OperatorTok{(}\DecValTok{1}\OperatorTok{,} \StringTok{"A+B"}\OperatorTok{,} \KeywordTok{new}\NormalTok{ TestCase}\OperatorTok{[]\{} \KeywordTok{new} \FunctionTok{TestCase}\OperatorTok{(}\StringTok{"1 2"}\OperatorTok{,}\StringTok{"3"}\OperatorTok{)} \OperatorTok{\}));}
\NormalTok{        bank}\OperatorTok{.}\FunctionTok{add}\OperatorTok{(}\KeywordTok{new} \FunctionTok{Problem}\OperatorTok{(}\DecValTok{2}\OperatorTok{,} \StringTok{"平方"}\OperatorTok{,} \KeywordTok{new}\NormalTok{ TestCase}\OperatorTok{[]\{} \KeywordTok{new} \FunctionTok{TestCase}\OperatorTok{(}\StringTok{"3"}\OperatorTok{,}\StringTok{"9"}\OperatorTok{),} \KeywordTok{new} \FunctionTok{TestCase}\OperatorTok{(}\StringTok{"4"}\OperatorTok{,}\StringTok{"16"}\OperatorTok{)} \OperatorTok{\}));}
        \ControlFlowTok{for} \OperatorTok{(}\DataTypeTok{int}\NormalTok{ id }\OperatorTok{:}\NormalTok{ bank}\OperatorTok{.}\FunctionTok{ids}\OperatorTok{())}                                  \CommentTok{// 遍历(类似 Iterator)}
            \BuiltInTok{System}\OperatorTok{.}\FunctionTok{out}\OperatorTok{.}\FunctionTok{println}\OperatorTok{(}\NormalTok{id }\OperatorTok{+} \StringTok{" =\textgreater{} "} \OperatorTok{+}\NormalTok{ bank}\OperatorTok{.}\FunctionTok{get}\OperatorTok{(}\NormalTok{id}\OperatorTok{).}\FunctionTok{getTitle}\OperatorTok{()} \OperatorTok{+} \StringTok{" ("} \OperatorTok{+}\NormalTok{ bank}\OperatorTok{.}\FunctionTok{get}\OperatorTok{(}\NormalTok{id}\OperatorTok{).}\FunctionTok{getCases}\OperatorTok{().}\FunctionTok{length} \OperatorTok{+} \StringTok{" 用例)"}\OperatorTok{);}
        \DataTypeTok{long}\NormalTok{ total }\OperatorTok{=}\NormalTok{ bank}\OperatorTok{.}\FunctionTok{ids}\OperatorTok{().}\FunctionTok{stream}\OperatorTok{().}\FunctionTok{mapToInt}\OperatorTok{(}\NormalTok{id }\OperatorTok{{-}\textgreater{}}\NormalTok{ bank}\OperatorTok{.}\FunctionTok{get}\OperatorTok{(}\NormalTok{id}\OperatorTok{).}\FunctionTok{getCases}\OperatorTok{().}\FunctionTok{length}\OperatorTok{).}\FunctionTok{sum}\OperatorTok{();}  \CommentTok{// Stream 统计}
        \BuiltInTok{System}\OperatorTok{.}\FunctionTok{out}\OperatorTok{.}\FunctionTok{println}\OperatorTok{(}\StringTok{"总用例数: "} \OperatorTok{+}\NormalTok{ total}\OperatorTok{);}
    \OperatorTok{\}}
\OperatorTok{\}}
\end{Highlighting}
\end{Shaded}

编译运行:

\begin{Shaded}
\begin{Highlighting}[]
\ExtensionTok{javac} \AttributeTok{{-}d}\NormalTok{ build src/oj/core/}\PreprocessorTok{*}\NormalTok{.java src/oj/io/Bank.java src/oj/Demo5a.java}
\ExtensionTok{java} \AttributeTok{{-}cp}\NormalTok{ build oj.Demo5a}
\CommentTok{\# 1 =\textgreater{} A+B (1 用例)}
\CommentTok{\# 2 =\textgreater{} 平方 (2 用例)}
\CommentTok{\# 总用例数: 3}
\end{Highlighting}
\end{Shaded}

\hypertarget{ux7b2cux4e8cux6b65-ux626bux76eeux5f55ux9898ux5e93-ux5916ux90e8-c-ux5224ux9898ux673a}{%
\subsubsection{\texorpdfstring{第二步 \GC --- 扫目录题库 + 外部 C++
判题机}{第二步 --- 扫目录题库 + 外部 C++ 判题机}}\label{ux7b2cux4e8cux6b65-ux626bux76eeux5f55ux9898ux5e93-ux5916ux90e8-c-ux5224ux9898ux673a}}

\texttt{ProblemLoader.cases(id)}(扫
\texttt{problems/\textless{}id\textgreater{}/N.in/N.out} -\textgreater{}
\texttt{List\textless{}TestCase\textgreater{}})、\texttt{ProblemRepository}(HashMap
缓存)、\texttt{MachineJudge}(\texttt{ProcessBuilder}
调判题机,代码见上文【新产物架构】)。先编译判题机,再跑 Java:

\begin{Shaded}
\begin{Highlighting}[]
\BuiltInTok{cd}\NormalTok{ judge }\KeywordTok{\&\&} \ExtensionTok{g++} \AttributeTok{{-}O2} \AttributeTok{{-}std}\OperatorTok{=}\NormalTok{c++17 }\AttributeTok{{-}o}\NormalTok{ judge judge.cpp }\KeywordTok{\&\&} \BuiltInTok{cd}\NormalTok{ ..}
\ExtensionTok{javac} \AttributeTok{{-}d}\NormalTok{ build }\VariableTok{$(}\FunctionTok{find}\NormalTok{ src }\AttributeTok{{-}name} \StringTok{\textquotesingle{}*.java\textquotesingle{}}\VariableTok{)}
\ExtensionTok{java} \AttributeTok{{-}cp} \StringTok{"build:lib/*"}\NormalTok{ oj.Main}
\end{Highlighting}
\end{Shaded}

\hypertarget{ux9a8cux6536ux6807ux51c6-5}{%
\subsection{验收标准}\label{ux9a8cux6536ux6807ux51c6-5}}

\begin{enumerate}
\def\labelenumi{\arabic{enumi}.}
\item
  \textbf{\texttt{ProblemLoader.cases(1001)}} 能从
  \texttt{problems/1001/} 扫到所有 \texttt{N.in}/\texttt{N.out}
  对，按编号升序返回
  \texttt{List\textless{}TestCase\textgreater{}}，编号不连续（如缺少
  \texttt{2.in}）时跳过不报错。
\item
  \textbf{\texttt{ProblemRepository.get(1001)}}
  第一次调用触发磁盘读取，第二次调用命中缓存（\texttt{cacheSize()}
  不增加），两次返回同一个对象引用。
\item
  \textbf{\texttt{SubmissionStore}} 连续 \texttt{save()}
  三条提交后，\texttt{loadAll()} 能完整读回三条，顺序一致，无
  \texttt{StreamCorruptedException}。
\item
  \textbf{\texttt{MachineJudge.judge()}} 在 judge 二进制返回合法 JSON
  时正确构造 \texttt{JudgeResult}；返回 \texttt{ERR} 状态时映射为
  \texttt{RE}；judge 进程超时未退出时强制销毁并返回 \texttt{RE}。
\item
  \textbf{Stream 统计}：对含 AC/WA/TLE 混合提交的
  \texttt{SubmissionStore} 执行 \texttt{StatsDemo}，输出各 Status
  计数正确，AC 用户名列表去重且排序。
\item
  \textbf{不出现任何 Socket/ServerSocket/网络相关代码}，不出现第一代
  \texttt{Solution} 接口的调用。
\end{enumerate}

\hypertarget{m5b-ux6570ux636eux5e93ux8854ux63a5ch11-jdbc-mysql}{%
\section{07 · M5b 数据库衔接（Ch11 JDBC +
MySQL）}\label{m5b-ux6570ux636eux5e93ux8854ux63a5ch11-jdbc-mysql}}

\begin{quote}
\textbf{考试相关度 ★★} ·
真题考点参考:\texttt{Connection/Statement/ResultSet} 查询入
List、更新。\textbf{{[}第一步{]}} \KS 用最简 JDBC(\texttt{Statement})把
OJ 的提交/题目存取 MySQL;\textbf{{[}第二步{]}}
\GC \texttt{ProblemDao}/事务/\texttt{PreparedStatement}/\texttt{ProblemService}
的 FS/DB 分工。
\end{quote}

\begin{quote}
文件能存代码、能存测试点，但它不会帮你在一秒内告诉你''过了 AC
的选手有几个''。当题目数量破百、提交记录过万，纯文件方案的代价才真正暴露出来。
\end{quote}

\hypertarget{ux6838ux5fc3ux75dbux70b9-2}{%
\subsection{【核心痛点】}\label{ux6838ux5fc3ux75dbux70b9-2}}

M5a 的 \texttt{ProblemRepository} 把全部
\texttt{ProblemMeta}（包括判题器类名、时间限制）缓存在内存的
\texttt{HashMap} 里，启动时从文件系统扫
\texttt{problems/\textless{}id\textgreater{}/config.txt}
逐一解析。\texttt{SubmissionStore} 则把 \texttt{Submission} 对象序列化进
\texttt{.ser} 文件，一个文件一条提交记录。

这两个方案在教学原型里够用，但三个场景会让它原形毕露：

\begin{enumerate}
\def\labelenumi{\arabic{enumi}.}
\tightlist
\item
  \textbf{统计}：想查某题通过率、某用户 AC 数，要把所有 \texttt{.ser}
  文件全部反序列化再在内存里过滤------O(N) 读 I/O，N 一大就超时。
\item
  \textbf{并发安全}：多线程判题队列同时往同一目录写
  \texttt{.ser}，没有任何互斥保障，文件名冲突或写到一半被读的情况真实发生。
\item
  \textbf{元数据变更}：改一道题的时间限制要手动编辑
  \texttt{config.txt}，没有原子性，改到一半系统崩溃就留下半截文件。
\end{enumerate}

本节引入 MySQL + JDBC，让 \textbf{元数据与历史记录进库、测试点文件留
FS}，形成 FS/DB 分工格局，同时完整实践 Ch11 的 JDBC 知识点。

\hypertarget{ux5f15ux5165ux8bfeux672cux77e5ux8bc6ux70b9-2}{%
\subsection{【引入课本知识点】}\label{ux5f15ux5165ux8bfeux672cux77e5ux8bc6ux70b9-2}}

\hypertarget{jdbc-ux56dbux6b65ux56faux5b9aux5957ux8def}{%
\subsubsection{JDBC
四步固定套路}\label{jdbc-ux56dbux6b65ux56faux5b9aux5957ux8def}}

\begin{verbatim}
DriverManager.getConnection(url, user, pwd)   // 建连
conn.prepareStatement(sql)                    // 预编译
ps.setXxx(index, value)                       // 绑参
ps.executeQuery() / ps.executeUpdate()        // 执行
\end{verbatim}

\texttt{PreparedStatement} 的核心价值是\textbf{防 SQL 注入}：参数占位符
\texttt{?} 由驱动负责转义，用户输入里的单引号、注释符永远不会被解释成
SQL 语法。

\hypertarget{resultset-ux6e38ux6807}{%
\subsubsection{ResultSet 游标}\label{resultset-ux6e38ux6807}}

\texttt{rs.next()} 把游标从''结果集首行之前''推进一行，返回
\texttt{false} 表示遍历完毕。列索引从 1 开始，或用列名
\texttt{rs.getString("title")}。

\hypertarget{ux4e8bux52a1}{%
\subsubsection{事务}\label{ux4e8bux52a1}}

MySQL 默认
\texttt{AUTO\_COMMIT=true}，每条语句独立提交。需要原子操作时：

\begin{Shaded}
\begin{Highlighting}[]
\NormalTok{conn}\OperatorTok{.}\FunctionTok{setAutoCommit}\OperatorTok{(}\KeywordTok{false}\OperatorTok{);}
\CommentTok{// ... 多条 SQL ...}
\NormalTok{conn}\OperatorTok{.}\FunctionTok{commit}\OperatorTok{();}          \CommentTok{// 全成功}
\CommentTok{// catch {-}\textgreater{} conn.rollback();}
\end{Highlighting}
\end{Shaded}

\texttt{SubmissionDao.save} 在单次提交里写两张表（\texttt{submissions}
插入 + 可选计数更新），必须在事务里完成，否则中途崩溃会留下孤儿记录。

\hypertarget{ux8fdeux63a5ux590dux7528}{%
\subsubsection{连接复用}\label{ux8fdeux63a5ux590dux7528}}

每次 \texttt{getConnection} 都要经历 TCP 握手 + 认证，开销约 10--30
ms。教学阶段用\textbf{单例连接}（\texttt{Db} 持有一个静态
\texttt{Connection}），生产级做法是连接池（HikariCP
等），了解即可，本节不实现。

\hypertarget{ux4e09ux4ee3ux6f14ux8fdbux5b9aux4f4d-2}{%
\subsection{【三代演进定位】}\label{ux4e09ux4ee3ux6f14ux8fdbux5b9aux4f4d-2}}

本节 \textbf{不属于判题机三代的演进节点}，而是为\textbf{第三代
MachineJudge 提供配置来源}：

\begin{longtable}[]{@{}
  >{\raggedright\arraybackslash}p{(\columnwidth - 4\tabcolsep) * \real{0.3333}}
  >{\raggedright\arraybackslash}p{(\columnwidth - 4\tabcolsep) * \real{0.3333}}
  >{\raggedright\arraybackslash}p{(\columnwidth - 4\tabcolsep) * \real{0.3333}}@{}}
\toprule\noalign{}
\begin{minipage}[b]{\linewidth}\raggedright
数据
\end{minipage} & \begin{minipage}[b]{\linewidth}\raggedright
来源（M5a 之前）
\end{minipage} & \begin{minipage}[b]{\linewidth}\raggedright
来源（M5b 之后）
\end{minipage} \\
\midrule\noalign{}
\endhead
\bottomrule\noalign{}
\endlastfoot
\texttt{timeLimitMs} / \texttt{memLimitMb} & \texttt{config.txt}
-\textgreater{} \texttt{ConfigFile.read} & \texttt{problems} 表
-\textgreater{} \texttt{ProblemDao.meta} \\
判题器类名 \texttt{judgeClass} & \texttt{config.txt} & \texttt{problems}
表 \\
提交历史 & \texttt{.ser} 文件 -\textgreater{} \texttt{SubmissionStore} &
\texttt{submissions} 表 -\textgreater{} \texttt{SubmissionDao.save} \\
测试点 \texttt{.in} / \texttt{.out} & 文件系统（保留） &
文件系统（保留，不变） \\
\end{longtable}

\texttt{MachineJudge} 调用外部 C++ 判题机时传入的
\texttt{timeMs}、\texttt{memMb} 由 \texttt{ProblemService.load} 从 DB
取出，路径不变、语义不变，只是数据源从 \texttt{config.txt} 换成了
MySQL。

\hypertarget{ux65b0ux4ea7ux7269ux67b6ux6784-2}{%
\subsection{【新产物架构】}\label{ux65b0ux4ea7ux7269ux67b6ux6784-2}}

\hypertarget{schema}{%
\subsubsection{Schema}\label{schema}}

\begin{Shaded}
\begin{Highlighting}[]
\KeywordTok{CREATE} \KeywordTok{TABLE} \ControlFlowTok{IF} \KeywordTok{NOT} \KeywordTok{EXISTS}\NormalTok{ users (}
    \KeywordTok{id}       \DataTypeTok{INT} \KeywordTok{PRIMARY} \KeywordTok{KEY}\NormalTok{ AUTO\_INCREMENT,}
\NormalTok{    username }\DataTypeTok{VARCHAR}\NormalTok{(}\DecValTok{64}\NormalTok{) }\KeywordTok{NOT} \KeywordTok{NULL} \KeywordTok{UNIQUE}\NormalTok{,}
\NormalTok{    pwd\_hash }\DataTypeTok{VARCHAR}\NormalTok{(}\DecValTok{128}\NormalTok{) }\KeywordTok{NOT} \KeywordTok{NULL}
\NormalTok{);}

\KeywordTok{CREATE} \KeywordTok{TABLE} \ControlFlowTok{IF} \KeywordTok{NOT} \KeywordTok{EXISTS}\NormalTok{ problems (}
    \KeywordTok{id}             \DataTypeTok{INT} \KeywordTok{PRIMARY} \KeywordTok{KEY}\NormalTok{,}
\NormalTok{    title          }\DataTypeTok{VARCHAR}\NormalTok{(}\DecValTok{256}\NormalTok{) }\KeywordTok{NOT} \KeywordTok{NULL}\NormalTok{,}
\NormalTok{    judge\_class    }\DataTypeTok{VARCHAR}\NormalTok{(}\DecValTok{256}\NormalTok{) }\KeywordTok{NOT} \KeywordTok{NULL}\NormalTok{,}
\NormalTok{    time\_limit\_ms  BIGINT }\KeywordTok{NOT} \KeywordTok{NULL} \KeywordTok{DEFAULT} \DecValTok{2000}\NormalTok{,}
\NormalTok{    mem\_limit\_mb   }\DataTypeTok{INT}    \KeywordTok{NOT} \KeywordTok{NULL} \KeywordTok{DEFAULT} \DecValTok{256}
\NormalTok{);}

\KeywordTok{CREATE} \KeywordTok{TABLE} \ControlFlowTok{IF} \KeywordTok{NOT} \KeywordTok{EXISTS}\NormalTok{ submissions (}
    \KeywordTok{id}          \DataTypeTok{INT} \KeywordTok{PRIMARY} \KeywordTok{KEY}\NormalTok{ AUTO\_INCREMENT,}
\NormalTok{    problem\_id  }\DataTypeTok{INT}         \KeywordTok{NOT} \KeywordTok{NULL}\NormalTok{,}
\NormalTok{    username    }\DataTypeTok{VARCHAR}\NormalTok{(}\DecValTok{64}\NormalTok{) }\KeywordTok{NOT} \KeywordTok{NULL}\NormalTok{,}
\NormalTok{    lang        }\DataTypeTok{VARCHAR}\NormalTok{(}\DecValTok{16}\NormalTok{) }\KeywordTok{NOT} \KeywordTok{NULL}\NormalTok{,}
\NormalTok{    status      }\DataTypeTok{VARCHAR}\NormalTok{(}\DecValTok{8}\NormalTok{)  }\KeywordTok{NOT} \KeywordTok{NULL}\NormalTok{,}
\NormalTok{    passed      }\DataTypeTok{INT}         \KeywordTok{NOT} \KeywordTok{NULL} \KeywordTok{DEFAULT} \DecValTok{0}\NormalTok{,}
\NormalTok{    total       }\DataTypeTok{INT}         \KeywordTok{NOT} \KeywordTok{NULL} \KeywordTok{DEFAULT} \DecValTok{0}\NormalTok{,}
\NormalTok{    time\_ms     BIGINT      }\KeywordTok{NOT} \KeywordTok{NULL} \KeywordTok{DEFAULT} \DecValTok{0}\NormalTok{,}
\NormalTok{    src\_path    }\DataTypeTok{VARCHAR}\NormalTok{(}\DecValTok{512}\NormalTok{),}
\NormalTok{    submitted\_at DATETIME   }\KeywordTok{NOT} \KeywordTok{NULL} \KeywordTok{DEFAULT} \FunctionTok{CURRENT\_TIMESTAMP}\NormalTok{,}
    \KeywordTok{FOREIGN} \KeywordTok{KEY}\NormalTok{ (problem\_id) }\KeywordTok{REFERENCES}\NormalTok{ problems(}\KeywordTok{id}\NormalTok{)}
\NormalTok{);}
\end{Highlighting}
\end{Shaded}

\texttt{src\_path}
存选手源码在文件系统上的绝对路径，文件本身不入库（大文本存 LONGTEXT 会让
\texttt{submissions} 表膨胀，且无法 \texttt{diff}）。

\begin{center}\rule{0.5\linewidth}{0.5pt}\end{center}

\hypertarget{oj.db.db-ux5355ux4f8bux8fdeux63a5ux7ba1ux7406}{%
\subsubsection{\texorpdfstring{\texttt{oj.db.Db} ---
单例连接管理}{oj.db.Db --- 单例连接管理}}\label{oj.db.db-ux5355ux4f8bux8fdeux63a5ux7ba1ux7406}}

\begin{Shaded}
\begin{Highlighting}[]
\KeywordTok{package}\ImportTok{ oj}\OperatorTok{.}\ImportTok{db}\OperatorTok{;}

\KeywordTok{import} \ImportTok{java}\OperatorTok{.}\ImportTok{sql}\OperatorTok{.}\ImportTok{Connection}\OperatorTok{;}
\KeywordTok{import} \ImportTok{java}\OperatorTok{.}\ImportTok{sql}\OperatorTok{.}\ImportTok{DriverManager}\OperatorTok{;}
\KeywordTok{import} \ImportTok{java}\OperatorTok{.}\ImportTok{sql}\OperatorTok{.}\ImportTok{SQLException}\OperatorTok{;}

\KeywordTok{public} \KeywordTok{class}\NormalTok{ Db }\OperatorTok{\{}
    \KeywordTok{private} \DataTypeTok{static} \DataTypeTok{final} \BuiltInTok{String} \BuiltInTok{URL}  \OperatorTok{=} \StringTok{"jdbc:mysql://localhost:3306/minioj?useSSL=false\&serverTimezone=UTC"}\OperatorTok{;}
    \KeywordTok{private} \DataTypeTok{static} \DataTypeTok{final} \BuiltInTok{String}\NormalTok{ USER }\OperatorTok{=} \StringTok{"oj"}\OperatorTok{;}
    \KeywordTok{private} \DataTypeTok{static} \DataTypeTok{final} \BuiltInTok{String}\NormalTok{ PASS }\OperatorTok{=} \StringTok{"ojpass"}\OperatorTok{;}

    \KeywordTok{private} \DataTypeTok{static} \BuiltInTok{Connection}\NormalTok{ conn}\OperatorTok{;}

    \KeywordTok{private} \FunctionTok{Db}\OperatorTok{()} \OperatorTok{\{\}}

    \KeywordTok{public} \DataTypeTok{static} \BuiltInTok{Connection} \FunctionTok{get}\OperatorTok{()} \KeywordTok{throws} \BuiltInTok{SQLException} \OperatorTok{\{}
        \ControlFlowTok{if} \OperatorTok{(}\NormalTok{conn }\OperatorTok{==} \KeywordTok{null} \OperatorTok{||}\NormalTok{ conn}\OperatorTok{.}\FunctionTok{isClosed}\OperatorTok{())} \OperatorTok{\{}
\NormalTok{            conn }\OperatorTok{=} \BuiltInTok{DriverManager}\OperatorTok{.}\FunctionTok{getConnection}\OperatorTok{(}\BuiltInTok{URL}\OperatorTok{,}\NormalTok{ USER}\OperatorTok{,}\NormalTok{ PASS}\OperatorTok{);}
        \OperatorTok{\}}
        \ControlFlowTok{return}\NormalTok{ conn}\OperatorTok{;}
    \OperatorTok{\}}

    \KeywordTok{public} \DataTypeTok{static} \DataTypeTok{void} \FunctionTok{close}\OperatorTok{()} \OperatorTok{\{}
        \ControlFlowTok{if} \OperatorTok{(}\NormalTok{conn }\OperatorTok{!=} \KeywordTok{null}\OperatorTok{)} \OperatorTok{\{}
            \ControlFlowTok{try} \OperatorTok{\{}
\NormalTok{                conn}\OperatorTok{.}\FunctionTok{close}\OperatorTok{();}
            \OperatorTok{\}} \ControlFlowTok{catch} \OperatorTok{(}\BuiltInTok{SQLException}\NormalTok{ e}\OperatorTok{)} \OperatorTok{\{}
\NormalTok{                e}\OperatorTok{.}\FunctionTok{printStackTrace}\OperatorTok{();}
            \OperatorTok{\}} \ControlFlowTok{finally} \OperatorTok{\{}
\NormalTok{                conn }\OperatorTok{=} \KeywordTok{null}\OperatorTok{;}
            \OperatorTok{\}}
        \OperatorTok{\}}
    \OperatorTok{\}}
\OperatorTok{\}}
\end{Highlighting}
\end{Shaded}

\textbf{教学要点}：\texttt{isClosed()}
保护让单例在连接被数据库服务端主动断开（如
\texttt{wait\_timeout}）后能自动重连，避免''幽灵连接''导致的
\texttt{Communications\ link\ failure}。

\begin{center}\rule{0.5\linewidth}{0.5pt}\end{center}

\hypertarget{oj.db.problemdao-ux5143ux6570ux636e-crud}{%
\subsubsection{\texorpdfstring{\texttt{oj.db.ProblemDao} --- 元数据
CRUD}{oj.db.ProblemDao --- 元数据 CRUD}}\label{oj.db.problemdao-ux5143ux6570ux636e-crud}}

\begin{Shaded}
\begin{Highlighting}[]
\KeywordTok{package}\ImportTok{ oj}\OperatorTok{.}\ImportTok{db}\OperatorTok{;}

\KeywordTok{import} \ImportTok{oj}\OperatorTok{.}\ImportTok{core}\OperatorTok{.}\ImportTok{ProblemMeta}\OperatorTok{;}

\KeywordTok{import} \ImportTok{java}\OperatorTok{.}\ImportTok{sql}\OperatorTok{.}\ImportTok{Connection}\OperatorTok{;}
\KeywordTok{import} \ImportTok{java}\OperatorTok{.}\ImportTok{sql}\OperatorTok{.}\ImportTok{PreparedStatement}\OperatorTok{;}
\KeywordTok{import} \ImportTok{java}\OperatorTok{.}\ImportTok{sql}\OperatorTok{.}\ImportTok{ResultSet}\OperatorTok{;}
\KeywordTok{import} \ImportTok{java}\OperatorTok{.}\ImportTok{sql}\OperatorTok{.}\ImportTok{SQLException}\OperatorTok{;}
\KeywordTok{import} \ImportTok{java}\OperatorTok{.}\ImportTok{util}\OperatorTok{.}\ImportTok{LinkedHashMap}\OperatorTok{;}
\KeywordTok{import} \ImportTok{java}\OperatorTok{.}\ImportTok{util}\OperatorTok{.}\ImportTok{Map}\OperatorTok{;}

\KeywordTok{public} \KeywordTok{class}\NormalTok{ ProblemDao }\OperatorTok{\{}

    \CommentTok{/**}
     \CommentTok{*} \CommentTok{按}\NormalTok{ id }\CommentTok{读取题目元数据，不含测试点。}
     \CommentTok{*} \CommentTok{测试点由}\NormalTok{ ProblemLoader }\CommentTok{从}\NormalTok{ FS }\CommentTok{读取，两者在}\NormalTok{ ProblemService }\CommentTok{里合并。}
     \CommentTok{*/}
    \KeywordTok{public}\NormalTok{ ProblemMeta }\FunctionTok{meta}\OperatorTok{(}\DataTypeTok{int}\NormalTok{ id}\OperatorTok{)} \KeywordTok{throws} \BuiltInTok{SQLException} \OperatorTok{\{}
        \BuiltInTok{Connection}\NormalTok{ conn }\OperatorTok{=}\NormalTok{ Db}\OperatorTok{.}\FunctionTok{get}\OperatorTok{();}
        \BuiltInTok{String}\NormalTok{ sql }\OperatorTok{=} \StringTok{"SELECT title, judge\_class, time\_limit\_ms FROM problems WHERE id = ?"}\OperatorTok{;}
        \ControlFlowTok{try} \OperatorTok{(}\BuiltInTok{PreparedStatement}\NormalTok{ ps }\OperatorTok{=}\NormalTok{ conn}\OperatorTok{.}\FunctionTok{prepareStatement}\OperatorTok{(}\NormalTok{sql}\OperatorTok{))} \OperatorTok{\{}
\NormalTok{            ps}\OperatorTok{.}\FunctionTok{setInt}\OperatorTok{(}\DecValTok{1}\OperatorTok{,}\NormalTok{ id}\OperatorTok{);}
            \ControlFlowTok{try} \OperatorTok{(}\BuiltInTok{ResultSet}\NormalTok{ rs }\OperatorTok{=}\NormalTok{ ps}\OperatorTok{.}\FunctionTok{executeQuery}\OperatorTok{())} \OperatorTok{\{}
                \ControlFlowTok{if} \OperatorTok{(!}\NormalTok{rs}\OperatorTok{.}\FunctionTok{next}\OperatorTok{())} \OperatorTok{\{}
                    \ControlFlowTok{throw} \KeywordTok{new} \BuiltInTok{SQLException}\OperatorTok{(}\StringTok{"Problem not found: "} \OperatorTok{+}\NormalTok{ id}\OperatorTok{);}
                \OperatorTok{\}}
                \ControlFlowTok{return} \KeywordTok{new} \FunctionTok{ProblemMeta}\OperatorTok{(}
\NormalTok{                    rs}\OperatorTok{.}\FunctionTok{getString}\OperatorTok{(}\StringTok{"title"}\OperatorTok{),}
\NormalTok{                    rs}\OperatorTok{.}\FunctionTok{getString}\OperatorTok{(}\StringTok{"judge\_class"}\OperatorTok{),}
\NormalTok{                    rs}\OperatorTok{.}\FunctionTok{getLong}\OperatorTok{(}\StringTok{"time\_limit\_ms"}\OperatorTok{)}
                \OperatorTok{);}
            \OperatorTok{\}}
        \OperatorTok{\}}
    \OperatorTok{\}}

    \CommentTok{/**}
     \CommentTok{*} \CommentTok{列出全部题目元数据（题目下拉用），按}\NormalTok{ id }\CommentTok{升序。}
     \CommentTok{*/}
    \KeywordTok{public} \BuiltInTok{Map}\OperatorTok{\textless{}}\BuiltInTok{Integer}\OperatorTok{,}\NormalTok{ ProblemMeta}\OperatorTok{\textgreater{}} \FunctionTok{listMeta}\OperatorTok{()} \KeywordTok{throws} \BuiltInTok{SQLException} \OperatorTok{\{}
        \BuiltInTok{Connection}\NormalTok{ conn }\OperatorTok{=}\NormalTok{ Db}\OperatorTok{.}\FunctionTok{get}\OperatorTok{();}
        \BuiltInTok{String}\NormalTok{ sql }\OperatorTok{=} \StringTok{"SELECT id, title, judge\_class, time\_limit\_ms FROM problems ORDER BY id"}\OperatorTok{;}
        \BuiltInTok{Map}\OperatorTok{\textless{}}\BuiltInTok{Integer}\OperatorTok{,}\NormalTok{ ProblemMeta}\OperatorTok{\textgreater{}}\NormalTok{ all }\OperatorTok{=} \KeywordTok{new} \BuiltInTok{LinkedHashMap}\OperatorTok{\textless{}\textgreater{}();}
        \ControlFlowTok{try} \OperatorTok{(}\BuiltInTok{PreparedStatement}\NormalTok{ ps }\OperatorTok{=}\NormalTok{ conn}\OperatorTok{.}\FunctionTok{prepareStatement}\OperatorTok{(}\NormalTok{sql}\OperatorTok{);}
             \BuiltInTok{ResultSet}\NormalTok{ rs }\OperatorTok{=}\NormalTok{ ps}\OperatorTok{.}\FunctionTok{executeQuery}\OperatorTok{())} \OperatorTok{\{}
            \ControlFlowTok{while} \OperatorTok{(}\NormalTok{rs}\OperatorTok{.}\FunctionTok{next}\OperatorTok{())} \OperatorTok{\{}
\NormalTok{                all}\OperatorTok{.}\FunctionTok{put}\OperatorTok{(}\NormalTok{rs}\OperatorTok{.}\FunctionTok{getInt}\OperatorTok{(}\StringTok{"id"}\OperatorTok{),} \KeywordTok{new} \FunctionTok{ProblemMeta}\OperatorTok{(}
\NormalTok{                    rs}\OperatorTok{.}\FunctionTok{getString}\OperatorTok{(}\StringTok{"title"}\OperatorTok{),}
\NormalTok{                    rs}\OperatorTok{.}\FunctionTok{getString}\OperatorTok{(}\StringTok{"judge\_class"}\OperatorTok{),}
\NormalTok{                    rs}\OperatorTok{.}\FunctionTok{getLong}\OperatorTok{(}\StringTok{"time\_limit\_ms"}\OperatorTok{)));}
            \OperatorTok{\}}
        \OperatorTok{\}}
        \ControlFlowTok{return}\NormalTok{ all}\OperatorTok{;}
    \OperatorTok{\}}

    \CommentTok{/**}
     \CommentTok{*} \CommentTok{插入或更新题目（管理端录题时调用）。}
     \CommentTok{*}\NormalTok{ INSERT }\CommentTok{... }\NormalTok{ON DUPLICATE KEY UPDATE }\CommentTok{保持幂等。}
     \CommentTok{*/}
    \KeywordTok{public} \DataTypeTok{void} \FunctionTok{upsert}\OperatorTok{(}\DataTypeTok{int}\NormalTok{ id}\OperatorTok{,}\NormalTok{ ProblemMeta meta}\OperatorTok{,} \DataTypeTok{int}\NormalTok{ memLimitMb}\OperatorTok{)} \KeywordTok{throws} \BuiltInTok{SQLException} \OperatorTok{\{}
        \BuiltInTok{Connection}\NormalTok{ conn }\OperatorTok{=}\NormalTok{ Db}\OperatorTok{.}\FunctionTok{get}\OperatorTok{();}
        \BuiltInTok{String}\NormalTok{ sql }\OperatorTok{=} \StringTok{"INSERT INTO problems (id, title, judge\_class, time\_limit\_ms, mem\_limit\_mb) "} \OperatorTok{+}
                     \StringTok{"VALUES (?, ?, ?, ?, ?) "} \OperatorTok{+}
                     \StringTok{"ON DUPLICATE KEY UPDATE "} \OperatorTok{+}
                     \StringTok{"title=VALUES(title), judge\_class=VALUES(judge\_class), "} \OperatorTok{+}
                     \StringTok{"time\_limit\_ms=VALUES(time\_limit\_ms), mem\_limit\_mb=VALUES(mem\_limit\_mb)"}\OperatorTok{;}
        \ControlFlowTok{try} \OperatorTok{(}\BuiltInTok{PreparedStatement}\NormalTok{ ps }\OperatorTok{=}\NormalTok{ conn}\OperatorTok{.}\FunctionTok{prepareStatement}\OperatorTok{(}\NormalTok{sql}\OperatorTok{))} \OperatorTok{\{}
\NormalTok{            ps}\OperatorTok{.}\FunctionTok{setInt}\OperatorTok{(}\DecValTok{1}\OperatorTok{,}\NormalTok{ id}\OperatorTok{);}
\NormalTok{            ps}\OperatorTok{.}\FunctionTok{setString}\OperatorTok{(}\DecValTok{2}\OperatorTok{,}\NormalTok{ meta}\OperatorTok{.}\FunctionTok{getTitle}\OperatorTok{());}
\NormalTok{            ps}\OperatorTok{.}\FunctionTok{setString}\OperatorTok{(}\DecValTok{3}\OperatorTok{,}\NormalTok{ meta}\OperatorTok{.}\FunctionTok{getJudgeClass}\OperatorTok{());}
\NormalTok{            ps}\OperatorTok{.}\FunctionTok{setLong}\OperatorTok{(}\DecValTok{4}\OperatorTok{,}\NormalTok{ meta}\OperatorTok{.}\FunctionTok{getTimeLimitMs}\OperatorTok{());}
\NormalTok{            ps}\OperatorTok{.}\FunctionTok{setInt}\OperatorTok{(}\DecValTok{5}\OperatorTok{,}\NormalTok{ memLimitMb}\OperatorTok{);}
\NormalTok{            ps}\OperatorTok{.}\FunctionTok{executeUpdate}\OperatorTok{();}
        \OperatorTok{\}}
    \OperatorTok{\}}
\OperatorTok{\}}
\end{Highlighting}
\end{Shaded}

注意 \texttt{meta(int)} \textbf{只返回 \texttt{ProblemMeta}}，不返回
\texttt{Problem}，因为 \texttt{Problem} 还需要测试点，而测试点来自
FS------这个组合动作交给 \texttt{ProblemService} 完成，职责分离。

\begin{center}\rule{0.5\linewidth}{0.5pt}\end{center}

\hypertarget{oj.db.submissiondao-ux63d0ux4ea4ux5386ux53f2ux6301ux4e45ux5316ux5e26ux4e8bux52a1}{%
\subsubsection{\texorpdfstring{\texttt{oj.db.SubmissionDao} ---
提交历史持久化（带事务）}{oj.db.SubmissionDao --- 提交历史持久化（带事务）}}\label{oj.db.submissiondao-ux63d0ux4ea4ux5386ux53f2ux6301ux4e45ux5316ux5e26ux4e8bux52a1}}

\begin{Shaded}
\begin{Highlighting}[]
\KeywordTok{package}\ImportTok{ oj}\OperatorTok{.}\ImportTok{db}\OperatorTok{;}

\KeywordTok{import} \ImportTok{oj}\OperatorTok{.}\ImportTok{core}\OperatorTok{.}\ImportTok{JudgeResult}\OperatorTok{;}
\KeywordTok{import} \ImportTok{oj}\OperatorTok{.}\ImportTok{core}\OperatorTok{.}\ImportTok{Submission}\OperatorTok{;}

\KeywordTok{import} \ImportTok{java}\OperatorTok{.}\ImportTok{sql}\OperatorTok{.}\ImportTok{Connection}\OperatorTok{;}
\KeywordTok{import} \ImportTok{java}\OperatorTok{.}\ImportTok{sql}\OperatorTok{.}\ImportTok{PreparedStatement}\OperatorTok{;}
\KeywordTok{import} \ImportTok{java}\OperatorTok{.}\ImportTok{sql}\OperatorTok{.}\ImportTok{ResultSet}\OperatorTok{;}
\KeywordTok{import} \ImportTok{java}\OperatorTok{.}\ImportTok{sql}\OperatorTok{.}\ImportTok{SQLException}\OperatorTok{;}
\KeywordTok{import} \ImportTok{java}\OperatorTok{.}\ImportTok{util}\OperatorTok{.}\ImportTok{ArrayList}\OperatorTok{;}
\KeywordTok{import} \ImportTok{java}\OperatorTok{.}\ImportTok{util}\OperatorTok{.}\ImportTok{List}\OperatorTok{;}

\KeywordTok{public} \KeywordTok{class}\NormalTok{ SubmissionDao }\OperatorTok{\{}

    \CommentTok{/**}
     \CommentTok{*} \CommentTok{将一条提交结果持久化进}\NormalTok{ submissions }\CommentTok{表。}
     \CommentTok{*} \CommentTok{整个操作在事务内：先}\NormalTok{ INSERT submissions}\CommentTok{，若判题已完成同步写}\NormalTok{ status}\CommentTok{/}\NormalTok{passed}\CommentTok{/}\NormalTok{total}\CommentTok{。}
     \CommentTok{*} \CommentTok{任何一步失败整体}\NormalTok{ rollback}\CommentTok{，保证不产生孤儿记录。}
     \CommentTok{*/}
    \KeywordTok{public} \DataTypeTok{void} \FunctionTok{save}\OperatorTok{(}\NormalTok{Submission sub}\OperatorTok{)} \KeywordTok{throws} \BuiltInTok{SQLException} \OperatorTok{\{}
        \BuiltInTok{Connection}\NormalTok{ conn }\OperatorTok{=}\NormalTok{ Db}\OperatorTok{.}\FunctionTok{get}\OperatorTok{();}
        \DataTypeTok{boolean}\NormalTok{ prev }\OperatorTok{=}\NormalTok{ conn}\OperatorTok{.}\FunctionTok{getAutoCommit}\OperatorTok{();}
\NormalTok{        conn}\OperatorTok{.}\FunctionTok{setAutoCommit}\OperatorTok{(}\KeywordTok{false}\OperatorTok{);}
        \ControlFlowTok{try} \OperatorTok{\{}
            \BuiltInTok{String}\NormalTok{ sql }\OperatorTok{=} \StringTok{"INSERT INTO submissions "} \OperatorTok{+}
                         \StringTok{"(problem\_id, username, lang, status, passed, total, time\_ms, src\_path) "} \OperatorTok{+}
                         \StringTok{"VALUES (?, ?, ?, ?, ?, ?, ?, ?)"}\OperatorTok{;}
            \ControlFlowTok{try} \OperatorTok{(}\BuiltInTok{PreparedStatement}\NormalTok{ ps }\OperatorTok{=}\NormalTok{ conn}\OperatorTok{.}\FunctionTok{prepareStatement}\OperatorTok{(}\NormalTok{sql}\OperatorTok{))} \OperatorTok{\{}
\NormalTok{                JudgeResult r }\OperatorTok{=}\NormalTok{ sub}\OperatorTok{.}\FunctionTok{getResult}\OperatorTok{();}
\NormalTok{                ps}\OperatorTok{.}\FunctionTok{setInt}\OperatorTok{(}\DecValTok{1}\OperatorTok{,}\NormalTok{ sub}\OperatorTok{.}\FunctionTok{getProblemId}\OperatorTok{());}
\NormalTok{                ps}\OperatorTok{.}\FunctionTok{setString}\OperatorTok{(}\DecValTok{2}\OperatorTok{,}\NormalTok{ sub}\OperatorTok{.}\FunctionTok{getUserName}\OperatorTok{());}
\NormalTok{                ps}\OperatorTok{.}\FunctionTok{setString}\OperatorTok{(}\DecValTok{3}\OperatorTok{,}\NormalTok{ sub}\OperatorTok{.}\FunctionTok{getLang}\OperatorTok{());}
\NormalTok{                ps}\OperatorTok{.}\FunctionTok{setString}\OperatorTok{(}\DecValTok{4}\OperatorTok{,}\NormalTok{ r }\OperatorTok{!=} \KeywordTok{null} \OperatorTok{?}\NormalTok{ r}\OperatorTok{.}\FunctionTok{getStatus}\OperatorTok{().}\FunctionTok{name}\OperatorTok{()} \OperatorTok{:} \StringTok{"PENDING"}\OperatorTok{);}
\NormalTok{                ps}\OperatorTok{.}\FunctionTok{setInt}\OperatorTok{(}\DecValTok{5}\OperatorTok{,}\NormalTok{ r }\OperatorTok{!=} \KeywordTok{null} \OperatorTok{?}\NormalTok{ r}\OperatorTok{.}\FunctionTok{getPassed}\OperatorTok{()} \OperatorTok{:} \DecValTok{0}\OperatorTok{);}
\NormalTok{                ps}\OperatorTok{.}\FunctionTok{setInt}\OperatorTok{(}\DecValTok{6}\OperatorTok{,}\NormalTok{ r }\OperatorTok{!=} \KeywordTok{null} \OperatorTok{?}\NormalTok{ r}\OperatorTok{.}\FunctionTok{getTotal}\OperatorTok{()} \OperatorTok{:} \DecValTok{0}\OperatorTok{);}
\NormalTok{                ps}\OperatorTok{.}\FunctionTok{setLong}\OperatorTok{(}\DecValTok{7}\OperatorTok{,}\NormalTok{ r }\OperatorTok{!=} \KeywordTok{null} \OperatorTok{?}\NormalTok{ r}\OperatorTok{.}\FunctionTok{getElapsedMs}\OperatorTok{()} \OperatorTok{:} \DecValTok{0L}\OperatorTok{);}
\NormalTok{                ps}\OperatorTok{.}\FunctionTok{setString}\OperatorTok{(}\DecValTok{8}\OperatorTok{,}\NormalTok{ sub}\OperatorTok{.}\FunctionTok{getSrcPath}\OperatorTok{());}
\NormalTok{                ps}\OperatorTok{.}\FunctionTok{executeUpdate}\OperatorTok{();}
            \OperatorTok{\}}
\NormalTok{            conn}\OperatorTok{.}\FunctionTok{commit}\OperatorTok{();}
        \OperatorTok{\}} \ControlFlowTok{catch} \OperatorTok{(}\BuiltInTok{SQLException}\NormalTok{ e}\OperatorTok{)} \OperatorTok{\{}
\NormalTok{            conn}\OperatorTok{.}\FunctionTok{rollback}\OperatorTok{();}
            \ControlFlowTok{throw}\NormalTok{ e}\OperatorTok{;}
        \OperatorTok{\}} \ControlFlowTok{finally} \OperatorTok{\{}
\NormalTok{            conn}\OperatorTok{.}\FunctionTok{setAutoCommit}\OperatorTok{(}\NormalTok{prev}\OperatorTok{);}
        \OperatorTok{\}}
    \OperatorTok{\}}

    \CommentTok{/**}
     \CommentTok{*} \CommentTok{查询某用户在某题上的所有提交，按时间倒序，最多返回}\NormalTok{ limit }\CommentTok{条。}
     \CommentTok{*/}
    \KeywordTok{public} \BuiltInTok{List}\OperatorTok{\textless{}}\NormalTok{Submission}\OperatorTok{\textgreater{}} \FunctionTok{listByUser}\OperatorTok{(}\BuiltInTok{String}\NormalTok{ username}\OperatorTok{,} \DataTypeTok{int}\NormalTok{ problemId}\OperatorTok{,} \DataTypeTok{int}\NormalTok{ limit}\OperatorTok{)} \KeywordTok{throws} \BuiltInTok{SQLException} \OperatorTok{\{}
        \BuiltInTok{Connection}\NormalTok{ conn }\OperatorTok{=}\NormalTok{ Db}\OperatorTok{.}\FunctionTok{get}\OperatorTok{();}
        \BuiltInTok{String}\NormalTok{ sql }\OperatorTok{=} \StringTok{"SELECT id, problem\_id, username, lang, status, passed, total, time\_ms, src\_path "} \OperatorTok{+}
                     \StringTok{"FROM submissions "} \OperatorTok{+}
                     \StringTok{"WHERE username = ? AND problem\_id = ? "} \OperatorTok{+}
                     \StringTok{"ORDER BY submitted\_at DESC LIMIT ?"}\OperatorTok{;}
        \BuiltInTok{List}\OperatorTok{\textless{}}\NormalTok{Submission}\OperatorTok{\textgreater{}}\NormalTok{ result }\OperatorTok{=} \KeywordTok{new} \BuiltInTok{ArrayList}\OperatorTok{\textless{}\textgreater{}();}
        \ControlFlowTok{try} \OperatorTok{(}\BuiltInTok{PreparedStatement}\NormalTok{ ps }\OperatorTok{=}\NormalTok{ conn}\OperatorTok{.}\FunctionTok{prepareStatement}\OperatorTok{(}\NormalTok{sql}\OperatorTok{))} \OperatorTok{\{}
\NormalTok{            ps}\OperatorTok{.}\FunctionTok{setString}\OperatorTok{(}\DecValTok{1}\OperatorTok{,}\NormalTok{ username}\OperatorTok{);}
\NormalTok{            ps}\OperatorTok{.}\FunctionTok{setInt}\OperatorTok{(}\DecValTok{2}\OperatorTok{,}\NormalTok{ problemId}\OperatorTok{);}
\NormalTok{            ps}\OperatorTok{.}\FunctionTok{setInt}\OperatorTok{(}\DecValTok{3}\OperatorTok{,}\NormalTok{ limit}\OperatorTok{);}
            \ControlFlowTok{try} \OperatorTok{(}\BuiltInTok{ResultSet}\NormalTok{ rs }\OperatorTok{=}\NormalTok{ ps}\OperatorTok{.}\FunctionTok{executeQuery}\OperatorTok{())} \OperatorTok{\{}
                \ControlFlowTok{while} \OperatorTok{(}\NormalTok{rs}\OperatorTok{.}\FunctionTok{next}\OperatorTok{())} \OperatorTok{\{}
\NormalTok{                    Submission sub }\OperatorTok{=} \KeywordTok{new} \FunctionTok{Submission}\OperatorTok{(}
\NormalTok{                        rs}\OperatorTok{.}\FunctionTok{getInt}\OperatorTok{(}\StringTok{"problem\_id"}\OperatorTok{),}
\NormalTok{                        rs}\OperatorTok{.}\FunctionTok{getString}\OperatorTok{(}\StringTok{"username"}\OperatorTok{),}
\NormalTok{                        rs}\OperatorTok{.}\FunctionTok{getString}\OperatorTok{(}\StringTok{"lang"}\OperatorTok{),}
\NormalTok{                        rs}\OperatorTok{.}\FunctionTok{getString}\OperatorTok{(}\StringTok{"src\_path"}\OperatorTok{)}
                    \OperatorTok{);}
\NormalTok{                    oj}\OperatorTok{.}\FunctionTok{core}\OperatorTok{.}\FunctionTok{Status}\NormalTok{ st }\OperatorTok{=}\NormalTok{ oj}\OperatorTok{.}\FunctionTok{core}\OperatorTok{.}\FunctionTok{Status}\OperatorTok{.}\FunctionTok{valueOf}\OperatorTok{(}\NormalTok{rs}\OperatorTok{.}\FunctionTok{getString}\OperatorTok{(}\StringTok{"status"}\OperatorTok{));}
\NormalTok{                    JudgeResult jr }\OperatorTok{=} \KeywordTok{new} \FunctionTok{JudgeResult}\OperatorTok{(}
\NormalTok{                        st}\OperatorTok{,}
\NormalTok{                        rs}\OperatorTok{.}\FunctionTok{getInt}\OperatorTok{(}\StringTok{"passed"}\OperatorTok{),}
\NormalTok{                        rs}\OperatorTok{.}\FunctionTok{getInt}\OperatorTok{(}\StringTok{"total"}\OperatorTok{),}
                        \StringTok{""}\OperatorTok{,}
\NormalTok{                        rs}\OperatorTok{.}\FunctionTok{getLong}\OperatorTok{(}\StringTok{"time\_ms"}\OperatorTok{)}
                    \OperatorTok{);}
\NormalTok{                    sub}\OperatorTok{.}\FunctionTok{setResult}\OperatorTok{(}\NormalTok{jr}\OperatorTok{);}
\NormalTok{                    result}\OperatorTok{.}\FunctionTok{add}\OperatorTok{(}\NormalTok{sub}\OperatorTok{);}
                \OperatorTok{\}}
            \OperatorTok{\}}
        \OperatorTok{\}}
        \ControlFlowTok{return}\NormalTok{ result}\OperatorTok{;}
    \OperatorTok{\}}
\OperatorTok{\}}
\end{Highlighting}
\end{Shaded}

\textbf{事务模式说明}：\texttt{save} 先保存原始
\texttt{conn.getAutoCommit()} 状态，\texttt{finally}
里恢复，这样单例连接被后续调用复用时不会携带残留的
\texttt{autoCommit=false} 状态。

\begin{center}\rule{0.5\linewidth}{0.5pt}\end{center}

\hypertarget{oj.service.problemservice-fsdb-ux5408ux5e76ux5165ux53e3}{%
\subsubsection{\texorpdfstring{\texttt{oj.service.ProblemService} ---
FS/DB
合并入口}{oj.service.ProblemService --- FS/DB 合并入口}}\label{oj.service.problemservice-fsdb-ux5408ux5e76ux5165ux53e3}}

\begin{Shaded}
\begin{Highlighting}[]
\KeywordTok{package}\ImportTok{ oj}\OperatorTok{.}\ImportTok{service}\OperatorTok{;}

\KeywordTok{import} \ImportTok{oj}\OperatorTok{.}\ImportTok{core}\OperatorTok{.}\ImportTok{Problem}\OperatorTok{;}
\KeywordTok{import} \ImportTok{oj}\OperatorTok{.}\ImportTok{core}\OperatorTok{.}\ImportTok{ProblemMeta}\OperatorTok{;}
\KeywordTok{import} \ImportTok{oj}\OperatorTok{.}\ImportTok{core}\OperatorTok{.}\ImportTok{TestCase}\OperatorTok{;}
\KeywordTok{import} \ImportTok{oj}\OperatorTok{.}\ImportTok{db}\OperatorTok{.}\ImportTok{ProblemDao}\OperatorTok{;}
\KeywordTok{import} \ImportTok{oj}\OperatorTok{.}\ImportTok{io}\OperatorTok{.}\ImportTok{ProblemLoader}\OperatorTok{;}

\KeywordTok{import} \ImportTok{java}\OperatorTok{.}\ImportTok{sql}\OperatorTok{.}\ImportTok{SQLException}\OperatorTok{;}
\KeywordTok{import} \ImportTok{java}\OperatorTok{.}\ImportTok{util}\OperatorTok{.}\ImportTok{List}\OperatorTok{;}
\KeywordTok{import} \ImportTok{java}\OperatorTok{.}\ImportTok{util}\OperatorTok{.}\ImportTok{Map}\OperatorTok{;}

\CommentTok{/**}
 \CommentTok{*} \CommentTok{统一加载入口：}\NormalTok{DB }\CommentTok{提供元数据，}\NormalTok{FS }\CommentTok{提供测试点，两者在此合并为完整}\NormalTok{ Problem}\CommentTok{。}
 \CommentTok{*} \CommentTok{调用方（}\NormalTok{OjFrame}\CommentTok{、}\NormalTok{JudgeWorker}\CommentTok{）只依赖本类，不直接接触}\NormalTok{ ProblemDao }\CommentTok{或}\NormalTok{ ProblemLoader}\CommentTok{。}
 \CommentTok{*/}
\KeywordTok{public} \KeywordTok{class}\NormalTok{ ProblemService }\OperatorTok{\{}

    \KeywordTok{private} \DataTypeTok{final}\NormalTok{ ProblemDao dao}\OperatorTok{;}
    \KeywordTok{private} \DataTypeTok{final}\NormalTok{ ProblemLoader loader}\OperatorTok{;}

    \KeywordTok{public} \FunctionTok{ProblemService}\OperatorTok{(}\NormalTok{ProblemDao dao}\OperatorTok{,}\NormalTok{ ProblemLoader loader}\OperatorTok{)} \OperatorTok{\{}
        \KeywordTok{this}\OperatorTok{.}\FunctionTok{dao} \OperatorTok{=}\NormalTok{ dao}\OperatorTok{;}
        \KeywordTok{this}\OperatorTok{.}\FunctionTok{loader} \OperatorTok{=}\NormalTok{ loader}\OperatorTok{;}
    \OperatorTok{\}}

    \CommentTok{/**}
     \CommentTok{*} \CommentTok{按题目}\NormalTok{ id }\CommentTok{返回含完整测试点的}\NormalTok{ Problem}\CommentTok{。}
     \CommentTok{*}
\CommentTok{     * @}\NormalTok{param id }\CommentTok{题目编号}
     \CommentTok{*} \CommentTok{@}\NormalTok{return }\CommentTok{合并了}\NormalTok{ DB }\CommentTok{元数据与}\NormalTok{ FS }\CommentTok{测试点的}\NormalTok{ Problem }\CommentTok{对象}
     \CommentTok{*} \CommentTok{@}\NormalTok{throws SQLException     }\CommentTok{数据库访问异常}
     \CommentTok{*} \CommentTok{@}\NormalTok{throws RuntimeException }\CommentTok{题目目录不存在或测试点格式错误}
     \CommentTok{*/}
    \KeywordTok{public}\NormalTok{ Problem }\FunctionTok{load}\OperatorTok{(}\DataTypeTok{int}\NormalTok{ id}\OperatorTok{)} \KeywordTok{throws} \BuiltInTok{SQLException} \OperatorTok{\{}
\NormalTok{        ProblemMeta meta }\OperatorTok{=}\NormalTok{ dao}\OperatorTok{.}\FunctionTok{meta}\OperatorTok{(}\NormalTok{id}\OperatorTok{);}             \CommentTok{// DB：元数据}
        \BuiltInTok{List}\OperatorTok{\textless{}}\NormalTok{TestCase}\OperatorTok{\textgreater{}}\NormalTok{ cases }\OperatorTok{=}\NormalTok{ loader}\OperatorTok{.}\FunctionTok{cases}\OperatorTok{(}\NormalTok{id}\OperatorTok{);}     \CommentTok{// FS：测试点}
        \ControlFlowTok{return} \KeywordTok{new} \FunctionTok{Problem}\OperatorTok{(}\NormalTok{id}\OperatorTok{,}\NormalTok{ meta}\OperatorTok{,}\NormalTok{ cases}\OperatorTok{);}
    \OperatorTok{\}}

    \CommentTok{/**} \CommentTok{只取元数据（}\NormalTok{GUI }\CommentTok{取时限/标题用，不读测试点）。} \CommentTok{*/}
    \KeywordTok{public}\NormalTok{ ProblemMeta }\FunctionTok{meta}\OperatorTok{(}\DataTypeTok{int}\NormalTok{ id}\OperatorTok{)} \KeywordTok{throws} \BuiltInTok{SQLException} \OperatorTok{\{}
        \ControlFlowTok{return}\NormalTok{ dao}\OperatorTok{.}\FunctionTok{meta}\OperatorTok{(}\NormalTok{id}\OperatorTok{);}
    \OperatorTok{\}}

    \CommentTok{/**} \CommentTok{列出全部题目元数据（题目下拉用）。} \CommentTok{*/}
    \KeywordTok{public} \BuiltInTok{Map}\OperatorTok{\textless{}}\BuiltInTok{Integer}\OperatorTok{,}\NormalTok{ ProblemMeta}\OperatorTok{\textgreater{}} \FunctionTok{listMeta}\OperatorTok{()} \KeywordTok{throws} \BuiltInTok{SQLException} \OperatorTok{\{}
        \ControlFlowTok{return}\NormalTok{ dao}\OperatorTok{.}\FunctionTok{listMeta}\OperatorTok{();}
    \OperatorTok{\}}
\OperatorTok{\}}
\end{Highlighting}
\end{Shaded}

\texttt{ProblemService}
是本节最关键的\textbf{组合点}：它让上层（\texttt{JudgeWorker}、\texttt{OjFrame}）只看到一个
\texttt{load(int)} 接口，完全不感知底层是 DB
还是文件。将来换成远程存储只改 \texttt{ProblemService}
内部，调用方零修改。

\begin{center}\rule{0.5\linewidth}{0.5pt}\end{center}

\hypertarget{ux6a21ux5757ux95f4ux6570ux636eux6d41m5b-ux5168ux666f}{%
\subsubsection{模块间数据流（M5b
全景）}\label{ux6a21ux5757ux95f4ux6570ux636eux6d41m5b-ux5168ux666f}}

\begin{verbatim}
OjFrame / JudgeWorker
        |
        v
 ProblemService.load(id)
   +- ProblemDao.meta(id)  -->  MySQL · problems 表
   +- ProblemLoader.cases(id) -> FS · problems/<id>/N.in,N.out
        |
        v
   Problem（含 ProblemMeta + List<TestCase>）
        |
        v
 MachineJudge.judge(...)  -->  C++ 外部判题机
        |
        v
 SubmissionDao.save(sub)  -->  MySQL · submissions 表（事务）
\end{verbatim}

测试点文件（\texttt{.in} / \texttt{.out}）与选手源码（\texttt{.java} /
\texttt{.cpp}）始终留在文件系统，MySQL 只存可查询的结构化字段。

\begin{quote}
自 M5b 起，提交历史改由 \texttt{SubmissionDao} 入库；M5a 的
\texttt{SubmissionStore}（\texttt{.ser} 序列化）停用，不再两套并存。
\end{quote}

\hypertarget{ux52a8ux624bux5b9eux73b0ux4e24ux6b65ux8d70javac-java-ux5b9eux64cd-4}{%
\subsection{动手实现:两步走(javac / java
实操)}\label{ux52a8ux624bux5b9eux73b0ux4e24ux6b65ux8d70javac-java-ux5b9eux64cd-4}}

\begin{quote}
需先装好 MySQL、建库 \texttt{mini\_oj},并把驱动
\texttt{mysql-connector-j-*.jar} 放进 \texttt{lib/}。
\end{quote}

\hypertarget{ux7b2cux4e00ux6b65-ux6700ux7b80-jdbcux5efaux8868ux63d2ux5165ux67e5ux8be2statement}{%
\subsubsection{\texorpdfstring{第一步 \KS --- 最简
JDBC:建表/插入/查询(Statement)}{第一步 --- 最简 JDBC:建表/插入/查询(Statement)}}\label{ux7b2cux4e00ux6b65-ux6700ux7b80-jdbcux5efaux8868ux63d2ux5165ux67e5ux8be2statement}}

\textbf{\texttt{src/oj/db/SimpleDb.java}}

\begin{Shaded}
\begin{Highlighting}[]
\KeywordTok{package}\ImportTok{ oj}\OperatorTok{.}\ImportTok{db}\OperatorTok{;}
\KeywordTok{import} \ImportTok{java}\OperatorTok{.}\ImportTok{sql}\OperatorTok{.*;}
\KeywordTok{import} \ImportTok{java}\OperatorTok{.}\ImportTok{util}\OperatorTok{.*;}
\KeywordTok{public} \KeywordTok{class}\NormalTok{ SimpleDb }\OperatorTok{\{}
    \DataTypeTok{static} \DataTypeTok{final} \BuiltInTok{String} \BuiltInTok{URL} \OperatorTok{=} \StringTok{"jdbc:mysql://localhost:3306/mini\_oj?serverTimezone=UTC"}\OperatorTok{;}
    \KeywordTok{public} \DataTypeTok{static} \DataTypeTok{void} \FunctionTok{main}\OperatorTok{(}\BuiltInTok{String}\OperatorTok{[]}\NormalTok{ args}\OperatorTok{)} \KeywordTok{throws} \BuiltInTok{Exception} \OperatorTok{\{}
        \ControlFlowTok{try} \OperatorTok{(}\BuiltInTok{Connection}\NormalTok{ c }\OperatorTok{=} \BuiltInTok{DriverManager}\OperatorTok{.}\FunctionTok{getConnection}\OperatorTok{(}\BuiltInTok{URL}\OperatorTok{,} \StringTok{"root"}\OperatorTok{,} \StringTok{"root"}\OperatorTok{);}
             \BuiltInTok{Statement}\NormalTok{ st }\OperatorTok{=}\NormalTok{ c}\OperatorTok{.}\FunctionTok{createStatement}\OperatorTok{())} \OperatorTok{\{}
\NormalTok{            st}\OperatorTok{.}\FunctionTok{executeUpdate}\OperatorTok{(}\StringTok{"CREATE TABLE IF NOT EXISTS sub(id INT, name VARCHAR(32), status VARCHAR(8))"}\OperatorTok{);}
\NormalTok{            st}\OperatorTok{.}\FunctionTok{executeUpdate}\OperatorTok{(}\StringTok{"INSERT INTO sub VALUES(1,\textquotesingle{}alice\textquotesingle{},\textquotesingle{}AC\textquotesingle{})"}\OperatorTok{);}      \CommentTok{// 增}
            \BuiltInTok{ResultSet}\NormalTok{ rs }\OperatorTok{=}\NormalTok{ st}\OperatorTok{.}\FunctionTok{executeQuery}\OperatorTok{(}\StringTok{"SELECT id,name,status FROM sub"}\OperatorTok{);}\CommentTok{// 查}
            \BuiltInTok{List}\OperatorTok{\textless{}}\BuiltInTok{String}\OperatorTok{\textgreater{}}\NormalTok{ rows }\OperatorTok{=} \KeywordTok{new} \BuiltInTok{ArrayList}\OperatorTok{\textless{}\textgreater{}();}
            \ControlFlowTok{while} \OperatorTok{(}\NormalTok{rs}\OperatorTok{.}\FunctionTok{next}\OperatorTok{())}                                                \CommentTok{// 游标遍历}
\NormalTok{                rows}\OperatorTok{.}\FunctionTok{add}\OperatorTok{(}\NormalTok{rs}\OperatorTok{.}\FunctionTok{getInt}\OperatorTok{(}\StringTok{"id"}\OperatorTok{)} \OperatorTok{+} \StringTok{" "} \OperatorTok{+}\NormalTok{ rs}\OperatorTok{.}\FunctionTok{getString}\OperatorTok{(}\StringTok{"name"}\OperatorTok{)} \OperatorTok{+} \StringTok{" "} \OperatorTok{+}\NormalTok{ rs}\OperatorTok{.}\FunctionTok{getString}\OperatorTok{(}\StringTok{"status"}\OperatorTok{));}
\NormalTok{            rows}\OperatorTok{.}\FunctionTok{forEach}\OperatorTok{(}\BuiltInTok{System}\OperatorTok{.}\FunctionTok{out}\OperatorTok{::}\NormalTok{println}\OperatorTok{);}
        \OperatorTok{\}}
    \OperatorTok{\}}
\OperatorTok{\}}
\end{Highlighting}
\end{Shaded}

讲解:\texttt{DriverManager.getConnection} -\textgreater{}
\texttt{Statement} -\textgreater{} \texttt{executeUpdate/executeQuery}
-\textgreater{} \texttt{ResultSet.next()} 游标------JDBC
四步固定套路,用在 OJ 自己的提交表上。

编译运行(驱动挂 classpath):

\begin{Shaded}
\begin{Highlighting}[]
\ExtensionTok{javac} \AttributeTok{{-}cp} \StringTok{"lib/*"} \AttributeTok{{-}d}\NormalTok{ build src/oj/db/SimpleDb.java}
\ExtensionTok{java} \AttributeTok{{-}cp} \StringTok{"build:lib/*"}\NormalTok{ oj.db.SimpleDb}
\CommentTok{\# 1 alice AC}
\end{Highlighting}
\end{Shaded}

\hypertarget{ux7b2cux4e8cux6b65-dao-ux4e8bux52a1-ux9884ux5904ux7406-fsdb-ux5206ux5de5}{%
\subsubsection{\texorpdfstring{第二步 \GC --- DAO + 事务 + 预处理 +
FS/DB
分工}{第二步 --- DAO + 事务 + 预处理 + FS/DB 分工}}\label{ux7b2cux4e8cux6b65-dao-ux4e8bux52a1-ux9884ux5904ux7406-fsdb-ux5206ux5de5}}

升级为 \texttt{Db.get()}
单例、\texttt{ProblemDao.meta/listMeta}、\texttt{SubmissionDao.save}(\texttt{PreparedStatement}
防注入 + \texttt{setAutoCommit(false)}/\texttt{commit}/\texttt{rollback}
事务)、\texttt{ProblemService} 合并 DB 元数据 + FS
测试点(代码见上文【新产物架构】):

\begin{Shaded}
\begin{Highlighting}[]
\ExtensionTok{javac} \AttributeTok{{-}cp} \StringTok{"lib/*"} \AttributeTok{{-}d}\NormalTok{ build }\VariableTok{$(}\FunctionTok{find}\NormalTok{ src }\AttributeTok{{-}name} \StringTok{\textquotesingle{}*.java\textquotesingle{}}\VariableTok{)}
\ExtensionTok{java} \AttributeTok{{-}cp} \StringTok{"build:lib/*"}\NormalTok{ oj.Main}
\end{Highlighting}
\end{Shaded}

\hypertarget{ux9a8cux6536ux6807ux51c6-6}{%
\subsection{验收标准}\label{ux9a8cux6536ux6807ux51c6-6}}

\begin{enumerate}
\def\labelenumi{\arabic{enumi}.}
\tightlist
\item
  执行 \texttt{schema.sql} 建库建表后，向 \texttt{problems}
  插入一条记录，运行 \texttt{ProblemDao.meta(id)} 能正确返回
  \texttt{ProblemMeta}，字段值与数据库一致。
\item
  \texttt{ProblemService.load(id)} 能同时取到 DB
  元数据（\texttt{getTitle()}、\texttt{getTimeLimitMs()}）和 FS
  测试点（\texttt{getCases().size()\ \textgreater{}\ 0}），两者均正确。
\item
  构造一个带 \texttt{JudgeResult}（status=AC, passed=3, total=3,
  elapsedMs=42）的 \texttt{Submission}，调用
  \texttt{SubmissionDao.save}，用 \texttt{SELECT\ *\ FROM\ submissions}
  验证数据库中出现对应行，status 字段值为 \texttt{"AC"}。
\item
  模拟 \texttt{save} 执行到一半时抛出异常（可在第二条 SQL 前手动
  \texttt{throw\ new\ SQLException("test")}），验证 \texttt{submissions}
  表中\textbf{没有}半截记录（事务 rollback 生效）。
\item
  \texttt{listByUser} 按 \texttt{submitted\_at\ DESC}
  返回，最新的提交排在 \texttt{get(0)}，结果条数不超过传入的
  \texttt{limit}。
\item
  在 \texttt{Db.get()} 关闭后再次调用，能自动重建连接，不抛
  \texttt{"Connection\ is\ closed"} 异常。
\item
  \texttt{ProblemDao.meta} 接受含单引号的 \texttt{title}（如
  \texttt{"It\textquotesingle{}s\ a\ Trap!"}），查询结果正确返回，不抛
  SQL 语法异常（验证 \texttt{PreparedStatement} 防注入有效）。
\end{enumerate}

\hypertarget{m5c-swing-ux684cux9762ux5ba2ux6237ux7aefch9ux5927ux524dux7aef}{%
\section{08 · M5c Swing
桌面客户端（Ch9，大前端）}\label{m5c-swing-ux684cux9762ux5ba2ux6237ux7aefch9ux5927ux524dux7aef}}

\begin{quote}
\textbf{考试相关度 ★★★(编程+填空双大题)} ·
真题考点参考:JFrame/布局/JButton/JTextField/JLabel/ActionListener。\textbf{{[}第一步{]}}
\KS 用考试级 Swing 写 OJ
提交窗口最简版(选题+输入+提交按钮+标签显示判题结果);\textbf{{[}第二步{]}}
\GC \texttt{OjFrame}/\texttt{OjController} MVC 解耦 + SwingWorker。
\end{quote}

\begin{quote}
判题逻辑已经具备：数据库存题、C++
判题机真正跑代码。现在缺少的是一扇窗户------用户还在用命令行粘贴代码、手动查库。本章用
Java Swing 为 Mini-OJ
装上可视化界面，把前三代的成果串联成一个完整的桌面应用。
\end{quote}

\hypertarget{ux6838ux5fc3ux75dbux70b9-3}{%
\subsection{【核心痛点】}\label{ux6838ux5fc3ux75dbux70b9-3}}

M5b 之后，系统后端已经相当完善：ProblemService
从数据库读取题目元数据，ProblemLoader 加载测试点，MachineJudge 调用外部
C++ 判题机真实编译运行。但所有交互都停留在命令行或单元测试层面。

具体痛点如下：

\begin{enumerate}
\def\labelenumi{\arabic{enumi}.}
\tightlist
\item
  \textbf{无题目浏览}：用户不知道有哪些题，必须手动查数据库或翻代码。
\item
  \textbf{无语言选择}：提交时语言类型硬编码在测试脚本里。
\item
  \textbf{无源码编辑区}：用户只能把源文件路径写死，无法在界面里直接粘贴或编写代码。
\item
  \textbf{无结果展示}：判题完成后结果只打印在控制台，没有弹窗或表格呈现给用户。
\end{enumerate}

本章目标：用 Swing 构建一个 MVC
结构的桌面客户端，让用户在图形界面里选题、选语言、贴代码、提交，然后看到
AC/WA/TLE 等结果。

\begin{center}\rule{0.5\linewidth}{0.5pt}\end{center}

\hypertarget{ux5f15ux5165ux8bfeux672cux77e5ux8bc6ux70b9-3}{%
\subsection{【引入课本知识点】}\label{ux5f15ux5165ux8bfeux672cux77e5ux8bc6ux70b9-3}}

本章对应 Ch9------Java 图形用户界面（Swing）。

\textbf{JFrame 与 BorderLayout}

\texttt{JFrame} 是 Swing 的顶层窗口容器。\texttt{BorderLayout}
是其默认布局管理器，将窗口分为 NORTH/SOUTH/CENTER/EAST/WEST
五个区域，适合''工具栏在上、内容在中、操作按钮在下''的经典布局。

\begin{Shaded}
\begin{Highlighting}[]
\BuiltInTok{JFrame}\NormalTok{ frame }\OperatorTok{=} \KeywordTok{new} \BuiltInTok{JFrame}\OperatorTok{(}\StringTok{"Mini{-}OJ"}\OperatorTok{);}
\NormalTok{frame}\OperatorTok{.}\FunctionTok{setLayout}\OperatorTok{(}\KeywordTok{new} \BuiltInTok{BorderLayout}\OperatorTok{());}
\NormalTok{frame}\OperatorTok{.}\FunctionTok{add}\OperatorTok{(}\NormalTok{toolbar}\OperatorTok{,} \BuiltInTok{BorderLayout}\OperatorTok{.}\FunctionTok{NORTH}\OperatorTok{);}
\NormalTok{frame}\OperatorTok{.}\FunctionTok{add}\OperatorTok{(}\NormalTok{scrollPane}\OperatorTok{,} \BuiltInTok{BorderLayout}\OperatorTok{.}\FunctionTok{CENTER}\OperatorTok{);}
\NormalTok{frame}\OperatorTok{.}\FunctionTok{add}\OperatorTok{(}\NormalTok{buttonPanel}\OperatorTok{,} \BuiltInTok{BorderLayout}\OperatorTok{.}\FunctionTok{SOUTH}\OperatorTok{);}
\end{Highlighting}
\end{Shaded}

\textbf{JComboBox --- 下拉选择框}

用于语言选择（cpp / python）和题目选择。泛型版本
\texttt{JComboBox\textless{}String\textgreater{}} 类型安全，通过
\texttt{getSelectedItem()} 获取当前选中值。题目下拉框统一用
\texttt{JComboBox\textless{}String\textgreater{}} 显示 ``id
标题''，\textbf{不要把它当成
\texttt{JComboBox\textless{}Integer\textgreater{}}}，取题号要靠 OjFrame
内部维护的 id 列表换算。

\begin{Shaded}
\begin{Highlighting}[]
\BuiltInTok{JComboBox}\OperatorTok{\textless{}}\BuiltInTok{String}\OperatorTok{\textgreater{}}\NormalTok{ langBox }\OperatorTok{=} \KeywordTok{new} \BuiltInTok{JComboBox}\OperatorTok{\textless{}\textgreater{}(}\KeywordTok{new} \BuiltInTok{String}\OperatorTok{[]\{}\StringTok{"cpp"}\OperatorTok{,} \StringTok{"python"}\OperatorTok{\});}
\BuiltInTok{String}\NormalTok{ lang }\OperatorTok{=} \OperatorTok{(}\BuiltInTok{String}\OperatorTok{)}\NormalTok{ langBox}\OperatorTok{.}\FunctionTok{getSelectedItem}\OperatorTok{();}
\end{Highlighting}
\end{Shaded}

\textbf{JTextArea --- 多行文本编辑区}

源码编辑区的核心组件。需要用 \texttt{JScrollPane}
包裹以支持滚动，设置等宽字体提升代码可读性。

\begin{Shaded}
\begin{Highlighting}[]
\BuiltInTok{JTextArea}\NormalTok{ codeArea }\OperatorTok{=} \KeywordTok{new} \BuiltInTok{JTextArea}\OperatorTok{(}\DecValTok{20}\OperatorTok{,} \DecValTok{60}\OperatorTok{);}
\NormalTok{codeArea}\OperatorTok{.}\FunctionTok{setFont}\OperatorTok{(}\KeywordTok{new} \BuiltInTok{Font}\OperatorTok{(}\StringTok{"Monospaced"}\OperatorTok{,} \BuiltInTok{Font}\OperatorTok{.}\FunctionTok{PLAIN}\OperatorTok{,} \DecValTok{13}\OperatorTok{));}
\BuiltInTok{JScrollPane}\NormalTok{ scroll }\OperatorTok{=} \KeywordTok{new} \BuiltInTok{JScrollPane}\OperatorTok{(}\NormalTok{codeArea}\OperatorTok{);}
\end{Highlighting}
\end{Shaded}

\textbf{JButton + Lambda 事件监听}

Ch9 的重点之一：用函数式接口 \texttt{ActionListener}（单方法接口）配合
Lambda 表达式绑定按钮事件，替代匿名内部类写法，代码更简洁。

\begin{Shaded}
\begin{Highlighting}[]
\BuiltInTok{JButton}\NormalTok{ submitBtn }\OperatorTok{=} \KeywordTok{new} \BuiltInTok{JButton}\OperatorTok{(}\StringTok{"提交"}\OperatorTok{);}
\NormalTok{submitBtn}\OperatorTok{.}\FunctionTok{addActionListener}\OperatorTok{(}\NormalTok{e }\OperatorTok{{-}\textgreater{}}\NormalTok{ controller}\OperatorTok{.}\FunctionTok{onSubmit}\OperatorTok{());}
\end{Highlighting}
\end{Shaded}

\texttt{ActionListener}
是函数式接口（\texttt{@FunctionalInterface}），\texttt{actionPerformed(ActionEvent\ e)}
是其唯一抽象方法，Lambda \texttt{e\ -\textgreater{}\ ...}
正是该方法的实现体。

\textbf{JLabel --- 状态标签}

工具栏里放一个
\texttt{JLabel}，用来实时显示当前状态（``就绪''、判题结果文本等），比每次都弹窗更轻量。本章
OjFrame 会真实持有这个状态标签字段并对外暴露
\texttt{getResultLabel()}，下一章（M6a）会直接复用它做''评测中\ldots``提示。

\textbf{JOptionPane --- 状态弹窗}

用于在判题完成后弹出结果对话框，无需手动创建 \texttt{JDialog}。

\begin{Shaded}
\begin{Highlighting}[]
\BuiltInTok{JOptionPane}\OperatorTok{.}\FunctionTok{showMessageDialog}\OperatorTok{(}\NormalTok{frame}\OperatorTok{,}\NormalTok{ result}\OperatorTok{.}\FunctionTok{toString}\OperatorTok{(),} \StringTok{"判题结果"}\OperatorTok{,}
    \BuiltInTok{JOptionPane}\OperatorTok{.}\FunctionTok{INFORMATION\_MESSAGE}\OperatorTok{);}
\end{Highlighting}
\end{Shaded}

\textbf{JTable --- 提交历史表格}

通过 \texttt{DefaultTableModel}
动态添加行，展示历史提交的题号、语言、状态、耗时。

\textbf{EDT（Event Dispatch Thread）规则}

Swing 不是线程安全的，所有 UI 操作必须在 EDT 上执行。长时间操作（如调用
MachineJudge）必须在后台线程（\texttt{SwingWorker}
或普通线程）里执行，完成后用 \texttt{SwingUtilities.invokeLater} 回到
EDT 更新界面。

\begin{center}\rule{0.5\linewidth}{0.5pt}\end{center}

\hypertarget{ux4e09ux4ee3ux6f14ux8fdbux5b9aux4f4d-3}{%
\subsection{【三代演进定位】}\label{ux4e09ux4ee3ux6f14ux8fdbux5b9aux4f4d-3}}

本章\textbf{不产生新判题机代}。Swing 客户端是''大前端''，其定位是：

\begin{verbatim}
用户 -> OjFrame（View）-> OjController（Controller）
         |                      |
    ProblemService（Model）   MachineJudge（第三代）
         |                      |
      数据库（M5b）          C++ 判题机进程
         |                      |
    SubmissionDao.save()     JudgeResult 返回
\end{verbatim}

\begin{itemize}
\tightlist
\item
  \textbf{第一代}（M1--M3，已完成）：Solution.solve 在 JVM
  内执行，题目硬编码，无 UI。
\item
  \textbf{第二代}（M4）：反射工厂 + 单文件配置，解耦题目元数据，仍无
  UI。
\item
  \textbf{第三代}（M5a）：ProcessBuilder 调用外部 C++
  判题机，真实编译运行，仍无 UI。
\item
  \textbf{M5c（本章）}：Swing 客户端直接调用第三代 MachineJudge 和 M5b
  的 DAO，不经过网络（\textbf{不使用
  Socket/ServerSocket}），所有调用均为本地方法调用。
\end{itemize}

删除说明：Ch13
网络内容已从课程中移除，因此客户端与判题机之间\textbf{不存在任何 Socket
通信}。OjController 直接持有 MachineJudge 和 SubmissionDao
的引用，在同一 JVM 进程内完成全部工作。

\begin{center}\rule{0.5\linewidth}{0.5pt}\end{center}

\hypertarget{ux65b0ux4ea7ux7269ux67b6ux6784-3}{%
\subsection{【新产物架构】}\label{ux65b0ux4ea7ux7269ux67b6ux6784-3}}

\hypertarget{ux6d89ux53caux7684ux7c7bux4e0eux63a5ux53e3}{%
\subsubsection{涉及的类与接口}\label{ux6d89ux53caux7684ux7c7bux4e0eux63a5ux53e3}}

\begin{longtable}[]{@{}
  >{\raggedright\arraybackslash}p{(\columnwidth - 4\tabcolsep) * \real{0.3333}}
  >{\raggedright\arraybackslash}p{(\columnwidth - 4\tabcolsep) * \real{0.3333}}
  >{\raggedright\arraybackslash}p{(\columnwidth - 4\tabcolsep) * \real{0.3333}}@{}}
\toprule\noalign{}
\begin{minipage}[b]{\linewidth}\raggedright
类/接口
\end{minipage} & \begin{minipage}[b]{\linewidth}\raggedright
包
\end{minipage} & \begin{minipage}[b]{\linewidth}\raggedright
职责
\end{minipage} \\
\midrule\noalign{}
\endhead
\bottomrule\noalign{}
\endlastfoot
\texttt{OjFrame} & \texttt{oj.gui} & View 层；持有全部 Swing
组件；不含业务逻辑 \\
\texttt{OjController} & \texttt{oj.gui} & Controller
层；处理提交事件；协调 Model \\
\texttt{ProblemService} & \texttt{oj.service} & Model 层；从 DB + FS
加载题目（M5b 已实现） \\
\texttt{MachineJudge} & \texttt{oj.judge} & 第三代判题机；ProcessBuilder
调用外部二进制 \\
\texttt{SubmissionDao} & \texttt{oj.db} & 持久化提交记录到数据库（M5b
已实现） \\
\texttt{Submission} & \texttt{oj.core} & 提交数据模型（M5a 已实现） \\
\texttt{JudgeResult} & \texttt{oj.core} & 判题结果数据模型 \\
\end{longtable}

\begin{center}\rule{0.5\linewidth}{0.5pt}\end{center}

\hypertarget{ojframeview-ux5c42}{%
\subsubsection{OjFrame（View 层）}\label{ojframeview-ux5c42}}

OjFrame 负责所有 Swing 组件的创建和布局，对外暴露 getter 让 Controller
读取输入、写入结果，不直接调用任何业务逻辑。题目下拉用
\texttt{JComboBox\textless{}String\textgreater{}}，内部用一个
\texttt{problemIds}
列表把下标映射回真正的题号，\texttt{getSelectedProblemId()}
在没有选中时返回 \texttt{-1}。工具栏里有一个真实的状态标签
\texttt{resultLabel}，并通过 \texttt{getResultLabel()} 暴露给
Controller（下一章 M6a 会复用它）。

\begin{Shaded}
\begin{Highlighting}[]
\KeywordTok{package}\ImportTok{ oj}\OperatorTok{.}\ImportTok{gui}\OperatorTok{;}

\KeywordTok{import} \ImportTok{javax}\OperatorTok{.}\ImportTok{swing}\OperatorTok{.*;}
\KeywordTok{import} \ImportTok{javax}\OperatorTok{.}\ImportTok{swing}\OperatorTok{.}\ImportTok{table}\OperatorTok{.}\ImportTok{DefaultTableModel}\OperatorTok{;}
\KeywordTok{import} \ImportTok{java}\OperatorTok{.}\ImportTok{awt}\OperatorTok{.*;}
\KeywordTok{import} \ImportTok{java}\OperatorTok{.}\ImportTok{util}\OperatorTok{.}\ImportTok{List}\OperatorTok{;}

\KeywordTok{public} \KeywordTok{class}\NormalTok{ OjFrame }\KeywordTok{extends} \BuiltInTok{JFrame} \OperatorTok{\{}

    \KeywordTok{private} \DataTypeTok{final} \BuiltInTok{JComboBox}\OperatorTok{\textless{}}\BuiltInTok{String}\OperatorTok{\textgreater{}}\NormalTok{ problemBox}\OperatorTok{;}
    \KeywordTok{private} \DataTypeTok{final} \BuiltInTok{JComboBox}\OperatorTok{\textless{}}\BuiltInTok{String}\OperatorTok{\textgreater{}}\NormalTok{ langBox}\OperatorTok{;}
    \KeywordTok{private} \DataTypeTok{final} \BuiltInTok{JTextArea}\NormalTok{ codeArea}\OperatorTok{;}
    \KeywordTok{private} \DataTypeTok{final} \BuiltInTok{JButton}\NormalTok{ submitBtn}\OperatorTok{;}
    \KeywordTok{private} \DataTypeTok{final} \BuiltInTok{JLabel}\NormalTok{ resultLabel}\OperatorTok{;}
    \KeywordTok{private} \DataTypeTok{final} \BuiltInTok{DefaultTableModel}\NormalTok{ historyModel}\OperatorTok{;}
    \KeywordTok{private} \DataTypeTok{final} \BuiltInTok{JTable}\NormalTok{ historyTable}\OperatorTok{;}

    \KeywordTok{private} \BuiltInTok{List}\OperatorTok{\textless{}}\BuiltInTok{Integer}\OperatorTok{\textgreater{}}\NormalTok{ problemIds}\OperatorTok{;}   \CommentTok{// 与 problemBox 下标对应}

    \KeywordTok{public} \FunctionTok{OjFrame}\OperatorTok{()} \OperatorTok{\{}
        \KeywordTok{super}\OperatorTok{(}\StringTok{"Mini{-}OJ 桌面客户端"}\OperatorTok{);}
        \FunctionTok{setDefaultCloseOperation}\OperatorTok{(}\BuiltInTok{JFrame}\OperatorTok{.}\FunctionTok{EXIT\_ON\_CLOSE}\OperatorTok{);}
        \FunctionTok{setSize}\OperatorTok{(}\DecValTok{900}\OperatorTok{,} \DecValTok{650}\OperatorTok{);}
        \FunctionTok{setLayout}\OperatorTok{(}\KeywordTok{new} \BuiltInTok{BorderLayout}\OperatorTok{(}\DecValTok{6}\OperatorTok{,} \DecValTok{6}\OperatorTok{));}

        \CommentTok{// {-}{-} 北部工具栏 {-}{-}{-}{-}{-}{-}{-}{-}{-}{-}{-}{-}{-}{-}{-}{-}{-}{-}{-}{-}{-}{-}{-}{-}{-}{-}{-}{-}{-}{-}{-}{-}{-}{-}{-}{-}{-}{-}{-}{-}{-}{-}}
        \BuiltInTok{JPanel}\NormalTok{ toolbar }\OperatorTok{=} \KeywordTok{new} \BuiltInTok{JPanel}\OperatorTok{(}\KeywordTok{new} \BuiltInTok{FlowLayout}\OperatorTok{(}\BuiltInTok{FlowLayout}\OperatorTok{.}\FunctionTok{LEFT}\OperatorTok{,} \DecValTok{8}\OperatorTok{,} \DecValTok{4}\OperatorTok{));}

\NormalTok{        toolbar}\OperatorTok{.}\FunctionTok{add}\OperatorTok{(}\KeywordTok{new} \BuiltInTok{JLabel}\OperatorTok{(}\StringTok{"题目："}\OperatorTok{));}
\NormalTok{        problemBox }\OperatorTok{=} \KeywordTok{new} \BuiltInTok{JComboBox}\OperatorTok{\textless{}\textgreater{}();}
\NormalTok{        problemBox}\OperatorTok{.}\FunctionTok{setPreferredSize}\OperatorTok{(}\KeywordTok{new} \BuiltInTok{Dimension}\OperatorTok{(}\DecValTok{260}\OperatorTok{,} \DecValTok{26}\OperatorTok{));}
\NormalTok{        toolbar}\OperatorTok{.}\FunctionTok{add}\OperatorTok{(}\NormalTok{problemBox}\OperatorTok{);}

\NormalTok{        toolbar}\OperatorTok{.}\FunctionTok{add}\OperatorTok{(}\KeywordTok{new} \BuiltInTok{JLabel}\OperatorTok{(}\StringTok{"语言："}\OperatorTok{));}
\NormalTok{        langBox }\OperatorTok{=} \KeywordTok{new} \BuiltInTok{JComboBox}\OperatorTok{\textless{}\textgreater{}(}\KeywordTok{new} \BuiltInTok{String}\OperatorTok{[]\{}\StringTok{"cpp"}\OperatorTok{,} \StringTok{"python"}\OperatorTok{\});}
\NormalTok{        toolbar}\OperatorTok{.}\FunctionTok{add}\OperatorTok{(}\NormalTok{langBox}\OperatorTok{);}

\NormalTok{        submitBtn }\OperatorTok{=} \KeywordTok{new} \BuiltInTok{JButton}\OperatorTok{(}\StringTok{"提交"}\OperatorTok{);}
\NormalTok{        toolbar}\OperatorTok{.}\FunctionTok{add}\OperatorTok{(}\NormalTok{submitBtn}\OperatorTok{);}

        \CommentTok{// 状态标签：实时显示判题状态，供 Controller 写入}
\NormalTok{        resultLabel }\OperatorTok{=} \KeywordTok{new} \BuiltInTok{JLabel}\OperatorTok{(}\StringTok{"就绪"}\OperatorTok{);}
\NormalTok{        resultLabel}\OperatorTok{.}\FunctionTok{setForeground}\OperatorTok{(}\KeywordTok{new} \BuiltInTok{Color}\OperatorTok{(}\BaseNTok{0x33}\OperatorTok{,} \BaseNTok{0x66}\OperatorTok{,} \BaseNTok{0x99}\OperatorTok{));}
\NormalTok{        toolbar}\OperatorTok{.}\FunctionTok{add}\OperatorTok{(}\NormalTok{resultLabel}\OperatorTok{);}

        \FunctionTok{add}\OperatorTok{(}\NormalTok{toolbar}\OperatorTok{,} \BuiltInTok{BorderLayout}\OperatorTok{.}\FunctionTok{NORTH}\OperatorTok{);}

        \CommentTok{// {-}{-} 中部源码编辑区 {-}{-}{-}{-}{-}{-}{-}{-}{-}{-}{-}{-}{-}{-}{-}{-}{-}{-}{-}{-}{-}{-}{-}{-}{-}{-}{-}{-}{-}{-}{-}{-}{-}{-}{-}{-}}
\NormalTok{        codeArea }\OperatorTok{=} \KeywordTok{new} \BuiltInTok{JTextArea}\OperatorTok{();}
\NormalTok{        codeArea}\OperatorTok{.}\FunctionTok{setFont}\OperatorTok{(}\KeywordTok{new} \BuiltInTok{Font}\OperatorTok{(}\StringTok{"Monospaced"}\OperatorTok{,} \BuiltInTok{Font}\OperatorTok{.}\FunctionTok{PLAIN}\OperatorTok{,} \DecValTok{13}\OperatorTok{));}
\NormalTok{        codeArea}\OperatorTok{.}\FunctionTok{setTabSize}\OperatorTok{(}\DecValTok{4}\OperatorTok{);}
        \BuiltInTok{JScrollPane}\NormalTok{ codeScroll }\OperatorTok{=} \KeywordTok{new} \BuiltInTok{JScrollPane}\OperatorTok{(}\NormalTok{codeArea}\OperatorTok{);}
\NormalTok{        codeScroll}\OperatorTok{.}\FunctionTok{setBorder}\OperatorTok{(}\BuiltInTok{BorderFactory}\OperatorTok{.}\FunctionTok{createTitledBorder}\OperatorTok{(}\StringTok{"源码"}\OperatorTok{));}
        \FunctionTok{add}\OperatorTok{(}\NormalTok{codeScroll}\OperatorTok{,} \BuiltInTok{BorderLayout}\OperatorTok{.}\FunctionTok{CENTER}\OperatorTok{);}

        \CommentTok{// {-}{-} 南部提交历史表格 {-}{-}{-}{-}{-}{-}{-}{-}{-}{-}{-}{-}{-}{-}{-}{-}{-}{-}{-}{-}{-}{-}{-}{-}{-}{-}{-}{-}{-}{-}{-}{-}{-}{-}}
        \BuiltInTok{String}\OperatorTok{[]}\NormalTok{ cols }\OperatorTok{=} \OperatorTok{\{}\StringTok{"提交 ID"}\OperatorTok{,} \StringTok{"题目 ID"}\OperatorTok{,} \StringTok{"语言"}\OperatorTok{,} \StringTok{"状态"}\OperatorTok{,} \StringTok{"通过/总计"}\OperatorTok{,} \StringTok{"耗时(ms)"}\OperatorTok{\};}
\NormalTok{        historyModel }\OperatorTok{=} \KeywordTok{new} \BuiltInTok{DefaultTableModel}\OperatorTok{(}\NormalTok{cols}\OperatorTok{,} \DecValTok{0}\OperatorTok{)} \OperatorTok{\{}
            \AttributeTok{@Override}
            \KeywordTok{public} \DataTypeTok{boolean} \FunctionTok{isCellEditable}\OperatorTok{(}\DataTypeTok{int}\NormalTok{ row}\OperatorTok{,} \DataTypeTok{int}\NormalTok{ col}\OperatorTok{)} \OperatorTok{\{}
                \ControlFlowTok{return} \KeywordTok{false}\OperatorTok{;}
            \OperatorTok{\}}
        \OperatorTok{\};}
\NormalTok{        historyTable }\OperatorTok{=} \KeywordTok{new} \BuiltInTok{JTable}\OperatorTok{(}\NormalTok{historyModel}\OperatorTok{);}
\NormalTok{        historyTable}\OperatorTok{.}\FunctionTok{setFillsViewportHeight}\OperatorTok{(}\KeywordTok{true}\OperatorTok{);}
        \BuiltInTok{JScrollPane}\NormalTok{ tableScroll }\OperatorTok{=} \KeywordTok{new} \BuiltInTok{JScrollPane}\OperatorTok{(}\NormalTok{historyTable}\OperatorTok{);}
\NormalTok{        tableScroll}\OperatorTok{.}\FunctionTok{setPreferredSize}\OperatorTok{(}\KeywordTok{new} \BuiltInTok{Dimension}\OperatorTok{(}\DecValTok{0}\OperatorTok{,} \DecValTok{180}\OperatorTok{));}
\NormalTok{        tableScroll}\OperatorTok{.}\FunctionTok{setBorder}\OperatorTok{(}\BuiltInTok{BorderFactory}\OperatorTok{.}\FunctionTok{createTitledBorder}\OperatorTok{(}\StringTok{"提交历史"}\OperatorTok{));}
        \FunctionTok{add}\OperatorTok{(}\NormalTok{tableScroll}\OperatorTok{,} \BuiltInTok{BorderLayout}\OperatorTok{.}\FunctionTok{SOUTH}\OperatorTok{);}

        \FunctionTok{setLocationRelativeTo}\OperatorTok{(}\KeywordTok{null}\OperatorTok{);}
    \OperatorTok{\}}

    \CommentTok{// {-}{-} 初始化题目下拉列表（由 Controller 在加载完数据后调用）{-}{-}}
    \KeywordTok{public} \DataTypeTok{void} \FunctionTok{loadProblems}\OperatorTok{(}\NormalTok{java}\OperatorTok{.}\FunctionTok{util}\OperatorTok{.}\FunctionTok{List}\OperatorTok{\textless{}}\BuiltInTok{Integer}\OperatorTok{\textgreater{}}\NormalTok{ ids}\OperatorTok{,}\NormalTok{ java}\OperatorTok{.}\FunctionTok{util}\OperatorTok{.}\FunctionTok{List}\OperatorTok{\textless{}}\BuiltInTok{String}\OperatorTok{\textgreater{}}\NormalTok{ titles}\OperatorTok{)} \OperatorTok{\{}
        \KeywordTok{this}\OperatorTok{.}\FunctionTok{problemIds} \OperatorTok{=}\NormalTok{ ids}\OperatorTok{;}
\NormalTok{        problemBox}\OperatorTok{.}\FunctionTok{removeAllItems}\OperatorTok{();}
        \ControlFlowTok{for} \OperatorTok{(}\DataTypeTok{int}\NormalTok{ i }\OperatorTok{=} \DecValTok{0}\OperatorTok{;}\NormalTok{ i }\OperatorTok{\textless{}}\NormalTok{ ids}\OperatorTok{.}\FunctionTok{size}\OperatorTok{();}\NormalTok{ i}\OperatorTok{++)} \OperatorTok{\{}
\NormalTok{            problemBox}\OperatorTok{.}\FunctionTok{addItem}\OperatorTok{(}\NormalTok{ids}\OperatorTok{.}\FunctionTok{get}\OperatorTok{(}\NormalTok{i}\OperatorTok{)} \OperatorTok{+} \StringTok{" "} \OperatorTok{+}\NormalTok{ titles}\OperatorTok{.}\FunctionTok{get}\OperatorTok{(}\NormalTok{i}\OperatorTok{));}
        \OperatorTok{\}}
    \OperatorTok{\}}

    \CommentTok{// {-}{-} 向历史表格追加一行 {-}{-}{-}{-}{-}{-}{-}{-}{-}{-}{-}{-}{-}{-}{-}{-}{-}{-}{-}{-}{-}{-}{-}{-}{-}{-}{-}{-}{-}{-}{-}{-}{-}{-}{-}{-}}
    \KeywordTok{public} \DataTypeTok{void} \FunctionTok{appendHistory}\OperatorTok{(}\DataTypeTok{int}\NormalTok{ subId}\OperatorTok{,} \DataTypeTok{int}\NormalTok{ probId}\OperatorTok{,} \BuiltInTok{String}\NormalTok{ lang}\OperatorTok{,}
                              \BuiltInTok{String}\NormalTok{ status}\OperatorTok{,} \BuiltInTok{String}\NormalTok{ passTotal}\OperatorTok{,} \DataTypeTok{long}\NormalTok{ ms}\OperatorTok{)} \OperatorTok{\{}
\NormalTok{        historyModel}\OperatorTok{.}\FunctionTok{addRow}\OperatorTok{(}\KeywordTok{new} \BuiltInTok{Object}\OperatorTok{[]\{}\NormalTok{subId}\OperatorTok{,}\NormalTok{ probId}\OperatorTok{,}\NormalTok{ lang}\OperatorTok{,}\NormalTok{ status}\OperatorTok{,}\NormalTok{ passTotal}\OperatorTok{,}\NormalTok{ ms}\OperatorTok{\});}
        \DataTypeTok{int}\NormalTok{ last }\OperatorTok{=}\NormalTok{ historyModel}\OperatorTok{.}\FunctionTok{getRowCount}\OperatorTok{()} \OperatorTok{{-}} \DecValTok{1}\OperatorTok{;}
\NormalTok{        historyTable}\OperatorTok{.}\FunctionTok{scrollRectToVisible}\OperatorTok{(}\NormalTok{historyTable}\OperatorTok{.}\FunctionTok{getCellRect}\OperatorTok{(}\NormalTok{last}\OperatorTok{,} \DecValTok{0}\OperatorTok{,} \KeywordTok{true}\OperatorTok{));}
    \OperatorTok{\}}

    \CommentTok{// {-}{-} Getters 供 Controller 读取 {-}{-}{-}{-}{-}{-}{-}{-}{-}{-}{-}{-}{-}{-}{-}{-}{-}{-}{-}{-}{-}{-}{-}{-}{-}{-}{-}}
    \KeywordTok{public} \DataTypeTok{int} \FunctionTok{getSelectedProblemId}\OperatorTok{()} \OperatorTok{\{}
        \DataTypeTok{int}\NormalTok{ idx }\OperatorTok{=}\NormalTok{ problemBox}\OperatorTok{.}\FunctionTok{getSelectedIndex}\OperatorTok{();}
        \ControlFlowTok{if} \OperatorTok{(}\NormalTok{idx }\OperatorTok{\textless{}} \DecValTok{0} \OperatorTok{||}\NormalTok{ problemIds }\OperatorTok{==} \KeywordTok{null}\OperatorTok{)} \ControlFlowTok{return} \OperatorTok{{-}}\DecValTok{1}\OperatorTok{;}
        \ControlFlowTok{return}\NormalTok{ problemIds}\OperatorTok{.}\FunctionTok{get}\OperatorTok{(}\NormalTok{idx}\OperatorTok{);}
    \OperatorTok{\}}

    \KeywordTok{public} \BuiltInTok{String} \FunctionTok{getSelectedLang}\OperatorTok{()} \OperatorTok{\{}
        \ControlFlowTok{return} \OperatorTok{(}\BuiltInTok{String}\OperatorTok{)}\NormalTok{ langBox}\OperatorTok{.}\FunctionTok{getSelectedItem}\OperatorTok{();}
    \OperatorTok{\}}

    \KeywordTok{public} \BuiltInTok{String} \FunctionTok{getSourceCode}\OperatorTok{()} \OperatorTok{\{}
        \ControlFlowTok{return}\NormalTok{ codeArea}\OperatorTok{.}\FunctionTok{getText}\OperatorTok{();}
    \OperatorTok{\}}

    \KeywordTok{public} \BuiltInTok{JButton} \FunctionTok{getSubmitBtn}\OperatorTok{()} \OperatorTok{\{}
        \ControlFlowTok{return}\NormalTok{ submitBtn}\OperatorTok{;}
    \OperatorTok{\}}

    \KeywordTok{public} \BuiltInTok{JLabel} \FunctionTok{getResultLabel}\OperatorTok{()} \OperatorTok{\{}
        \ControlFlowTok{return}\NormalTok{ resultLabel}\OperatorTok{;}
    \OperatorTok{\}}
\OperatorTok{\}}
\end{Highlighting}
\end{Shaded}

\begin{center}\rule{0.5\linewidth}{0.5pt}\end{center}

\hypertarget{ojcontrollercontroller-ux5c42}{%
\subsubsection{OjController（Controller
层）}\label{ojcontrollercontroller-ux5c42}}

OjController 持有 Model 的全部引用，构造签名为
\texttt{(OjFrame\ view,\ ProblemService\ svc,\ MachineJudge\ judge,\ SubmissionDao\ dao)}。构造时用
\texttt{svc.listMeta()} 把题目填进下拉框，并绑定提交按钮的 Lambda
事件监听器。提交操作在后台线程执行，避免阻塞 EDT。

注意几个与全局契约对齐的关键点：

\begin{itemize}
\tightlist
\item
  ProblemService 用 \texttt{listMeta()} / \texttt{meta(int)}，不是
  \texttt{listAll()}。
\item
  \texttt{SubmissionDao.save(Submission)} 返回 \texttt{void}，所以拿提交
  ID 要先 \texttt{dao.save(sub)} 再 \texttt{sub.getId()}（id 由
  Submission 静态计数自增）。
\item
  临时源文件写好后\textbf{不要立即删除}------它的路径要存进 Submission
  的 \texttt{srcPath}，留作记录。
\end{itemize}

\begin{Shaded}
\begin{Highlighting}[]
\KeywordTok{package}\ImportTok{ oj}\OperatorTok{.}\ImportTok{gui}\OperatorTok{;}

\KeywordTok{import} \ImportTok{oj}\OperatorTok{.}\ImportTok{core}\OperatorTok{.}\ImportTok{JudgeResult}\OperatorTok{;}
\KeywordTok{import} \ImportTok{oj}\OperatorTok{.}\ImportTok{core}\OperatorTok{.}\ImportTok{ProblemMeta}\OperatorTok{;}
\KeywordTok{import} \ImportTok{oj}\OperatorTok{.}\ImportTok{core}\OperatorTok{.}\ImportTok{Submission}\OperatorTok{;}
\KeywordTok{import} \ImportTok{oj}\OperatorTok{.}\ImportTok{db}\OperatorTok{.}\ImportTok{SubmissionDao}\OperatorTok{;}
\KeywordTok{import} \ImportTok{oj}\OperatorTok{.}\ImportTok{judge}\OperatorTok{.}\ImportTok{MachineJudge}\OperatorTok{;}
\KeywordTok{import} \ImportTok{oj}\OperatorTok{.}\ImportTok{service}\OperatorTok{.}\ImportTok{ProblemService}\OperatorTok{;}

\KeywordTok{import} \ImportTok{javax}\OperatorTok{.}\ImportTok{swing}\OperatorTok{.*;}
\KeywordTok{import} \ImportTok{java}\OperatorTok{.}\ImportTok{io}\OperatorTok{.}\ImportTok{IOException}\OperatorTok{;}
\KeywordTok{import} \ImportTok{java}\OperatorTok{.}\ImportTok{nio}\OperatorTok{.}\ImportTok{charset}\OperatorTok{.}\ImportTok{StandardCharsets}\OperatorTok{;}
\KeywordTok{import} \ImportTok{java}\OperatorTok{.}\ImportTok{nio}\OperatorTok{.}\ImportTok{file}\OperatorTok{.}\ImportTok{Files}\OperatorTok{;}
\KeywordTok{import} \ImportTok{java}\OperatorTok{.}\ImportTok{nio}\OperatorTok{.}\ImportTok{file}\OperatorTok{.}\ImportTok{Path}\OperatorTok{;}
\KeywordTok{import} \ImportTok{java}\OperatorTok{.}\ImportTok{util}\OperatorTok{.}\ImportTok{ArrayList}\OperatorTok{;}
\KeywordTok{import} \ImportTok{java}\OperatorTok{.}\ImportTok{util}\OperatorTok{.}\ImportTok{List}\OperatorTok{;}
\KeywordTok{import} \ImportTok{java}\OperatorTok{.}\ImportTok{util}\OperatorTok{.}\ImportTok{Map}\OperatorTok{;}

\KeywordTok{public} \KeywordTok{class}\NormalTok{ OjController }\OperatorTok{\{}

    \KeywordTok{private} \DataTypeTok{final}\NormalTok{ OjFrame view}\OperatorTok{;}
    \KeywordTok{private} \DataTypeTok{final}\NormalTok{ ProblemService svc}\OperatorTok{;}
    \KeywordTok{private} \DataTypeTok{final}\NormalTok{ MachineJudge judge}\OperatorTok{;}
    \KeywordTok{private} \DataTypeTok{final}\NormalTok{ SubmissionDao dao}\OperatorTok{;}

    \KeywordTok{private} \DataTypeTok{static} \DataTypeTok{final} \BuiltInTok{String}\NormalTok{ PROBLEMS\_DIR }\OperatorTok{=} \StringTok{"problems"}\OperatorTok{;}
    \KeywordTok{private} \DataTypeTok{static} \DataTypeTok{final} \BuiltInTok{String}\NormalTok{ USER\_NAME    }\OperatorTok{=} \StringTok{"student"}\OperatorTok{;}   \CommentTok{// 演示用固定用户}

    \KeywordTok{public} \FunctionTok{OjController}\OperatorTok{(}\NormalTok{OjFrame view}\OperatorTok{,}
\NormalTok{                        ProblemService svc}\OperatorTok{,}
\NormalTok{                        MachineJudge judge}\OperatorTok{,}
\NormalTok{                        SubmissionDao dao}\OperatorTok{)} \OperatorTok{\{}
        \KeywordTok{this}\OperatorTok{.}\FunctionTok{view}  \OperatorTok{=}\NormalTok{ view}\OperatorTok{;}
        \KeywordTok{this}\OperatorTok{.}\FunctionTok{svc}   \OperatorTok{=}\NormalTok{ svc}\OperatorTok{;}
        \KeywordTok{this}\OperatorTok{.}\FunctionTok{judge} \OperatorTok{=}\NormalTok{ judge}\OperatorTok{;}
        \KeywordTok{this}\OperatorTok{.}\FunctionTok{dao}   \OperatorTok{=}\NormalTok{ dao}\OperatorTok{;}

        \FunctionTok{initProblems}\OperatorTok{();}
        \FunctionTok{bindEvents}\OperatorTok{();}
    \OperatorTok{\}}

    \CommentTok{// 从数据库加载题目列表，填充下拉框}
    \KeywordTok{private} \DataTypeTok{void} \FunctionTok{initProblems}\OperatorTok{()} \OperatorTok{\{}
        \ControlFlowTok{try} \OperatorTok{\{}
            \BuiltInTok{Map}\OperatorTok{\textless{}}\BuiltInTok{Integer}\OperatorTok{,}\NormalTok{ ProblemMeta}\OperatorTok{\textgreater{}}\NormalTok{ all }\OperatorTok{=}\NormalTok{ svc}\OperatorTok{.}\FunctionTok{listMeta}\OperatorTok{();}
            \BuiltInTok{List}\OperatorTok{\textless{}}\BuiltInTok{Integer}\OperatorTok{\textgreater{}}\NormalTok{ ids    }\OperatorTok{=} \KeywordTok{new} \BuiltInTok{ArrayList}\OperatorTok{\textless{}\textgreater{}(}\NormalTok{all}\OperatorTok{.}\FunctionTok{keySet}\OperatorTok{());}
            \BuiltInTok{List}\OperatorTok{\textless{}}\BuiltInTok{String}\OperatorTok{\textgreater{}}\NormalTok{  titles }\OperatorTok{=} \KeywordTok{new} \BuiltInTok{ArrayList}\OperatorTok{\textless{}\textgreater{}();}
\NormalTok{            ids}\OperatorTok{.}\FunctionTok{forEach}\OperatorTok{(}\NormalTok{id }\OperatorTok{{-}\textgreater{}}\NormalTok{ titles}\OperatorTok{.}\FunctionTok{add}\OperatorTok{(}\NormalTok{all}\OperatorTok{.}\FunctionTok{get}\OperatorTok{(}\NormalTok{id}\OperatorTok{).}\FunctionTok{getTitle}\OperatorTok{()));}
\NormalTok{            view}\OperatorTok{.}\FunctionTok{loadProblems}\OperatorTok{(}\NormalTok{ids}\OperatorTok{,}\NormalTok{ titles}\OperatorTok{);}
        \OperatorTok{\}} \ControlFlowTok{catch} \OperatorTok{(}\BuiltInTok{Exception}\NormalTok{ ex}\OperatorTok{)} \OperatorTok{\{}
            \BuiltInTok{JOptionPane}\OperatorTok{.}\FunctionTok{showMessageDialog}\OperatorTok{(}\NormalTok{view}\OperatorTok{,}
                \StringTok{"加载题目失败："} \OperatorTok{+}\NormalTok{ ex}\OperatorTok{.}\FunctionTok{getMessage}\OperatorTok{(),} \StringTok{"错误"}\OperatorTok{,}
                \BuiltInTok{JOptionPane}\OperatorTok{.}\FunctionTok{ERROR\_MESSAGE}\OperatorTok{);}
        \OperatorTok{\}}
    \OperatorTok{\}}

    \CommentTok{// 绑定提交按钮：Lambda {-}\textgreater{} 后台线程执行判题，SwingUtilities 回 EDT 更新 UI}
    \KeywordTok{private} \DataTypeTok{void} \FunctionTok{bindEvents}\OperatorTok{()} \OperatorTok{\{}
\NormalTok{        view}\OperatorTok{.}\FunctionTok{getSubmitBtn}\OperatorTok{().}\FunctionTok{addActionListener}\OperatorTok{(}\NormalTok{e }\OperatorTok{{-}\textgreater{}} \FunctionTok{onSubmit}\OperatorTok{());}
    \OperatorTok{\}}

    \DataTypeTok{void} \FunctionTok{onSubmit}\OperatorTok{()} \OperatorTok{\{}
        \DataTypeTok{int}\NormalTok{ problemId }\OperatorTok{=}\NormalTok{ view}\OperatorTok{.}\FunctionTok{getSelectedProblemId}\OperatorTok{();}
        \ControlFlowTok{if} \OperatorTok{(}\NormalTok{problemId }\OperatorTok{\textless{}} \DecValTok{0}\OperatorTok{)} \OperatorTok{\{}
            \BuiltInTok{JOptionPane}\OperatorTok{.}\FunctionTok{showMessageDialog}\OperatorTok{(}\NormalTok{view}\OperatorTok{,} \StringTok{"请先选择一道题目。"}\OperatorTok{,}
                \StringTok{"提示"}\OperatorTok{,} \BuiltInTok{JOptionPane}\OperatorTok{.}\FunctionTok{WARNING\_MESSAGE}\OperatorTok{);}
            \ControlFlowTok{return}\OperatorTok{;}
        \OperatorTok{\}}

        \BuiltInTok{String}\NormalTok{ src }\OperatorTok{=}\NormalTok{ view}\OperatorTok{.}\FunctionTok{getSourceCode}\OperatorTok{().}\FunctionTok{trim}\OperatorTok{();}
        \ControlFlowTok{if} \OperatorTok{(}\NormalTok{src}\OperatorTok{.}\FunctionTok{isEmpty}\OperatorTok{())} \OperatorTok{\{}
            \BuiltInTok{JOptionPane}\OperatorTok{.}\FunctionTok{showMessageDialog}\OperatorTok{(}\NormalTok{view}\OperatorTok{,} \StringTok{"源码不能为空。"}\OperatorTok{,}
                \StringTok{"提示"}\OperatorTok{,} \BuiltInTok{JOptionPane}\OperatorTok{.}\FunctionTok{WARNING\_MESSAGE}\OperatorTok{);}
            \ControlFlowTok{return}\OperatorTok{;}
        \OperatorTok{\}}

        \BuiltInTok{String}\NormalTok{ lang }\OperatorTok{=}\NormalTok{ view}\OperatorTok{.}\FunctionTok{getSelectedLang}\OperatorTok{();}
\NormalTok{        view}\OperatorTok{.}\FunctionTok{getSubmitBtn}\OperatorTok{().}\FunctionTok{setEnabled}\OperatorTok{(}\KeywordTok{false}\OperatorTok{);}
\NormalTok{        view}\OperatorTok{.}\FunctionTok{getResultLabel}\OperatorTok{().}\FunctionTok{setText}\OperatorTok{(}\StringTok{"评测中…"}\OperatorTok{);}

        \CommentTok{// 后台线程：写源文件 {-}\textgreater{} 调 MachineJudge {-}\textgreater{} 入库 {-}\textgreater{} 回 EDT 更新 UI}
        \KeywordTok{new} \BuiltInTok{Thread}\OperatorTok{(()} \OperatorTok{{-}\textgreater{}} \OperatorTok{\{}
\NormalTok{            JudgeResult result }\OperatorTok{=} \KeywordTok{null}\OperatorTok{;}
            \DataTypeTok{int}\NormalTok{ subId }\OperatorTok{=} \OperatorTok{{-}}\DecValTok{1}\OperatorTok{;}
            \ControlFlowTok{try} \OperatorTok{\{}
                \CommentTok{// 1. 落源文件（保留路径，存进 Submission.srcPath，不立即删除）}
                \BuiltInTok{String}\NormalTok{ ext }\OperatorTok{=}\NormalTok{ lang}\OperatorTok{.}\FunctionTok{equals}\OperatorTok{(}\StringTok{"cpp"}\OperatorTok{)} \OperatorTok{?} \StringTok{".cpp"} \OperatorTok{:} \StringTok{".py"}\OperatorTok{;}
\NormalTok{                Path   tmp }\OperatorTok{=}\NormalTok{ Files}\OperatorTok{.}\FunctionTok{createTempFile}\OperatorTok{(}\StringTok{"oj\_submit\_"}\OperatorTok{,}\NormalTok{ ext}\OperatorTok{);}
\NormalTok{                Files}\OperatorTok{.}\FunctionTok{write}\OperatorTok{(}\NormalTok{tmp}\OperatorTok{,}\NormalTok{ src}\OperatorTok{.}\FunctionTok{getBytes}\OperatorTok{(}\NormalTok{StandardCharsets}\OperatorTok{.}\FunctionTok{UTF\_8}\OperatorTok{));}

                \CommentTok{// 2. 获取题目元数据（时间限制）}
\NormalTok{                ProblemMeta meta }\OperatorTok{=}\NormalTok{ svc}\OperatorTok{.}\FunctionTok{meta}\OperatorTok{(}\NormalTok{problemId}\OperatorTok{);}
                \DataTypeTok{long}\NormalTok{ timeMs }\OperatorTok{=}\NormalTok{ meta}\OperatorTok{.}\FunctionTok{getTimeLimitMs}\OperatorTok{();}
                \DataTypeTok{int}\NormalTok{  memMb  }\OperatorTok{=} \DecValTok{256}\OperatorTok{;}   \CommentTok{// 默认内存限制}

                \CommentTok{// 3. 调用第三代 MachineJudge（本地进程，无网络）}
                \BuiltInTok{String}\NormalTok{ problemDir }\OperatorTok{=}\NormalTok{ PROBLEMS\_DIR }\OperatorTok{+} \StringTok{"/"} \OperatorTok{+}\NormalTok{ problemId}\OperatorTok{;}
\NormalTok{                result }\OperatorTok{=}\NormalTok{ judge}\OperatorTok{.}\FunctionTok{judge}\OperatorTok{(}\NormalTok{problemDir}\OperatorTok{,}\NormalTok{ tmp}\OperatorTok{.}\FunctionTok{toString}\OperatorTok{(),}
\NormalTok{                    lang}\OperatorTok{,}\NormalTok{ timeMs}\OperatorTok{,}\NormalTok{ memMb}\OperatorTok{);}

                \CommentTok{// 4. 构造提交记录并持久化到数据库}
                \CommentTok{//    save() 返回 void，提交 ID 由 Submission 自增，存库后用 getId() 取回}
\NormalTok{                Submission sub }\OperatorTok{=} \KeywordTok{new} \FunctionTok{Submission}\OperatorTok{(}
\NormalTok{                    problemId}\OperatorTok{,}\NormalTok{ USER\_NAME}\OperatorTok{,}\NormalTok{ lang}\OperatorTok{,}\NormalTok{ tmp}\OperatorTok{.}\FunctionTok{toString}\OperatorTok{());}
\NormalTok{                sub}\OperatorTok{.}\FunctionTok{setResult}\OperatorTok{(}\NormalTok{result}\OperatorTok{);}
\NormalTok{                dao}\OperatorTok{.}\FunctionTok{save}\OperatorTok{(}\NormalTok{sub}\OperatorTok{);}
\NormalTok{                subId }\OperatorTok{=}\NormalTok{ sub}\OperatorTok{.}\FunctionTok{getId}\OperatorTok{();}

            \OperatorTok{\}} \ControlFlowTok{catch} \OperatorTok{(}\BuiltInTok{IOException} \OperatorTok{|} \BuiltInTok{InterruptedException}\NormalTok{ ex}\OperatorTok{)} \OperatorTok{\{}
                \BuiltInTok{Thread}\OperatorTok{.}\FunctionTok{currentThread}\OperatorTok{().}\FunctionTok{interrupt}\OperatorTok{();}
                \DataTypeTok{final} \BuiltInTok{String}\NormalTok{ msg }\OperatorTok{=}\NormalTok{ ex}\OperatorTok{.}\FunctionTok{getMessage}\OperatorTok{();}
                \BuiltInTok{SwingUtilities}\OperatorTok{.}\FunctionTok{invokeLater}\OperatorTok{(()} \OperatorTok{{-}\textgreater{}} \OperatorTok{\{}
\NormalTok{                    view}\OperatorTok{.}\FunctionTok{getResultLabel}\OperatorTok{().}\FunctionTok{setText}\OperatorTok{(}\StringTok{"提交出错"}\OperatorTok{);}
\NormalTok{                    view}\OperatorTok{.}\FunctionTok{getSubmitBtn}\OperatorTok{().}\FunctionTok{setEnabled}\OperatorTok{(}\KeywordTok{true}\OperatorTok{);}
                    \BuiltInTok{JOptionPane}\OperatorTok{.}\FunctionTok{showMessageDialog}\OperatorTok{(}\NormalTok{view}\OperatorTok{,}
                        \StringTok{"提交出错："} \OperatorTok{+}\NormalTok{ msg}\OperatorTok{,} \StringTok{"错误"}\OperatorTok{,}
                        \BuiltInTok{JOptionPane}\OperatorTok{.}\FunctionTok{ERROR\_MESSAGE}\OperatorTok{);}
                \OperatorTok{\});}
                \ControlFlowTok{return}\OperatorTok{;}
            \OperatorTok{\}}

            \CommentTok{// 回到 EDT 更新界面}
            \DataTypeTok{final}\NormalTok{ JudgeResult finalResult }\OperatorTok{=}\NormalTok{ result}\OperatorTok{;}
            \DataTypeTok{final} \DataTypeTok{int}\NormalTok{         finalSubId  }\OperatorTok{=}\NormalTok{ subId}\OperatorTok{;}
            \BuiltInTok{SwingUtilities}\OperatorTok{.}\FunctionTok{invokeLater}\OperatorTok{(()} \OperatorTok{{-}\textgreater{}} \OperatorTok{\{}
\NormalTok{                view}\OperatorTok{.}\FunctionTok{getSubmitBtn}\OperatorTok{().}\FunctionTok{setEnabled}\OperatorTok{(}\KeywordTok{true}\OperatorTok{);}
\NormalTok{                view}\OperatorTok{.}\FunctionTok{getResultLabel}\OperatorTok{().}\FunctionTok{setText}\OperatorTok{(}\NormalTok{finalResult}\OperatorTok{.}\FunctionTok{toString}\OperatorTok{());}

                \CommentTok{// 追加历史表格一行}
\NormalTok{                view}\OperatorTok{.}\FunctionTok{appendHistory}\OperatorTok{(}
\NormalTok{                    finalSubId}\OperatorTok{,}\NormalTok{ problemId}\OperatorTok{,}\NormalTok{ lang}\OperatorTok{,}
\NormalTok{                    finalResult}\OperatorTok{.}\FunctionTok{getStatus}\OperatorTok{().}\FunctionTok{name}\OperatorTok{(),}
\NormalTok{                    finalResult}\OperatorTok{.}\FunctionTok{getPassed}\OperatorTok{()} \OperatorTok{+} \StringTok{"/"} \OperatorTok{+}\NormalTok{ finalResult}\OperatorTok{.}\FunctionTok{getTotal}\OperatorTok{(),}
\NormalTok{                    finalResult}\OperatorTok{.}\FunctionTok{getElapsedMs}\OperatorTok{());}

                \CommentTok{// 弹窗显示判题结果}
                \DataTypeTok{int}\NormalTok{ msgType }\OperatorTok{=}\NormalTok{ finalResult}\OperatorTok{.}\FunctionTok{isAccepted}\OperatorTok{()}
                    \OperatorTok{?} \BuiltInTok{JOptionPane}\OperatorTok{.}\FunctionTok{INFORMATION\_MESSAGE}
                    \OperatorTok{:} \BuiltInTok{JOptionPane}\OperatorTok{.}\FunctionTok{WARNING\_MESSAGE}\OperatorTok{;}
                \BuiltInTok{JOptionPane}\OperatorTok{.}\FunctionTok{showMessageDialog}\OperatorTok{(}\NormalTok{view}\OperatorTok{,}
\NormalTok{                    finalResult}\OperatorTok{.}\FunctionTok{toString}\OperatorTok{()} \OperatorTok{+} \StringTok{"}\SpecialCharTok{\textbackslash{}n}\StringTok{"} \OperatorTok{+}\NormalTok{ finalResult}\OperatorTok{.}\FunctionTok{getDetail}\OperatorTok{(),}
                    \StringTok{"判题结果"}\OperatorTok{,}\NormalTok{ msgType}\OperatorTok{);}
            \OperatorTok{\});}
        \OperatorTok{\}).}\FunctionTok{start}\OperatorTok{();}
    \OperatorTok{\}}
\OperatorTok{\}}
\end{Highlighting}
\end{Shaded}

\begin{center}\rule{0.5\linewidth}{0.5pt}\end{center}

\hypertarget{ux7a0bux5e8fux5165ux53e3main-ux7c7b}{%
\subsubsection{程序入口（Main
类）}\label{ux7a0bux5e8fux5165ux53e3main-ux7c7b}}

Main 按全局契约组装 Model 层：\texttt{ProblemDao} /
\texttt{SubmissionDao} 都是\textbf{无参构造}（内部用 \texttt{Db.get()}
取单例连接，没有
\texttt{Db.connect()}），\texttt{ProblemService(ProblemDao,\ ProblemLoader)}、\texttt{MachineJudge(String\ judgeBinary)}
都按契约构造。所有 Swing 操作放进 \texttt{SwingUtilities.invokeLater}。

\begin{Shaded}
\begin{Highlighting}[]
\KeywordTok{package}\ImportTok{ oj}\OperatorTok{.}\ImportTok{gui}\OperatorTok{;}

\KeywordTok{import} \ImportTok{oj}\OperatorTok{.}\ImportTok{db}\OperatorTok{.}\ImportTok{ProblemDao}\OperatorTok{;}
\KeywordTok{import} \ImportTok{oj}\OperatorTok{.}\ImportTok{db}\OperatorTok{.}\ImportTok{SubmissionDao}\OperatorTok{;}
\KeywordTok{import} \ImportTok{oj}\OperatorTok{.}\ImportTok{io}\OperatorTok{.}\ImportTok{ProblemLoader}\OperatorTok{;}
\KeywordTok{import} \ImportTok{oj}\OperatorTok{.}\ImportTok{judge}\OperatorTok{.}\ImportTok{MachineJudge}\OperatorTok{;}
\KeywordTok{import} \ImportTok{oj}\OperatorTok{.}\ImportTok{service}\OperatorTok{.}\ImportTok{ProblemService}\OperatorTok{;}

\KeywordTok{import} \ImportTok{javax}\OperatorTok{.}\ImportTok{swing}\OperatorTok{.*;}

\KeywordTok{public} \KeywordTok{class}\NormalTok{ Main }\OperatorTok{\{}

    \KeywordTok{public} \DataTypeTok{static} \DataTypeTok{void} \FunctionTok{main}\OperatorTok{(}\BuiltInTok{String}\OperatorTok{[]}\NormalTok{ args}\OperatorTok{)} \OperatorTok{\{}
        \CommentTok{// 初始化 Model 层（DAO 无参构造，内部用 Db.get() 单例连接）}
\NormalTok{        ProblemService svc   }\OperatorTok{=} \KeywordTok{new} \FunctionTok{ProblemService}\OperatorTok{(}\KeywordTok{new} \FunctionTok{ProblemDao}\OperatorTok{(),}
                                                  \KeywordTok{new} \FunctionTok{ProblemLoader}\OperatorTok{(}\StringTok{"problems"}\OperatorTok{));}
\NormalTok{        MachineJudge   judge }\OperatorTok{=} \KeywordTok{new} \FunctionTok{MachineJudge}\OperatorTok{(}\StringTok{"judge/judge"}\OperatorTok{);}
\NormalTok{        SubmissionDao  dao   }\OperatorTok{=} \KeywordTok{new} \FunctionTok{SubmissionDao}\OperatorTok{();}

        \CommentTok{// 所有 Swing 操作必须在 EDT 上执行}
        \BuiltInTok{SwingUtilities}\OperatorTok{.}\FunctionTok{invokeLater}\OperatorTok{(()} \OperatorTok{{-}\textgreater{}} \OperatorTok{\{}
\NormalTok{            OjFrame f }\OperatorTok{=} \KeywordTok{new} \FunctionTok{OjFrame}\OperatorTok{();}
            \KeywordTok{new} \FunctionTok{OjController}\OperatorTok{(}\NormalTok{f}\OperatorTok{,}\NormalTok{ svc}\OperatorTok{,}\NormalTok{ judge}\OperatorTok{,}\NormalTok{ dao}\OperatorTok{);}
\NormalTok{            f}\OperatorTok{.}\FunctionTok{setVisible}\OperatorTok{(}\KeywordTok{true}\OperatorTok{);}
        \OperatorTok{\});}
    \OperatorTok{\}}
\OperatorTok{\}}
\end{Highlighting}
\end{Shaded}

\begin{center}\rule{0.5\linewidth}{0.5pt}\end{center}

\hypertarget{ux5173ux952eux4ea4ux4e92ux65f6ux5e8fux6587ux5b57ux7248}{%
\subsubsection{关键交互时序（文字版）}\label{ux5173ux952eux4ea4ux4e92ux65f6ux5e8fux6587ux5b57ux7248}}

\begin{verbatim}
[EDT] 用户点击"提交"
  -> OjController.onSubmit()
      -> 校验输入（EDT）
      -> 禁用按钮 + 状态标签显示"评测中…"（EDT）
      -> 启动新线程（非 EDT）
          -> Files.createTempFile / Files.write  <- 写源文件（路径留作 srcPath）
          -> svc.meta(problemId)                 <- 取题目时限
          -> judge.judge(...)                    <- ProcessBuilder，第三代
          -> new Submission(...) + setResult     <- 构造提交记录
          -> dao.save(sub); subId = sub.getId()  <- JDBC 持久化（save 返回 void）
          -> SwingUtilities.invokeLater(...)     <- 回 EDT
              -> 状态标签写入结果文本
              -> view.appendHistory(...)         <- 更新表格
              -> JOptionPane.showMessageDialog   <- 弹窗
              -> 恢复按钮
\end{verbatim}

\textbf{没有任何 Socket/ServerSocket/网络调用}。MachineJudge 通过
\texttt{ProcessBuilder} 启动本地可执行文件（系统进程），SubmissionDao
通过 JDBC 访问本地数据库，两者均为进程内的本地调用。

\begin{center}\rule{0.5\linewidth}{0.5pt}\end{center}

\hypertarget{ux52a8ux624bux5b9eux73b0ux4e24ux6b65ux8d70javac-java-ux5b9eux64cd-5}{%
\subsection{动手实现:两步走(javac / java
实操)}\label{ux52a8ux624bux5b9eux73b0ux4e24ux6b65ux8d70javac-java-ux5b9eux64cd-5}}

\hypertarget{ux7b2cux4e00ux6b65-ux6700ux7b80-oj-ux63d0ux4ea4ux7a97ux53e3ux5355ux6587ux4ef6}{%
\subsubsection{\texorpdfstring{第一步 \KS --- 最简 OJ
提交窗口(单文件)}{第一步 --- 最简 OJ 提交窗口(单文件)}}\label{ux7b2cux4e00ux6b65-ux6700ux7b80-oj-ux63d0ux4ea4ux7a97ux53e3ux5355ux6587ux4ef6}}

\textbf{\texttt{src/oj/gui/SimpleOj.java}}

\begin{Shaded}
\begin{Highlighting}[]
\KeywordTok{package}\ImportTok{ oj}\OperatorTok{.}\ImportTok{gui}\OperatorTok{;}
\KeywordTok{import} \ImportTok{javax}\OperatorTok{.}\ImportTok{swing}\OperatorTok{.*;}
\KeywordTok{import} \ImportTok{java}\OperatorTok{.}\ImportTok{awt}\OperatorTok{.*;}
\KeywordTok{public} \KeywordTok{class}\NormalTok{ SimpleOj }\OperatorTok{\{}
    \KeywordTok{public} \DataTypeTok{static} \DataTypeTok{void} \FunctionTok{main}\OperatorTok{(}\BuiltInTok{String}\OperatorTok{[]}\NormalTok{ args}\OperatorTok{)} \OperatorTok{\{}
        \BuiltInTok{SwingUtilities}\OperatorTok{.}\FunctionTok{invokeLater}\OperatorTok{(()} \OperatorTok{{-}\textgreater{}} \OperatorTok{\{}               \CommentTok{// 在 EDT 上建界面}
            \BuiltInTok{JFrame}\NormalTok{ f }\OperatorTok{=} \KeywordTok{new} \BuiltInTok{JFrame}\OperatorTok{(}\StringTok{"Mini{-}OJ 提交"}\OperatorTok{);}
\NormalTok{            f}\OperatorTok{.}\FunctionTok{setLayout}\OperatorTok{(}\KeywordTok{new} \BuiltInTok{FlowLayout}\OperatorTok{());}
\NormalTok{            f}\OperatorTok{.}\FunctionTok{setSize}\OperatorTok{(}\DecValTok{380}\OperatorTok{,} \DecValTok{150}\OperatorTok{);}
\NormalTok{            f}\OperatorTok{.}\FunctionTok{setDefaultCloseOperation}\OperatorTok{(}\BuiltInTok{JFrame}\OperatorTok{.}\FunctionTok{EXIT\_ON\_CLOSE}\OperatorTok{);}
            \BuiltInTok{JComboBox}\OperatorTok{\textless{}}\BuiltInTok{String}\OperatorTok{\textgreater{}}\NormalTok{ prob }\OperatorTok{=} \KeywordTok{new} \BuiltInTok{JComboBox}\OperatorTok{\textless{}\textgreater{}(}\KeywordTok{new} \BuiltInTok{String}\OperatorTok{[]\{} \StringTok{"1 A+B"}\OperatorTok{,} \StringTok{"2 平方"} \OperatorTok{\});}
            \BuiltInTok{JTextField}\NormalTok{ input }\OperatorTok{=} \KeywordTok{new} \BuiltInTok{JTextField}\OperatorTok{(}\DecValTok{12}\OperatorTok{);}
            \BuiltInTok{JButton}\NormalTok{ submit }\OperatorTok{=} \KeywordTok{new} \BuiltInTok{JButton}\OperatorTok{(}\StringTok{"提交"}\OperatorTok{);}
            \BuiltInTok{JLabel}\NormalTok{ result }\OperatorTok{=} \KeywordTok{new} \BuiltInTok{JLabel}\OperatorTok{(}\StringTok{"等待提交"}\OperatorTok{);}
\NormalTok{            submit}\OperatorTok{.}\FunctionTok{addActionListener}\OperatorTok{(}\NormalTok{e }\OperatorTok{{-}\textgreater{}}                \CommentTok{// Lambda 监听 ActionEvent}
\NormalTok{                result}\OperatorTok{.}\FunctionTok{setText}\OperatorTok{(}\StringTok{"已提交["} \OperatorTok{+}\NormalTok{ prob}\OperatorTok{.}\FunctionTok{getSelectedItem}\OperatorTok{()} \OperatorTok{+} \StringTok{"] 输出="} \OperatorTok{+}\NormalTok{ input}\OperatorTok{.}\FunctionTok{getText}\OperatorTok{()));}
\NormalTok{            f}\OperatorTok{.}\FunctionTok{add}\OperatorTok{(}\KeywordTok{new} \BuiltInTok{JLabel}\OperatorTok{(}\StringTok{"题目"}\OperatorTok{));}\NormalTok{ f}\OperatorTok{.}\FunctionTok{add}\OperatorTok{(}\NormalTok{prob}\OperatorTok{);}
\NormalTok{            f}\OperatorTok{.}\FunctionTok{add}\OperatorTok{(}\KeywordTok{new} \BuiltInTok{JLabel}\OperatorTok{(}\StringTok{"输出"}\OperatorTok{));}\NormalTok{ f}\OperatorTok{.}\FunctionTok{add}\OperatorTok{(}\NormalTok{input}\OperatorTok{);}
\NormalTok{            f}\OperatorTok{.}\FunctionTok{add}\OperatorTok{(}\NormalTok{submit}\OperatorTok{);}\NormalTok{ f}\OperatorTok{.}\FunctionTok{add}\OperatorTok{(}\NormalTok{result}\OperatorTok{);}
\NormalTok{            f}\OperatorTok{.}\FunctionTok{setVisible}\OperatorTok{(}\KeywordTok{true}\OperatorTok{);}
        \OperatorTok{\});}
    \OperatorTok{\}}
\OperatorTok{\}}
\end{Highlighting}
\end{Shaded}

讲解:\texttt{JFrame}+\texttt{FlowLayout}+\texttt{JComboBox}(选题)+\texttt{JTextField}(输入)+\texttt{JButton}+\texttt{JLabel}+\texttt{ActionListener}
Lambda------Ch9 全部考点,组成 OJ 的提交窗口。

编译运行(需图形界面环境):

\begin{Shaded}
\begin{Highlighting}[]
\ExtensionTok{javac} \AttributeTok{{-}d}\NormalTok{ build src/oj/gui/SimpleOj.java}
\ExtensionTok{java} \AttributeTok{{-}cp}\NormalTok{ build oj.gui.SimpleOj}
\end{Highlighting}
\end{Shaded}

\hypertarget{ux7b2cux4e8cux6b65-mvc-ux89e3ux8026-ux540eux53f0ux5224ux9898}{%
\subsubsection{\texorpdfstring{第二步 \GC --- MVC 解耦 +
后台判题}{第二步 --- MVC 解耦 + 后台判题}}\label{ux7b2cux4e8cux6b65-mvc-ux89e3ux8026-ux540eux53f0ux5224ux9898}}

拆成
\texttt{OjFrame}(View)/\texttt{OjController}(Controller)/\texttt{ProblemService}+\texttt{MachineJudge}(Model),点提交后用
\texttt{SwingWorker} 在后台调判题机、回 EDT
刷新,界面不卡(代码见上文【新产物架构】):

\begin{Shaded}
\begin{Highlighting}[]
\ExtensionTok{javac} \AttributeTok{{-}cp} \StringTok{"lib/*"} \AttributeTok{{-}d}\NormalTok{ build }\VariableTok{$(}\FunctionTok{find}\NormalTok{ src }\AttributeTok{{-}name} \StringTok{\textquotesingle{}*.java\textquotesingle{}}\VariableTok{)}
\ExtensionTok{java} \AttributeTok{{-}cp} \StringTok{"build:lib/*"}\NormalTok{ oj.gui.Main}
\end{Highlighting}
\end{Shaded}

\hypertarget{ux9a8cux6536ux6807ux51c6-7}{%
\subsection{验收标准}\label{ux9a8cux6536ux6807ux51c6-7}}

\begin{enumerate}
\def\labelenumi{\arabic{enumi}.}
\tightlist
\item
  \textbf{题目下拉}：启动后
  \texttt{JComboBox\textless{}String\textgreater{}}
  自动填充数据库中全部题目（``id 标题''），按 id
  升序，无需手动输入题号；取题号一律走
  \texttt{getSelectedProblemId()}，未选中返回 \texttt{-1}。
\item
  \textbf{语言选择}：\texttt{JComboBox} 包含 \texttt{cpp} 和
  \texttt{python} 两项，选中值传递给 MachineJudge。
\item
  \textbf{源码编辑}：\texttt{JTextArea} 使用等宽字体，支持 Tab
  缩进，可粘贴任意长度源码。
\item
  \textbf{提交不冻结
  UI}：点击提交后按钮立即变灰、状态标签显示''评测中\ldots``，判题期间界面保持可响应；判题完成后按钮恢复。
\item
  \textbf{状态标签}：OjFrame 工具栏里有一个真实的
  \texttt{JLabel}（\texttt{getResultLabel()}
  暴露），提交过程中显示''评测中\ldots``，完成后显示
  \texttt{JudgeResult.toString()}（如 \texttt{AC\ 3/3\ (47ms)}）。
\item
  \textbf{结果弹窗}：判题完成后弹出 \texttt{JOptionPane}，AC 时显示
  \texttt{INFORMATION\_MESSAGE}，非 AC 时显示
  \texttt{WARNING\_MESSAGE}，内容为 \texttt{JudgeResult.toString()} 加
  detail 字段。
\item
  \textbf{历史表格}：每次提交后向 \texttt{JTable} 追加一行，包含提交
  ID、题目 ID、语言、状态、通过/总计、耗时，表格自动滚动到最新行。
\item
  \textbf{持久化}：提交记录通过 \texttt{dao.save(sub)}（返回
  void）写入数据库，提交 ID 用 \texttt{sub.getId()} 取回；源文件路径作为
  \texttt{srcPath} 一并入库（写文件后不立即删除）。
\item
  \textbf{构造契约}：DAO 无参构造（内部
  \texttt{Db.get()}），\texttt{ProblemService(ProblemDao,\ ProblemLoader)}、\texttt{MachineJudge(String)}
  按契约组装；OjController 构造签名为
  \texttt{(OjFrame,\ ProblemService,\ MachineJudge,\ SubmissionDao)}。
\item
  \textbf{无网络依赖}：全程无
  \texttt{Socket}、\texttt{ServerSocket}、\texttt{URL}、\texttt{HttpURLConnection}
  等任何网络 API；MachineJudge 和 SubmissionDao 均为本地调用。
\item
  \textbf{MVC 分离}：\texttt{OjFrame}
  不包含任何业务逻辑，\texttt{OjController} 不直接操作 Swing
  组件样式，\texttt{ProblemService}/\texttt{MachineJudge}/\texttt{SubmissionDao}
  不感知 UI 层。
\item
  \textbf{EDT 合规}：所有 \texttt{view.*} 调用均在 EDT
  上执行；MachineJudge 调用在独立线程；通过
  \texttt{SwingUtilities.invokeLater} 回到 EDT，无线程安全问题。
\end{enumerate}

\hypertarget{m6a-ux591aux7ebfux7a0bux5e76ux53d1ux5224ux9898ch12ux6536ux5b98ux603bux51b3ux8d5b}{%
\section{09 · M6a
多线程并发判题（Ch12，收官总决赛）}\label{m6a-ux591aux7ebfux7a0bux5e76ux53d1ux5224ux9898ch12ux6536ux5b98ux603bux51b3ux8d5b}}

\begin{quote}
\textbf{考试相关度 ★★} ·
真题考点参考:\texttt{extends\ Thread}/run/start、join、synchronized。\textbf{{[}第一步{]}}
\KS 给 OJ 判题加最简并发(\texttt{Thread/Runnable} +
\texttt{synchronized/join});\textbf{{[}第二步{]}} \GC 泛型阻塞队列 +
工作线程池 + 调 C++ 判题机。
\end{quote}

\begin{quote}
至此，Mini-OJ
终于能同时接待多名同学提交代码，而不必让后来者苦苦等待------这一章是整个项目的收官之战。
\end{quote}

\begin{center}\rule{0.5\linewidth}{0.5pt}\end{center}

\hypertarget{ux6838ux5fc3ux75dbux70b9-4}{%
\subsection{【核心痛点】}\label{ux6838ux5fc3ux75dbux70b9-4}}

在 M5 完成之后，系统已经能调用真实的 C++
判题机对外部源码进行编译、运行和评测。然而有两个致命问题依然悬而未决：

\textbf{问题一：串行阻塞。}
如果两位同学同时点击''提交''，第二位提交会被完全阻塞，直到第一份代码评测完毕才能开始。课堂演示时只要稍微多几个人同时操作，界面就会卡住无响应。

\textbf{问题二：死循环挂死系统。} 提交一份 \texttt{while(true)\{\}} 的
Java 代码，在 M3 的纯 JVM 方案里会把整个 JVM 进程卡死。在 M5
虽然调用了外部判题机，但如果 Java 主线程同步等待
\texttt{Process.waitFor()}，同样会造成 GUI 无响应。

解决方案的核心是两步：第一步，把判题任务放进一个\textbf{阻塞队列}，让提交操作立即返回；第二步，用\textbf{工作线程}异步消费队列，每个工作线程调起
C++ 判题机------TLE 由 C++ 端的墙钟看门狗加 \texttt{setrlimit}
掐断死循环并返回 \texttt{TLE} 状态，Java 这侧不再需要外套超时计时器。

\begin{center}\rule{0.5\linewidth}{0.5pt}\end{center}

\hypertarget{ux5f15ux5165ux8bfeux672cux77e5ux8bc6ux70b9-4}{%
\subsection{【引入课本知识点】}\label{ux5f15ux5165ux8bfeux672cux77e5ux8bc6ux70b9-4}}

\hypertarget{ch12thread-ux4e0e-runnable}{%
\subsubsection{Ch12：Thread 与
Runnable}\label{ch12thread-ux4e0e-runnable}}

Java 并发的两种启动方式：

\begin{Shaded}
\begin{Highlighting}[]
\CommentTok{// 方式一：继承 Thread}
\KeywordTok{class}\NormalTok{ MyThread }\KeywordTok{extends} \BuiltInTok{Thread} \OperatorTok{\{}
    \KeywordTok{public} \DataTypeTok{void} \FunctionTok{run}\OperatorTok{()} \OperatorTok{\{} \CommentTok{/* ... */} \OperatorTok{\}}
\OperatorTok{\}}

\CommentTok{// 方式二：实现 Runnable（推荐，保留继承槽位）}
\BuiltInTok{Thread}\NormalTok{ t }\OperatorTok{=} \KeywordTok{new} \BuiltInTok{Thread}\OperatorTok{(}\KeywordTok{new} \FunctionTok{MyRunnable}\OperatorTok{());}
\NormalTok{t}\OperatorTok{.}\FunctionTok{setDaemon}\OperatorTok{(}\KeywordTok{true}\OperatorTok{);} \CommentTok{// 守护线程：主程序退出时不阻止 JVM 关闭}
\NormalTok{t}\OperatorTok{.}\FunctionTok{start}\OperatorTok{();}
\end{Highlighting}
\end{Shaded}

\hypertarget{ch12synchronized-waitnotify}{%
\subsubsection{Ch12：synchronized +
wait/notify}\label{ch12synchronized-waitnotify}}

\texttt{synchronized}
锁住对象监视器，保证同一时刻只有一个线程进入临界区。\texttt{wait()}
让当前线程\textbf{释放锁并挂起}，\texttt{notify()}
唤醒一个等待中的线程（\texttt{notifyAll()} 唤醒全部）。

\begin{Shaded}
\begin{Highlighting}[]
\KeywordTok{synchronized} \OperatorTok{(}\NormalTok{lock}\OperatorTok{)} \OperatorTok{\{}
    \ControlFlowTok{while} \OperatorTok{(}\NormalTok{条件不满足}\OperatorTok{)} \OperatorTok{\{}
\NormalTok{        lock}\OperatorTok{.}\FunctionTok{wait}\OperatorTok{();}   \CommentTok{// 释放锁，挂起}
    \OperatorTok{\}}
    \CommentTok{// 操作共享数据}
\NormalTok{    lock}\OperatorTok{.}\FunctionTok{notifyAll}\OperatorTok{();} \CommentTok{// 唤醒等待方}
\OperatorTok{\}}
\end{Highlighting}
\end{Shaded}

\texttt{wait()} 必须在 \texttt{while} 循环里检查条件，而不是
\texttt{if}------这是防范\textbf{虚假唤醒（spurious
wakeup）}的标准写法。

\hypertarget{ch15ux6cdbux578bux963bux585eux961fux5217}{%
\subsubsection{Ch15：泛型阻塞队列}\label{ch15ux6cdbux578bux963bux585eux961fux5217}}

教材 Ch15 讲泛型容器。我们用泛型让
\texttt{JudgeQueue\textless{}T\textgreater{}} 既能装
\texttt{JudgeTask}，将来也能换其他任务类型，而不需要修改队列本身------这正是\textbf{开闭原则}的实践：对扩展开放，对修改关闭。

\hypertarget{swingworkeredt-ux4e0eux540eux53f0ux7ebfux7a0b}{%
\subsubsection{SwingWorker：EDT
与后台线程}\label{swingworkeredt-ux4e0eux540eux53f0ux7ebfux7a0b}}

Swing
的所有界面操作必须在\textbf{事件派发线程（EDT）}上执行；而判题调用（ProcessBuilder、文件
IO）必须在\textbf{后台线程}上执行，否则界面冻结。\texttt{SwingWorker\textless{}T,\ V\textgreater{}}
提供了标准的两阶段模型：

\begin{itemize}
\tightlist
\item
  \texttt{doInBackground()}：后台线程执行耗时操作，返回结果。
\item
  \texttt{done()}：EDT 线程取回结果并刷新界面。
\end{itemize}

\begin{center}\rule{0.5\linewidth}{0.5pt}\end{center}

\hypertarget{ux4e09ux4ee3ux6f14ux8fdbux5b9aux4f4d-4}{%
\subsection{【三代演进定位】}\label{ux4e09ux4ee3ux6f14ux8fdbux5b9aux4f4d-4}}

\begin{longtable}[]{@{}
  >{\raggedright\arraybackslash}p{(\columnwidth - 8\tabcolsep) * \real{0.2000}}
  >{\raggedright\arraybackslash}p{(\columnwidth - 8\tabcolsep) * \real{0.2000}}
  >{\raggedright\arraybackslash}p{(\columnwidth - 8\tabcolsep) * \real{0.2000}}
  >{\raggedright\arraybackslash}p{(\columnwidth - 8\tabcolsep) * \real{0.2000}}
  >{\raggedright\arraybackslash}p{(\columnwidth - 8\tabcolsep) * \real{0.2000}}@{}}
\toprule\noalign{}
\begin{minipage}[b]{\linewidth}\raggedright
代次
\end{minipage} & \begin{minipage}[b]{\linewidth}\raggedright
所在里程碑
\end{minipage} & \begin{minipage}[b]{\linewidth}\raggedright
判题执行方式
\end{minipage} & \begin{minipage}[b]{\linewidth}\raggedright
并发能力
\end{minipage} & \begin{minipage}[b]{\linewidth}\raggedright
超时控制
\end{minipage} \\
\midrule\noalign{}
\endhead
\bottomrule\noalign{}
\endlastfoot
第一代 & M1--M3 & Java 对象 \texttt{Solution.solve()} 在 JVM 内直接调用
& 无，全程串行 & 无法控制死循环 \\
第二代 & M4 & 反射工厂 \texttt{JudgeFactory} 动态加载 \texttt{Judge}
实现 & 无，全程串行 & 无法控制死循环 \\
第三代（串行） & M5a & \texttt{MachineJudge} 通过
\texttt{ProcessBuilder} 调外部 C++ 判题机 & 无，\texttt{waitFor()}
同步阻塞 & C++ 侧 \texttt{setrlimit} 控制 \\
\textbf{第三代（并发化）} & \textbf{M6a（本章）} & \textbf{同
M5a，但由工作线程异步消费 \texttt{JudgeQueue}} & \textbf{N
个守护工作线程并行} & \textbf{C++ 墙钟看门狗 + \texttt{setrlimit}，Java
不外套超时} \\
\end{longtable}

本章是第三代的\textbf{并发化收官}：判题执行逻辑不变，变的是调度层------从单线程同步调用，升级为生产者-消费者队列驱动的多线程并发调用。

\begin{center}\rule{0.5\linewidth}{0.5pt}\end{center}

\hypertarget{ux65b0ux4ea7ux7269ux67b6ux6784-4}{%
\subsection{【新产物架构】}\label{ux65b0ux4ea7ux7269ux67b6ux6784-4}}

本章新增或改造以下类，均位于全局契约规定的包路径下。

\hypertarget{ux7c7bux5173ux7cfbux603bux89c8-1}{%
\subsubsection{类关系总览}\label{ux7c7bux5173ux7cfbux603bux89c8-1}}

\begin{verbatim}
OjController (生产者)
    | onSubmit()
    |  +- SwingWorker.doInBackground()
    |       +- JudgeQueue.put(task)      <--- 返回立即，不阻塞 EDT
    |
JudgeQueue<JudgeTask>  (阻塞队列，synchronized + wait/notify)
    |
JudgeWorker (消费者，Runnable，守护线程)
    |  run() 死循环
    |   +- JudgeQueue.take()             <--- 空时挂起
    |   +- 把源码字节写到持久文件
    |   +- MachineJudge.judge(...)       <--- 调 C++ 判题机
    |   +- new Submission(...) + save    <--- 入库
    |   +- JudgeTask.complete(result)    <--- 回填，唤醒等待方
    |
JudgeTask (工作项，携带源码 + synchronized 的 complete/await)
    +- SwingWorker.done()  <--- EDT 刷新界面
\end{verbatim}

\begin{center}\rule{0.5\linewidth}{0.5pt}\end{center}

\hypertarget{oj.core.judgetask}{%
\subsubsection{oj.core.JudgeTask}\label{oj.core.judgetask}}

携带一次提交的全部信息，同时充当\textbf{Future}的简化版：\texttt{await()}
挂起调用者直到评测完成，\texttt{complete()} 回填结果并唤醒。

\begin{Shaded}
\begin{Highlighting}[]
\KeywordTok{package}\ImportTok{ oj}\OperatorTok{.}\ImportTok{core}\OperatorTok{;}

\KeywordTok{import} \ImportTok{java}\OperatorTok{.}\ImportTok{io}\OperatorTok{.}\ImportTok{Serializable}\OperatorTok{;}

\KeywordTok{public} \KeywordTok{class}\NormalTok{ JudgeTask }\KeywordTok{implements} \BuiltInTok{Serializable} \OperatorTok{\{}
    \KeywordTok{private} \DataTypeTok{static} \DataTypeTok{final} \DataTypeTok{long}\NormalTok{ serialVersionUID }\OperatorTok{=} \DecValTok{1L}\OperatorTok{;}

    \KeywordTok{private} \DataTypeTok{final} \DataTypeTok{int}\NormalTok{ submissionId}\OperatorTok{;}
    \KeywordTok{private} \DataTypeTok{final} \DataTypeTok{int}\NormalTok{ problemId}\OperatorTok{;}
    \KeywordTok{private} \DataTypeTok{final} \BuiltInTok{String}\NormalTok{ userName}\OperatorTok{;}
    \KeywordTok{private} \DataTypeTok{final} \BuiltInTok{String}\NormalTok{ lang}\OperatorTok{;}
    \KeywordTok{private} \DataTypeTok{final} \DataTypeTok{byte}\OperatorTok{[]}\NormalTok{ source}\OperatorTok{;}      \CommentTok{// 源码字节，避免依赖临时文件路径}
    \KeywordTok{private}\NormalTok{ JudgeResult result}\OperatorTok{;}
    \KeywordTok{private} \DataTypeTok{boolean}\NormalTok{ done}\OperatorTok{;}

    \KeywordTok{public} \FunctionTok{JudgeTask}\OperatorTok{(}\DataTypeTok{int}\NormalTok{ submissionId}\OperatorTok{,} \DataTypeTok{int}\NormalTok{ problemId}\OperatorTok{,}
                     \BuiltInTok{String}\NormalTok{ userName}\OperatorTok{,} \BuiltInTok{String}\NormalTok{ lang}\OperatorTok{,} \DataTypeTok{byte}\OperatorTok{[]}\NormalTok{ source}\OperatorTok{)} \OperatorTok{\{}
        \KeywordTok{this}\OperatorTok{.}\FunctionTok{submissionId} \OperatorTok{=}\NormalTok{ submissionId}\OperatorTok{;}
        \KeywordTok{this}\OperatorTok{.}\FunctionTok{problemId}    \OperatorTok{=}\NormalTok{ problemId}\OperatorTok{;}
        \KeywordTok{this}\OperatorTok{.}\FunctionTok{userName}     \OperatorTok{=}\NormalTok{ userName}\OperatorTok{;}
        \KeywordTok{this}\OperatorTok{.}\FunctionTok{lang}         \OperatorTok{=}\NormalTok{ lang}\OperatorTok{;}
        \KeywordTok{this}\OperatorTok{.}\FunctionTok{source}       \OperatorTok{=}\NormalTok{ source}\OperatorTok{;}
        \KeywordTok{this}\OperatorTok{.}\FunctionTok{done}         \OperatorTok{=} \KeywordTok{false}\OperatorTok{;}
    \OperatorTok{\}}

    \CommentTok{/**} \CommentTok{工作线程调用：回填结果，唤醒所有}\NormalTok{ await}\CommentTok{()} \CommentTok{调用者} \CommentTok{*/}
    \KeywordTok{public} \KeywordTok{synchronized} \DataTypeTok{void} \FunctionTok{complete}\OperatorTok{(}\NormalTok{JudgeResult r}\OperatorTok{)} \OperatorTok{\{}
        \KeywordTok{this}\OperatorTok{.}\FunctionTok{result} \OperatorTok{=}\NormalTok{ r}\OperatorTok{;}
        \KeywordTok{this}\OperatorTok{.}\FunctionTok{done}   \OperatorTok{=} \KeywordTok{true}\OperatorTok{;}
        \FunctionTok{notifyAll}\OperatorTok{();}
    \OperatorTok{\}}

    \CommentTok{/**} \CommentTok{提交方调用：阻塞到评测完成后返回结果} \CommentTok{*/}
    \KeywordTok{public} \KeywordTok{synchronized}\NormalTok{ JudgeResult }\FunctionTok{await}\OperatorTok{()} \KeywordTok{throws} \BuiltInTok{InterruptedException} \OperatorTok{\{}
        \ControlFlowTok{while} \OperatorTok{(!}\NormalTok{done}\OperatorTok{)} \OperatorTok{\{}
            \FunctionTok{wait}\OperatorTok{();}
        \OperatorTok{\}}
        \ControlFlowTok{return}\NormalTok{ result}\OperatorTok{;}
    \OperatorTok{\}}

    \KeywordTok{public} \DataTypeTok{int} \FunctionTok{getSubmissionId}\OperatorTok{()} \OperatorTok{\{} \ControlFlowTok{return}\NormalTok{ submissionId}\OperatorTok{;} \OperatorTok{\}}
    \KeywordTok{public} \DataTypeTok{int} \FunctionTok{getProblemId}\OperatorTok{()}    \OperatorTok{\{} \ControlFlowTok{return}\NormalTok{ problemId}\OperatorTok{;} \OperatorTok{\}}
    \KeywordTok{public} \BuiltInTok{String} \FunctionTok{getUserName}\OperatorTok{()}  \OperatorTok{\{} \ControlFlowTok{return}\NormalTok{ userName}\OperatorTok{;} \OperatorTok{\}}
    \KeywordTok{public} \BuiltInTok{String} \FunctionTok{getLang}\OperatorTok{()}      \OperatorTok{\{} \ControlFlowTok{return}\NormalTok{ lang}\OperatorTok{;} \OperatorTok{\}}
    \KeywordTok{public} \DataTypeTok{byte}\OperatorTok{[]} \FunctionTok{getSource}\OperatorTok{()}    \OperatorTok{\{} \ControlFlowTok{return}\NormalTok{ source}\OperatorTok{;} \OperatorTok{\}}
    \KeywordTok{public}\NormalTok{ JudgeResult }\FunctionTok{getResult}\OperatorTok{()} \OperatorTok{\{} \ControlFlowTok{return}\NormalTok{ result}\OperatorTok{;} \OperatorTok{\}}
\OperatorTok{\}}
\end{Highlighting}
\end{Shaded}

\begin{center}\rule{0.5\linewidth}{0.5pt}\end{center}

\hypertarget{oj.judge.queue.judgequeue}{%
\subsubsection{oj.judge.queue.JudgeQueue}\label{oj.judge.queue.judgequeue}}

泛型阻塞队列，用 \texttt{LinkedList} 作底层存储，通过
\texttt{synchronized\ +\ wait/notify} 实现生产者-消费者协调。

\begin{Shaded}
\begin{Highlighting}[]
\KeywordTok{package}\ImportTok{ oj}\OperatorTok{.}\ImportTok{judge}\OperatorTok{.}\ImportTok{queue}\OperatorTok{;}

\KeywordTok{import} \ImportTok{java}\OperatorTok{.}\ImportTok{util}\OperatorTok{.}\ImportTok{LinkedList}\OperatorTok{;}
\KeywordTok{import} \ImportTok{java}\OperatorTok{.}\ImportTok{util}\OperatorTok{.}\ImportTok{Queue}\OperatorTok{;}

\KeywordTok{public} \KeywordTok{class}\NormalTok{ JudgeQueue}\OperatorTok{\textless{}}\NormalTok{T}\OperatorTok{\textgreater{}} \OperatorTok{\{}
    \KeywordTok{private} \DataTypeTok{final} \BuiltInTok{Queue}\OperatorTok{\textless{}}\NormalTok{T}\OperatorTok{\textgreater{}}\NormalTok{ queue }\OperatorTok{=} \KeywordTok{new} \BuiltInTok{LinkedList}\OperatorTok{\textless{}\textgreater{}();}
    \KeywordTok{private} \DataTypeTok{final} \DataTypeTok{int}\NormalTok{ capacity}\OperatorTok{;}

    \KeywordTok{public} \FunctionTok{JudgeQueue}\OperatorTok{(}\DataTypeTok{int}\NormalTok{ capacity}\OperatorTok{)} \OperatorTok{\{}
        \KeywordTok{this}\OperatorTok{.}\FunctionTok{capacity} \OperatorTok{=}\NormalTok{ capacity}\OperatorTok{;}
    \OperatorTok{\}}

    \CommentTok{/**} \CommentTok{生产者：队列满时挂起，有空位后入队并唤醒消费者} \CommentTok{*/}
    \KeywordTok{public} \KeywordTok{synchronized} \DataTypeTok{void} \FunctionTok{put}\OperatorTok{(}\NormalTok{T item}\OperatorTok{)} \KeywordTok{throws} \BuiltInTok{InterruptedException} \OperatorTok{\{}
        \ControlFlowTok{while} \OperatorTok{(}\NormalTok{queue}\OperatorTok{.}\FunctionTok{size}\OperatorTok{()} \OperatorTok{\textgreater{}=}\NormalTok{ capacity}\OperatorTok{)} \OperatorTok{\{}
            \FunctionTok{wait}\OperatorTok{();}
        \OperatorTok{\}}
\NormalTok{        queue}\OperatorTok{.}\FunctionTok{offer}\OperatorTok{(}\NormalTok{item}\OperatorTok{);}
        \FunctionTok{notifyAll}\OperatorTok{();}
    \OperatorTok{\}}

    \CommentTok{/**} \CommentTok{消费者：队列空时挂起，有元素后出队并唤醒生产者} \CommentTok{*/}
    \KeywordTok{public} \KeywordTok{synchronized}\NormalTok{ T }\FunctionTok{take}\OperatorTok{()} \KeywordTok{throws} \BuiltInTok{InterruptedException} \OperatorTok{\{}
        \ControlFlowTok{while} \OperatorTok{(}\NormalTok{queue}\OperatorTok{.}\FunctionTok{isEmpty}\OperatorTok{())} \OperatorTok{\{}
            \FunctionTok{wait}\OperatorTok{();}
        \OperatorTok{\}}
\NormalTok{        T item }\OperatorTok{=}\NormalTok{ queue}\OperatorTok{.}\FunctionTok{poll}\OperatorTok{();}
        \FunctionTok{notifyAll}\OperatorTok{();}
        \ControlFlowTok{return}\NormalTok{ item}\OperatorTok{;}
    \OperatorTok{\}}

    \KeywordTok{public} \KeywordTok{synchronized} \DataTypeTok{int} \FunctionTok{size}\OperatorTok{()} \OperatorTok{\{}
        \ControlFlowTok{return}\NormalTok{ queue}\OperatorTok{.}\FunctionTok{size}\OperatorTok{();}
    \OperatorTok{\}}
\OperatorTok{\}}
\end{Highlighting}
\end{Shaded}

\textbf{设计说明：} \texttt{capacity}
参数让队列在请求洪峰时有背压（back-pressure）能力，防止内存无限增长。课堂场景设
64 即可。

\begin{center}\rule{0.5\linewidth}{0.5pt}\end{center}

\hypertarget{oj.judge.queue.judgeworker}{%
\subsubsection{oj.judge.queue.JudgeWorker}\label{oj.judge.queue.judgeworker}}

消费者，实现
\texttt{Runnable}，以守护线程方式运行。核心流程：从队列取任务
-\textgreater{} 把源码字节\textbf{写到持久文件}（如
\texttt{submissions/sub\textless{}id\textgreater{}.\textless{}ext\textgreater{}}，\textbf{不要写完立刻删}，因为这个路径要存进
Submission 的 \texttt{srcPath}）-\textgreater{} 调 \texttt{MachineJudge}
-\textgreater{} \texttt{new\ Submission(...)} + \texttt{setResult} +
\texttt{dao.save(sub)} -\textgreater{} 回填 \texttt{JudgeTask}。

关键契约点：

\begin{itemize}
\tightlist
\item
  不要 \texttt{new\ MachineJudge()} /
  \texttt{new\ SubmissionDao()}------它们从构造器传入。
\item
  \texttt{SubmissionDao.save(Submission)} 返回
  \texttt{void}，\textbf{没有} \texttt{save(int,\ JudgeResult)}
  这种重载；要先构造 \texttt{Submission} 再 \texttt{save}。
\item
  任何异常都要 catch 住，绝不让工作线程崩溃退出；异常时给
  \texttt{task.complete()} 一个 \texttt{RE} 结果。
\end{itemize}

\begin{Shaded}
\begin{Highlighting}[]
\KeywordTok{package}\ImportTok{ oj}\OperatorTok{.}\ImportTok{judge}\OperatorTok{.}\ImportTok{queue}\OperatorTok{;}

\KeywordTok{import} \ImportTok{oj}\OperatorTok{.}\ImportTok{core}\OperatorTok{.}\ImportTok{JudgeResult}\OperatorTok{;}
\KeywordTok{import} \ImportTok{oj}\OperatorTok{.}\ImportTok{core}\OperatorTok{.}\ImportTok{JudgeTask}\OperatorTok{;}
\KeywordTok{import} \ImportTok{oj}\OperatorTok{.}\ImportTok{core}\OperatorTok{.}\ImportTok{Status}\OperatorTok{;}
\KeywordTok{import} \ImportTok{oj}\OperatorTok{.}\ImportTok{core}\OperatorTok{.}\ImportTok{Submission}\OperatorTok{;}
\KeywordTok{import} \ImportTok{oj}\OperatorTok{.}\ImportTok{db}\OperatorTok{.}\ImportTok{SubmissionDao}\OperatorTok{;}
\KeywordTok{import} \ImportTok{oj}\OperatorTok{.}\ImportTok{judge}\OperatorTok{.}\ImportTok{MachineJudge}\OperatorTok{;}

\KeywordTok{import} \ImportTok{java}\OperatorTok{.}\ImportTok{io}\OperatorTok{.}\ImportTok{IOException}\OperatorTok{;}
\KeywordTok{import} \ImportTok{java}\OperatorTok{.}\ImportTok{io}\OperatorTok{.}\ImportTok{OutputStream}\OperatorTok{;}
\KeywordTok{import} \ImportTok{java}\OperatorTok{.}\ImportTok{nio}\OperatorTok{.}\ImportTok{file}\OperatorTok{.}\ImportTok{Files}\OperatorTok{;}
\KeywordTok{import} \ImportTok{java}\OperatorTok{.}\ImportTok{nio}\OperatorTok{.}\ImportTok{file}\OperatorTok{.}\ImportTok{Path}\OperatorTok{;}
\KeywordTok{import} \ImportTok{java}\OperatorTok{.}\ImportTok{nio}\OperatorTok{.}\ImportTok{file}\OperatorTok{.}\ImportTok{Paths}\OperatorTok{;}

\KeywordTok{public} \KeywordTok{class}\NormalTok{ JudgeWorker }\KeywordTok{implements} \BuiltInTok{Runnable} \OperatorTok{\{}
    \KeywordTok{private} \DataTypeTok{static} \DataTypeTok{final} \BuiltInTok{String}\NormalTok{ PROBLEMS\_DIR    }\OperatorTok{=} \StringTok{"problems"}\OperatorTok{;}
    \KeywordTok{private} \DataTypeTok{static} \DataTypeTok{final} \BuiltInTok{String}\NormalTok{ SUBMISSIONS\_DIR }\OperatorTok{=} \StringTok{"submissions"}\OperatorTok{;}

    \KeywordTok{private} \DataTypeTok{final}\NormalTok{ JudgeQueue}\OperatorTok{\textless{}}\NormalTok{JudgeTask}\OperatorTok{\textgreater{}}\NormalTok{ queue}\OperatorTok{;}
    \KeywordTok{private} \DataTypeTok{final}\NormalTok{ SubmissionDao dao}\OperatorTok{;}
    \KeywordTok{private} \DataTypeTok{final}\NormalTok{ MachineJudge machine}\OperatorTok{;}

    \KeywordTok{public} \FunctionTok{JudgeWorker}\OperatorTok{(}\NormalTok{JudgeQueue}\OperatorTok{\textless{}}\NormalTok{JudgeTask}\OperatorTok{\textgreater{}}\NormalTok{ queue}\OperatorTok{,}
\NormalTok{                       SubmissionDao dao}\OperatorTok{,}
\NormalTok{                       MachineJudge machine}\OperatorTok{)} \OperatorTok{\{}
        \KeywordTok{this}\OperatorTok{.}\FunctionTok{queue}   \OperatorTok{=}\NormalTok{ queue}\OperatorTok{;}
        \KeywordTok{this}\OperatorTok{.}\FunctionTok{dao}     \OperatorTok{=}\NormalTok{ dao}\OperatorTok{;}
        \KeywordTok{this}\OperatorTok{.}\FunctionTok{machine} \OperatorTok{=}\NormalTok{ machine}\OperatorTok{;}
    \OperatorTok{\}}

    \AttributeTok{@Override}
    \KeywordTok{public} \DataTypeTok{void} \FunctionTok{run}\OperatorTok{()} \OperatorTok{\{}
        \ControlFlowTok{while} \OperatorTok{(}\KeywordTok{true}\OperatorTok{)} \OperatorTok{\{}                          \CommentTok{// 守护线程，随 JVM 退出}
\NormalTok{            JudgeTask task }\OperatorTok{=} \KeywordTok{null}\OperatorTok{;}
            \ControlFlowTok{try} \OperatorTok{\{}
\NormalTok{                task }\OperatorTok{=}\NormalTok{ queue}\OperatorTok{.}\FunctionTok{take}\OperatorTok{();}            \CommentTok{// 队列空时挂起，不自旋}
\NormalTok{                JudgeResult result }\OperatorTok{=} \FunctionTok{evaluate}\OperatorTok{(}\NormalTok{task}\OperatorTok{);}

                \CommentTok{// 构造提交记录并持久化（save 返回 void，没有 save(id,result) 重载）}
\NormalTok{                Submission sub }\OperatorTok{=} \KeywordTok{new} \FunctionTok{Submission}\OperatorTok{(}
\NormalTok{                    task}\OperatorTok{.}\FunctionTok{getProblemId}\OperatorTok{(),}\NormalTok{ task}\OperatorTok{.}\FunctionTok{getUserName}\OperatorTok{(),}
\NormalTok{                    task}\OperatorTok{.}\FunctionTok{getLang}\OperatorTok{(),} \FunctionTok{srcPathFor}\OperatorTok{(}\NormalTok{task}\OperatorTok{));}
\NormalTok{                sub}\OperatorTok{.}\FunctionTok{setResult}\OperatorTok{(}\NormalTok{result}\OperatorTok{);}
\NormalTok{                dao}\OperatorTok{.}\FunctionTok{save}\OperatorTok{(}\NormalTok{sub}\OperatorTok{);}

\NormalTok{                task}\OperatorTok{.}\FunctionTok{complete}\OperatorTok{(}\NormalTok{result}\OperatorTok{);}
            \OperatorTok{\}} \ControlFlowTok{catch} \OperatorTok{(}\BuiltInTok{InterruptedException}\NormalTok{ e}\OperatorTok{)} \OperatorTok{\{}
                \BuiltInTok{Thread}\OperatorTok{.}\FunctionTok{currentThread}\OperatorTok{().}\FunctionTok{interrupt}\OperatorTok{();}
                \ControlFlowTok{break}\OperatorTok{;}
            \OperatorTok{\}} \ControlFlowTok{catch} \OperatorTok{(}\BuiltInTok{Exception}\NormalTok{ e}\OperatorTok{)} \OperatorTok{\{}
                \CommentTok{// 任何异常都不能让工作线程崩溃退出；回填一个 RE}
\NormalTok{                JudgeResult err }\OperatorTok{=} \KeywordTok{new} \FunctionTok{JudgeResult}\OperatorTok{(}
\NormalTok{                    Status}\OperatorTok{.}\FunctionTok{RE}\OperatorTok{,} \DecValTok{0}\OperatorTok{,} \DecValTok{0}\OperatorTok{,} \StringTok{"Worker error: "} \OperatorTok{+}\NormalTok{ e}\OperatorTok{.}\FunctionTok{getMessage}\OperatorTok{(),} \DecValTok{0L}\OperatorTok{);}
                \ControlFlowTok{if} \OperatorTok{(}\NormalTok{task }\OperatorTok{!=} \KeywordTok{null}\OperatorTok{)} \OperatorTok{\{}
\NormalTok{                    task}\OperatorTok{.}\FunctionTok{complete}\OperatorTok{(}\NormalTok{err}\OperatorTok{);}
                \OperatorTok{\}}
            \OperatorTok{\}}
        \OperatorTok{\}}
    \OperatorTok{\}}

    \KeywordTok{private}\NormalTok{ JudgeResult }\FunctionTok{evaluate}\OperatorTok{(}\NormalTok{JudgeTask task}\OperatorTok{)} \KeywordTok{throws} \BuiltInTok{IOException} \OperatorTok{\{}
        \CommentTok{// 1. 把源码字节写到持久文件（保留，不删除；路径要存进库）}
\NormalTok{        Path dir }\OperatorTok{=}\NormalTok{ Paths}\OperatorTok{.}\FunctionTok{get}\OperatorTok{(}\NormalTok{SUBMISSIONS\_DIR}\OperatorTok{);}
\NormalTok{        Files}\OperatorTok{.}\FunctionTok{createDirectories}\OperatorTok{(}\NormalTok{dir}\OperatorTok{);}
\NormalTok{        Path src }\OperatorTok{=}\NormalTok{ Paths}\OperatorTok{.}\FunctionTok{get}\OperatorTok{(}\FunctionTok{srcPathFor}\OperatorTok{(}\NormalTok{task}\OperatorTok{));}
        \ControlFlowTok{try} \OperatorTok{(}\BuiltInTok{OutputStream}\NormalTok{ out }\OperatorTok{=}\NormalTok{ Files}\OperatorTok{.}\FunctionTok{newOutputStream}\OperatorTok{(}\NormalTok{src}\OperatorTok{))} \OperatorTok{\{}
\NormalTok{            out}\OperatorTok{.}\FunctionTok{write}\OperatorTok{(}\NormalTok{task}\OperatorTok{.}\FunctionTok{getSource}\OperatorTok{());}
        \OperatorTok{\}}

        \CommentTok{// 2. 题目测试点目录}
        \BuiltInTok{String}\NormalTok{ problemDir }\OperatorTok{=}\NormalTok{ PROBLEMS\_DIR }\OperatorTok{+} \StringTok{"/"} \OperatorTok{+}\NormalTok{ task}\OperatorTok{.}\FunctionTok{getProblemId}\OperatorTok{();}

        \CommentTok{// 3. 调外部 C++ 判题机（TLE / MLE 由 C++ 侧 setrlimit 控制）}
        \ControlFlowTok{return}\NormalTok{ machine}\OperatorTok{.}\FunctionTok{judge}\OperatorTok{(}
\NormalTok{            problemDir}\OperatorTok{,}
\NormalTok{            src}\OperatorTok{.}\FunctionTok{toString}\OperatorTok{(),}
\NormalTok{            task}\OperatorTok{.}\FunctionTok{getLang}\OperatorTok{(),}
            \DecValTok{2000L}\OperatorTok{,}   \CommentTok{// 最大墙钟 ms，C++ 判题机内部会按题目配置截断}
            \DecValTok{256}      \CommentTok{// 最大内存 MB}
        \OperatorTok{);}
    \OperatorTok{\}}

    \CommentTok{// 持久化源文件路径：submissions/sub\textless{}id\textgreater{}.\textless{}ext\textgreater{}}
    \KeywordTok{private} \BuiltInTok{String} \FunctionTok{srcPathFor}\OperatorTok{(}\NormalTok{JudgeTask task}\OperatorTok{)} \OperatorTok{\{}
        \ControlFlowTok{return}\NormalTok{ SUBMISSIONS\_DIR }\OperatorTok{+} \StringTok{"/sub"} \OperatorTok{+}\NormalTok{ task}\OperatorTok{.}\FunctionTok{getSubmissionId}\OperatorTok{()}
             \OperatorTok{+} \FunctionTok{langToExt}\OperatorTok{(}\NormalTok{task}\OperatorTok{.}\FunctionTok{getLang}\OperatorTok{());}
    \OperatorTok{\}}

    \KeywordTok{private} \BuiltInTok{String} \FunctionTok{langToExt}\OperatorTok{(}\BuiltInTok{String}\NormalTok{ lang}\OperatorTok{)} \OperatorTok{\{}
        \ControlFlowTok{switch} \OperatorTok{(}\NormalTok{lang}\OperatorTok{.}\FunctionTok{toLowerCase}\OperatorTok{())} \OperatorTok{\{}
            \ControlFlowTok{case} \StringTok{"java"}\OperatorTok{:}   \ControlFlowTok{return} \StringTok{".java"}\OperatorTok{;}
            \ControlFlowTok{case} \StringTok{"python"}\OperatorTok{:} \ControlFlowTok{return} \StringTok{".py"}\OperatorTok{;}
            \ControlFlowTok{case} \StringTok{"c"}\OperatorTok{:}      \ControlFlowTok{return} \StringTok{".c"}\OperatorTok{;}
            \KeywordTok{default}\OperatorTok{:}       \ControlFlowTok{return} \StringTok{".cpp"}\OperatorTok{;}
        \OperatorTok{\}}
    \OperatorTok{\}}
\OperatorTok{\}}
\end{Highlighting}
\end{Shaded}

\begin{center}\rule{0.5\linewidth}{0.5pt}\end{center}

\hypertarget{oj.gui.ojcontroller-swingworker-ux63d0ux4ea4}{%
\subsubsection{oj.gui.OjController --- SwingWorker
提交}\label{oj.gui.ojcontroller-swingworker-ux63d0ux4ea4}}

\texttt{OjController}
是本章改造最重的类。构造签名仍按全局契约：\texttt{(OjFrame\ view,\ ProblemService\ svc,\ MachineJudge\ judge,\ SubmissionDao\ dao)}。构造里做三件事：建
\texttt{JudgeQueue\textless{}JudgeTask\textgreater{}}、启动 N 个守护
\texttt{JudgeWorker}（\textbf{复用传入的
\texttt{dao}、\texttt{judge}，不要 \texttt{new\ MachineJudge()} /
\texttt{new\ SubmissionDao()}}）、用 \texttt{svc.listMeta()}
初始化题目下拉并绑定提交。

提交按钮的响应逻辑从直接调用 \texttt{MachineJudge}，改为：在
\texttt{SwingWorker.doInBackground()} 里把任务放进队列、调
\texttt{await()} 等结果，在 \texttt{done()} 里回到 EDT 刷新
\texttt{getResultLabel()}、追加历史、非 AC 弹窗。题号一律走
\texttt{view.getSelectedProblemId()}，不要把下拉框当
\texttt{JComboBox\textless{}Integer\textgreater{}} 强转。

\begin{Shaded}
\begin{Highlighting}[]
\KeywordTok{package}\ImportTok{ oj}\OperatorTok{.}\ImportTok{gui}\OperatorTok{;}

\KeywordTok{import} \ImportTok{oj}\OperatorTok{.}\ImportTok{core}\OperatorTok{.}\ImportTok{JudgeResult}\OperatorTok{;}
\KeywordTok{import} \ImportTok{oj}\OperatorTok{.}\ImportTok{core}\OperatorTok{.}\ImportTok{JudgeTask}\OperatorTok{;}
\KeywordTok{import} \ImportTok{oj}\OperatorTok{.}\ImportTok{core}\OperatorTok{.}\ImportTok{ProblemMeta}\OperatorTok{;}
\KeywordTok{import} \ImportTok{oj}\OperatorTok{.}\ImportTok{db}\OperatorTok{.}\ImportTok{SubmissionDao}\OperatorTok{;}
\KeywordTok{import} \ImportTok{oj}\OperatorTok{.}\ImportTok{judge}\OperatorTok{.}\ImportTok{MachineJudge}\OperatorTok{;}
\KeywordTok{import} \ImportTok{oj}\OperatorTok{.}\ImportTok{judge}\OperatorTok{.}\ImportTok{queue}\OperatorTok{.}\ImportTok{JudgeQueue}\OperatorTok{;}
\KeywordTok{import} \ImportTok{oj}\OperatorTok{.}\ImportTok{judge}\OperatorTok{.}\ImportTok{queue}\OperatorTok{.}\ImportTok{JudgeWorker}\OperatorTok{;}
\KeywordTok{import} \ImportTok{oj}\OperatorTok{.}\ImportTok{service}\OperatorTok{.}\ImportTok{ProblemService}\OperatorTok{;}

\KeywordTok{import} \ImportTok{javax}\OperatorTok{.}\ImportTok{swing}\OperatorTok{.*;}
\KeywordTok{import} \ImportTok{java}\OperatorTok{.}\ImportTok{nio}\OperatorTok{.}\ImportTok{charset}\OperatorTok{.}\ImportTok{StandardCharsets}\OperatorTok{;}
\KeywordTok{import} \ImportTok{java}\OperatorTok{.}\ImportTok{util}\OperatorTok{.}\ImportTok{ArrayList}\OperatorTok{;}
\KeywordTok{import} \ImportTok{java}\OperatorTok{.}\ImportTok{util}\OperatorTok{.}\ImportTok{List}\OperatorTok{;}
\KeywordTok{import} \ImportTok{java}\OperatorTok{.}\ImportTok{util}\OperatorTok{.}\ImportTok{Map}\OperatorTok{;}
\KeywordTok{import} \ImportTok{java}\OperatorTok{.}\ImportTok{util}\OperatorTok{.}\ImportTok{concurrent}\OperatorTok{.}\ImportTok{atomic}\OperatorTok{.}\ImportTok{AtomicInteger}\OperatorTok{;}

\KeywordTok{public} \KeywordTok{class}\NormalTok{ OjController }\OperatorTok{\{}
    \KeywordTok{private} \DataTypeTok{static} \DataTypeTok{final} \DataTypeTok{int}\NormalTok{ WORKER\_COUNT   }\OperatorTok{=} \DecValTok{4}\OperatorTok{;}
    \KeywordTok{private} \DataTypeTok{static} \DataTypeTok{final} \DataTypeTok{int}\NormalTok{ QUEUE\_CAPACITY }\OperatorTok{=} \DecValTok{64}\OperatorTok{;}
    \KeywordTok{private} \DataTypeTok{static} \DataTypeTok{final} \BuiltInTok{String}\NormalTok{ USER\_NAME   }\OperatorTok{=} \StringTok{"student"}\OperatorTok{;}

    \KeywordTok{private} \DataTypeTok{final}\NormalTok{ OjFrame view}\OperatorTok{;}
    \KeywordTok{private} \DataTypeTok{final}\NormalTok{ ProblemService svc}\OperatorTok{;}
    \KeywordTok{private} \DataTypeTok{final}\NormalTok{ MachineJudge judge}\OperatorTok{;}
    \KeywordTok{private} \DataTypeTok{final}\NormalTok{ SubmissionDao dao}\OperatorTok{;}

    \KeywordTok{private} \DataTypeTok{final}\NormalTok{ JudgeQueue}\OperatorTok{\textless{}}\NormalTok{JudgeTask}\OperatorTok{\textgreater{}}\NormalTok{ queue}\OperatorTok{;}
    \KeywordTok{private} \DataTypeTok{final} \BuiltInTok{AtomicInteger}\NormalTok{ submissionIdGen }\OperatorTok{=} \KeywordTok{new} \BuiltInTok{AtomicInteger}\OperatorTok{(}\DecValTok{1}\OperatorTok{);}

    \KeywordTok{public} \FunctionTok{OjController}\OperatorTok{(}\NormalTok{OjFrame view}\OperatorTok{,}
\NormalTok{                        ProblemService svc}\OperatorTok{,}
\NormalTok{                        MachineJudge judge}\OperatorTok{,}
\NormalTok{                        SubmissionDao dao}\OperatorTok{)} \OperatorTok{\{}
        \KeywordTok{this}\OperatorTok{.}\FunctionTok{view}  \OperatorTok{=}\NormalTok{ view}\OperatorTok{;}
        \KeywordTok{this}\OperatorTok{.}\FunctionTok{svc}   \OperatorTok{=}\NormalTok{ svc}\OperatorTok{;}
        \KeywordTok{this}\OperatorTok{.}\FunctionTok{judge} \OperatorTok{=}\NormalTok{ judge}\OperatorTok{;}
        \KeywordTok{this}\OperatorTok{.}\FunctionTok{dao}   \OperatorTok{=}\NormalTok{ dao}\OperatorTok{;}
        \KeywordTok{this}\OperatorTok{.}\FunctionTok{queue} \OperatorTok{=} \KeywordTok{new}\NormalTok{ JudgeQueue}\OperatorTok{\textless{}\textgreater{}(}\NormalTok{QUEUE\_CAPACITY}\OperatorTok{);}

        \FunctionTok{startWorkers}\OperatorTok{();}
        \FunctionTok{initProblems}\OperatorTok{();}
        \FunctionTok{bindEvents}\OperatorTok{();}
    \OperatorTok{\}}

    \CommentTok{// 启动 N 个守护工作线程，复用传入的 dao / judge（不要 new）}
    \KeywordTok{private} \DataTypeTok{void} \FunctionTok{startWorkers}\OperatorTok{()} \OperatorTok{\{}
        \ControlFlowTok{for} \OperatorTok{(}\DataTypeTok{int}\NormalTok{ i }\OperatorTok{=} \DecValTok{0}\OperatorTok{;}\NormalTok{ i }\OperatorTok{\textless{}}\NormalTok{ WORKER\_COUNT}\OperatorTok{;}\NormalTok{ i}\OperatorTok{++)} \OperatorTok{\{}
            \BuiltInTok{Thread}\NormalTok{ t }\OperatorTok{=} \KeywordTok{new} \BuiltInTok{Thread}\OperatorTok{(}
                \KeywordTok{new} \FunctionTok{JudgeWorker}\OperatorTok{(}\NormalTok{queue}\OperatorTok{,}\NormalTok{ dao}\OperatorTok{,}\NormalTok{ judge}\OperatorTok{),}
                \StringTok{"judge{-}worker{-}"} \OperatorTok{+}\NormalTok{ i}\OperatorTok{);}
\NormalTok{            t}\OperatorTok{.}\FunctionTok{setDaemon}\OperatorTok{(}\KeywordTok{true}\OperatorTok{);}   \CommentTok{// JVM 退出时不等待工作线程}
\NormalTok{            t}\OperatorTok{.}\FunctionTok{start}\OperatorTok{();}
        \OperatorTok{\}}
    \OperatorTok{\}}

    \CommentTok{// 用 svc.listMeta() 填充题目下拉框}
    \KeywordTok{private} \DataTypeTok{void} \FunctionTok{initProblems}\OperatorTok{()} \OperatorTok{\{}
        \ControlFlowTok{try} \OperatorTok{\{}
            \BuiltInTok{Map}\OperatorTok{\textless{}}\BuiltInTok{Integer}\OperatorTok{,}\NormalTok{ ProblemMeta}\OperatorTok{\textgreater{}}\NormalTok{ all }\OperatorTok{=}\NormalTok{ svc}\OperatorTok{.}\FunctionTok{listMeta}\OperatorTok{();}
            \BuiltInTok{List}\OperatorTok{\textless{}}\BuiltInTok{Integer}\OperatorTok{\textgreater{}}\NormalTok{ ids    }\OperatorTok{=} \KeywordTok{new} \BuiltInTok{ArrayList}\OperatorTok{\textless{}\textgreater{}(}\NormalTok{all}\OperatorTok{.}\FunctionTok{keySet}\OperatorTok{());}
            \BuiltInTok{List}\OperatorTok{\textless{}}\BuiltInTok{String}\OperatorTok{\textgreater{}}\NormalTok{  titles }\OperatorTok{=} \KeywordTok{new} \BuiltInTok{ArrayList}\OperatorTok{\textless{}\textgreater{}();}
\NormalTok{            ids}\OperatorTok{.}\FunctionTok{forEach}\OperatorTok{(}\NormalTok{id }\OperatorTok{{-}\textgreater{}}\NormalTok{ titles}\OperatorTok{.}\FunctionTok{add}\OperatorTok{(}\NormalTok{all}\OperatorTok{.}\FunctionTok{get}\OperatorTok{(}\NormalTok{id}\OperatorTok{).}\FunctionTok{getTitle}\OperatorTok{()));}
\NormalTok{            view}\OperatorTok{.}\FunctionTok{loadProblems}\OperatorTok{(}\NormalTok{ids}\OperatorTok{,}\NormalTok{ titles}\OperatorTok{);}
        \OperatorTok{\}} \ControlFlowTok{catch} \OperatorTok{(}\BuiltInTok{Exception}\NormalTok{ ex}\OperatorTok{)} \OperatorTok{\{}
            \BuiltInTok{JOptionPane}\OperatorTok{.}\FunctionTok{showMessageDialog}\OperatorTok{(}\NormalTok{view}\OperatorTok{,}
                \StringTok{"加载题目失败："} \OperatorTok{+}\NormalTok{ ex}\OperatorTok{.}\FunctionTok{getMessage}\OperatorTok{(),} \StringTok{"错误"}\OperatorTok{,}
                \BuiltInTok{JOptionPane}\OperatorTok{.}\FunctionTok{ERROR\_MESSAGE}\OperatorTok{);}
        \OperatorTok{\}}
    \OperatorTok{\}}

    \KeywordTok{private} \DataTypeTok{void} \FunctionTok{bindEvents}\OperatorTok{()} \OperatorTok{\{}
\NormalTok{        view}\OperatorTok{.}\FunctionTok{getSubmitBtn}\OperatorTok{().}\FunctionTok{addActionListener}\OperatorTok{(}\NormalTok{e }\OperatorTok{{-}\textgreater{}} \FunctionTok{onSubmit}\OperatorTok{());}
    \OperatorTok{\}}

    \KeywordTok{private} \DataTypeTok{void} \FunctionTok{onSubmit}\OperatorTok{()} \OperatorTok{\{}
        \CommentTok{// 1. 从界面读取输入（在 EDT 上执行，安全）}
        \DataTypeTok{int}\NormalTok{ problemId }\OperatorTok{=}\NormalTok{ view}\OperatorTok{.}\FunctionTok{getSelectedProblemId}\OperatorTok{();}
        \ControlFlowTok{if} \OperatorTok{(}\NormalTok{problemId }\OperatorTok{\textless{}} \DecValTok{0}\OperatorTok{)} \OperatorTok{\{}
            \BuiltInTok{JOptionPane}\OperatorTok{.}\FunctionTok{showMessageDialog}\OperatorTok{(}\NormalTok{view}\OperatorTok{,} \StringTok{"请先选择一道题目。"}\OperatorTok{,}
                \StringTok{"提示"}\OperatorTok{,} \BuiltInTok{JOptionPane}\OperatorTok{.}\FunctionTok{WARNING\_MESSAGE}\OperatorTok{);}
            \ControlFlowTok{return}\OperatorTok{;}
        \OperatorTok{\}}

        \BuiltInTok{String}\NormalTok{ lang }\OperatorTok{=}\NormalTok{ view}\OperatorTok{.}\FunctionTok{getSelectedLang}\OperatorTok{();}
        \BuiltInTok{String}\NormalTok{ src  }\OperatorTok{=}\NormalTok{ view}\OperatorTok{.}\FunctionTok{getSourceCode}\OperatorTok{();}
        \ControlFlowTok{if} \OperatorTok{(}\NormalTok{src}\OperatorTok{.}\FunctionTok{isBlank}\OperatorTok{())} \OperatorTok{\{}
            \BuiltInTok{JOptionPane}\OperatorTok{.}\FunctionTok{showMessageDialog}\OperatorTok{(}\NormalTok{view}\OperatorTok{,} \StringTok{"源码不能为空。"}\OperatorTok{,}
                \StringTok{"提示"}\OperatorTok{,} \BuiltInTok{JOptionPane}\OperatorTok{.}\FunctionTok{WARNING\_MESSAGE}\OperatorTok{);}
            \ControlFlowTok{return}\OperatorTok{;}
        \OperatorTok{\}}

        \DataTypeTok{int}\NormalTok{ subId }\OperatorTok{=}\NormalTok{ submissionIdGen}\OperatorTok{.}\FunctionTok{getAndIncrement}\OperatorTok{();}
        \DataTypeTok{byte}\OperatorTok{[]}\NormalTok{ srcBytes }\OperatorTok{=}\NormalTok{ src}\OperatorTok{.}\FunctionTok{getBytes}\OperatorTok{(}\NormalTok{StandardCharsets}\OperatorTok{.}\FunctionTok{UTF\_8}\OperatorTok{);}
        \DataTypeTok{final}\NormalTok{ JudgeTask task }\OperatorTok{=} \KeywordTok{new} \FunctionTok{JudgeTask}\OperatorTok{(}
\NormalTok{            subId}\OperatorTok{,}\NormalTok{ problemId}\OperatorTok{,}\NormalTok{ USER\_NAME}\OperatorTok{,}\NormalTok{ lang}\OperatorTok{,}\NormalTok{ srcBytes}\OperatorTok{);}

        \CommentTok{// 2. 禁用提交按钮，显示等待状态}
\NormalTok{        view}\OperatorTok{.}\FunctionTok{getSubmitBtn}\OperatorTok{().}\FunctionTok{setEnabled}\OperatorTok{(}\KeywordTok{false}\OperatorTok{);}
\NormalTok{        view}\OperatorTok{.}\FunctionTok{getResultLabel}\OperatorTok{().}\FunctionTok{setText}\OperatorTok{(}\StringTok{"评测中…"}\OperatorTok{);}

        \CommentTok{// 3. 后台线程：入队 + 等待结果}
        \BuiltInTok{SwingWorker}\OperatorTok{\textless{}}\NormalTok{JudgeResult}\OperatorTok{,} \BuiltInTok{Void}\OperatorTok{\textgreater{}}\NormalTok{ worker }\OperatorTok{=} \KeywordTok{new} \BuiltInTok{SwingWorker}\OperatorTok{\textless{}\textgreater{}()} \OperatorTok{\{}
            \AttributeTok{@Override}
            \KeywordTok{protected}\NormalTok{ JudgeResult }\FunctionTok{doInBackground}\OperatorTok{()} \KeywordTok{throws} \BuiltInTok{Exception} \OperatorTok{\{}
\NormalTok{                queue}\OperatorTok{.}\FunctionTok{put}\OperatorTok{(}\NormalTok{task}\OperatorTok{);}     \CommentTok{// 若队列满则挂起（背压）}
                \ControlFlowTok{return}\NormalTok{ task}\OperatorTok{.}\FunctionTok{await}\OperatorTok{();} \CommentTok{// 等待工作线程回填结果}
            \OperatorTok{\}}

            \AttributeTok{@Override}
            \KeywordTok{protected} \DataTypeTok{void} \FunctionTok{done}\OperatorTok{()} \OperatorTok{\{}  \CommentTok{// 回到 EDT}
\NormalTok{                view}\OperatorTok{.}\FunctionTok{getSubmitBtn}\OperatorTok{().}\FunctionTok{setEnabled}\OperatorTok{(}\KeywordTok{true}\OperatorTok{);}
                \ControlFlowTok{try} \OperatorTok{\{}
\NormalTok{                    JudgeResult r }\OperatorTok{=} \FunctionTok{get}\OperatorTok{();}
\NormalTok{                    view}\OperatorTok{.}\FunctionTok{getResultLabel}\OperatorTok{().}\FunctionTok{setText}\OperatorTok{(}\NormalTok{r}\OperatorTok{.}\FunctionTok{toString}\OperatorTok{());}
\NormalTok{                    view}\OperatorTok{.}\FunctionTok{appendHistory}\OperatorTok{(}
\NormalTok{                        task}\OperatorTok{.}\FunctionTok{getSubmissionId}\OperatorTok{(),}\NormalTok{ task}\OperatorTok{.}\FunctionTok{getProblemId}\OperatorTok{(),}
\NormalTok{                        task}\OperatorTok{.}\FunctionTok{getLang}\OperatorTok{(),}\NormalTok{ r}\OperatorTok{.}\FunctionTok{getStatus}\OperatorTok{().}\FunctionTok{name}\OperatorTok{(),}
\NormalTok{                        r}\OperatorTok{.}\FunctionTok{getPassed}\OperatorTok{()} \OperatorTok{+} \StringTok{"/"} \OperatorTok{+}\NormalTok{ r}\OperatorTok{.}\FunctionTok{getTotal}\OperatorTok{(),}
\NormalTok{                        r}\OperatorTok{.}\FunctionTok{getElapsedMs}\OperatorTok{());}
                    \ControlFlowTok{if} \OperatorTok{(!}\NormalTok{r}\OperatorTok{.}\FunctionTok{isAccepted}\OperatorTok{())} \OperatorTok{\{}
                        \BuiltInTok{JOptionPane}\OperatorTok{.}\FunctionTok{showMessageDialog}\OperatorTok{(}
\NormalTok{                            view}\OperatorTok{,}
\NormalTok{                            r}\OperatorTok{.}\FunctionTok{getDetail}\OperatorTok{(),}
\NormalTok{                            r}\OperatorTok{.}\FunctionTok{getStatus}\OperatorTok{().}\FunctionTok{name}\OperatorTok{(),}
                            \BuiltInTok{JOptionPane}\OperatorTok{.}\FunctionTok{WARNING\_MESSAGE}\OperatorTok{);}
                    \OperatorTok{\}}
                \OperatorTok{\}} \ControlFlowTok{catch} \OperatorTok{(}\BuiltInTok{Exception}\NormalTok{ ex}\OperatorTok{)} \OperatorTok{\{}
\NormalTok{                    view}\OperatorTok{.}\FunctionTok{getResultLabel}\OperatorTok{().}\FunctionTok{setText}\OperatorTok{(}\StringTok{"提交失败："} \OperatorTok{+}\NormalTok{ ex}\OperatorTok{.}\FunctionTok{getMessage}\OperatorTok{());}
                \OperatorTok{\}}
            \OperatorTok{\}}
        \OperatorTok{\};}
\NormalTok{        worker}\OperatorTok{.}\FunctionTok{execute}\OperatorTok{();}
    \OperatorTok{\}}
\OperatorTok{\}}
\end{Highlighting}
\end{Shaded}

\begin{center}\rule{0.5\linewidth}{0.5pt}\end{center}

\hypertarget{oj.gui.ojframe-ux754cux9762ux9aa8ux67b6}{%
\subsubsection{oj.gui.OjFrame ---
界面骨架}\label{oj.gui.ojframe-ux754cux9762ux9aa8ux67b6}}

OjFrame 沿用 08 的 API（已含
\texttt{getResultLabel()}），本章不新增字段、不重命名
getter。控制器需要的方法就是 08
那一套：\texttt{getSubmitBtn()}、\texttt{getSelectedProblemId()}、\texttt{getSelectedLang()}、\texttt{getSourceCode()}、\texttt{getResultLabel()}、\texttt{loadProblems(...)}、\texttt{appendHistory(...)}。

\begin{center}\rule{0.5\linewidth}{0.5pt}\end{center}

\hypertarget{ux52a8ux624bux5b9eux73b0ux4e24ux6b65ux8d70javac-java-ux5b9eux64cd-6}{%
\subsection{动手实现:两步走(javac / java
实操)}\label{ux52a8ux624bux5b9eux73b0ux4e24ux6b65ux8d70javac-java-ux5b9eux64cd-6}}

\hypertarget{ux7b2cux4e00ux6b65-ux6700ux7b80ux5e76ux53d1ux5224ux9898thread-synchronized-join}{%
\subsubsection{\texorpdfstring{第一步 \KS --- 最简并发判题:Thread +
synchronized +
join}{第一步 --- 最简并发判题:Thread + synchronized + join}}\label{ux7b2cux4e00ux6b65-ux6700ux7b80ux5e76ux53d1ux5224ux9898thread-synchronized-join}}

\textbf{\texttt{src/oj/Demo6a.java}}

\begin{Shaded}
\begin{Highlighting}[]
\KeywordTok{package}\ImportTok{ oj}\OperatorTok{;}
\KeywordTok{public} \KeywordTok{class}\NormalTok{ Demo6a }\OperatorTok{\{}
    \DataTypeTok{static} \DataTypeTok{int}\NormalTok{ done }\OperatorTok{=} \DecValTok{0}\OperatorTok{;}
    \DataTypeTok{static} \KeywordTok{synchronized} \DataTypeTok{void} \FunctionTok{finish}\OperatorTok{(}\BuiltInTok{String}\NormalTok{ who}\OperatorTok{)} \OperatorTok{\{}     \CommentTok{// 同步:保护共享计数}
\NormalTok{        done}\OperatorTok{++;}
        \BuiltInTok{System}\OperatorTok{.}\FunctionTok{out}\OperatorTok{.}\FunctionTok{println}\OperatorTok{(}\NormalTok{who }\OperatorTok{+} \StringTok{" 判完,累计 "} \OperatorTok{+}\NormalTok{ done}\OperatorTok{);}
    \OperatorTok{\}}
    \KeywordTok{public} \DataTypeTok{static} \DataTypeTok{void} \FunctionTok{main}\OperatorTok{(}\BuiltInTok{String}\OperatorTok{[]}\NormalTok{ args}\OperatorTok{)} \KeywordTok{throws} \BuiltInTok{InterruptedException} \OperatorTok{\{}
        \BuiltInTok{Runnable}\NormalTok{ job }\OperatorTok{=} \OperatorTok{()} \OperatorTok{{-}\textgreater{}} \OperatorTok{\{}
            \ControlFlowTok{try} \OperatorTok{\{} \BuiltInTok{Thread}\OperatorTok{.}\FunctionTok{sleep}\OperatorTok{(}\DecValTok{50}\OperatorTok{);} \OperatorTok{\}} \ControlFlowTok{catch} \OperatorTok{(}\BuiltInTok{InterruptedException}\NormalTok{ e}\OperatorTok{)} \OperatorTok{\{} \OperatorTok{\}}
            \FunctionTok{finish}\OperatorTok{(}\BuiltInTok{Thread}\OperatorTok{.}\FunctionTok{currentThread}\OperatorTok{().}\FunctionTok{getName}\OperatorTok{());}
        \OperatorTok{\};}
        \BuiltInTok{Thread}\NormalTok{ t1 }\OperatorTok{=} \KeywordTok{new} \BuiltInTok{Thread}\OperatorTok{(}\NormalTok{job}\OperatorTok{,} \StringTok{"worker{-}1"}\OperatorTok{);}
        \BuiltInTok{Thread}\NormalTok{ t2 }\OperatorTok{=} \KeywordTok{new} \BuiltInTok{Thread}\OperatorTok{(}\NormalTok{job}\OperatorTok{,} \StringTok{"worker{-}2"}\OperatorTok{);}
\NormalTok{        t1}\OperatorTok{.}\FunctionTok{start}\OperatorTok{();}\NormalTok{ t2}\OperatorTok{.}\FunctionTok{start}\OperatorTok{();}                        \CommentTok{// 两个判题线程并发}
\NormalTok{        t1}\OperatorTok{.}\FunctionTok{join}\OperatorTok{();}\NormalTok{  t2}\OperatorTok{.}\FunctionTok{join}\OperatorTok{();}                          \CommentTok{// 主线程等它们判完}
        \BuiltInTok{System}\OperatorTok{.}\FunctionTok{out}\OperatorTok{.}\FunctionTok{println}\OperatorTok{(}\StringTok{"全部完成: "} \OperatorTok{+}\NormalTok{ done}\OperatorTok{);}
    \OperatorTok{\}}
\OperatorTok{\}}
\end{Highlighting}
\end{Shaded}

讲解:\texttt{Runnable} + \texttt{Thread.start()}
并发、\texttt{synchronized} 保护共享变量、\texttt{join()}
让主线程等待------Ch12 考点,模拟 OJ 多份提交并发判题。

编译运行:

\begin{Shaded}
\begin{Highlighting}[]
\ExtensionTok{javac} \AttributeTok{{-}d}\NormalTok{ build src/oj/Demo6a.java}
\ExtensionTok{java} \AttributeTok{{-}cp}\NormalTok{ build oj.Demo6a}
\CommentTok{\# worker{-}1 判完,累计 1}
\CommentTok{\# worker{-}2 判完,累计 2   (顺序可能不同)}
\CommentTok{\# 全部完成: 2}
\end{Highlighting}
\end{Shaded}

\hypertarget{ux7b2cux4e8cux6b65-ux6cdbux578bux963bux585eux961fux5217-ux5de5ux4f5cux7ebfux7a0bux6c60-ux8c03-c-ux5224ux9898ux673a}{%
\subsubsection{\texorpdfstring{第二步 \GC --- 泛型阻塞队列 + 工作线程池
+ 调 C++
判题机}{第二步 --- 泛型阻塞队列 + 工作线程池 + 调 C++ 判题机}}\label{ux7b2cux4e8cux6b65-ux6cdbux578bux963bux585eux961fux5217-ux5de5ux4f5cux7ebfux7a0bux6c60-ux8c03-c-ux5224ux9898ux673a}}

\texttt{JudgeTask}(\texttt{wait/notify}
回填结果)、\texttt{JudgeQueue\textless{}T\textgreater{}}(泛型阻塞队列)、\texttt{JudgeWorker}(N
个守护线程消费队列、落临时文件、调 \texttt{MachineJudge})、GUI 用
\texttt{SwingWorker}(代码见上文【新产物架构】):

\begin{Shaded}
\begin{Highlighting}[]
\BuiltInTok{cd}\NormalTok{ judge }\KeywordTok{\&\&} \ExtensionTok{g++} \AttributeTok{{-}O2} \AttributeTok{{-}std}\OperatorTok{=}\NormalTok{c++17 }\AttributeTok{{-}o}\NormalTok{ judge judge.cpp }\KeywordTok{\&\&} \BuiltInTok{cd}\NormalTok{ ..}
\ExtensionTok{javac} \AttributeTok{{-}cp} \StringTok{"lib/*"} \AttributeTok{{-}d}\NormalTok{ build }\VariableTok{$(}\FunctionTok{find}\NormalTok{ src }\AttributeTok{{-}name} \StringTok{\textquotesingle{}*.java\textquotesingle{}}\VariableTok{)}
\ExtensionTok{java} \AttributeTok{{-}cp} \StringTok{"build:lib/*"}\NormalTok{ oj.gui.Main}
\end{Highlighting}
\end{Shaded}

\hypertarget{ux9a8cux6536ux6807ux51c6-8}{%
\subsection{验收标准}\label{ux9a8cux6536ux6807ux51c6-8}}

\begin{enumerate}
\def\labelenumi{\arabic{enumi}.}
\tightlist
\item
  \textbf{并发提交不阻塞
  EDT}：同时点击提交按钮两次，两份代码均能在队列中排队，状态标签''评测中\ldots``正常显示，界面不卡顿。
\item
  \textbf{死循环源码返回 TLE}：提交 \texttt{while(true)\{\}} 的 Java 或
  C++ 代码，C++ 判题机在配置的时间限制后返回 \texttt{TLE}，Java
  侧正常展示结果，不挂死。
\item
  \textbf{工作线程数量正确}：用调试输出或线程 dump 确认恰好有
  \texttt{WORKER\_COUNT}（默认 4）个名为 \texttt{judge-worker-N}
  的守护线程在运行，且它们复用构造器传入的同一个 \texttt{dao} 与
  \texttt{judge}，没有任何 \texttt{new\ MachineJudge()} /
  \texttt{new\ SubmissionDao()}。
\item
  \textbf{结果最终一致}：任意一次提交的 \texttt{JudgeResult}，在
  \texttt{OjFrame} 展示的内容与 \texttt{SubmissionDao}
  持久化到数据库的内容完全一致；源文件落在
  \texttt{submissions/sub\textless{}id\textgreater{}.\textless{}ext\textgreater{}}
  且未被删除，路径存进 Submission 的 \texttt{srcPath}。
\item
  \textbf{异常不崩溃工作线程}：向工作线程注入一个不存在的题目
  ID，验证捕获 \texttt{Exception}
  后工作线程仍然存活并继续处理后续任务，对应任务返回 \texttt{RE} 状态。
\item
  \textbf{背压生效}：将 \texttt{QUEUE\_CAPACITY} 设为 2，同时发起 5
  次提交，第 3--5 次的 \texttt{SwingWorker.doInBackground()} 应阻塞在
  \texttt{queue.put()}，不丢任务也不抛异常。
\item
  \textbf{getter 与 08 一致}：OjController 全程使用 08 定义的 OjFrame
  API（\texttt{getSubmitBtn} / \texttt{getSelectedProblemId} /
  \texttt{getSelectedLang} / \texttt{getSourceCode} /
  \texttt{getResultLabel}），题号走
  \texttt{getSelectedProblemId()}，不把下拉当
  \texttt{JComboBox\textless{}Integer\textgreater{}} 强转；OjController
  构造签名为
  \texttt{(OjFrame,\ ProblemService,\ MachineJudge,\ SubmissionDao)}。
\item
  \textbf{项目整体可运行}：从 M1 到 M6a 的全部测试用例（含 M4
  单文件配置、M5a 外部判题机、M5b 数据库持久化、M5c Swing
  界面）在同一个项目内编译通过并正确运行，不存在孤立分支或未接入的模块。
\end{enumerate}

\begin{center}\rule{0.5\linewidth}{0.5pt}\end{center}

\begin{quote}
\textbf{项目收官说明} M6a 是 Mini-OJ
的最后一个里程碑。至此，系统完整地串联了面向对象设计（M1--M3）、反射与配置（M4）、外部进程与文件
IO（M5a）、数据库持久化（M5b）、Swing
GUI（M5c），以及本章的多线程并发（M6a）。Ch13
网络编程不纳入本项目，\texttt{JudgeServer}/\texttt{JudgeClient}/\texttt{Socket}/\texttt{ServerSocket}
均不出现在任何章节中。同学们掌握了这条主线，便已综合运用了一学期 Java
课程的核心知识。
\end{quote}

\hypertarget{mini-oj-ux7b2cux4e09ux4ee3ux5224ux9898ux673a-judge}{%
\section{Mini-OJ 第三代判题机
(judge)}\label{mini-oj-ux7b2cux4e09ux4ee3ux5224ux9898ux673a-judge}}

单文件、编译型的判题机:\textbf{真编译、真运行、真限资源},跑全部测试点,比对输出,输出一行
\textbf{JSON} 判定。 是
Mini-OJ「判题机三代演进」里的\textbf{第三代};Java
OJ(\texttt{MachineJudge})只负责 \texttt{ProcessBuilder} 调用它 +
正则解析结果。

\hypertarget{ux80fdux529bux4e00ux89c8}{%
\subsection{能力一览}\label{ux80fdux529bux4e00ux89c8}}

\begin{itemize}
\tightlist
\item
  \textbf{多语言}:\texttt{cpp} / \texttt{c} / \texttt{python} /
  \texttt{java}(自动编译 / 语法检查 / 构造运行命令)。
\item
  \textbf{完整判定}:\texttt{AC\ /\ WA\ /\ PE\ /\ TLE\ /\ MLE\ /\ RE\ /\ CE\ /\ ERR}。
\item
  \textbf{资源限制}:\texttt{setrlimit}(CPU/AS/输出/栈)+ 父进程墙钟看门狗
  + \textbf{进程组整组 SIGKILL}(连同子孙进程一起掐);死循环必出
  \texttt{TLE}、爆内存必出 \texttt{MLE}。
\item
  \textbf{编译保护}:编译也有墙钟上限(g++ 10s / javac
  15s),防模板炸弹挂死。
\item
  \textbf{比对}:逐行去尾空白 + 去末尾空行;\texttt{-\/-special}
  开浮点容差(1e-6,逐 token);仅空白排布不同判 \texttt{PE}。
\item
  \textbf{健壮}:缺文件 / 无用例 / 编译器缺失 / 异常输出
  都有明确状态,绝不崩。
\end{itemize}

\hypertarget{ux7f16ux8bd1}{%
\subsection{编译}\label{ux7f16ux8bd1}}

\begin{Shaded}
\begin{Highlighting}[]
\FunctionTok{make}            \CommentTok{\# 等价于 g++ {-}O2 {-}std=c++17 {-}o judge judge.cpp}
\end{Highlighting}
\end{Shaded}

\hypertarget{ux7528ux6cd5}{%
\subsection{用法}\label{ux7528ux6cd5}}

\begin{Shaded}
\begin{Highlighting}[]
\ExtensionTok{./judge} \AttributeTok{{-}{-}problem} \OperatorTok{\textless{}}\NormalTok{题目目录}\OperatorTok{\textgreater{}}\NormalTok{ {-}{-}src }\OperatorTok{\textless{}}\NormalTok{提交源码}\OperatorTok{\textgreater{}}\NormalTok{ {-}{-}lang }\OperatorTok{\textless{}}\NormalTok{cpp}\KeywordTok{|}\ExtensionTok{c}\KeywordTok{|}\ExtensionTok{python}\KeywordTok{|}\ExtensionTok{java}\OperatorTok{\textgreater{}} \DataTypeTok{\textbackslash{}}
\NormalTok{        [{-}{-}time{-}ms 1000] [{-}{-}mem{-}mb 256] }\PreprocessorTok{[{-}{-}}\SpecialStringTok{special}\PreprocessorTok{]}
\CommentTok{\# {-}{-}time / {-}{-}mem 是 {-}{-}time{-}ms / {-}{-}mem{-}mb 的别名}
\end{Highlighting}
\end{Shaded}

题目目录约定(成对、按数字排序;非 \texttt{.in/.out} 文件忽略):

\begin{verbatim}
problems/<id>/
+-- 1.in   1.out
+-- 2.in   2.out
+-- 3.in   3.out
\end{verbatim}

\hypertarget{ux8f93ux51fastdout-ux4e00ux884c-json}{%
\subsection{输出(stdout 一行
JSON)}\label{ux8f93ux51fastdout-ux4e00ux884c-json}}

\begin{Shaded}
\begin{Highlighting}[]
\FunctionTok{\{}\DataTypeTok{"status"}\FunctionTok{:}\StringTok{"AC"}\FunctionTok{,}\DataTypeTok{"passed"}\FunctionTok{:}\DecValTok{3}\FunctionTok{,}\DataTypeTok{"total"}\FunctionTok{:}\DecValTok{3}\FunctionTok{,}\DataTypeTok{"time\_ms"}\FunctionTok{:}\DecValTok{3}\FunctionTok{,}\DataTypeTok{"mem\_kb"}\FunctionTok{:}\DecValTok{1840}\FunctionTok{,}\DataTypeTok{"detail"}\FunctionTok{:}\StringTok{""}\FunctionTok{\}}
\end{Highlighting}
\end{Shaded}

\begin{longtable}[]{@{}ll@{}}
\toprule\noalign{}
字段 & 含义 \\
\midrule\noalign{}
\endhead
\bottomrule\noalign{}
\endlastfoot
\texttt{status} &
\texttt{AC}/\texttt{WA}/\texttt{PE}/\texttt{TLE}/\texttt{MLE}/\texttt{RE}/\texttt{CE}/\texttt{ERR} \\
\texttt{passed} / \texttt{total} & 通过测试点数 / 总数 \\
\texttt{time\_ms} & 最大测试点 CPU 耗时 \\
\texttt{mem\_kb} & 峰值内存(\texttt{ru\_maxrss}) \\
\texttt{detail} & 失败说明(WA/PE/RE 的 case 号、CE 的编译错误等) \\
\end{longtable}

状态:AC 通过 · WA 答案错 · \textbf{PE 格式错(仅空白差异)} · TLE 超时 ·
MLE 超内存 · RE 运行错(崩溃/非零退出) · CE 编译/语法错 · ERR
判题机内部错。

\hypertarget{ux4f8bux5b50}{%
\subsection{例子}\label{ux4f8bux5b50}}

\begin{Shaded}
\begin{Highlighting}[]
\FunctionTok{make}\NormalTok{ test       }\CommentTok{\# 跑全部示例:AC/WA/TLE/RE/CE/MLE/PE + python + 浮点 special}
\ExtensionTok{./judge} \AttributeTok{{-}{-}problem}\NormalTok{ examples/problems/aplusb }\AttributeTok{{-}{-}src}\NormalTok{ examples/solutions/ac.cpp  }\AttributeTok{{-}{-}lang}\NormalTok{ cpp}
\ExtensionTok{./judge} \AttributeTok{{-}{-}problem}\NormalTok{ examples/problems/avg    }\AttributeTok{{-}{-}src}\NormalTok{ examples/solutions/avg.cpp }\AttributeTok{{-}{-}lang}\NormalTok{ cpp }\AttributeTok{{-}{-}special}
\end{Highlighting}
\end{Shaded}

\hypertarget{ux653eux7f6eux4f4dux7f6eux96c6ux6210ux8fdb-mini-oj-ux9879ux76ee}{%
\subsection{放置位置(集成进 Mini-OJ
项目)}\label{ux653eux7f6eux4f4dux7f6eux96c6ux6210ux8fdb-mini-oj-ux9879ux76ee}}

\begin{verbatim}
mini-oj/
+-- src/oj/...           # Java 代码(MachineJudge 在 oj.judge)
+-- problems/<id>/       # 题目:config.txt + N.in/N.out
+-- judge/
    +-- judge.cpp
    +-- judge            # make 出来的二进制(.gitignore)
\end{verbatim}

\hypertarget{java-oj-ux5982ux4f55ux8c03ux7528ux7b2cux4e09ux4ee3ux6838ux5fc3processbuilder-ux6b63ux5219ux89e3ux6790}{%
\subsection{Java OJ 如何调用(第三代核心:ProcessBuilder +
正则解析)}\label{java-oj-ux5982ux4f55ux8c03ux7528ux7b2cux4e09ux4ee3ux6838ux5fc3processbuilder-ux6b63ux5219ux89e3ux6790}}

\begin{Shaded}
\begin{Highlighting}[]
\CommentTok{// oj.judge.MachineJudge}
\BuiltInTok{Process}\NormalTok{ p }\OperatorTok{=} \KeywordTok{new} \BuiltInTok{ProcessBuilder}\OperatorTok{(}
    \StringTok{"judge/judge"}\OperatorTok{,} \StringTok{"{-}{-}problem"}\OperatorTok{,} \StringTok{"problems/1"}\OperatorTok{,} \StringTok{"{-}{-}src"}\OperatorTok{,} \StringTok{"sub.cpp"}\OperatorTok{,} \StringTok{"{-}{-}lang"}\OperatorTok{,} \StringTok{"cpp"}\OperatorTok{,}
    \StringTok{"{-}{-}time{-}ms"}\OperatorTok{,} \StringTok{"1000"}\OperatorTok{,} \StringTok{"{-}{-}mem{-}mb"}\OperatorTok{,} \StringTok{"256"}
\OperatorTok{).}\FunctionTok{start}\OperatorTok{();}
\BuiltInTok{String}\NormalTok{ json }\OperatorTok{=} \KeywordTok{new} \BuiltInTok{String}\OperatorTok{(}\NormalTok{p}\OperatorTok{.}\FunctionTok{getInputStream}\OperatorTok{().}\FunctionTok{readAllBytes}\OperatorTok{(),}\NormalTok{ StandardCharsets}\OperatorTok{.}\FunctionTok{UTF\_8}\OperatorTok{);}
\NormalTok{p}\OperatorTok{.}\FunctionTok{waitFor}\OperatorTok{();}
\CommentTok{// 用正则从 json 抽 status/passed/total/time\_ms,构造 oj.core.JudgeResult}
\end{Highlighting}
\end{Shaded}

完整封装见教程 \textbf{M5a · 第三代判题机} 一章的
\texttt{MachineJudge}。

\hypertarget{ux9650ux5236ux4e0eux8fb9ux754cux6559ux5b66ux5355ux673aux7248}{%
\subsection{限制与边界(教学/单机版)}\label{ux9650ux5236ux4e0eux8fb9ux754cux6559ux5b66ux5355ux673aux7248}}

\begin{itemize}
\tightlist
\item
  用 \texttt{setrlimit} + 进程组隔离,但\textbf{未做完整沙箱}(无
  seccomp/namespace/chroot、不禁系统调用、不隔离文件系统)。\textbf{只在你信任的本地/实验室环境运行提交。}
\item
  内存判定以 \texttt{ru\_maxrss} 为准;\texttt{RLIMIT\_AS} 只对
  \texttt{c/cpp} 设硬上限(\texttt{java}/\texttt{python}
  解释器预留巨量虚拟内存,设 AS 会误杀)。
\item
  CPU 时限到点可能由内核以 \texttt{SIGXCPU} 或硬限 \texttt{SIGKILL}
  触发,判题机据''CPU 已超限''统一归 \texttt{TLE}。
\item
  \texttt{java} 提交的 public 类名须为
  \texttt{Main}(判题机会把源码复制成 \texttt{Main.java} 再
  \texttt{javac})。
\end{itemize}

\end{document}
